\documentclass[11pt]{report}
\usepackage[utf8]{inputenc}

% Colors
\usepackage[dvipsnames]{xcolor}

\definecolor{C0}{HTML}{1d1d1d}
\definecolor{C1}{HTML}{1e3668}
\definecolor{C2}{HTML}{199d8b}
\definecolor{C3}{HTML}{d52f4c}
\definecolor{C4}{HTML}{5ab2d6}
\definecolor{C5}{HTML}{ffb268}

\color{C0}

% fonts
\usepackage[no-math]{fontspec}

\emergencystretch=8pt
\hyphenpenalty=1000 % default 50
\tolerance=800      % default 200
%\righthyphenmin=4
%\lefthyphenmin=4

\setmainfont[
    BoldFont = Vollkorn Bold,
    ItalicFont = Vollkorn Italic,
    BoldItalicFont={Vollkorn Bold Italic},
    RawFeature=+lnum,
]{Vollkorn}

\setsansfont[
    BoldFont = Lato Bold,
    FontFace={l}{n}{*-Light},
    FontFace={l}{it}{*-Light Italic},
]{Lato}

\usepackage[scale=.78]{luatexja-fontspec}
\setmainjfont{SourceHanSerifSC-Bold}%{Source Han Serif CN VF Regular}


\usepackage{amsmath,amssymb,amsthm}
\usepackage{mathtools}
\newcommand\smop\smashoperator

\usepackage{unicode-math}
\mathitalicsmode=1

\setmathfont[
mathit = sym,
mathup = sym,
mathbf = sym,
math-style = TeX, 
bold-style = TeX
]{Wholegrain Math}
\everydisplay{\Umathoperatorsize\displaystyle=4ex}
\AtBeginDocument{\renewcommand\setminus{\smallsetminus}}
%\setmathfont{TeX Gyre Pagella Math}[range=\sum,Scale=1.5]


% section format
\usepackage{sectsty}
%\allsectionsfont{\normalfont\Vollkorn\bfseries}
\allsectionsfont{\normalfont\sffamily\bfseries}

\usepackage{titlesec}

\titleformat{\chapter}[display]{\sffamily\bfseries}{}{0pt}{\fontsize{22}{28}\selectfont {\fontseries{l}\sffamily Chapter \thechapter}\\\vspace{-.2em}}
\titleformat{\part}[display]{\sffamily\bfseries}{}{0pt}{\fontsize{22}{28}\selectfont {\fontseries{l}\sffamily Part \thepart}\\\vspace{10cm}}

\usepackage{tocloft}
\renewcommand{\cfttoctitlefont}{\fontsize{22}{22}\selectfont\sffamily\bfseries}


\def\myraggedright{%
  \rightskip=0pt plus 1fil
  \spaceskip .3333em \xspaceskip .5em\relax
}


% Bibliography
%\usepackage[backend=biber,style=alphabetic,maxalphanames=5,maxcitenames=99,maxbibnames=99,firstinits=true,url=false,eprint=false]{biblatex}
\usepackage[backend=biber,backref,style=alphabetic,maxcitenames=99,maxbibnames=99,firstinits=true,url=false]{biblatex}
\addbibresource{thesis.bib}
\defcounter{biburlnumpenalty}{3000}
\defcounter{biburlucpenalty}{6000}
\defcounter{biburllcpenalty}{9000}

\DefineBibliographyStrings{english}{%
  backrefpage = {cited on page},%
  backrefpages = {cited on pages}%
}

% Page Formatting
\usepackage[
    paper=b5paper,
    inner=22mm,         % Inner margin
    outer=22mm,         % Outer margin
    bindingoffset=0mm, % Binding offset
    top=28mm,           % Top margin
    bottom=22mm,        % Bottom margin
    %showframe,         % show how the type block is set on the page
]{geometry}

\setlength{\parindent}{0em}
\setlength{\parskip}{.7em}


\usepackage{afterpage}
\usepackage{pdflscape}

\usepackage{enumitem}
\renewcommand{\labelitemi}{\normalfont\bfseries\textendash}

\usepackage{appendix}
\usepackage{setspace}

% header
\usepackage{fancyhdr}
%\renewcommand{\headrulewidth}{0pt}

\fancypagestyle{body}{%
    \fancyhf{}
    \fancyhead[R]{{\sffamily page \textbf{\thepage}}}
    \fancyhead[L]{{\sffamily chapter \textbf{\thechapter}}}%~section \thesection}}
}

\fancypagestyle{front}{%
    \fancyhf{}
    \fancyhead[R]{{\sffamily page \textbf{\thepage}}}
}
\setlength{\headheight}{14pt}

\usepackage{etoolbox}
\patchcmd{\chapter}{\thispagestyle{plain}}{\thispagestyle{front}}{}{}

\AtBeginDocument{\addtocontents{toc}{\protect\thispagestyle{front}}}

\makeatletter
%\setlength{\cftbeforefigskip}{16pt}
% \patchcmd{<cmd>}{<search>}{<replace>}{<succes>}{<failure>}
\patchcmd{\@chapter}{\addtocontents{lof}{\protect\addvspace{10\p@}}}{}{}{}% LoF
\patchcmd{\@chapter}{\addtocontents{lot}{\protect\addvspace{10\p@}}}{}{}{}% LoT
\makeatother

\setlength{\cftbeforesecskip}{6pt}
\setlength{\cftbeforesubsecskip}{6pt}

% Algorithms
\usepackage{algorithm}
\AtBeginEnvironment{algorithmic}{\setstretch{1.25}}
\counterwithin{algorithm}{chapter}
\usepackage{algorithmicx}
\usepackage[noend]{algpseudocode}

\algrenewcommand\algorithmicindent{1.0em}
\let\Algorithm\algorithm
\renewcommand\algorithm[1][]{\Algorithm[#1]}%\fontsize{11}{16}\selectfont}

\newenvironment{labelalgorithm}[4][t]{%
\begin{algorithm}[#1]
%\newcommand{\thealgorithmlabel}{#2}
\newcommand{\thealgorithmname}{#3}
%\newcommand{\thealgorithmcap}{#4}
\customlabel{alg:name:#2}{\textproc{#3}}
%\customlabel{alg:cap:#2}{#4}
\caption{#4}\label{alg:#2}
}{\end{algorithm}}


\makeatletter
\newcommand{\customlabel}[2]{%
   \protected@write \@auxout {}{\string \newlabel {#1}{{#2}{\thepage}{#2}{#1}{}} }%
   \hypertarget{#1}{}%
}
\makeatother


%\algdef{SE}[FUNCTION]{Procedure}{EndProcedure}%
%   [2]{\algorithmicclass\ \textproc{#1}\ifthenelse{\equal{#2}{}}{}{$($#2$)$}}%
%   {\algorithmicend\ \algorithmicclass}%

\algnewcommand\algorithmicclass{\textbf{class}}
\algdef{SE}[FUNCTION]{Class}{EndClass}%
   [2]{\algorithmicclass\ \textproc{#1}\ifthenelse{\equal{#2}{}}{}{$($#2$)$}}%
   {\algorithmicend\ \algorithmicclass}%

% Tells algorithmicx not to print an empty line if `noend' is set 
\makeatletter
\ifthenelse{\equal{\ALG@noend}{t}}%
  {\algtext*{EndClass}}
  {}%
\makeatother

%\usepackage{showlabels}

% Math 

\newtheorem{theorem}{Theorem}[chapter]
\newtheorem{assumption}[theorem]{Assumption}
\newtheorem{lemma}[theorem]{Lemma}
\newtheorem{corollary}[theorem]{Corollary}
\newtheorem{proposition}[theorem]{Proposition}

\newtheorem{definition}[theorem]{Definition}

\theoremstyle{definition}
\newtheorem{example}[theorem]{Example}
\newtheorem{remarkx}[theorem]{Remark}
\newenvironment{remark}
  {\pushQED{\qed}\renewcommand{\qedsymbol}{$\triangle$}\remarkx}
  {\popQED\endremarkx}


\usepackage{thmtools}
\usepackage{thm-restate}



% Graphics, tables, and captions

\usepackage{pgffor}
\usepackage{graphicx}
\usepackage{pgfplots}
\usepackage{subcaption}
\captionsetup{labelfont=bf}
\captionsetup{width=.9\textwidth}

\usepackage{tabularx}
\usepackage{ltablex}
\usepackage{booktabs}
\usepackage{array}

\usepackage{hyperref}
\hypersetup{
    colorlinks,
    linkcolor=C1,
    citecolor=C1,
    urlcolor=C1
}

\usepackage{pdfpages}
\usepackage[nameinlink,noabbrev,capitalize]{cleveref}

\crefname{ineq}{}{}
\crefname{equation}{}{}
\creflabelformat{ineq}{#2{\textup{(1)}}#3}
\creflabelformat{equation}{#2\textup{(#1)}#3}

\usepackage{lipsum}  


\allowdisplaybreaks
\newcommand{\note}[1]{{\color{Red}#1}}
\newcommand{\TODO}[1]{{\color{blue}TODO: #1}}

\newcommand{\tableemph}{\textbf}
%\newcommand{\tableemph}{*}
%\newcommand{\tableemph}{\color{Red}}


%\usepackage{dsfont}
\newcommand{\bOne}{\mathbb{1}}

%\usepackage{bm}
%\renewcommand{\vec}{\mathbf}
\renewcommand{\vec}{\symbfup}
\newcommand{\ms}{\texttt}

\DeclareMathOperator*{\argmin}{argmin}


\newcommand{\cT}{\textsf{\textup{\fontsize{7}{7}\selectfont{H}}}}
\newcommand{\rT}{\textsf{\textup{\fontsize{7}{7}\selectfont{T}}}}
\newcommand{\T}{\textsf{\textup{\color{red}T}}}

\renewcommand{\d}{\textup{d}}
\newcommand{\tr}{\operatorname{tr}}

\newcommand{\R}{\mathbb{R}}
\newcommand{\F}{{\textsf{\textup{F}}}}

\newcommand{\EE}{\mathbb{E}}
\newcommand{\PP}{\mathbb{P}}

\newcommand{\W}{d_{\mathrm{W}}}
\newcommand{\KS}{d_{\mathrm{KS}}}
\newcommand{\TV}{d_{\mathrm{TV}}}

\newcommand{\qq}[3][]{%
    \ifstrempty{#1}{%
        [#2]_{#3}^{\circ\textup{q}}
    }{%
        [#2]_{#3}^{\textup{{#1}q}}
    }%
}

\newcommand{\ff}[3][]{%
    \ifstrempty{#1}{%
        [#2]_{#3}^{\circ\textup{p}}
    }{%
        [#2]_{#3}^{\textup{{#1}p}}
    }%
}


\newcommand{\nv}{n_{\textup{v}}}
\newcommand{\lmin}{\lambda_{\textup{min}}}
\newcommand{\lmax}{\lambda_{\textup{max}}}

\newcommand{\ii}{\symbfit{i}}

\newcommand{\samp}[2][]{#1\langle #2 #1\rangle}

\newcommand{\mf}[1]{(#1)}
\newcommand{\A}{\mf{\vec{A}}}
\newcommand{\fA}{f\A}
\newcommand{\pA}{p\A}
\newcommand{\fT}{f\mf{\vec{T}}}

\newcommand{\Qhat}{\mkern.5mu\widehat{\mkern-.5mu\vec{Q}\mkern-.5mu}\mkern.5mu}
    
\newcommand{\lan}{\textup{\textsf{lan-FA}}}
\newcommand{\err}{\textup{\textsf{err}}}
\newcommand{\res}{\textup{\textsf{res}}}
\newcommand{\Res}{\textup{\textsf{res}}}
\newcommand{\Opt}{\textup{\textsf{opt}}}

\newcommand{\lanopt}{\textup{\textsf{lan-OR}}}

\AtBeginDocument{\renewcommand{\Re}{\operatorname{Re}}}
\AtBeginDocument{\renewcommand{\Im}{\operatorname{Im}}}

% Document information
\usepackage{wallpaper}


\let\svthefootnote\thefootnote
\newcommand\freefootnote[1]{%
  \let\thefootnote\relax%
  \footnotetext{#1}%
  \let\thefootnote\svthefootnote%
}
\interfootnotelinepenalty=10000



\begin{document}

\fontsize{11}{16}\selectfont

\pagestyle{empty}
\renewcommand{\thepage}{Copyright}
\begin{center}
\textcopyright\,Copyright 2022
\\
Tyler Chen
\thispagestyle{empty}
\end{center}


\clearpage
\renewcommand{\thepage}{UW title page}
\begin{center}

\textbf{Lanczos-based  methods for matrix functions}

\vspace{3em}
Tyler Chen 

\vfill

\textit{A dissertation}
\\
submitted in partial fulfillment of the
\\
requirements for the degree of

\vspace{3em}
\textbf{Doctor of Philosophy}

\vspace{3em}
University of Washington
\\
2022

\vspace{3em}
Reading Committee:
\\Anne Greenbaum, Chair
\\Thomas Trogdon, Chair
\\Aleksandr Aravkin

\vfill
Program Authorized to Offer Degree:
\\
Applied Mathematics

\end{center}

\thispagestyle{empty}

\clearpage
\renewcommand{\thepage}{UW abstract}
\begin{center}
University of Washington

\vspace{3em}
    \textit{Abstract}

\vspace{3em}
\textbf{Lanczos-based  methods for matrix functions}

\vspace{3em}
Tyler Chen

\vspace{3em}

Chairs of the Supervisory Committee:
\\Anne Greenbaum
\\Thomas Trogdon

\vspace{3em}
Department of Applied Mathematics

\end{center}

\vfill
{\fontsize{11}{24}\selectfont
We study Lanczos-based methods for tasks involving matrix functions.
We begin by resurfacing a range of ideas regarding matrix-free quadrature which, to the best of our knowledge, have not been treated simultaneously. 
This enables the development of a unified perspective from which a number of commonly used randomized methods for spectrum and spectral sum approximation can be understood.
We proceed to develop optimal Krylov subspace methods for approximating the product of a rational matrix function with a fixed vector.
Finally, we show how the optimality of such methods can be used to obtain fine-grained spectrum dependent bounds for standard Lanczos-based methods for approximating a wide class of matrix functions applied to a vector.
}
\thispagestyle{empty}



\ThisULCornerWallPaper{1}{special_pages/UW_logo.pdf}
\begin{titlepage}
\renewcommand{\thepage}{Title page}
~\vfill~
\sffamily

\fontsize{19.2}{22}\selectfont
Lanczos-based  methods for matrix functions

\vspace{1.5cm}
\fontsize{26}{22}\selectfont
\textbf{TYLER CHEN}

\vspace{5cm}

\fontsize{10}{12}\selectfont
A dissertation submitted in partial fulfillment 
\\
of the requirements for the degree of

\textbf{Doctor of Philosophy}

\vspace{3cm}
2022
\\
University of Washington
\\
Department of Applied Mathematics

\end{titlepage}


\pagenumbering{roman}
\pagestyle{front}

\iffalse
\section*{Preface}

The sun has set over the Olympics on what is undoubtedly one of the last clear days of the year.
Seattle is beautiful in the summer, but I've always preferred the rain -- the persistent grey sidesteps the feeling of dissonance that accompanies sunny weather.
It is only fitting, then, that this document, which is the culmination of my time here, will be written over the long nights of a Pacific northwest winter.  

\vspace{.5em}
October 30, 2021
\vspace{-.7em}

Chinatown/International District
\vspace{-.7em}

Seattle, Washington

\begin{verse}
\textbf{Promise}

Long nights, short years. Forgiving
\\
silence.

When morning comes, and pain --

no one is a stranger, this whole world is your home. 
\end{verse}
\fi

\includepdf[]{special_pages/dedication_plain.pdf}

\section*{Preface}

At night, my apartment looks out to a thousand illuminated windows.
I'm drawn such views because they make me feel an isolating sense of closeness; behind each window is a person---a family enjoying dinner, a student working on their homework, a cleaning person ending their shift.
I do not know them and they do not know me, yet we are all connected in this moment of existence. 
This is \emph{sonder}, a concept for which we find a definition in the Dictionary of Obscure Sorrows:
\begin{quote}
\textbf{sonder}

{\textit n.} the realization that each random passerby is living a life as vivid and complex as your own%---populated with their own ambitions, friends, routines, worries and inherited craziness---an epic story that continues invisibly around you like an anthill sprawling deep underground, with elaborate passageways to thousands of other lives that you’ll never know existed, in which you might appear only once, as an extra sipping coffee in the background, as a blur of traffic passing on the highway, as a lighted window at dusk.
\end{quote}

Sonder, even with the accompanying melancholy, has been the single most consistent force driving my success throughout my tertiary education.
It only fitting, then, that it receives mention in my dissertation, the symbolic culmination of my formal education.


%In a sense, this thesis is the culmination of my formal education.

%Surprisingly minimal requirements regarding the formatting of 

\iffalse
I continually find myself chasing certain aesthetic emotions.
After hours, sonder is perhaps the 


emotional aesthetic,

I hope this thesis invokes such 

editorial control.

This is expressed in two main ways. First, exposition

Second typesetting

but also in the typesetting and 
\fi

\begin{flushright}
\vspace{5em}
Chinatown/International District
\vspace{-.7em}

Seattle, Washington %2022
\end{flushright}




\clearpage
\section*{Acknowledgements}


First and foremost, I would like to thank my thesis advisors, Anne Greenbaum and Thomas Trogdon, whose support and guidance throughout my PhD have shaped me into the researcher I am today.
Many PhD students struggle to find a good advisor, but I had two excellent ones.
Anne, thank you for sharing with me your vast wisdom on numerical linear algebra and for always making time for me when I stopped by your office unannounced.
Tom, thank you for pushing me to produce work of the highest quality and for answering my many, perhaps repetitive, career advice questions.
I feel truly fortunate to have crossed path's with each of you, and I look forward to ongoing collaboration in the future.

Many others have contributed directly to my academic development.
I'm particularly grateful to my undergraduate advisor, Roger Tobin, for his mentorship throughout my time at Tufts.
This was a formative time in my life during which much of my current perspective on academia was developed.
%as well as to Megan Callow and everyone in the IWP. 
%Alexandra Elbakyan
I'm also thankful for the collaborators I've had over the past several years, Raghu Bollapragada, Erin Carson, Yu-Chen Cheng, Eric Hallman, Hexuan Liu, Cameron Musco, Christopher Musco, Shashanka Ubaru, Rachel Ward, Natalie Wellen, and Qichen Xu, each of who contributed, in one way or another, to the content of this thesis.
I especially want to thank Yu-Chen for teaching me much of what I know about probability during the first year of my PhD.
Finally, thanks to my thesis committee, Sasha Aravkin and Maryam Fazel, for contributing their time and energy to my thesis and defense.

I'm indebted to many folks at UW who supported me throughout my time here.
This includes everyone at Hall Health, upstairs and downstairs, the staff and faculty in the Applied Math department, and the many students who I regularly interacted with.
When I first vested the department as a prospective student, it was immediately clear that there was a wonderful community of students. 
This was the deciding factor in my decision to come to UW, and I have been continually grateful for this community throughout my time in the program. 
%I'd especially like to thank Andreas, Brian, and Jeremy for modeling what PhD student life could look like and for helping to foster the amath community during my first few years in the program.
I was also extremely blessed to have a wonderful cohort without whom the last five years would have been far less enjoyable.
%Thanks especially to Diya, Jorge, Katherine, Micah, Nora, Roman, and Ying-Jen for the many memories we made.


Finally, I want to acknowledge my family and friends.
This thesis is dedicated to my grandparents, 陈祖浩 and 鲁月娥, the greatest supporters of my education.
To my parents, Ming-Guang (陈明光) and Peggy, thank you for providing for me since before I was even born.
This achievement is as much yours as it is mine.
To my partner, 黃羿綺, I'm so happy our lives intersected when they did.
Thank you for the constant love and support, and for quarantining with me during the pandemic. 
I cannot imagine what the last few years would have looked like without you.
Last, a shout out to  Jesse, Matthew, Max, Tailong, and Vaibhav, my ``internet friends'', who have each been with me through many stages of my life, despite large geographic distances.

\vfill

This material is based on work supported by the National Science Foundation under grant number DGE-1762114.
Any opinion, findings, and conclusions or recommendations expressed in this material are those of the author and do not necessarily reflect the views of the National Science Foundation.






{
\fontsize{11}{13}\selectfont
\tableofcontents
}

\clearpage
\pagenumbering{arabic}
\pagestyle{body}

\chapter{Introduction}
\label{chap:intro}

\freefootnote{A reference sheet containing common notation and useful factscan be found in \cref{chap:notation}.}

Computational approaches to today's most pressing and world-changing questions are reliant on subroutines for fundamental linear algebraic tasks.
The focus of this thesis is on the design and analysis of algorithms for an increasingly prevalent subset of such tasks: those involving matrix functions of Hermitian (or real symmetric) matrices.
For the duration of this thesis, \( \vec{A} \) will be a \( n\times n \) Hermitian matrix with eigenvalues \( \Lambda := \{ \lambda_i \}_{i=0}^{n-1} \) and (orthonormal) eigenvectors \( \{ \vec{u}_i\}_{i=0}^{n-1} \); i.e.,
\begin{equation}
    \label{eqn:A}
    \vec{A} = \smop{\sum_{i=0}^{n-1}} \lambda_i \vec{u}_i \vec{u}_i^\cT.
\end{equation}
A matrix function transforms the eigenvalues of a Hermitian (or symmetric) matrix according to some scalar function, while leaving the eigenvectors untouched.
\begin{definition}
\label{def:fA}
The matrix function \( \fA \), induced by \( f:\R\to\R \) and \( \vec{A} \), is defined as 
\begin{equation*}
    \fA := \smop{\sum_{i=0}^{n-1}} f(\lambda_i) \vec{u}_i \vec{u}_i^\cT.
\end{equation*}
   
\end{definition}
Perhaps the most well known example of a matrix function is the matrix inverse \( \vec{A}^{-1} \), which corresponds to the inverse function \( f(x) = x^{-1} \). 
Other common matrix functions including the matrix sign, logarithm, exponential, square root, and inverse square root, each of which has many applications throughout the mathematical sciences.

A common task involving matrix functions is computing the product \( \fA \vec{v} \) of a matrix function \( \fA \) with a fixed vector \( \vec{v} \); for instance, the matrix inverse applied to a vector corresponds to the solution of a linear system of equations.
Beyond the multitude of applications of linear systems, matrix functions applied to vectors are used for 
computing the overlap operator in quantum chromodynamics \cite{eshof_frommer_lippert_schilling_van_der_vorst_02}, 
solving differential equations in applied math \cite{saad_92,hochbruck_lubich_97}, 
Gaussian process sampling in statistics \cite{pleiss_jankowiak_eriksson_damle_gardner_20},
principle component projection and regression in data science \cite{jin_sidford_19},
and a range of other applications \cite{higham_08}.


Another related and especially interesting task involving matrix functions is estimating the \emph{spectral sum},
\begin{equation}
    \label{eqn:spectral_sum}
    \tr(\fA) = \smop{\sum_{i=0}^{n-1}} f(\lambda_i).
\end{equation}
Applications of spectral sums include
characterizing the degree of protein folding in biology \cite{estrada_00},
studying the thermodynamics of spin systems in quantum physics and chemistry \cite{weisse_wellein_alvermann_fehske_06,schnalle_schnack_10,schnack_richter_steinigeweg_20,jin_willsch_willsch_lagemann_michielsen_deraedt_21},
benchmarking quantum devices in quantum information theory \cite{jozsa_94},
maximum likelihood estimation in statistics \cite{barry_pace_99,pace_lesage_04},
designing better public transit in urban planning \cite{bergermann_stoll_22,wang_sun_musco_bao_21}, and 
finding triangle counts and other structure in network science \cite{avron_10,dong_benson_bindel_19,benzi_boito_20}.

The trace of matrix functions is intimately related to the spectral measure of \( \vec{A} \) which encodes the eigenvalues of \( \vec{A} \).
\begin{definition}
\label{def:CESM}
The cumulative empirical spectral measure (CESM) \( \Phi : \R \to [0,1] \), induced by \( \vec{A} \), is defined by
\begin{equation*}
    \Phi(x)
    = \Phi_{\vec{A}}
    := \smop{\sum_{i=0}^{n-1}} n^{-1} \bOne[ \lambda_i \leq x].
\end{equation*}
\end{definition}
Here \( \bOne[ \texttt{true} ] =1 \) and \( \bOne[\texttt{false}] = 0 \).

Not only is \( \Phi(x) \) itself a spectral sum for each \( x\in\R \), but 
\begin{equation*}
    \tr(\fA) = n \int f\,\d\Phi.
\end{equation*}
In this sense, approximating the CESM \( \Phi \) is equivalent to approximating spectral sums.
However, approximations to \( \Phi \) are also useful in that they provide a global picture of the spectrum of \( \vec{A} \).
Such coarse grained approximations are used in 
electronic structure computations %\cite{ducastelle_cryotlackmann_70,wheeler_blumstein_72,haydock_heine_kelly_72,alben_blume_krakauer_schwartz_75,girard_87,deraedt_devries_89,skilling_89,hutchinson_89} 
and other tasks in physics\footnote{In physics, the ``density'' \( \d\Phi/\d x \) is often called the density of states (DOS).} \cite{weisse_wellein_alvermann_fehske_06,jin_willsch_willsch_lagemann_michielsen_deraedt_21},
probing the behavior of neural networks in machine learning \cite{ghorbani_krishnan_xiao_19,papyan_19,granziol_wan_garipov_19,yao_gholami_keutzer_mahoney_20}, load balancing modern parallel eigensolvers in numerical linear algebra \cite{polizzi_09,evsl_19}, and computing the product of matrix functions with vectors \cite{fan_shuman_ubaru_saad_19}.

The simplest, and arguably most elegant, approach to spectrum and spectral sum approximation involves computing quadratic forms \( \vec{v}^\cT \fA \vec{v} \) for suitably chosen random vectors \( \vec{v} \).
For any fixed \( \vec{v} \), the task of computing \( \vec{v}^\cT \fA \vec{v} \) is intimately related to quadrature \cite{golub_meurant_94,golub_meurant_09} and, besides the many applications of spectrum and spectral sum approximation, is used for estimating the error of Krylov subspace methods \cite{dahlquist_eisenstat_golub_72,golub_strakos_94,golub_meurant_09}.



%\subsection{Krylov subspace methods}
\section{Lanczos-based methods}
\label{sec:lanczos}

The algorithms we study in this thesis fall into a general class of algorithms called Krylov subspace methods (KSMs).
KSMs produce approximations using information from the set of low-degree polynomials in \( \vec{A} \) applied to a vector \( \vec{v} \); i.e. from the so-called Krylov subspace generated by \( \vec{A} \) and \( \vec{v} \).
\begin{definition}
    \label{def:Kk}
    The dimension \( k \) Krylov subspace \( \mathcal{K}_k \) generated by \( \vec{A} \) and \( \vec{v} \) is defined as
\begin{equation*}
    \mathcal{K}_{k} 
    =\mathcal{K}_{k}(\vec{A},\vec{v}) 
    := \operatorname{span}\{\vec{v}, \vec{A}\vec{v}, \ldots, \vec{A}^{k-1}\vec{v} \}
    = \{ \pA\vec{v} : \deg(p) < k \}.
\end{equation*}
\end{definition}

The information from a given Krylov subspace can be used to approximate \( \fA \vec{v} \) and \( \vec{v}^\cT \fA \vec{v} \).
In particular, a natural approach is to use the approximations 
\begin{equation*}
    \fA \vec{v} \approx  \ff{f}{s}\mf{\vec{A}} \vec{v}
    ,\qquad
    \vec{v}^\cT \fA \vec{v} \approx  \vec{v}^\cT \ff{f}{s}\mf{\vec{A}} \vec{v},
\end{equation*}
where \( \ff{f}{s} : \R\to\R \) is a degree \( s \) polynomial chosen to approximate \( f \).

Throughout this thesis, the symbol ``\( \circ \)'' should be interpreted as a parameter encompassing any other parameters which impact how \( \ff{f}{s}\mf{\vec{A}}  \) is determined for \( f \). 
For instance, once choice of \( \circ \) may correspond to the interpolating polynomial to \( f \) at some set of nodes while another choice of \( \circ \) may correspond to the Chebyshev approximation to \( f \).
Specific choices of \( \circ \) corresponding to widely used algorithms will be defined as they come up.
Here the use of ``p'' stands for polynomial, and will be used to differentiate between polynomial approximations of a function and quadrature approximations of a distribution function, which will be defined later.

\begin{remark}
When \( \vec{A} \) is Hermitian, Krylov subspace methods are, in one way or another, related to the Lanczos algorithm \cite{lanczos_50} described in \cref{alg:lanczos}.
    Even so, we use the term \emph{Lanczos-based methods} to refer to algorithms which make use of the information generated by the Lanczos algorithm in some non-trivial way. 
This is in contrast to methods, such as those based on explicit polynomial approximation, which can easily be constructed directly.
\end{remark}

\begin{assumption}
    From this point onwards, we will assume \( \|\vec{v}\|_2=1 \).
\end{assumption}


The Lanczos algorithm (\cref{alg:lanczos}) \cite{lanczos_50} produces an orthonormal basis \( \{ \vec{q}_i \}_{i=0}^{k} \) for the Krylov subspace \( \mathcal{K}_{k+1} \) such that, for all \( i = 0, 1, \ldots, k \),
\begin{equation*}
    \operatorname{span}\{ \vec{q}_0, \vec{q}_1, \ldots, \vec{q}_{i} \} = 
    \mathcal{K}_{i+1}.
%    ,\qquad \forall i\in0:k+1.%,1,\ldots, k.
\end{equation*}
These basis vectors satisfy a three term recurrence, for all \( i = 0,1, \ldots,k-1 \),
\begin{equation*}
    \vec{A} \vec{q}_i = \beta_{i-1}\vec{q}_{i-1} + \alpha_i \vec{q}_i + \beta_i \vec{q}_{i+1}
    %,\qquad \forall i \in 0:k%i = 0,1,\ldots,k-1
\end{equation*}
with initial conditions \( \vec{q}_{-1} = \vec{0} \) and \( \beta_{-1} = 0 \).
The coefficients \( \{ \alpha_i \}_{i=0}^{k-1} \) and \( \{ \beta_i \}_{i=0}^{k-1} \) defining the three term recurrence are also generated by the algorithm.
This recurrence can be written in matrix form as
\begin{equation}
    \label{eqn:lanczos_three_term}
    \vec{A} \vec{Q} = \vec{Q} \vec{T} + \beta_{k-1} \vec{q}_k \vec{e}_{k-1}^\rT
\end{equation}
where 
\begin{equation*}
    \vec{Q} := 
    \begin{bmatrix}
        |&|&&|\\
        \vec{q}_0 & \vec{q}_1 & \cdots & \vec{q}_{k-1}\\
        |&|&&|\\
    \end{bmatrix}
    ,\quad
    \vec{T} := 
    \begin{bmatrix}
        \alpha_0 & \beta_0 \\
        \beta_0 & \alpha_1 & \ddots & \phantom{\ddots}\\
        &\ddots & \ddots & \beta_{k-2} \\
        &&\beta_{k-2} & \alpha_{k-1} \\
    \end{bmatrix}.
\end{equation*}

%\begin{remark}
%Comparing the Stieltjes procedure \cref{alg:stieltjes} to the Lanczos algorithm \cref{alg:lanczos}, it is clear that they are intimately related.
%In particular, \( \vec{q}_i = p_i\mf{\vec{A}} \vec{v} \), and \( \vec{T} = \vec{M} \).
%Indeed, the Lanczos algorithm is simply the Stieltjes procedure applied to a certain weighted spectral measure which we shall discuss in the next chapter.
%\end{remark}
\begin{labelalgorithm}[H]{lanczos}{Lanczos}{Lanczos algorithm}
\begin{algorithmic}[1]
    \Procedure{\thealgorithmname}{$\vec{A}, \vec{v}, k$}
    \State \( \vec{q}_0 = \vec{v} \), \( \beta_{-1} = 0 \), \( \vec{q}_{-1} = \vec{0} \)%, \( \beta_{-1} = 0 \)
    \For {\( i=0,1,\ldots, k-1 \)}
        \State \( \tilde{\vec{q}}_{i+1} = \vec{A} \vec{q}_{i} - \beta_{i-1} \vec{q}_{i-1} \)
        \State \( \alpha_{i} = \vec{q}_i^\cT \tilde{\vec{q}}_{i+1} \)
        \State \( \hat{\vec{q}}_{i+1} = \tilde{\vec{q}}_{i+1} - \alpha_{i} \vec{q}_i \)
        \State optionally, reorthogonalize,  \( \hat{\vec{q}}_{i+1} \) against \( \{\vec{q}_j\}_{j=0}^{i} \)
        \State \( \beta_{i} = \| \hat{\vec{q}}_{i+1} \| \)
        \State \( \vec{q}_{i+1} = \hat{\vec{q}}_{i+1} / \beta_{i} \)
    \EndFor
    \State \Return \( \{ \vec{q}_i \}_{i=0}^{k} \), \( \{ \alpha_i \}_{i=0}^{k-1} \), \( \{ \beta_i \}_{i=0}^{k-1} \).
\EndProcedure
\end{algorithmic}
\end{labelalgorithm}



\begin{remark}
    It is not uncommon for the  matrices which we call \( \vec{Q} \) and \( \vec{T} \) to be denoted by \( \vec{Q}_{\mkern2mu k} \) and \( \vec{T}_k \).
    We omit these subscripts for legibility, as the number of iterations \( k \) can be treated as fixed throughout this thesis.
    Note also that we begin indexing at zero so that indices match the degree of the corresponding polynomial.
\end{remark}

%\begin{assumption}
%    Throughout this thesis, we will assume \( \|\vec{v}\|_2 = 1 \). 
%    This is essentially without loss of generality as all algorithms we discuss have an obvious dependence on  \( \| \vec{v} \| \).
%\end{assumption}

\subsection{The insufficiency of interval-based bounds}
\label{sec:poly_insuff}

As we previously noted, this thesis is concerned with polynomial approximations to \( \fA \vec{v} \) and \( \vec{v}^\cT \fA \vec{v} \).
The error of such methods is often closely related to problems in scalar polynomial approximation theory.
In particular, note that 
\begin{equation*}
    \| g \A \|_2
    = \max_{x \in \Lambda} |g(x)|
    =: \| g \|_\Lambda, 
\end{equation*}
where \( \Lambda \) is the set of eigenvalues of \( \vec{A} \).
Here we have introduced the notation \( \|g\|_S = \sup_{x \in S}|g(x)| \) for \( g:\mathbb{C}\to\mathbb{C} \) and \( S\subset \mathbb{C} \). \label{def:norm_sup}
Let \( \| \cdot \| \) be any norm induced by a positive definite matrix with the same eigenvectors as \( \vec{A} \). \label{def:norm}
Then, a simple application of the sub-multiplicative property of matrix norms (see \cref{thm:norm_exchange}) implies 
\begin{equation}
    \label{eqn:linear_form}
    \frac{\| \fA \vec{v} -  \ff{f}{s}\mf{\vec{A}} \vec{v} \|}{\|\vec{v}\|} 
    \leq \| \fA - \ff{f}{s}\mf{\vec{A}}  \|_2  
   = \| f-\ff{f}{s} \|_{\Lambda}.
\end{equation}
Recalling our assumption \( \|\vec{v}\|_2 = 1 \), we also have
\begin{equation}
    \label{eqn:quadratic_form}
    | \vec{v}^\cT \fA \vec{v} - \vec{v}^\cT \ff{f}{s}\mf{\vec{A}} \vec{v} | 
    \leq \| \fA - \ff{f}{s}\mf{\vec{A}} \|_2 \| \vec{v} \|_2^2
    = \| f-\ff{f}{s} \|_{\Lambda}.
\end{equation}
Thus, we see that the quality of our approximations can be studied in terms of the quality of the polynomial approximation \( \ff{f}{s} \) to \( f \) on the eigenvalues of \( \vec{A} \).

Since polynomial approximation on an interval is well understood, it is common to bound the quality of Krylov subspace methods in terms of the best polynomial approximations on an interval.
It is always true that 
\begin{equation*}
    \| f-p \|_{\Lambda} \leq \| f-p \|_{\mathcal{I}},
\end{equation*}
where \( \mathcal{I} := [\lmin,\lmax] \) is the smallest interval containing all of the eigenvalues.
Thus, it is common to bound the left hand sides of \cref{eqn:linear_form,eqn:quadratic_form} by an expression like
\begin{equation}
    \label{eqn:unif_bound}
    2 \min_{\deg(p) < s} \| f - p \|_{\mathcal{I}}.
\end{equation}
However, while such bounds are useful in some situations, they often provide a large overestimate of the true behavior of Lanczos-based methods and are therefore unsuitable for use as practical stopping criteria.

\subsection{The effect of finite precision arithmetic}

While reorthogonalization in the Lanczos algorithm is unnecessary in exact arithmetic, omitting it often results in drastically different behavior when using finite precision arithmetic.
Specifically, the Lanczos basis \( \vec{Q} \) may be far from orthogonal, and the tridiagonal matrix \( \vec{T} \) may be far from what would have been obtained in exact arithmetic.
Because the Lanzcos algorithm is ostensibly unstable, there has been a widespread hesitance towards Lanczos-based approaches for problems involving matrix functions, at least without reorthogonalization \cite{jaklic_prelovsek_94,silver_roeder_voter_kress_96,aichhorn_daghofer_evertz_vondelinden_03,weisse_wellein_alvermann_fehske_06,ubaru_chen_saad_17,granziol_wan_garipov_19}.

The two primary effects of finite precision arithmetic on Lanczos-based methods when run without reorthogonalization are (i) a delay of convergence (increase in the number of iterations to reach a given level of accuracy) and (ii) a reduction in the maximal attainable accuracy.
However, while both effects are easily noticeable on most problems, they do not imply that reorthogonalization is needed.
In fact, throughout this thesis, we argue that Lanczos-based methods are highly effective even without reorthogonalization.

%In fact, \emph{the first major theme of this thesis is that Lanczos-based methods for matrix functions are often the best choice of Krylov subspace methods, even when accounting for the effects of finite precision arithmetic.}

%These effects are reasonably well understood in the context of linear systems \cite{greenbaum_89,greenbaum_97a}, and for some other functions such as the matrix exponential \cite{druskin_greenbaum_knizhnerman_98}. 





\subsection{A motivating example}

We now provide a simple and familiar example chosen to illustrate the themes introduced in this section.
The conjugate gradient algorithm (CG) \cite{hestenes_stiefel_52} is used to solve positive definite linear systems of equations, and is perhaps the most well-known Lanczos-based KSM.
When applied to a positive-definite linear system \( \vec{A} \vec{x} = \vec{v} \), CG produces iterates \( \textsf{cg}_k\in\mathcal{K}_k \) optimal in the \( \vec{A} \)-norm.
This optimality implies the error bounds\footnote{Throughout, we will occasionally use symbol ``\( x \)'' for the identity function \( x: \lambda\mapsto \lambda \) rather than an unspecified real value. Thus, expressions  like \( x p \) and \( x^{-1} - p \) should respectively be interpreted to mean the functions \( \lambda\mapsto \lambda p(\lambda) \) and \( \lambda\mapsto \lambda^{-1} - p(\lambda) \).}
\begin{align}
    \frac{\| \vec{A}^{-1} \vec{v} - \textsf{cg}_k \|_{\vec{A}}}{\| \vec{v} \|_{\vec{A}}}
    &\stackrel{\text{(a)}}{\leq} \min_{\deg(p)<k} \| x^{-1} - p \|_{\Lambda} %\label{eqn:spectrum_bd}
    \stackrel{\text{(b)}}{\leq} \min_{\deg(p)<k} \| x^{-1} - p \|_{\mathcal{I}}. \label{eqn:spec_int_bd}
    %&\stackrel{\text{(a)}}{\leq} \min_{\deg(p)<k} \max_{\lambda\in\Lambda}| \lambda^{-1} - p(\lambda) | %\label{eqn:spectrum_bd}
    %\stackrel{\text{(b)}}{\leq} \min_{\deg(p)<k} \max_{\lambda\in\mathcal{I}} | \lambda^{-1} - p(\lambda) |. \label{eqn:spec_int_bd}
\end{align}

Given \( \mathcal{I} = [\lmin,\lmax] \), \href{eqn:spec_int_bd}{(\ref{eqn:spec_int_bd}b)} can be computed analytically and, roughly speaking, it decreases linearly at a rate proportional to \( (1 - 1/\sqrt{\kappa}) \). 
In other words, to reach accuracy \( \epsilon \), CG requires at most \( O(\sqrt{\kappa}\log(\epsilon^{-1})) \) iterations, where \( \kappa = \lmax / \lmin \) is the condition number of \( \vec{A} \).
This is the well known \emph{root condition number bound} for CG.
On the other hand, \href{eqn:spec_int_bd}{(\ref{eqn:spec_int_bd}a)} may be significantly better than the latter bound involving \( \mathcal{I} \) and provide a more realistic picture of the convergence of CG.
However, since \href{eqn:spec_int_bd}{(\ref{eqn:spec_int_bd}a)} depends on the spectrum of \( \vec{A} \), which is typically unknown, the bound's use is in that it provides intuition into the theoretical behavior of CG rather than as a practical stopping critera.

\begin{figure}
    \begin{center}
        \includegraphics{imgs/ch1_unif_vs_spec.pdf} 
    \end{center}
    \caption[CG convergence (with and without reorthgonalization) and error bounds.]{%
    CG convergence (with and without reorthgonalization) and error bounds.
    \hspace{.25em}\emph{Legend}:
    CG error \( \| \vec{A}^{-1} \vec{v} - \textsf{cg}_k \|_{\vec{A}} / \| \vec{v} \|_{\vec{A}} \) with 
    ({\protect\raisebox{0mm}{\protect\includegraphics[]{imgs/legend/l2.pdf}}}) 
    and without
    ({\protect\raisebox{0mm}{\protect\includegraphics[]{imgs/legend/l1.pdf}}})
    reorthgonalization, 
    bound \href{eqn:spec_int_bd}{(\ref{eqn:spec_int_bd}a)} on \( \Lambda \) 
    ({\protect\raisebox{0mm}{\protect\includegraphics[]{imgs/legend/l4.pdf}}}),
    and bound \href{eqn:spec_int_bd}{(\ref{eqn:spec_int_bd}b)} on \( \mathcal{I} \) 
    ({\protect\raisebox{0mm}{\protect\includegraphics[]{imgs/legend/l3.pdf}}}).
    }
    \label{fig:ch1_unif_vs_spec}
\end{figure}

In \cref{fig:ch1_unif_vs_spec}, we plot the bounds \href{eqn:spec_int_bd}{(\ref{eqn:spec_int_bd}a)} and \href{eqn:spec_int_bd}{(\ref{eqn:spec_int_bd}b)} as well as the actual errors in a Lanczos-based implementation of CG run with and without reorthogonalization for a spectrum with exponentially spaced eigenvalues (the precise details are not important at this point).
Observe that the bound \href{eqn:spec_int_bd}{(\ref{eqn:spec_int_bd}a)} decreases significantly faster than \href{eqn:spec_int_bd}{(\ref{eqn:spec_int_bd}b)} for the given spectrum.
Note also that, even without reorthogonalization, CG converges significantly faster than \href{eqn:spec_int_bd}{(\ref{eqn:spec_int_bd}b)}.


An alternative to CG is the Chebyshev semi-iterative method \cite{glanders_shortley_50,gutknecht_rollin_02}.
This method can be implemented without using the Lanczos algorithm by directly constructing an explicit polynomial approximation to \( x^{-1} \) on \( \mathcal{I} \). 
However, as a result, the algorithm is unable to adapt to the spectrum of \( \vec{A} \) and usually converges very similarly to \href{eqn:spec_int_bd}{(\ref{eqn:spec_int_bd}b)}.
Moreover, if \( \mathcal{I} \) is estimated inaccurately, then the algorithm may become unstable. 
Interestingly, even in finite precision arithmetic, an iterate very close to what would be produced by the Chebyshev method can be obtained from the Lanczos method \cite{greenbaum_89,druskin_knizhnerman_91,musco_musco_sidford_18}.
Since the extreme eigenvalues of \( \vec{T} \) typically provide a very good estimate for \( \mathcal{I} \), this means that a Lanczos-based implementation of the Chebyshev method avoids the need for a priori parameter selection.

\begin{remark}
    Optimization algorithms such as accelerated gradient descent also attain a root condition number iteration complexity on any smooth and strongly convex function (such as \( \vec{x} \mapsto \frac{1}{2} \vec{x}^\cT \vec{A} \vec{x} - \vec{x}^\cT \vec{v} \) for positive definite linear systems).
Since this rate is optimal among first-order methods for smooth and strongly convex functions, accelerated gradient descent is often referred to as ``optimal''.
The fact that CG has a similar convergence guarantee often leads CG to be introduced as an alternative to accelerated gradient descent for linear systems.
However, accelerated gradient descent is essentially equivalent to the Chebyshev semi-iterative method when applied to the above objective and therefore is not typically competitive with CG in terms of the number of matrix-vector products.
\end{remark}

%\section{Some historical context}

\section{Context and contributions}
\iffalse
The Lanczos method was introduced in 1950 as a eigenvalue algorithm \cite{lanczos_50}
The conjugate gradient algorithm was also introduced in the early 1950s, and is typically attributed the joint paper of Hestenes and Stiefel \cite{hestenes_stiefel_52}, although Hestenes, Stiefel, and Lanczos \cite{lanczos_50,lanczos_52} all had earlier work involving the algorithm.
For a more comprehensive overview on the early history of the algorithms, we turn readers to \cite{golub_oleary_89}


Based on earlier works by 


National Bureau of Standards (now the National Institute of Standards and Technology).

\fi
%There are a number of books, theses, and monographs which discuss the use of Krylov subspace methods for tasks involving matrix functions.
Essentially any numerical linear algebra textbook will have at least one chapter on KSMs.
In fact, there are a number of texts which focus on the classical uses of KSMs: eigenvalue problems and solving linear systems of equations \cite{paige_71,greenbaum_97,saad_11,meurant_strakos_06,meurant_06,liesen_strakos_13}; see also \cite{golub_oleary_89} for a historical overview of early developments.
The large number of such resources means that treatments of topics such as the behavior of algorithms in finite precision arithmetic can be found at a range of levels of detail. 

Resources dealing with general matrix functions are less plentiful.
While many textbooks might have a chapter on functions such as the exponential or square root, the only recent book I am aware of which is dedicated specifically to matrix functions is \cite{higham_08}.
However, this book is focused primarily on the case of computing \( \fA \), with only one chapter devoted to the task of computing \( \fA \vec{v} \) and no discussion on the task of \( \vec{v}^\cT \fA \vec{v} \).
There are a number of more specialized texts on these topics.
The topic of \( \vec{v}^\cT \fA \vec{v} \) is covered thoroughly in \cite{golub_meurant_94,golub_meurant_09} from a theoretical quadrature perspective.
A widely used practical quadrature algorithm for estimating spectral sums and spectral densities called the kernel polynomial method is discussed in \cite{weisse_wellein_alvermann_fehske_06} but not analyzed theoretically.
The thesis \cite{schweitzer_16} discusses practical error bounds for methods for computing \( \fA\vec{v} \) for symmetric and non-symmetric \( \vec{A} \) as well as restarting techniques for non-symmetric problems, and the thesis \cite{cortinovis_22} discusses low-rank approximation of matrix functions as well as stochastic trace estimation of matrix functions.
None of the above texts discuss thoroughly the impacts of finite precision arithmetic.

During my PhD studies, it became strikingly clear that the state of knowledge surrounding Lanczos-based methods for matrix functions is fragmented.
For instance, there are several lines of work within the quantum physics literature which contain results not discovered in applied math until decades later.
Conversely, practitioners in physics, data science, and machine learning, often lack knowledge regarding the practical behavior of Lanczos-based methods in finite precision arithmetic. 
While this can be partially attributed to a lack of due-diligence in studying background material, a larger problem is that the requisite background is not easily accessible to non-specialists.
In fact, some of the most relevant work on Lanczos-based methods in finite precision arithmetic is not even well known within the numerical analysis community.


This thesis aims to fill some of the gaps in the presentation of Lanczos-based methods for matrix functions by providing a comprehensive background on the topic.
Indeed, several chapters are primarily expository, with the express goal of providing a more thorough context for the other chapters. 
%Of particular interested is a chapter on the impacts of finite precision arithmetic on Lanczos-based methods for matrix functions which we hope will serve as a useful reference going forwards. 
While there are a number of technical contributions, arguably the most significant contributions of this thesis are the following two themes:
\begin{itemize}
    \item 
        Bounds based on polynomial approximation on a single interval are insufficient to describe the true behavior of Lanczos-based methods for matrix functions. 
    Instead, one should seek bounds based on the spectrum which are able to take advantage of more fine-grained spectral structure such as gaps and outlying eigenvalues.
    \item 
        While Lanczos-based methods may behave differently in finite precision arithmetic than exact arithmetic, they still outperform Krylov subspace methods based on explicit polynomial approximations. 
        Moreover, the hyper-parameters in explicit polynomial methods can be determined effectively through the use of Lanczos-based implementations.
\end{itemize}
It is my hope that the contributions of this thesis are presented in a way which will promote understanding of and intuition for Lanczos-based methods for matrix functions outside of the numerical analysis community.



This thesis contains primarily work which appears in the following papers:

\begin{refsection}
    \nocite{
    chen_trogdon_ubaru_22,
    chen_trogdon_ubaru_21,
    chen_greenbaum_musco_musco_22a,
    chen_greenbaum_musco_musco_22b,
    }
    \printbibliography[heading=none]
\end{refsection}

Rather than stapling these papers together, I have arranged the content in accordance with my broader goals for this thesis.
In addition, a large amount of new exposition has been included in order to tie things together more cleanly and to provide additional context for non-experts.
Towards this end, many of the numerical examples from the above papers have been modified for consistency with the rest of the thesis.
The files needed to generate the figures in this thesis (as well as the thesis itself) are freely available online.


%\cite{chen_cheng_22}

%\cite{chen_hallman_22}

 

\chapter{Scalar polynomials}


In this chapter, we introduce some basic theory regarding scalar polynomials%
\footnote{Polynomials were perhaps the first abstract mathematical object I encountered --- I distinctly remember grappling with the meaning of a variable \(x\), which is ``somehow a number but also not'', in the later years of elementary school. It's amazing, then, that the theory of polynomials of a single variable underlies so much of my PhD thesis.}
which will come in handy throughout the rest of this thesis.
A scalar polynomial of degree \( k \) is a function of the form 
\begin{equation*}
    x\mapsto c_0 + c_1 x + \cdots + c_k x^k,
%    ,\qquad 
%    c_0,c_1, \ldots, c_k\in \R.
\end{equation*}
where \( c_0, c_1, \ldots, c_k \) are fixed scalars.
Such a polynomial is naturally extended to matrices as
\begin{equation*}
    \vec{A} \mapsto c_0 \vec{I} + c_1 \vec{A} + \cdots + c_k \vec{A}^k
\end{equation*}
in a manner compatible with our definition of matrix functions.
Thus, matrix polynomials are intimately related to Krylov subspace methods.

%The precise relationship with Lanczos-based methods is discussed in \cref{sec:poly_insuff} where we provide more detail on our assertion that bounds based on best polynomial approximation on a single interval are often insufficient.


\section{Basic definitions}

Our discussion on quadrature centers around approximating distributions and integrals against distribution functions. 
Several examples of such functions are illustrated in \cref{fig:ch2_sample_dist}.

\begin{figure}[htb]
    \begin{center}
        \includegraphics{imgs/ch2_sample_dist.pdf} 
    \end{center}
    \caption[Sample unit mass distribution functions.]{%
    Sample unit mass distribution functions.
    \hspace{.25em}\emph{Legend}:
    continuous increasing distribution function ({\protect\raisebox{0mm}{\protect\includegraphics[]{imgs/legend/l1nm.pdf}}}),
    continuous distribution function which is not weakly-increasing ({\protect\raisebox{0mm}{\protect\includegraphics[]{imgs/legend/l2nm.pdf}}}),
    discrete weakly-increasing distribution function ({\protect\raisebox{0mm}{\protect\includegraphics[]{imgs/legend/l3nm.pdf}}}),
    discrete distribution function which is not weakly-increasing ({\protect\raisebox{0mm}{\protect\includegraphics[]{imgs/legend/l4nm.pdf}}}).
    }
    \label{fig:ch2_sample_dist}
\end{figure}


\begin{definition}
A (signed) unit mass distribution function \( \Upsilon \) is a right continuous function \( \Upsilon:\R\to\R \) such that \( \lim_{x\to\infty} \Upsilon(x) = 1 \) and \( \lim_{x\to-\infty} \Upsilon(x) = 0 \).
If \( \Upsilon \) is also weakly increasing, we say it is a probability distribution function.
\end{definition}

\begin{remark}
    If \( \Upsilon \) is differentiable, then the derivative \( \d\Upsilon / \d{x} = \Upsilon' \) is the usual probability density function.
    Likewise, in the sense of distributions, \( \d\bOne[a \leq x] /\d{x} = \delta(x-a) \) where \( \delta(x-a) \) is a unit mass Dirac delta function centered at \( a \).
    Thus, if \( \Upsilon \) is piecewise constant, then \( \d\Upsilon/\d{x} \) can be expressed in terms of the sum of weighted Dirac delta functions, where a delta function is located at each discontinuity and weighted by the size of the corresponding jump.
\end{remark}

We now introduce several definitions which we will use throughout the next several chapters.
\begin{definition}
    Given a function \( f \) and distribution function \( \Upsilon \) we denote by 
        \( 
        \int_{a}^{b} f \,\d \Upsilon
        \) the standard Riemann--Stieltjes integral% defined by
        \begin{equation*}
            \int_a^b f \,\d \Upsilon
            := \lim_{\|P\|\to 0} \smop[r]{\sum_{i=0}^{p-1}} f(c_i) \big( \Upsilon(x^{(i+1)}) - \Upsilon(x^{(i)}) \big),
        \end{equation*}
        where \( P = \{ a = x^{(0)} < \cdots < x^{(p)} = b \} \) is a partition of \( [a,b] \), \( \| P \| = \max_i | x^{(i+1)} - x^{(i)}| \), and \( c_i \in [x_i,x_{i+1}] \).

    For notational clarity, we will often write 
    \(
    \int f \,\d \Upsilon
    \)
    in which case \( a,b \) can be taken as \( \pm \infty \).%\in \mathbb{R} \cup \{\pm \infty\} \) are understood to be such that \( \Upsilon \) is constant on \( (-\infty,a) \) and \( (b,\infty) \).
\end{definition}


\begin{definition}
    In the setting of the previous definition,
    \begin{equation*}
        \int_{a}^{b} f \, |\d \Upsilon|
        := \lim_{\|P\|\to 0} \smop[r]{\sum_{i=0}^{p-1}} f(c_i) \big| \Upsilon(x^{(i+1)}) - \Upsilon(x^{(i)}) \big|.
    \end{equation*}
\end{definition}

\begin{definition}
    Let \( \Upsilon \) be a (distribution) function.
The total variation (TV) of \( \Upsilon \), denoted \( \TV(\Upsilon) \), is defined by
\begin{equation*}
    \TV(\Upsilon) := \int |\d \Upsilon|.
\end{equation*}
\end{definition}
\begin{remark}
If \( \Psi \) is a weakly-increasing unit-mass distribution function, then \( \TV(\Psi) = 1 \). 
\end{remark}

To measure the similarity of two distribution functions, we will typically use the Wasserstein (earth mover) distance.
\begin{definition}
\label[definition]{def:wasserstein}
Let \( \Upsilon_1 \) and \( \Upsilon_2 \) be two probability distribution functions.
The Wasserstein distance between \( \Upsilon_1 \) and \( \Upsilon_2 \), denoted \( \W(\Upsilon_1,\Upsilon_2) \), is defined by
\begin{equation*}
    \W(\Upsilon_1,\Upsilon_2) 
    := \int | \Upsilon_1 - \Upsilon_2 | \,\d{x}.
\end{equation*}
\end{definition}

It is well known that the Wasserstein distance between two distribution functions has a dual form involving 1-Lipshitz functions. 
\begin{definition}
    We say that \( f\in \mathsf{Lip}(L,S) \) if \( |f(x) - f(y)| \leq L|x-y| \) for all \( x,y \in S \subseteq \R \).
\end{definition}

\begin{lemma}
\label[lemma]{thm:wasserstein_1lip}
    Suppose \( \Upsilon_1 \) and \( \Upsilon_2 \) are two probability distribution functions of bounded total variation each constant on \( (-\infty,a) \) and \( (b,\infty) \).
Then,
\begin{equation*}
    \W(\Upsilon_1,\Upsilon_2) = \sup \left\{ \int f \d \big( \Upsilon_1 - \Upsilon_2\big) : f\in\textup{\textsf{Lip}}(1,[a,b]) \right\}.
\end{equation*}
\end{lemma}
\begin{remark}
    In some situations, other metrics may be more meaningful. 
    For instance, if it is important for two distribution functions to agree to very high precision in a certain region, but only to moderate accuracy in others, then the Wasserstein distance may be unsuitable.
\end{remark}


\section{Orthogonal polynomials}
\label{sec:OP}


Throughout this thesis, \( \mu \) will be a non-negative unit mass distribution function. 
Associated with \( \mu \) is the inner product \( \langle \cdot, \cdot \rangle_\mu \) defined by
\label{def:mu}
\begin{equation}
    \label{eqn:mu_ip}
    \langle f, g\rangle_\mu\: := \int fg \,\d\mu.
\end{equation}

The set of \( \mu \)-square-integrable functions forms a Hilbert space with respect to this inner product, so we may hope to find an orthonormal basis \( \{ p_i \}_{i=0}^{\infty} \) of polynomials with \( \deg(p_i) = i \). \label{def:OP}
Such a basis is easily produced by a simple modification of the Gram-Schmidt algorithm which results in a naive implementation of the so-called Stieltjes algorithm.
An implementation is described in \cref{alg:stieltjes_naive}.

\begin{labelalgorithm}[H]{stieltjes_naive}{Stieltjes}{Stieltjes algorithm (naive)}
\begin{algorithmic}[1]
    \Procedure{\thealgorithmname}{$\mu,k$}
    \State \( p_0 = 1 \)
    \For {\( i=0,1,\ldots, k-1 \)}
        \State \( \tilde{p}_{i+1} = x p_i \)
        \State \( \hat{p}_{i+1} = \tilde{p}_{i+1} - \big(\langle p_0, \tilde{p}_{i+1} \rangle_\mu\: p_0  +  \cdots + \langle p_{i}, \tilde{p}_{i+1} \rangle_\mu\: p_i \big) \)
        \State \( p_{i+1} = \hat{p}_{i+1} / \| \hat{p}_{i+1} \|_\mu\: \)
    \EndFor
    \State \Return \( \{ p_i \}_{i=0}^{k} \) 
\EndProcedure
\end{algorithmic}
\end{labelalgorithm}

Note that the polynomials satisfy, for all \( i\geq 0 \),
\begin{equation}
    \label{eqn:all_term_recurrence}
    x p_i = \| \hat{p}_{i+1} \|_\mu\: p_{i+1} + \langle p_0, \tilde{p}_{i+1} \rangle_\mu\: p_0  +  \cdots + \langle p_{i}, \tilde{p}_{i+1} \rangle_\mu\: p_i.
\end{equation}
This can be written in matrix form as
\begin{equation*}
x\, [p_0, p_1, \ldots ] = [p_0, p_1, \ldots ] \vec{H},
\end{equation*}
where \( [p_0, p_1, \ldots ] \) is a \emph{quasi-matrix} whose columns are the polynomials \( \{ p_i \}_{i=0}^{\infty} \) and \( \vec{H} \) is a semi-infinite upper-Hessenberg matrix.
Moreover, for all \( i,j \geq 0 \),
\begin{equation*}
    [\vec{H}]_{i,j} = \langle p_i, x p_j \rangle_\mu  = \langle p_j, x p_i \rangle_\mu = [\vec{H}]_{j,i};
\end{equation*}
that is, \( \vec{H} \) is symmetric. 
Since, by construction, \( \vec{H} \) is upper-Hessenberg, this implies that \( \vec{H} \) is symmetric tridiagonal!

Therefore, for all \( i\geq0 \), \cref{eqn:all_term_recurrence} becomes the symmetric three-term recurrence
\begin{equation}
    \label{eqn:poly_three_term}
    x p_i =  \beta_{i-1} p_{i-1} + \alpha_i p_i + \beta_{i} p_{i+1}
\end{equation}
with initial conditions \( p_0 = 1 \), \( p_{-1} = 0 \), and \( \beta_{-1} = 0 \), where \( \{ \alpha_i \}_{i\geq0} \) and \( \{ \beta_i \}_{i\geq 0} \) are chosen to enforce orthogonality.
In particular, \cref{alg:stieltjes_naive} can be modified to take advantage of this short-recurrence, resulting in the standard implementation of the Stieltjes algorithm, \cref{alg:stieltjes}.
In the case that \( \beta_i = 0 \), the algorithm should be terminated as the dimension of Krylov subspaces does not continue to grow.

\begin{labelalgorithm}[htb]{stieltjes}{Stieltjes}{Stieltjes algorithm}
\begin{algorithmic}[1]
    \Procedure{\thealgorithmname}{$\mu,k$}
    \State \( p_0 = 1 \)
    \For {\( i=0,1,\ldots, k-1 \)}
        \State \( \tilde{p}_{i+1} = x p_i \)
        \State \( \alpha_i = \langle p_i, \tilde{p}_{i+1} \rangle_\mu \)
        \State \( \hat{p}_{i+1} = \tilde{p}_{i+1} - \alpha_i p_i \)
        \State \( \beta_{i} = \| \hat{p}_{i+1} \|_\mu\: \) 
        \State \( p_{i+1} = \hat{p}_{i+1} / \beta_i \)
    \EndFor
    \State \Return \( \{ p_i \}_{i=0}^{k} \), \( \{ \alpha_i \}_{i=0}^{k-1} \), \( \{\beta_i\}_{i=0}^{k-1} \) 
\EndProcedure
\end{algorithmic}
\end{labelalgorithm}

\begin{definition}
    \label{def:jacobi}
    The, possibly semi-infinite, tridiagonal matrix \( \vec{M} = \vec{M}(\mu) \) giving the three-term recurrence coefficients for the orthogonal polynomials of \( \mu \) is called the Jacobi matrix corresponding to \( \mu \). 
    Unless specified otherwise, the coefficients are 
    \begin{equation*}
    \vec{M} = 
    \begin{bmatrix}
        \alpha_0 & \beta_0 \\
        \beta_0 & \alpha_1 & \beta_1 & \phantom{\ddots}\\
        &\beta_1 & \alpha_2 & \ddots \\
        &&\ddots & \ddots \\
    \end{bmatrix}.
\end{equation*}
%and for a distribution \( \nu \), different from \( \mu \), we will denote the corresponding Jacobi matrix by \( \vec{M}(\nu) \).   
\end{definition}

An important property of a distribution function \( \Upsilon \) are it's polynomial moments.
We are particularly interested in those induced by \( \mu \).
\begin{definition}
\label{def:modified_moments}
For each \( i \geq 0 \), the modified moments of \( \Upsilon \) (with respect to \( \mu \)) are 
\begin{equation*}
    m_i = m_i(\Upsilon,\mu) := \int p_i \,\d\Upsilon.
\end{equation*}
If \( m_0, \ldots, m_{s} < \infty \), we say \( \Upsilon \) has finite moments through degree \( s \).
\end{definition}
 

Jacobi matrices have many interesting properties, several of which we review here.

\begin{lemma}
    The upper-leftmost \( k\times k \) submatrix of a Jacobi matrix is entirely determined by the moments through degree \( 2k-1 \) of the associated distribution function.
\end{lemma}

\begin{proof}
    This is a direct consequence of the fact that the \( k \)-point (degree \( 2k-1 \)) Gaussian quadrature rule for a distribution function can be determined from the upper-leftmost \( k\times k \) submatrix of the associated Jacobi matrix.
    This argument will be made whole in \cref{sec:gq}.
\end{proof}

\begin{lemma}
    \label{thm:jacobi_eigen}
    Denote the zeros of \( p_{k} \) by \( \{ \theta_j^{(k)} \}_{j=0}^{k-1} \).
Then, for any \( j=0,1,\ldots, k-1 \),
\begin{equation*}
    \label{eqn:jacobi_evec}
    \big[ p_0(\theta_{j}^{(k)}), p_1(\theta_{j}^{(k)}), \ldots, p_{k-1}(\theta_{j}^{(k)}) \big]^\cT
\end{equation*}
is an eigenvector of \( [\vec{M}]_{:k,:k} \) with eigenvalue \( \theta_{j}^{(k)} \).
Moreover, all eigenvectors are obtained in this way.
\end{lemma}

\begin{proof}
    In matrix form, \cref{eqn:poly_three_term} becomes
    \begin{equation*}
        x\, [p_0, p_1, \ldots, p_{k-1} ] = [p_0, p_1, \ldots, p_{k-1}] \vec{M} + \beta_{k-1} p_{k} \vec{e}_{k-1}^\rT.
    \end{equation*}
    Evaluating each side of the above equality at \( \theta_j^{(k)} \) gives the first part of the result.

    To show all eigenvectors are obtained in this way, it suffices to show that \( \{ \theta_j^{(k)} \}_{0\leq j<k} \) are distinct.
    Let \( \{ t_j \}_{j=0}^{k'-1} \) be the points at which \( p_{k} \) changes signs. 
    Then, 
    \begin{equation*}
        \int p_k {\prod_{j=0}^{k'-1}} (x-t_j) \,\d \mu \neq 0
    \end{equation*}
    since the integrand does not change signs.
    This implies \( k' = k \) since \( p_k \) is orthogonal to all polynomials of lower degree.
\end{proof}





\subsection{Chebyshev polynomials}

Owing to the deep connection between Chebyshev polynomials and approximation theory \cite{trefethen_19}, one particularly important choice of \( \mu \) is the distribution function corresponding to the Chebyshev polynomials of the first kind. 
We will often treat this case with special care.
%While we treat the case of the Chebyshev polynomials of the first kind, similar results can usually be expected for other classical orthogonal polynomials.
\begin{definition}
    \label{def:muT}
    The Chebyshev distribution function of the first kind, \( \mu_{a,b}^T : [a,b] \to [0,1] \), is defined as
    \begin{equation*}
        \mu_{a,b}^T 
        := \frac{1}{2} + \frac{1}{\pi}\arcsin\left(\frac{2}{b-a}x - \frac{b+a}{b-a}\right).
    \end{equation*}
    Thus, for \( x\in[a,b] \),
    \begin{equation*}
        \frac{\d\mu_{a,b}^T}{\d{x}}
        = \frac{2}{\pi(b-a)} \Big(1-\big(\frac{2}{b-a}x - \frac{b+a}{b-a}\big)^2\Big)^{-1/2}.
    \end{equation*}    
\end{definition}

\begin{definition}
    The Chebyshev polynomials of the first kind, denoted \( \{ T_i \}_{i=0}^{\infty} \), are defined by the recurrence \( T_0 = 1 \), \( T_1 = x \), and, for all \( i \geq 1 \),
    \begin{equation*}
        T_{i+1} := 2 x T_i - T_{i-1}.
    \end{equation*}
\end{definition}

    It can be verified that the orthogonal polynomials \( \{p_i\}_{i=0}^{\infty} \) with respect to \( \mu_{a,b}^T \) are given by \( p_0 = T_0 = 1 \) and, for all \( i \geq 1 \),
    \begin{equation*}
        p_i = \sqrt{2} T_i \left( \frac{2}{b-a} x + \frac{b+a}{b-a} \right).
    \end{equation*}

    Therefore, the Jacobi matrix \( \vec{M}( \mu_{a,b}^T ) \) has diagonal and off diagonals entries given by
    \begin{equation*}
        \left[ \frac{a+b}{2}, \frac{a+b}{2}, \ldots \right]
        \quad\text{and}\quad
        \left[ \frac{b-a}{2\sqrt{2}} , \frac{b-a}{4}, \frac{b-a}{4} , \ldots  \right],
    \end{equation*}
    respectively.



\section{Polynomial approximations and bounds}
\label{sec:uniform_poly_bounds}

As noted in the introduction, we use the notation \( \ff{f}{s} \) to denote a degree \( s \) polynomial obtained from \( f \) by some algorithm parameterized by \( \circ \).
In particular, we will make the following definitions.
\begin{definition}
    \label{def:fffs}
    Given a non-negative unit mass distribution function \( \mu \) with degree \( s+1 \) orthogonal polynomial \( p_{s+1} \) we define,
    \begin{equation*}
        \ff{f}{s} := 
        \begin{cases}
            \circ = \text{i} & \text{degree $s$ interpolant to \( f \) at roots of $p_{s+1}$} \\
            \circ = \text{a} & \text{degree $s$ truncated series for \( f \) in \( \langle \cdot, \cdot \rangle_\mu \)} \\
%            \circ = \text{d-i} & \text{damped degree $s$ interpolant to \( f \) at roots of $p_{s+1}$} \\
%            \circ = \text{d-a} & \text{damped degree $s$ approximation to \( f \) in \( \langle \cdot, \cdot \rangle_\mu \)} \\
        \end{cases}
    \end{equation*}.
\end{definition}

The damped projection \( \ff[d-a]{f}{s} \) and damped interpolant \( \ff[d-i]{f}{s} \) are respectively obtained by scaling each of the coefficients of \( \ff[a]{f}{s} \) and \( \ff[i]{f}{s} \), when represented as a linear combination of the orthogonal polynomials \( \{ p_i \}_{i=0}^{s} \), by constants \( \rho_i \) for each \( i = 0,1,\ldots, s\).
\begin{definition}
\label{def:damping}
    Write \( \ff{f}{s} \) in a polynomial series with respect to \( \langle \cdot, \cdot \rangle_\mu \)\,; i.e. as
    \begin{equation*}
        \ff{f}{s} = \: \smop{\sum_{i=0}^{s}} c_i p_i.
    \end{equation*}
    Then, given damping coefficients \( \{ \rho_i \}_{i=0}^{s} \) with \( 0\leq \rho_i \leq 1 \) for all \( i \),
    \begin{equation*}
        \ff[d-$\circ$]{f}{s} := \: \smop{\sum_{i=0}^{s}} \rho_i c_i p_i.
    \end{equation*}
\end{definition}



We now review several classical results from approximation theory which we will use throughout this thesis.
These are constructive bounds for the case \( \mu = \mu_{-1,1}^T \) which provide upper bounds for the quality of the \emph{best} polynomial approximation to \( f \).
In fact, both  \( \ff[a]{f}{s} \) and \( \ff[i]{f}{s} \) provide nearly optimal approximations in many settings \cite{trefethen_19}.

A full treatment requires at least a textbook, and we refer readers to \cite{trefethen_19} for an excellent such book.
The following theorems are summarized from Theorems 7.2 and 8.2 in \cite{trefethen_19}.

\begin{definition}
    We say that \( f\in \mathsf{BV}(d,V,S) \) if, on \( S\subseteq \R \), \( f \) is \( d \) times differentiable, its derivatives through \( f^{(d-1)} \) are absolutely continuous, and the \( d \)-th derivative \( f^{(d)} \) has total variation bounded above by some constant \( V \) on \( S \).
\end{definition}


\begin{definition}
    For \( \rho \geq 1 \) the Bernstein ellipse \( E_\rho(a,b) \) is the ellipse centered at \( \frac{a+b}{2} \) with semi-axis lengths \( \frac{b-a}{2} \frac{1}{2} (\rho + \rho^{-1}) \) and \( \frac{b-a}{2}\frac{1}{2} (\rho+\rho^{-1}) \) along the real and imaginary directions; i.e
    \begin{equation*}
        E_\rho(a,b) = \left\{ z \in\mathbb{C} : z = \frac{b-a}{2} \frac{1}{2} (u + u^{-1}) + \frac{a+b}{2}, \: u = \rho \exp(i\theta), \: \theta\in[0,2\pi) \right\}. 
    \end{equation*}
\end{definition}


\begin{definition}
    We say that \( f\in \mathsf{Anl}(\rho,M,[a,b]) \) if \( f \) is analytic on the region enclosed by \( E_\rho(a,b) \) where it satisfies \( |f(x)|\leq M \).
\end{definition}

\begin{theorem}
    \label{thm:bv_cheb}
    For an integer \( d \geq 0 \), suppose \( f \) is \( d \) times differentiable, its derivatives through \( f^{(d-1)} \) are absolutely continuous, and the \( d \)-th derivative \( f^{(d)} \) has total variation bounded above by some constant \( V \) on \( [-1,-1] \); i.e., suppose \( f \in \textup{\textsf{BV}}(d,V,[-1,1]) \).
    Then, with \( \mu = \mu_{-1,1}^T \), for any \( s > d \), 
    \begin{equation*}
        \| f - \ff[a]{f}{s} \|_{[-1,1]} \leq \frac{2 V}{\pi d(s-d)^d}
        ,\qquad
        \| f - \ff[i]{f}{s} \|_{[-1,1]} \leq \frac{4 V}{\pi d(s-d)^d}.
    \end{equation*}
\end{theorem}

\begin{theorem}
    \label{thm:analytic_cheb}
    Suppose \( f \) is analytic on the region enclosed by the Bernstein ellipse \( E_\rho(-1,1) \) where it satisfies \( \|f\|_{E_{\rho(-1,1)}} \leq M \); i.e., suppose \( f\in \textup{\textsf{Anl}}(\rho,M,[-1,1]) \).
    %Let a function \( f \) analytic in \( [-1,1] \) be analytically continuance to the open Bernstein ellipse \( E_{\rho} \), where it satisfies \( |f(z)| \leq M \) for some \( M \).
    Then, with \( \mu = \mu_{-1,1}^T \), for any \( s \geq 0 \), 
    \begin{equation*}
        \| f - \ff[a]{f}{s} \|_{[-1,1]} \leq \frac{2 M \rho^{-k}}{\rho - 1}
        ,\qquad
        \| f - \ff[i]{f}{s} \|_{[-1.1]} \leq \frac{4 M \rho^{-k}}{\rho - 1}.
    \end{equation*}
\end{theorem}

\begin{lemma}
    \label[lemma]{thm:cheb_rates}
    Set \( c_{\textup{i}} = 2 \) and \( c_{\textup{a}} = 1 \).
    Then, for \( \circ \in \{ \textup{i}, \textup{a} \} \), \( \| f - \ff[]{f}{s} \|_\infty < \epsilon / 2 \) provided
        \begin{equation*}
        s \geq 
         \begin{cases}
             \frac{1}{\ln(\rho)} \ln \left( \frac{4 c_\circ M}{\rho-1} \right) + \frac{1}{\ln(\rho)} \ln \left( \varepsilon^{-1} \right) & f \in \mathsf{Anl}(\rho,M,[a,b]), \\
             d + \left( \frac{4 c_\circ V}{\pi d} \right)^{1/d} \varepsilon^{-1/d} & f \in \mathsf{BV}(d,V,[a,b]) .
        \end{cases}   
        \end{equation*}
\end{lemma}
\begin{proof}
    Define \( \tilde{f} : \R\to\R\) by \( \tilde{f}(x) = f( \frac{b-a}{2} x + \frac{a+b}{2} ) \).
    Then we have that,
    \begin{equation*}
        \min_{\deg(p) \leq s} \| f - p \|_{[a,b]}
        = \min_{\deg(p) \leq s} \| \tilde{f} - p \|_{[-1,1]}.
    \end{equation*}
    
    Note that if \( f\in \mathsf{Anl}(\rho,M,[a,b]) \) then \( \tilde{f} \in \mathsf{Anl}(\rho,M,[-1,1]) \) and if \( f\in \mathsf{BV}(d,V,[a,b]) \) then \( \tilde{f} \in \mathsf{BV}(d,V,[-1,1]) \).

The result then follows by setting the upper bounds in \cref{thm:bv_cheb,thm:analytic_cheb} to \( \epsilon/2 \) and solving for \( s \).
\end{proof}

Next, we consider polynomial approximations to 1-Lipshitz functions.
Note that there exist 1-Lipshitz functions whose derivatives are not of bounded variation. 
Therefore we cannot simply use \cref{thm:bv_cheb}.
Fortunately, the best approximation of differentiable functions is well studied. 
In particular, we have the following theorem due to Jackson; see \cite[Section 87]{achieser_92} and \cite[Section 6]{cheney_00} for details.
\begin{theorem} % see also cheney ch 4 sec 6
\label{thm:jackson}
    Suppose \( f \) is 1-Lipshitz on \( [-1,1] \); i.e., suppose that \( f \in \textup{\textsf{Lip}}(1,[-1,1]) \).
Then,
\begin{equation*}
    \min_{\deg(p) \leq s} \| f - p \|_{[-1,1]} \leq \frac{\pi}{2} (s+1)^{-1}.
\end{equation*}
\end{theorem}
In fact, the constant \( \pi/2 \) is the best possible under the stated conditions.

While the ``vanilla'' Chebyshev projection and interpolation do not attain this rate for all 1-Lipshitz functions, we can constructively obtain polynomial approximations which do attain this rate by damping.
\begin{definition}
\label{def:jackson_coeffs}
For \( i=0,1,\ldots, s \), the degree \( s \) Jackson's damping coefficients are
\begin{equation*}
    \rho_{i}^{J}
    = \frac{(s-i+2)\cos \left( \frac{i \pi}{s+2} \right) + \sin \left( \frac{i \pi}{s+2} \right)\cot \left( \frac{\pi}{s+2} \right)}{s+2}.
\end{equation*}
\end{definition}
The damped projection and interpolant then satisfy a similar bound to \cref{thm:jackson}.
\begin{restatable}{theorem}{dampedcheb}
    \label{thm:damped_cheb}
    Suppose \( f \in \textup{\textsf{Lip}}(1,[-1,1]) \), \( \mu  = \mu_{-1,1}^T \), and we use Jackson's damping coefficients as in \cref{def:jackson_coeffs}.
    Then for \( \circ \in \{ \textup{d-i}, \textup{d-a} \} \),
    \begin{equation*}
        \| f - \ff[]{f}{s} \|_{[-1,1]} \leq \frac{\pi^2}{2} (s+2)^{-1}.
    \end{equation*}
\end{restatable}
We provide a proof of this statement in \cref{sec:damped_cheb}.
While our proof is based closely on \cite{rivlin_81}, the exact constant we obtain is sharper than other bounds we know of for the quality of the damped projection \( \ff[d-a]{f}{s} \).
Moreover, our version of the proof works for the damped interpolant \( \ff[d-i]{f}{s} \). 
We were unable to find a similar result in the literature, although we suspect such a result is known.
%We remark that it is somewhat curious that the bound for the damped projection and interpolant have the same constant; indeed most bounds for the Chebyshev inter


\begin{subappendices}
\section{Proof of Jackson's theorem}
\label{sec:damped_cheb}

In this section we prove \cref{thm:damped_cheb}.
%\dampedcheb*
We follow Chapter 1 of \cite{rivlin_81} closely, starting with trigonometric polynomials on \( [-\pi,\pi] \) and then mapping to algebraic polynomials on \( [-1,1] \).
Throughout this section we maintain the notation of \cite{rivlin_81}, so the constants in this section do not necessarily have the same meaning as the rest of the paper.
In particular, \( n \) is the degree of the trigonometric polynomials used.



Given \( g:\R\to\R \), 1-Lipshitz and \( 2\pi \)-periodic, for \( \circ \in \{ \textup{i}, \textup{a} \} \), define
\begin{equation*}
    s_n^\circ(\theta) 
    := \frac{a_0^\circ}{2} 
    + \smop{\sum_{k=1}^{n}} \left( a_k^\circ \cos(k\theta) + b_k^\circ \sin(k\theta) \right)
\end{equation*}
where, for \( k=0,1,\ldots, n \)
\begin{equation*}
    a_k^\circ := \frac{1}{\pi}\int_{-\pi}^{\pi^-} g(\phi) \cos(k\phi) M_n^\circ(\phi) \,\d\phi
    ,\qquad
    b_k^\circ := \frac{1}{\pi}\int_{-\pi}^{\pi^-} g(\phi) \sin(k\phi) M_n^\circ(\phi) \,\d\phi.
\end{equation*}
Here \( M_n^{\textup{a}}(\phi) := 1 \) and 
\begin{equation*}
    M_n^{\textup{i}}(\phi) 
    :=  \frac{\pi}{n} \sum_{i\in\mathbb{Z}}\delta(\phi-\phi_i)
    ,\qquad 
    \phi_i 
    := \frac{2\pi(i-1/2)}{2n} - \pi
\end{equation*}
where \( \delta(\phi) \) is a Dirac delta distribution centered at zero.
Thus, \( s_n^\textup{a} \) is the truncation of the Fourier series of \( g \) while \( s_n^{\textup{i}} \) is the interpolant to \( g \) at the equally spaced nodes \( \{ \phi_i \}_{i-0}^{2n} \).

\begin{remark}
    Note that \( \int_{-\pi}^{\pi^-} \) means an integral over \( [-\pi,\pi) \); i.e. the upper endpoint of integration is excluded. 
    This is important for integrals involving \( M_n^{\textup{i}} \) which can have nonzero integral at a single point.
\end{remark}


Finally, define the damped interpolant/approximant
\begin{equation*}
    q_n^\circ(\theta) 
    := \frac{a^\circ_0}{2} + \smop{\sum_{k=1}^{n}} \rho_k \left( a_k^\circ \cos(k \theta) + b_k^\circ \sin(k\theta) \right)
\end{equation*}
where the damping coefficients \( \{ \rho_{k} \}_{k=1}^{n} \) are arbitrary real numbers.
Our aim is to bound \( \| q_n^\circ - g \|_{[-\pi,\pi]} \).

\begin{lemma}
    \label[lemma]{thm:Dn_zero}
    For all \( k=0,1,\ldots, n \),
    \begin{equation*}
        \frac{1}{\pi} \int_{-\pi}^{\pi^-} \cos(k \phi) M_n^\circ(\phi) \,\d\phi
        = \begin{cases}
            2 & k = 0 \\%\in 4n\mathbb{Z} \\
            0 & k \in 1,2,\ldots, n
        \end{cases}
    \end{equation*}
    and
    \begin{equation*}
        \frac{1}{\pi} \int_{-\pi}^{\pi^-} \sin(k x) M_n^\circ(\phi) \,\d\phi
        = 0.
    \end{equation*}
\end{lemma}

\begin{proof} 
    Clearly \( \frac{1}{\pi} \int_{-\pi}^{\pi^-} \cos(0 \phi) M_n^{\textup{a}} \d\phi = 2 \), and for \( k > 0 \), we have,
    \begin{equation*}
        \frac{1}{\pi} \int_{-\pi}^{\pi^-} \cos(k\phi) M_n^{\textup{a}}(\phi) \,\d\phi
%        = \frac{\sin(-k\pi) - \sin(k\pi)}{k \pi} 
        = 0.
    \end{equation*}
    By definition, 
    \begin{equation*}
        \frac{1}{\pi} \int_{-\pi}^{\pi^-} \cos(k\phi) M_n^{\textup{i}}(\phi) \,\d\phi
        = \frac{1}{n} \smop{\sum_{j=1}^{2n}} \cos(k \phi_j)
        = \frac{1}{n} \operatorname{Re} \smop{\sum_{j=1}^{2n}} \exp( \ii k \phi_j).
    \end{equation*}
    The case for \( k=0 \) is clear.
    Assume \( k > 0 \).
    Then, using that \( \phi_j = \frac{\pi}{n}(j-\frac{1}{2}) - \pi \) we have
    \begin{equation*}
        \frac{1}{n} \operatorname{Re} \smop{\sum_{j=1}^{2n}} \exp( \ii k \phi_j) 
        = \frac{1}{n} \operatorname{Re}\left[ \exp \left(\ii k\left(\frac{\pi}{2n}- \pi\right) \right) \smop{\sum_{j=0}^{2n-1}} \exp \left( \ii \frac{k\pi}{n} j \right) \right].
    \end{equation*}
    The result follows by observing that 
    \begin{equation*}
        \smop{\sum_{j=0}^{2n-1}} \exp\left( \ii \frac{k \pi}{n} j \right) 
        =  \frac{\exp(2\ii k \pi)-1}{\exp(\ii k \pi/n) - 1} 
        = 0.
    \end{equation*}
    Finally, since \( \sin(k\phi) \) is odd and \( M_n^\circ \) is symmetric about zero, the corresponding integrals are zero.
\end{proof}


We now introduce a generalized version of \cite[Lemma 1.4]{rivlin_81}.
\begin{lemma}
    Define
    \begin{equation*}
        u_n(\phi) := \frac{1}{2} + \smop{\sum_{k=1}^{n}} \rho_{k} \cos(k\phi).
    \end{equation*} 
    Then,
    \begin{equation*}
        q_n^\circ(\theta) = \frac{1}{\pi} \int_{-\pi}^{\pi^-}g(\phi+\theta) u_n(\phi) M_n^\circ(\phi+\theta) \,\d\phi.
    \end{equation*}
\end{lemma}

\begin{proof}
    First, note that
    \begin{align*}
        \pi q_n^\circ(\theta)
        &= \frac{1}{2} \bigg( \int_{-\pi}^{\pi^-} g(\phi) M_n^\circ(\phi) \,\d\phi  \bigg) 
        + \smop{\sum_{k=1}^{n}} \rho_k 
        \bigg( 
        \bigg( \int_{-\pi}^{\pi^-} g(\phi) \cos(k\phi)  M_n^\circ(\phi) \,\d\phi \bigg) \cos(k\theta)
        \\&\hspace{4em} + \bigg( \int_{-\pi}^{\pi^-} g(\phi) \sin(k\phi)  M_n^\circ(\phi) \,\d\phi \bigg) \sin(k\theta)
        \bigg)
        \\&= \int_{-\pi}^{\pi^-} g(\phi) \bigg( \frac{1}{2} + \smop{\sum_{k=1}^{n}} \rho_{k}( \cos(k\phi) \cos(k\theta) + \sin(k\phi) \sin(k\theta)) \bigg) M_n^\circ(\phi) \,\d\phi.
    \end{align*}
    % 
    Thus, using the identity \( \cos(\alpha)\cos(\beta) + \sin(\alpha)\sin(\beta) = \cos(\alpha-\beta) \) and the definition of \( u_n \)
    \begin{align*}
        q_n^\circ(\theta) 
        &= \frac{1}{\pi} \int_{-\pi}^{\pi^-} g(\phi) \left( \frac{1}{2} + \smop{\sum_{k=1}^{n}} \rho_{k} \cos(k (\phi-\theta)) \right) M_n^\circ(\phi) \,\d\phi
        \\&= \frac{1}{\pi} \int_{-\pi}^{\pi^-} g(\phi) u_n(\phi-\theta) M_n^\circ(\phi) \,\d\phi
    \end{align*} 
    Now, note that  \( g \), \( u_n \), and \( M_n^\circ \) are \( 2\pi \)-periodic so by a change of variables,
    \begin{align*}
        \frac{1}{\pi} \int_{-\pi}^{\pi^-} g(\phi) u_n(\phi-\theta) M_n^\circ(\phi) \,\d\phi
        &=\frac{1}{\pi} \int_{-\pi-\theta}^{\pi^--\theta} g(\phi+\theta) u_n(\phi) M_n^\circ(\phi+\theta) \,\d\phi
        \\&=\frac{1}{\pi} \int_{-\pi}^{\pi^-} g(\phi + \theta) u_n(\phi) M_n^\circ(\phi+\theta) \,\d\phi.
        \tag*{\qedhere}
    \end{align*}  
\end{proof}


Next, we prove a result similar to \cite[Lemma 1.7]{rivlin_81}, but by assuming that \( g \) is 1-Lipshitz we obtain a slightly better constant.
\begin{lemma}
    \label[lemma]{thm:damped_trig}
    Suppose \( u_n(\phi) \geq 0 \) for all \( \phi \).
    Then, if \( g \) is 1-Lipshitz,
    \begin{equation*}
        \|g - q_n^\circ \|_{[-\pi,\pi]} \leq \frac{\pi}{\sqrt{2}} ( 1 - \rho_{1})^{1/2} .
    \end{equation*} 
\end{lemma}


\begin{proof}
    Fix any \( \theta \in [-\pi,\pi] \).
    Recall that \( g \) is \( 1 \)-Lipshitz so that \( |g(\theta) - g(\phi+\theta)| \leq |\phi| \).
    Using this and the fact that \( u_n \) is non-negative,
    \begin{align*}
        |g(\theta) - q_n^\circ(\theta)|
        &= \bigg| \frac{1}{\pi} \int_{-\pi}^{\pi^-} \left( g(\theta) - g(\phi+\theta) \right)u_n(\phi) M_n^\circ(\phi+\theta) \,\d\phi \bigg|
        \\&\leq \frac{1}{\pi} \int_{-\pi}^{\pi^-} |\phi| u_n(\phi) M_n^\circ(\phi+\theta) \,\d\phi .
    \end{align*}
    Next, note \( M_n^\circ \) and \( u_n \) are \( 2\pi \)-periodic.
    Using this followed by the fact that \( \cos(k(\phi-\theta)) = \cos(k\phi) \cos(k\theta) - \sin(k\phi)\sin(k\theta) \), the definition of \( u_n \), and \cref{thm:Dn_zero}, we have
    \begin{equation*}
        \frac{1}{\pi} \int_{-\pi}^{\pi^-} M_n^\circ(\phi+\theta) u_n(\phi) \,\d\phi 
        = \frac{1}{\pi} \int_{-\pi}^{\pi^-} M_n^\circ(\phi) u_n(\phi-\theta) \,\d\phi 
        = 1.
    \end{equation*}
    Therefore, by the Cauchy--Schwarz inequality, 
    \begin{align*}
        \hspace{7em}&\hspace{-7em}\bigg(\frac{1}{\pi} \int_{-\pi}^{\pi^-} |\phi| u_n(\phi) M_n^\circ(\phi+\theta) \,\d\phi \bigg)^2 
        \\&=\bigg(\frac{1}{\pi} \int_{-\pi}^{\pi^-} |\phi| u_n(\phi) \cdot u_n(\phi) M_n^\circ(\phi+\theta)  \d\phi \bigg)^2
        \\&\leq \bigg( \frac{1}{\pi} \int_{-\pi}^{\pi^-} \phi^2 u_n(\phi) M_n^\circ(\phi+\theta) \,\d\phi \bigg)
        \bigg( \frac{1}{\pi} \int_{-\pi}^{\pi^-} u_n(\phi) M_n^\circ(\phi+\theta) \,\d\phi \bigg)
        \\&=\frac{1}{\pi} \int_{-\pi}^{\pi^-} \phi^2 u_n(\phi) M_n^\circ(\phi+\theta) \,\d\phi.
    \end{align*}
    Using the fact that  \( \phi^2 \leq \frac{\pi^2}{2} (1 -\cos(\phi)) \) we have 
    \begin{equation*}
        \frac{1}{\pi} \int_{-\pi}^{\pi^-} \phi^2 u_n(\phi) D_n(\phi+\theta) \,\d\phi
        \leq \frac{\pi^2}{2} \frac{1}{\pi} \int_{-\pi}^{\pi^-} (1-\cos(\phi)) u_n(\phi) M_n^\circ(\phi+\theta) \,\d\phi.
    \end{equation*}
    Next, we use that \( \cos(\phi) \cos(k\phi) = \frac{1}{2}( \cos((k-1)\phi) + \cos((k+1)\phi)) \) and \cref{thm:Dn_zero} to obtain
    \begin{equation*} 
        \frac{\pi^2}{2} \frac{1}{\pi} \int_{-\pi}^{\pi^-} (1-\cos(\phi)) u_n(\phi) M_n^\circ(\phi+\theta) \,\d\phi
        = \frac{\pi^2}{2} (1 - \rho_{1}).
    \end{equation*}  
    Combining this sequence of inequalities we find that
    \begin{equation*}
        |g(\theta) - q_n^\circ(\theta)| \leq \frac{\pi}{\sqrt{2}} ( 1 - \rho_{1})^{1/2}. \tag*{\qedhere}
    \end{equation*} 
\end{proof}


\iffalse
    Thus, from the definition of modulus of continuity,
    \begin{align*}
        |g(\theta) - q_n^\circ(\theta)|
        &= \left| \frac{1}{\pi} \int_{-\pi}^{\pi^-} \left( g(\theta) - g(\phi+\theta) \right)u_n(\phi) D_n(\phi+\theta) \,\d\phi \right|
        \\&= \left| \frac{1}{\pi} \int_{-\pi}^{\pi^-} \omega(|\phi|) u_n(\phi) M_n^\circ(\phi+\theta) \,\d\phi \right|
    \end{align*}
    By \note{todo}, 
    \begin{align*}
        \omega(|\phi|) \leq \left( 1 + n |\phi| \right) \omega \left( \frac{1}{n} \right)
    \end{align*}
    Thus,
    \begin{align*}
        |g(\theta) - q_n^\circ(\theta)|
        \leq \omega \left( \frac{1}{n} \right) \left( 1 + \frac{n}{\pi} \int_{-\pi}^{\pi^-} |\phi| u_n(\phi) M_n^\circ(\phi+\theta) \,\d\phi \right).
    \end{align*}
    Next, by Cauchy--Schwarz inequality,
    \begin{align*}
        \frac{1}{\pi} \int_{-\pi}^{\pi^-} |\phi| u_n(\phi) M_n^\circ(\phi+\theta) \,\d\phi 
        &=\frac{1}{\pi} \int_{-\pi}^{\pi^-} (|\phi| u_n(\phi)M_n^\circ(\phi+\theta))^{1/2} (u_n(\phi)M_n^\circ(\phi+\theta))^{1/2}  \d\phi 
        \\&\leq \left( \frac{1}{\pi} \int_{-\pi}^{\pi^-} \phi^2 u_n(\phi) M_n^\circ(\phi+\theta) \,\d\phi \right)^{1/2} 
        \left( \frac{1}{\pi} \int_{-\pi}^{\pi^-} u_n(\phi) M_n^\circ(\phi+\theta) \,\d\phi \right)^{1/2}
        \\&=\left( \frac{1}{\pi} \int_{-\pi}^{\pi^-} \phi^2 u_n(\phi) M_n^\circ(\phi+\theta) \,\d\phi \right)^{1/2} 
        \tag*{\note{why does proof have $\leq$?}}
    \end{align*}
    Using the fact that  \( \phi^2 \leq \frac{\pi^2}{2} (1 -\cos(\phi)) \) we have 
    \begin{align*}
        \frac{1}{\pi}\int_{[-\pi,\pi)} \phi^2 u_n(\phi) D_n(\phi+\theta) \,\d\phi
        &\leq \frac{\pi^2}{2} \frac{1}{\pi}\int_{[-\pi,\pi)} (1-\cos(\phi)) u_n(\phi) M_n^\circ(\phi+\theta) \,\d\phi
    \end{align*}
    Next, we use that \( \cos(\phi) \cos(k\phi) = \frac{1}{2}( \cos((k-1)\phi) + \cos((s+1)\phi)) \) and \cref{thm:Dn_zero} to obtain
    \begin{align*} 
        \frac{\pi^2}{2} \frac{1}{\pi}\int_{[-\pi,\pi)} (1-\cos(\phi)) u_n(\phi) M_n^\circ(\phi+\theta) \,\d\phi
        &= \frac{\pi^2}{2} (1 - g_{1,n}).
    \end{align*}  
    
    Combining this sequence of inequalities we find that
    \begin{align*}
        |g(\theta) - q_n^\circ(\theta)| \leq \omega \left( \frac{1}{n} \right) \left[ 1 + \frac{n \pi}{\sqrt{2}} ( 1 - g_{1,n})^{1/2} \right].
    \end{align*} 
\fi

    \begin{lemma}
        \label[lemma]{thm:damped_weights}
        If we use Jackson's damping coefficents from \cref{def:jackson_coeffs}, then \( u_n \) is positive and
        \begin{equation*}    
        \frac{\pi}{\sqrt{2}} ( 1 - \rho_{1})^{1/2}
            \leq \frac{\pi^2}{2}(n+2)^{-1}.
        \end{equation*}        
    \end{lemma}
    
\begin{proof}
    Let \( \{ c_\ell \}_{\ell=0}^{n} \) be any real numbers. 
    Then 
    \begin{equation*}
        \bigg( \smop[r]{\sum_{\ell=0}^{n}} c_\ell \exp( \ii \ell \theta) \bigg)
        \bigg( \smop[r]{\sum_{\ell=0}^{n}} c_\ell \exp( - \ii \ell \theta) \bigg)
        =
        \bigg| \smop[r]{\sum_{\ell=0}^{n}} c_\ell \exp( \ii \ell \theta) \bigg|^2
        \geq 0.
    \end{equation*}
    Expanding and using that \( \exp(\ii k \theta) + \exp(-\ii k \theta) = 2\cos(k\theta) \) we find
    \begin{equation*}
        \bigg( \smop[r]{\sum_{\ell=0}^{n}} c_\ell \exp( \ii \ell  \theta) \bigg)
        \bigg( \smop[r]{\sum_{\ell=0}^{n}} c_\ell \exp( - \ii \ell \theta) \bigg)
        = \smop{\sum_{k=0}^{n}} c_k^2 
        + 2\sum_{p=1}^{n} \sum_{k=0}^{n-p} c_k c_{k+p} \cos(p\theta). 
    \end{equation*}
    Because \( u_n \) must have the constant term equal to \( 1/2 \) we require \( c_0^2 + \ldots + c_n^2 = 1/2 \).

    For \( \ell = 0,1,\ldots, n \), let
    \begin{equation*}
        c_\ell = c \sin \left( \frac{\ell+1}{n+2} \pi \right)
    \end{equation*}
    where
    \begin{equation*}
        c^2 = \bigg( \smop[r]{\sum\limits_{\ell=0}^{n}} 2\sin^2 \left( \frac{\ell+1}{n+2} \pi \right) \bigg)^{-1}
        = \frac{1}{n+2}.
    \end{equation*}
    Then setting \( \rho_0 = 1 \) and 
    \begin{equation*}
        \rho_k = 2 \smop{\sum_{\ell=0}^{n-k}} c_\ell c_{c+k}
    \end{equation*}
    we obtain \( u_n \).

    Next, we show that these damping coefficients are equal to those described above. 
    Following \cite{weisse_wellein_alvermann_fehske_06} we have that
    \begin{align*}
        2 \smop{\sum_{\ell=0}^{n-k}} c_\ell c_{\ell+k}
        &= 2 c^2 \smop{\sum_{\ell=0}^{n-k}} \sin\left( \frac{\ell+1}{n+2} \pi \right) \sin \left( \frac{\ell+k+1}{n+2}\pi \right)
        \\&= 2 c^2 \smop{\sum_{\ell=1}^{n-k+1}} \sin\left( \frac{\ell}{n+2} \pi \right) \sin \left( \frac{\ell+k}{n+2}\pi \right)
        \\&= c^2 \smop{\sum_{\ell=1}^{n-k+1}} \left( \cos\left( \frac{k}{n+2} \pi \right) - \cos \left( \frac{2\ell+k}{n+2}\pi \right) \right)
    % \sin(\alkha)\sin(\beta) = \frac{1}{2} \left( \cos( \alkha -\beta ) - \cos(\alkha + \beta) \right)  
        \\&= c^2 \left( (n-k) \cos\left( \frac{k}{n+2} \pi \right) 
             - \Re \smop{\sum_{\ell=1}^{n-k+1}} \exp\left(\ii \frac{2\ell+k}{n+2}\pi \right) \right)
        \\&= c^2 \left( (n-k+1) \cos\left( \frac{k}{n+2} \pi \right) 
             - \sin \left( \frac{\ell}{n+2}\pi \right) \cot \left( \frac{\pi}{n+2} \right) \right).
    \end{align*}
    These are exactly Jackson's damping coefficients.

    Using this expression, it's easy to verify that \( \rho_1 = \cos(\pi / (n+2)) \).
    Thus
    \begin{equation*}
        (1-\rho_{1})^{1/2} 
        = \left( 1 - \cos \left( \frac{\pi}{n+2} \right) \right)^{1/2}
        = \sqrt{2} \sin \left( \frac{\pi}{2n+4} \right)
        \leq \sqrt{2} \frac{\pi}{2n+4}
    \end{equation*}
    so
    \begin{equation*}
%        \label{eqn:1-rho1}
        \frac{\pi}{\sqrt{2}} ( 1 - \rho_{1})^{1/2}
        \leq \frac{\pi^2}{2n+4}
%        \leq \frac{\pi^2}{2}n^{-1}.
        \tag*{\qedhere}
    \end{equation*} 
\end{proof}

Finally, we prove the desired theorem.
\begin{proof}[Proof of \cref{thm:damped_cheb}]

    Without loss of generality, we can consider the case that \( f \) is 1-Lipshitz.
%    \note{this needs to be done more carefully, but I didn't want to dump too much time into it until I'm more confident in the rest of the results in this section. The techniques are also basically standard and always deferred to exercises in books. I will clean up later.}
    For \( \theta \in [-\pi,\pi) \) define \( g \) by \( g(\theta) = f(\cos(\theta)) \).
    Then \( g \) is 1-Lipshitz, \( 2\pi \)-periodic, and even. % since \( \cos \) is \( 2\pi \)-periodic and even.
    Next define the inverse mapping of the damped trigonometric polynomial \( q_n^\circ \) for \( g \) as
    \begin{equation*}
        p_n^\circ(t) = q_n^\circ(\arccos(t)).
    \end{equation*}
    For any \( t\in [-1,1] \), setting \( \theta = \arccos(t) \in [0,\pi] \) we use \cref{thm:damped_trig,thm:damped_weights} to obtain the bound
    \begin{equation*}
        | p_n^\circ(t) - f(t) |
        = | p_n^\circ(\cos(\theta)) - f(\cos(\theta)) |
        = | q_n^\circ(\theta) - g(\theta) |
        \leq \frac{\pi^2}{2} n^{-1}.
    \end{equation*}

    We will now show that \( p_n^\circ(t) = [f]_n^{\textup{d-}\circ} \).
    The mapping \( \theta = \arccos(t) \) gives the Chebyshev polynomials; indeed, it is well known that
    \begin{equation*}
        T_k(t) = \cos(k\arccos(t)).
    \end{equation*}
    Since \( g \) is even we have \( b_k^\circ = 0 \) so 
    \begin{equation*}
        p_n^\circ(t) 
        = q_n^\circ(\arccos(t)) 
        = \frac{a_0^\circ}{2} + \smop{\sum_{k=1}^{n}} \rho_k a_k^\circ T_k(t).
    \end{equation*}
    Thus, our goal is to show that \( a_k^\circ \) are the coefficients for the Chebyshev approximation/interpolation series.

    Towards this end, recall that
    \begin{equation*}
        a_k^\circ 
        = \frac{1}{\pi} \int_{-\pi}^{\pi^-}g(\phi) \cos(k\phi)M_n^\circ(\phi) \,\d\phi.
    \end{equation*}
    Since \( g \) is even we can replace the integral on \( [-\pi,\pi) \) an integral on \( (0,\pi) \) and an integral on \( [0,\pi) \).
    We first consider the case \( M_n^{\textup{a}}(\phi) = 1 \).
    Noting that \( - \pi^{-1}  \arccos(t) = \mu_{-1,1}^T(t) \), we find
    \begin{equation*}
        a_k^\textup{a} = 2 \int_{-1}^{1} f T_k \,\d \mu_{-1,1}^T
    \end{equation*}
    as desired.
    For \( j=1,2,\ldots, 2n \), we have the Chebyshev nodes
    \begin{equation*}
        \cos(\phi_j) 
        = \cos \left( \frac{2\pi (j-1/2)}{2n} - \pi \right)
        = - \cos \left( \frac{\pi (j-1/2)}{n}\right).
    \end{equation*}
    Thus,
    \begin{equation*}
        M_n^\textup{i}(t) = \frac{\pi}{n} \sum_{i\in \mathbb{Z}}  \delta(x-x_i)
        ,\qquad x_i = - \cos \left( \frac{\pi (j-1/2)}{n}\right)
    \end{equation*}
    so
    \begin{equation*}
        a_k^\textup{i} = 2 \smop{\sum_{i=1}^{2n}} \frac{\pi}{n} f(x_i) T_k(x_i)
        = 2 \int_{-1}^{1} f T_k \,\d \qq[g]{\mu_{-1,1}^T}{n}
    \end{equation*}
    as desired.    
    The result follows by renaming \( n \) to \( s \).
\end{proof}



\end{subappendices}









\chapter{Matrix-free quadrature}
\label{chap:QF}

    This chapter focuses primarily on quadrature rules for the \emph{weighted} CESM induced by \( \vec{A} \) and a unit vector \( \vec{v} \).
    \begin{definition}
    \label{def:wCESM}
        The weighted CESM \( \Psi : \R\to[0,1] \), induced by \( \vec{A} \) and a unit vector \( \vec{v} \), is defined by
    \begin{equation*}
        \Psi(x)
        = \Psi_{\vec{A},\vec{v}}(x)
        := \vec{v}^\cT \bOne[\vec{A}\leq x] \vec{v}.
%        = x\mapsto \smop{\sum_{i=0}^{n-1}} |\vec{v}^\cT \vec{u}_i|^2 \bOne[ \lambda_i \leq x].
    \end{equation*}
    \end{definition}
    This definition implies that
    \begin{equation*}
        \int f \,\d\Psi = \vec{v}^\cT \fA \vec{v},
    \end{equation*}
    so it is clear that \( \Psi \) is closely related to the task of approximating \( \vec{v}^\cT \fA \vec{v} \). 
    In fact, in this chapter, we take the perspective that Krylov subspace methods for \( \vec{v}^\cT \fA \vec{v} \) are in correspondence with quadrature rules for \( \Psi \).
    Such a perspective was popularized by \cite{golub_meurant_94,golub_meurant_09}

    \begin{remark}
        It is now clear that the Lanczos algorithm \cref{alg:lanczos} is simply the Stieltjes procedure \cref{alg:stieltjes} applied to the weighted CESM \( \Psi \).
        Specifically, \( \vec{q}_i \propto p_i\A\vec{v} \) for \( i=0,1,\ldots, k \) and the tridiagonal matrix \( \vec{T} \) generated by Lanczos is equal to the Jacobi matrix \( \vec{M}(\Psi) \).
    \end{remark}

    %This chapter builds towards developing quadrature approximations to the CESM \( \Phi \) defined in \cref{eqn:CESM} which we discuss in \cref{chap:random_quadrature}

    The chapter title, \emph{matrix-free quadrature} refers to the fact that our approach to constructing quadrature approximations to \( \Psi \) involve matrix-free algorithms; i.e. algorithms which access \( \vec{A} \) only through matrix-vector products.
    While the quadrature rules we study are standard in approximation theory, it is worth noting several critical differences between classical quadrature methods and the algorithm studied in this chapter.
    First, the costs of the algorithms in this paper are determined primarily by the number of matrix-vector products.
    This is because we typically only want to approximate \( \int f\,\d\Psi \) for a single, or perhaps a few, functions. 
    On the other hand, the weight functions which classical quadrature rules approximate never change, so nodes and weights can be precomputed and the dominant cost becomes the cost to evaluate \( f \) at the quadrature nodes.
    Second, while classical weight functions, such as the weight functions for Jacobi or Hermite polynomials, are typically relatively uniform in the interior of the interval of integration, \( \Psi \) may vary wildly from application to application. 
    In some cases \( \Psi \) might resemble the distribution function of a classical weight function whereas in others it might have large gaps, jumps, and other oddities.
    These distinctions are hinted at throughout the chapter and illustrated explicitly in the numerical examples at the end of this chapter.


    Moment based methods for estimating the weighted CESM\footnote{In physics, the ``density'' \( \d\Psi/\d{x} \) is often called the local density of states (local DOS).} have been used in physics for at least half a century.
    Early approaches were based on monomial moments \cite{cyrotlackmann_67,cyrotlackmann_69,ducastelle_cryotlackmann_70,ducastelle_cryotlackmann_71,cryotlackmann_ducastelle_71}, but the use of modified Chebyshev moments \cite{wheeler_blumstein_72} and Lanczos-based approaches \cite{haydock_heine_kelly_pendry_72,haydock_heine_kelly_72,haydock_heine_kelly_75} were soon introduced.


    %\cite{nex_78} integrals


    \section{Extracting moments from a Krylov subspace}
    \label{sec:krylov_moments}


    Bases for the Krylov subspace \( \mathcal{K}_{k+1} = \operatorname{span}\{ \vec{v}, \vec{A}\vec{v}, \ldots, \vec{A}^k \vec{v} \} \) can be computed using \( k \) matrix-vector products with \( \vec{A} \) and contain a wealth of information about the interaction of \( \vec{A} \) with \( \vec{v} \); in particular, they contain the information necessary to compute the moments of \( \Psi \) through degree \( 2k \).
    Indeed, for all \( i,j \geq 0 \),
    \begin{equation*}
        (\vec{A}^j\vec{v})^\cT (\vec{A}^j\vec{v}) = \vec{v}^\cT \vec{A}^{i+j} \vec{v} = \int x^{i+j} \,\d\Psi.
    \end{equation*}
    Note, however, that it is sometimes more straightforward to obtain the moments through degree \( s \), for some \( s\leq 2k \).
    Thus, we will use \( s \) to denote the degree of the maximum moment we compute and \( k \) to denote the number of matrix-vector products used.



    \subsection{Computing modified moments directly}
    \label{sec:moments_direct}
   
    Perhaps the most obvious approach to computing modified moments is to construct the basis \( [ p_0\A\vec{v}, \ldots, p_{k}\A\vec{v}] \) for \( \mathcal{K}_{k+1} \) and then compute 
    \begin{equation*}
        \vec{v}^\cT [ p_0\A\vec{v}, \ldots, p_{k}\A\vec{v}].
    \end{equation*}
    This can be done using \( k \) matrix-vector products and \( O(n) \) storage using the matrix recurrence version of \cref{eqn:poly_three_term}.
    Indeed for all \( i=0,1,\ldots, k-1 \) we have that
    \begin{equation*}
        \vec{A} p_i\A \vec{v} =  \beta_{i-1} p_{i-1}\A\vec{v} + \alpha_i p_i\A\vec{v} + \beta_{i} p_{i+1}\A\vec{v}
    \end{equation*}
    from which we can implement an efficient matrix-free algorithm to compute the modified moments \( \{ m_i \}_{i=0}^{k} \), as shown in \cref{alg:moments}.

    \begin{labelalgorithm}[H]{moments}{get-moments}{Get modified moments of \( \Psi \) wrt. \( \mu \)}
    \begin{algorithmic}[1]
        \Procedure{\thealgorithmname}{$\vec{A}, \vec{v}, k, \mu$}
        \State \( \vec{q}_0 = \vec{v} \), \( m_0 = \vec{v}^\cT\vec{v} \), \( \vec{q}_{-1} = \vec{0} \), \( \beta_{-1} = 0 \)
        \For {\( i=0,1,\ldots, k-1 \)}
            \State \( \vec{q}_{i+1} = \frac{1}{\beta_{i}} \left( \vec{A} \vec{q}_i - \alpha_i \vec{q}_i - \beta_{i-1} \vec{q}_{i-1} \right) \)
            \State \( m_{i+1}  = \vec{v}^\cT \vec{q}_{i+1} \)
        \EndFor
        \State \Return \( \{ m_i \}_{i=0}^{k} \)
    \EndProcedure
    \end{algorithmic}
    \end{labelalgorithm}

    If we instead compute 
    \begin{equation*}
        [ p_0\A\vec{v}, \ldots, p_{k}\A\vec{v}]^\cT
        [ p_0\A\vec{v}, \ldots, p_{k}\A\vec{v}],
    \end{equation*}
    then we have the information required to compute the modified moments though degree \( 2k \).
    However, it is not immediately clear how to do this without the \( O(kn) \) memory required to store a basis for \( \mathcal{K}_{k+1} \).
    It turns out it is indeed generally possible to compute these moments without storing the entire basis, and we discuss a principled approach for doing so using connection coefficients in \cref{sec:connection_coeffs}.

    One case where extracting the moments to degree \( 2k \), without storing a basis for Krylov subspace, is straightforward is when \( \mu = \mu_{a,b}^T \).
    This is because, for all \( i \geq 0 \), the Chebyshev polynomials satisfy the identities
    \begin{equation*}
        T_{2i} = 2 (T_i)^2 - 1 
        ,\qquad
        T_{2i+1} = 2 T_i T_{i+1} - x. 
    \end{equation*}
    Thus, using the recurrence for the Chebyshev polynomials and their relation to the orthogonal polynomials with respect to \( \mu_{a,b}^T \), we obtain \cref{alg:cheb_moments}.
    This algorithm is well-known in papers on the kernel polynomial method are variants; see for instance \cite{skilling_89,silver_roder_94,weisse_wellein_alvermann_fehske_06,hallman_21}.% and is based on related ideas from \cite{}.

    \begin{remark}   
    The Chebyshev polynomials grow rapidly outside of the interval \( [-1,1] \).
    Therefore, if the spectrum of \( \vec{A} \) extends beyond this interval, then computing the Chebyshev polynomials in \( \vec{A} \) may suffer from numerical instabilities.
        Instead, the distribution function \( \mu_{a,b}^T \) and corresponding orthogonal polynomials should be used for some choice of \( a \) and \( b \) with \( \mathcal{I} \subset [a,b] \).

    %then \( \vec{A} \) can be scaled and shifted so the eigenvalues are in \( [-1,1] \) using the transform
    %\begin{equation*}
    %    \vec{A} \mapsto \frac{2}{b-a} \left( \vec{A} - \frac{a+b}{2}\vec{I} \right).
    %\end{equation*}
    %    This transformed matrix doesn't need to be computed explicitly since matrix-vector products with a vector \( \vec{x} \) can be written as a linear combination of \( \vec{A}\vec{x} \) and \( \vec{x} \); e.g. \( \frac{2}{b-a} ( \vec{A} - \frac{a+b}{2}\vec{I} ) \vec{x} = \frac{2}{b-a}( \vec{A}\vec{x} - \frac{a+b}{2} \vec{x} )\).
    \end{remark}


    \begin{labelalgorithm}[H]{cheb_moments}{get-Chebyshev-moments}{Get modified moments of \( \Psi \) wrt. \( \mu_{a,b}^T \)}
    \begin{algorithmic}[1]
        \Procedure{\thealgorithmname}{$\vec{A}, \vec{v}, k, a,b$}
        \State \( \vec{q}_0 = \vec{v} \), \( m_0 = \vec{q}_0^\cT \vec{q}_0 \)
        \State \( \vec{q}_1 = \frac{2}{b-a}(\vec{A} \vec{q}_0 - \frac{a+b}{2} \vec{q}_0) \), \( m_1 =  \sqrt{2} \vec{q}_0^\cT \vec{q}_1 \)
        \For {\( i=1,2,\ldots, k-1 \)}
            \State \( m_{2i}  = \sqrt{2} (2\vec{q}_i^\cT \vec{q}_{i} - m_0) \)
            \State \( \vec{q}_{i+1} = 2 \frac{2}{b-a} (\vec{A} \vec{q}_i - \frac{a+b}{2} \vec{q}_i) - \vec{q}_{i-1} \)
            \State \( m_{2i+1}  = \sqrt{2} (2\vec{q}_i^\cT \vec{q}_{i+1}) - m_1 \)
        \EndFor
        \State \( m_{2k} = \sqrt{2}( 2 \vec{q}_k^\cT\vec{q}_k - m_0) \)
        \State \Return \( \{ m_i \}_{i=0}^{2k} \)
    \EndProcedure
    \end{algorithmic}
    \end{labelalgorithm}


    \subsection{Connection coefficients to compute more modified moments}
    \label{sec:connection_coeffs}

    We now discuss how to use connection coefficients to compute the modified moments of \( \Psi \) with respect to \( \mu \) given knowledge of either
    (i) the modified moments of \( \Psi \) with respect to some distribution \( \nu \) or 
    (ii) the tridiagonal matrix computed using \cref{alg:lanczos}.
    Much of our discussion on connection coefficients is based on \cite{webb_olver_21}; see also \cite{fischer_golub_91}.

    \begin{definition}
        \label[definition]{def:connection_coeffs}
        The connection coefficient matrix \( \vec{C} = \vec{C}_{\mu\to\nu} \) is the upper triangular matrix representing a change of basis between the orthogonal polynomials \( \{ p_i \}_{i=0}^{\infty} \) with respect to \( \mu \) and the orthogonal polynomials \( \{ q_i \}_{i=0}^{\infty} \) with respect to \( \nu \), whose entries satisfy,
        \begin{equation*}
            p_s = [\vec{C}]_{0,s} q_0 +  [\vec{C}]_{1,s} q_1 +  \cdots + [\vec{C}]_{s,s} q_s.
        \end{equation*}
    \end{definition}

    \Cref{def:connection_coeffs} implies that, for all \( i = 0,1, \ldots, s \),
    \begin{equation*}
        m_i =
        \int p_i \,\d\Psi
        = 
        \smop{\sum_{j=0}^{i}} [\vec{C}]_{j,i} \int q_j \,\d\Psi
        =
        \smop{\sum_{j=0}^{i}} [\vec{C}]_{j,i} n_j
    \end{equation*}
    where \( \{n_i\}_{i=0}^{s} \) are the modified moments of \( \Psi \) with respect to \( \nu \).
    Thus, we can easily obtain the modified moments \( \{ m_i \}_{i=0}^{s} \) of \( \Psi \) with respect to \( \mu \) from the modified moments of \( \Psi \) with respect to \( \nu \).
    In particular, if \( \vec{m} \) and \( \vec{n} \) denote the vectors of modified moments, then \( \vec{m} = \vec{C}^\rT \vec{n} \).

    Moreover, in the special case that \( \nu \) has the same moments as \( \Psi \) through degree \( s \), so, for any \( j\geq 0 \), 
    \begin{equation*}
        n_j = \int q_j \,\d\Psi 
        = \int q_j \,\d\nu
        = \int q_0 q_j \,\d\nu 
        = \bOne[j = 0] .
    \end{equation*}
    Therefore, the modified moments of \( \Psi \) (with respect to \( \mu \)) through degree \( s \) can be computed by
    \begin{equation*}
        m_i=
        \int p_i \,\d\Psi = 
        [ \vec{C}]_{0,i}.
    \end{equation*}

    In order to use the above expressions, we must compute the connection coefficient matrix.
    \Cref{def:connection_coeffs} implies that for all \( i\leq j \), the entries of the connection coefficient matrix are given by
    \begin{equation*}
        [\vec{C}]_{i,j} = \int q_i p_j \,\d\nu.
    \end{equation*}    
    Unsurprisingly, then, the entries of the connection coefficient matrix \( \vec{C} = \vec{C}_{\mu\to\nu} \) can be obtained by a recurrence relation.
    \begin{proposition}[{\cite[Corollary 3.3]{webb_olver_21}}]
    \label[proposition]{thm:connection_coefficient_recurrence}
        Suppose the Jacobi matrices for \( \mu \) and \( \nu \) are respectively given by
        \begin{equation*}
            \vec{M}(\mu) =
        \begin{bmatrix}
            \alpha_0 & \beta_0 \\
            \beta_0 & \alpha_1 & \beta_1 & \phantom{\ddots}\\
            &\beta_1 & \alpha_2 & \ddots \\
            &&\ddots & \ddots \\
        \end{bmatrix}
        ,\qquad
            \vec{M}(\nu) =
        \begin{bmatrix}
            \gamma_0 & \delta_0 \\
            \delta_0 & \gamma_1 & \delta_1 & \phantom{\ddots}\\
            &\delta_1 & \gamma_2 & \ddots \\
            &&\ddots & \ddots \\
        \end{bmatrix}.
        \end{equation*}

        Then the entries of \( \vec{C} = \vec{C}_{\mu\to\nu} \) satisfy, for \( i,j \geq 0 \), the following recurrence:
    \begin{align*}
        [\vec{C}]_{0,0} &= 1 
        \\
        [\vec{C}]_{0,1} &= (\gamma_0 - \alpha_0)/\beta_0
        \\
        [\vec{C}]_{1,1} &= \delta_0 / \beta_0 
        \\
        [\vec{C}]_{0,j} &= ((\gamma_0 - \alpha_{j-1}) [\vec{C}]_{0,j-1} + \delta_0 [\vec{C}]_{2,j-1} - \beta_{j-2} [\vec{C}]_{0,j-2} ) / \beta_{j-1}
        \\
        [\vec{C}]_{i,j} &= (\delta_{i-1} [\vec{C}]_{i-1,j-1} + (\gamma_i - \alpha_{j-1}) [\vec{C}]_{i,j-1} + \delta_i[\vec{C}]_{i+1,j-1} - \beta_{j-2} [ \vec{C}]_{i,j-2}) / \beta_{j-1}.
    \end{align*}
    \end{proposition}

    \Cref{thm:connection_coefficient_recurrence} yields a natural algorithm for computing the connection coefficient matrix \( \vec{C}_{\mu\to\nu} \).
    This algorithm is shown as \cref{alg:connection_coeffs}.
    Note that \( \vec{C} \) is, by definition, upper triangular, so \( [\vec{C}]_{i,j}  = 0 \) whenever \( i > j \).
    We remark that for certain cases, particularly transforms between the Jacobi matrices of classical orthogonal polynomials, faster algorithms are known \cite{townsend_webb_olver_17}.
    We do not focus on such details in this paper as the cost of products with \( \vec{A} \) is typically far larger than the cost of computing \( \vec{C}_{\mu\to\nu} \).

    \begin{labelalgorithm}[H]{connection_coeffs}{get-connection-coeffs}{Get connection coefficients}
    \begin{algorithmic}[1]
        %\Procedure{\thealgorithmname}{$[\vec{M}(\mu)]_{:k_{m_1},:k_{m_2}},[\vec{M}(\nu)]_{:k_{n_1},:k_{n_2}}$}
        \Procedure{\thealgorithmname}{$\mu,k_{\mu},k_{\mu}',\nu,k_{\nu},k_{\nu}'$}
        \State \( [\vec{C}]_{0,0} = 1 \), \( [\vec{C}]_{i',j'} = 0 \) if \( i' > j' \) or \( j' = -1 \)
        \For {$j=1,2,\ldots,\min(k_{\mu}',k_{\nu}+k_{\nu}')$}
        \For {$i=0,1,\ldots,\min(j,k_{\nu}+k_{\nu}'-j)$}
%        \State {$[\vec{C}]_{i,j} = \frac{1}{\beta_{j-1}}(\delta_{i-1} [\vec{C}]_{i-1,j-1} + (\gamma_i - \alpha_{j-1}) [\vec{C}]_{i,j-1} + \delta_i[\vec{C}]_{i+1,j-1} - \beta_{j-2} [ \vec{C}]_{i,j-2}) $}
        \State 
        \(
        \begin{aligned}%
            \relax[\vec{C}]_{i,j} = (&\delta_{i-1} [\vec{C}]_{i-1,j-1} + (\gamma_i - \alpha_{j-1}) [\vec{C}]_{i,j-1} 
            \\[-.6em]&\hspace{3em}+ \delta_i[\vec{C}]_{i+1,j-1} 
            - \beta_{j-2} [ \vec{C}]_{i,j-2}) / \beta_{j-1} 
        \end{aligned}\)
        \EndFor
        \EndFor
        \State \Return \( \vec{C} = \vec{C}_{\mu\to\nu} \)
    \EndProcedure
    \end{algorithmic}
    \end{labelalgorithm}

    \begin{remark}
        From \cref{thm:connection_coefficient_recurrence} it is not hard to see that \( [\vec{C}]_{:k,:k} \) can be computed using \( [\vec{M}(\nu)]_{:k,:k} \) and \( [\vec{M}(\mu)]_{:k,:k} \). 
        Moreover, \( [\vec{C}]_{0,:2k+1} \) can be computed using \( [\vec{M}(\nu)]_{:k+1,:k} \) and \( [\vec{M}(\mu)]_{:2k,:2k} \). 
        In general \( \vec{M}(\mu) \) will be known fully, and in such cases, the modified moments through degree \( 2k \) can be computed using the information generated by Lanczos run for \( k \) iterations.
    \end{remark}

    We can use \cref{alg:connection_coeffs} in conjunction with \cref{alg:cheb_moments} and \cref{alg:lanczos} to compute modified moments with respect to \( \mu \).
    This is shown in \cref{alg:moments_cheb} and \cref{alg:moments_lanc} respectively.

    \begin{labelalgorithm}[H]{moments_cheb}{get-moments-from-Cheb}{Get modified moments wrt. \( \mu \) of weighed CESM (via Chebyshev moments)}
    \begin{algorithmic}[1]
        \Procedure{\thealgorithmname}{$\vec{A}, \vec{v}, s, \mu, a, b$}
        \State \( k = \lceil s/2 \rceil \)
        \State \( \{ n_i \}_{i=0}^{2k} = \text{\ref{alg:name:cheb_moments}}(\vec{A},\vec{v}, k,a,b)  \)
        \State \( \vec{C} = \text{\ref{alg:name:connection_coeffs}}([\vec{M}(\mu)]_{:2k,:2k},[\vec{M}(\mu_{a,b}^T)]_{:2k,:2k}) \)
        \For{\( i=0,1,\ldots, s \)}
            \State \( m_i = \sum_{j=0}^{i} [\vec{C}]_{j,i} n_j \)
        \EndFor
        \State \Return \( \{ m_i \}_{i=0}^{s} \)
    \EndProcedure
    \end{algorithmic}
    \end{labelalgorithm}


    \begin{labelalgorithm}[H]{moments_lanc}{get-moments-from-Lanczos}{Get modified moments wrt. \( \mu \) of weighed CESM (via Lanczos)}
    \begin{algorithmic}[1]
        \Procedure{\thealgorithmname}{$\vec{A}, \vec{v}, s, \mu$}
        \State \( k = \lceil s/2 \rceil \)
        \State \( [\vec{T}]_{i:k+1,:k} = \text{\ref{alg:name:lanczos}}(\vec{A},\vec{v}, k)  \)
        \State \( \vec{C} = \text{\ref{alg:name:connection_coeffs}}([\vec{M}(\mu)]_{:2k,:2k},[\vec{T}]_{:k+1,:k}) \)
        \For{\( i=0,1,\ldots, s \)}
            \State \( m_i = [\vec{C}]_{0,i} \)
        \EndFor
        \State \Return \( \{ m_i \}_{i=0}^{s} \)
    \EndProcedure
    \end{algorithmic}
    \end{labelalgorithm}
    
    \begin{remark}
    Given the modified moments of \( \Psi \) with respect to \( \mu \), the tridiagonal matrix produced by \cref{alg:lanczos} can itself be obtained \cite{sack_donovan_71}.
    This is quite similar to \( s \)-step Lanczos methods designed to reduce communication on distributed memory computers.
    However, if implemented naively, such methods can be even more susceptible to the effects of finite precision arithmetic than the regular Lanczos method \cite{carson_demmel_15,carson_20}, so special care must be taken when implementing such an algorithm.
    \end{remark}


\section{Quadrature approximations for weighted spectral measures}
\label{sec:quadrature}

We now discuss how to use the information extracted by the algorithms in the previous section to obtain quadrature rules for the weighted CESM \( \Psi \).
We begin with a discussion on quadrature by interpolation in \cref{sec:iq} followed by a discussion on Guassian quadrature in \cref{sec:gq} and quadrature by approximation in \cref{sec:aq}.
Finally, in \cref{sec:damped}, we describe how damping can be used to ensure the positivity of quadrature approximations.

\subsection{Quadrature by interpolation}
\label{sec:iq}

Our first class of quadrature approximations for \( \Psi \) is the degree \( s \) quadrature by interpolation \( \qq[i]{\Psi}{s} \) (i.e. $\circ = \textrm{i}$) which is defined by the relation
\begin{equation}
    \label{eqn:interp_f}
    \int f \,\d\qq[i]{\Psi}{s}
    := \int \ff[i]{f}{s} \,\d\Psi,
\end{equation}
where \( \ff[i]{f}{s} \) is the degree \( s \) polynomial interpolating a function \( f \) at the zeros \( \{\theta_{j}^{(s+1)}\}_{j=0}^{s} \) of \( p_{s+1} \), the degree $s+1$ orthogonal polynomial with respect to \( \mu \).
\Cref{eqn:interp_f} implies that
\begin{equation*}
    \qq[i]{\Psi}{s} = \smop{\sum_{j=0}^{s}} \omega_{j}^{} \bOne [\theta_{j}^{(s+1)} \leq x ]
\end{equation*}
where the weights \( \{ \omega_{j}^{} \}_{j=0}^{s} \) are chosen such that the moments of \( \qq[i]{\Psi}{s} \) agree with those of \( \Psi \) through degree \( s \).


One approach to doing this is by solving the Vandermonde-like linear system of equations%
\begin{equation}
    \label{eqn:vandermonde}
    \begin{bmatrix}
        \vphantom{\Big(} 
        p_0(\theta_{0}^{(s+1)}) & \cdots & p_0 (\theta_{s}^{(s+1)}) \\
        \vphantom{\Big(} 
        \vdots & & \vdots \\
        \vphantom{\Big(} 
        p_{s}(\theta_{0}^{(s+1)}) & \cdots & p_{s}(\theta_{s}^{(s+1)}) \\
    \end{bmatrix}
    \begin{bmatrix}
        \vphantom{\Big(}\omega_{0}^{} \\ 
        \vphantom{\Big(} 
        \vdots \\ \vphantom{\Big(} \omega_{s}^{}
    \end{bmatrix}
    =
    \begin{bmatrix}
        \vphantom{\Big(} \int p_0 \,\d\Psi \\ 
        \vphantom{\Big(} 
        \vdots \\ \vphantom{\Big(} \int p_{s} \,\d\Psi
    \end{bmatrix}.
\end{equation}
which we denote by \( \vec{P} \vec{\omega} = \vec{m} \). 
This will ensure that polynomials of degree up to \( s \) are integrated exactly. 

While it is not necessary to restrict the test polynomials to be the orthogonal polynomials \( \{ p_i \}_{i=0}^{\infty} \) with respect to \( \mu \) nor the interpolation nodes to be the zeros \( \{ \theta_j^{(s+1)} \}_{j=0}^{s} \) of \( p_{s+1} \), doing so has several advantages.
If arbitrary polynomials are used, the matrix \( \vec{P} \) may be exponentially ill-conditioned; i.e. the condition number of the matrix could grow exponentially in \( s \).
This can cause numerical issues with solving \( \vec{P} \vec{\omega} = \vec{m} \).
If orthogonal polynomials are used, then as in \cref{eqn:jacobi_evec} we see that the columns of \( \vec{P} \) are eigenvectors of the Jacobi matrix \( \vec{M} \).
Since \( \vec{M} \) is symmetric, this implies that \( \vec{P} \) has orthogonal columns; i.e. \( \vec{P}^\cT\vec{P} \) is diagonal.
%This means that the columns of the matrix will not be near linearly independent. 
Therefore, we can easily apply \( \vec{P}^{-1} \) through a product with \( \vec{P}^\cT \) and an appropriately chosen diagonal matrix.
In particular, if \( \vec{S} \) is the orthonormal matrix of eigenvectors, then \( \vec{P} = \vec{S} \operatorname{diag}([\vec{S}]_{0,:})^{-1} \) so that \( \vec{P}^{-1} = \operatorname{diag}([\vec{S}]_{0,:})\vec{S}^\cT \).
This yields \cref{alg:iq}.

\begin{labelalgorithm}[H]{iq}{get-IQ}{Quadrature by interpolation}
\begin{algorithmic}[1]
    \Procedure{\thealgorithmname}{$\{m_i\}_{i=0}^{s},\mu$}
    \State \( \vec{\theta},\vec{S} = \Call{eig}{[\vec{M}(\mu)]_{:s+1,:s+1}} \) \Comment{Eigenvectors normalized to unit length}
    \State \( \vec{\omega} = \operatorname{diag}([\vec{S}]_{0,:})\vec{S}^\cT \vec{m} \) 
%    \Comment{Eigenvectors normalized so the first row is 1}
%    \State \( \bm{\omega} = \Call{solve}{\vec{P},\vec{m}} \) 
    \State \Return \( \qq[i]{\Psi}{s} = \sum_{j=0}^{s} [\vec{\omega}]_{j} \bOne[ [\vec{\theta}]_{j} \leq x] \)
\EndProcedure
\end{algorithmic}
\end{labelalgorithm}

\begin{remark}   
In certain cases, such as \( \mu = \mu_{a,b}^T \), \( \vec{P}^{-1} \) can be applied quickly and stably using fast transforms, such as the discrete cosine transform, without ever constructing \( \vec{P} \).
\end{remark}



\subsection{Gaussian quadrature}
\label{sec:gq}

While interpolation-based quadrature rules supported on \( k \) nodes do not, in general, integrate polynomials of degree higher than \( k-1 \) exactly, if we allow the nodes to be chosen adaptively we can do better.
The degree \( 2k-1 \) Gaussian quadrature rule \( \qq[g]{\Psi}{2k-1} \) for \( \Psi \) is obtained by constructing an quadrature by interpolation rule at the roots \( \{ \theta_{i}^{(k)} \}_{i=1}^{k} \) of the degree \( k \) orthogonal polynomial \( p_{k} \) of \( \Psi \) (i.e. by taking \( \mu = \Psi \)).
\begin{theorem}
If \( p \) is any polynomial of degree at most \( 2k-1 \), then 
\begin{equation*}
    \int p\,\d\Psi = \int p \,\d\qq[g]{\Psi}{2k-1}.
\end{equation*}
I.e., the Gaussian quadrature rule integrates polynomials of degree \( 2k-1 \) exactly.
\end{theorem}

\begin{proof}
We can decompose \( p \) as
\begin{equation*}
    p = q  p_k + r
\end{equation*}
where \( q \) and \( r \) are each polynomials of degree at most \( k-1 \).
Since \( p_k \) is the \( k \)-th orthogonal polynomial with respect to \( \mu = \Psi \), it is orthogonal to all polynomials of lower degree, including \( q \).
Thus,
\begin{equation*}
    \int p \,\d\Psi
    = \int q p_k \,\d\Psi + \int r\,\d \Psi
    = \int r \,\d\Psi.
\end{equation*}
    On the other hand, since the interpolation nodes \( \{\theta_j^{(k)}\}_{j=0}^{k-1} \) are the roots of  \( p_k \),
\begin{equation*}
    \int p \,\d\qq[i]{\Psi}{k-1}
    = \smop{\sum_{j=0}^{k-1}} \omega_{j} \big( q(\theta_{j}^{(k)} ) p_k(\theta_{j}^{(k)} ) + r(\theta_{j}^{(k)} )\big)
    = \smop{\sum_{j=0}^{k-1}} \omega_{j} r(\theta_{j}^{(k)} )
    = \int r \,\d\qq[i]{\Psi}{k-1}.
\end{equation*}
Since the quadrature rule \( \qq[i]{\Psi}{k} \) is interpolatory of degree \( k \), this implies
\begin{equation*}
    \int p \,\d\Psi
    = \int p \,\d\qq[i]{\Psi}{k-1}.
    \qedhere
\end{equation*}
\end{proof}

Because the polynomials \( \{ p_i \}_{i=0}^{\infty} \) are orthogonal with respect to the probability distribution \( \Psi \) function,
we have that, for all \( i\geq 0 \),
\begin{equation*}
    m_i = \vec{v}^\cT p_i(\vec{A})\vec{v} 
    =
    \int p_i p_0 \,\d\Psi(\vec{A},\vec{v})
    = \bOne[i=0].
    %\begin{cases}
    %    1 & i=0 \\
    %    0 & i>0
    %\end{cases}.
\end{equation*}
This means the right hand side \( \vec{m} \) of \cref{eqn:vandermonde} is the first canonical unit vector \( \vec{e}_0 = [1,0,\ldots, 0]^\cT \).
Thus, as in \cref{alg:iq}, \( \vec{\omega} = \operatorname{diag}([\vec{S}]_{0,:}) [\vec{S}]_{0,:}\); that is, the quadrature weights are the squares of the first components of the unit length eigenvectors of \( [\vec{T}]_{:k,:k} \).
We then arrive at \cref{alg:gq} for obtaining a Gaussian quadrature rule for \( \Psi(\vec{A},\vec{v}) \) from the tridiagonal matrix \( [\vec{T}]_{:k,:k} \) generated by \cref{alg:lanczos}.

\begin{labelalgorithm}[H]{gq}{get-GQ}{Gaussian quadrature}
\begin{algorithmic}[1]
    \Procedure{\thealgorithmname}{$[\vec{T}]_{:k,:k}$}
%    \State \( [\vec{T}]_{:s,:s} = \textsc{Lanczos}(\vec{A},\vec{v},s) \)
    \State \( \vec{\theta}, \vec{S} = \textsc{eig}([\vec{T}]_{:k,:k}) \) \Comment{Eigenvectors normalized to unit length}
    \State \( \vec{\omega} = \operatorname{diag}([\vec{S}]_{0,:}) [\vec{S}]_{0,:} \)
    \State \Return \( \qq[g]{\Psi}{2k-1} = \sum_{j=0}^{k-1} [\vec{\omega}]_j \bOne[[\vec{\theta}]_j \leq x] \)
\EndProcedure
\end{algorithmic}
\end{labelalgorithm}

\begin{remark}
    To construct a Gaussian quadrature rule, the three term recurrence for the orthogonal polynomials of \( \Psi \) must be determined. 
    Thus, the main computational cost is computing the tridiagonal matrix giving this recurrence. 
    However, due to orthogonality, we know that all but the degree zero modified moments are zero and do not need to compute the moments.
    This is in contrast to other schemes where the polynomial recurrence is known but the modified moments must be computed.
\end{remark}



\subsection{Quadrature by approximation}
\label{sec:aq}

Rather than defining a quadrature approximation using an interpolating polynomial, we might instead define an approximation \( \qq[a]{\Psi}{s} \) by the relation
\begin{equation*}
    %\label{eqn:approx_f}
    \int f\d\qq[a]{\Psi}{s}
    := \int \ff[a]{f}{s} \,\d\Psi
\end{equation*}
where \( \ff[a]{f}{s} \) is the projection of \( f \) onto the orthogonal polynomials with respect to \( \mu \) through degree \( s \) in the inner product \( \langle \cdot, \cdot\rangle_\mu\: \).
That is,
%Explicitly, \( \ff[a]{f}{s} \) is given by 
\begin{equation*}
    \ff[a]{f}{s}
    := \:\smop{\sum_{i=0}^{s}} \langle f, p_i \rangle_\mu\: p_i
    = \:\smop{\sum_{i=0}^{s}} \left(\int f p_i\d\mu \right) p_i.
\end{equation*}

Expanding the integral of \(  \ff[a]{f}{s} \) against \( \Psi \), 
\begin{equation*}
    \int \ff[a]{f}{s} \,\d\Psi
%    &= \int \sum_{i=0}^{s} \langle f, p_i \rangle_\mu\: p_i \,\d\Psi
%    = \sum_{i=0}^{s} \langle f, p_i \rangle_\mu\: \int p_i \,\d\Psi
    = \int \smop{\sum_{i=0}^{s}} \left(\int f p_i \,\d\mu\right) p_i \,\d\Psi
    = \int f  \bigg( \smop[r]{\sum_{i=0}^{s}} \left(\int p_i \,\d\Psi\right) p_i \bigg) \,\d\mu.
\end{equation*}
This implies 
\begin{equation*}
%    \label{eqn:kpm_expanded}
    \frac{\d\qq[a]{\Psi}{s}}{\d\mu} 
    = \:\smop{\sum_{i=0}^{s}} \left( \int p_i \,\d\Psi \right) p_i 
    = \:\smop{\sum_{i=0}^{s}} m_i p_i 
\end{equation*}
where \( \d\qq[a]{\Psi}{s}/\d\mu \) is the Radon--Nikodym derivative of \( \qq[a]{\Psi}{s} \) with respect to \( \mu \).


Supposing%
\footnote{If \( \Psi = \Psi(\vec{A},\vec{v}) \) then \( \Psi \) is not absolutely continuous with respect to the Lebesgue measure (or any equivalent measure) so the Radon--Nikodym derivative does not exist. 
However, there are absolutely continuous distribution distributions with the same modified moments as \( \Psi \) up to arbitrary degree, so conceptually one can use such a distribution instead.}
that the Radon--Nikodym derivative \( \d\Psi/\d\mu \) exists, we observe 
\begin{equation*}
    m_i = \int p_i \,\d\Psi
    = \int p_i \frac{\d\Psi}{\d\mu} \,\d\mu
\end{equation*}
is the \( \mu \)-projection of \( \d\Psi/\d\mu \) onto \( p_i \) for \( i=0,1,\ldots, s \). 
Thus \( \d\Psi/\d\mu \) is approximated in a truncated orthogonal polynomial series as
\begin{equation*}
    \smop{\sum_{i=0}^{s}} \left( \int p_i \frac{\d\Psi}{\d\mu} \,\d\mu \right) \: p_i
    =
    \smop{\sum_{i=0}^{s}} \left( \int p_i \,\d\Psi \right) \: p_i.
\end{equation*}
This means the density \( \d\qq[a]{\Psi}{s}/\d{x} \) is, at least formally, the polynomial approximation to the Radon--Nikodym derivative \( \d\Psi/\d\mu \) times the density \( \d\mu/\d{x} \); i.e. \( \d\qq[a]{\Psi}{s}/\d{x} = \ff[a]{\d\Psi/\d\mu}{s} \).

We can obtain an approximation to the density \( \d\Psi/\d{x} \) by using the density \( \d\mu/\d{x} \) and the definition of the Radon--Nikodym derivative: 
\begin{equation*}
    \frac{\d\qq[a]{\Psi}{s}}{\d{x}} 
    =  \frac{\d\qq[a]{\Psi}{s}}{\d\mu}  \frac{\d\mu}{\d{x}}
    = \frac{\d\mu}{\d{x}} \smop{\sum_{i=0}^{s}} m_i p_i.
\end{equation*}
Converting this ``density'' to a distribution gives the approximation \( \qq[a]{\Psi}{s} \) shown in \cref{alg:aq}. 
%is the well sometimes called the KPM approximation to the empirical spectral measure\footnote{The empirical spectral measure is commonly called ``density of states'' in the physical sciences.} \( \d\Phi/\d{x} \).

\begin{labelalgorithm}[H]{aq}{get-AQ}{Quadrature by approximation}
\begin{algorithmic}[1]
    \Procedure{\thealgorithmname}{$\{ m_i \}_{i=0}^{s}, \mu$}%[\vec{M}]_{:s+1,:s+1}$}
    \State \Return \( \qq[a]{\Psi}{s} = \left( x \mapsto \sum_{i=0}^{s} m_i \int_{-\infty}^{x} p_i \,\d\mu \right) \)
\EndProcedure
\end{algorithmic}
\end{labelalgorithm}

\begin{remark}
    When \( \mu = \mu_{a,b}^T \), \( {\d\qq[a]{\Psi}{s}} / {\d\mu} \) can be evaluated quickly at Chebyshev nodes by means of the discrete cosine transform.
    This allows the density \( {\d\qq[a]{\Psi}{s}} / {\d{x}} \) to be evaluated quickly at these points.
\end{remark}


\subsubsection{Evaluating spectral sums and the relation to quadrature by interpolation}

We have written the output of \cref{alg:aq} as a distribution function for consistency.
However, note that 
\begin{equation*}
    \int f \,\d \qq[a]{\Psi}{s}
    = \smop{\sum_{i=0}^{s}} m_i \int f  p_i \,\d\mu.
\end{equation*}
Thus, if used for the task of spectral sum approximation, the distribution function \( \qq[a]{\Psi}{s} \) need not be computed. 
Rather, the \( \mu \)-projections of \( f \) onto the orthogonal polynomials with respect to \( \mu \) can be used instead.
In many cases, the values of these projections are known analytically, and even if they are unknown, computing them is a scalar problem independent of the matrix size \( n \).


A natural approach to computing the \( \mu \)-projections of \( f \) numerically is to use a quadrature approximation for \( \mu \).
Specifically, we might use the \( d \)-point Gaussian quadrature rule \( \qq[g]{\mu}{2d-1} \) for \( \mu \) to approximate \( \int f p_i \,\d\mu \).
This gives us the approximation 
\begin{equation*}
    \smop{\sum_{i=0}^{s}} m_i \int f p_i \,\d\mu
    \approx \smop{\sum_{i=0}^{s}}  m_i \int f p_i \,\d \qq[g]{\mu}{2s+1}
    = \smop{\sum_{i=0}^{s}}  m_i \smop{\sum_{j=0}^{d-1}} \omega_{i}^{(d)} f(\theta_{j}^{(d)}) p_i(\theta_j^{(d)})
\end{equation*}
where \( \{ \omega_i^{(d)} \}_{i=0}^{d-1} \) are the Gaussian quadrature weights for \( \mu \).

Similar to above, denote by \( \vec{P} \) the \( d\times d \) Vandermonde-like matrix of orthogonal polynomials with respect to \( \mu \) evaluated at the zeros of \( p_{d+1} \).
If \( \vec{S} \) is the orthonormal matrix of eigenvectors of the \( d\times d \) Jacobi matrix \( [\vec{M}]_{:d,:d} \) for \( \mu \), then recall that the Gaussian quadrature weights \( \vec{\omega}^{(d)} \) are given by \( \operatorname{diag}([\vec{S}]_{0,:s+1})([\vec{S}]_{:s+1,:})^\cT \).
This yields \cref{alg:aaq}.

\begin{remark}
In the case \( d=s+1 \), \cref{alg:aaq} is equivalent to \cref{alg:iq}.
\end{remark}

\begin{labelalgorithm}[H]{aaq}{get-aAQ}{Approximate quadrature by approximation}
\begin{algorithmic}[1]
    \Procedure{\thealgorithmname}{$\{m_i\}_{i=0}^{s},d,\mu$}%[\vec{M}]_{:d,:d}$}
    \State \( \vec{\theta},\vec{S} = \Call{eig}{[\vec{M}(\mu)]_{:d,:d}} \) \Comment{Eigenvectors normalized to unit length}
    \State \( \vec{\omega} = \operatorname{diag}([\vec{S}]_{0,:s+1})([\vec{S}]_{:s+1,:})^\cT \vec{m} \) 
%    \Comment{Eigenvectors normalized so the first row is 1}
%    \State \( \bm{\omega} = \Call{solve}{\vec{P},\vec{m}} \) 
    \State \Return \( \qq[i]{\Psi}{s} = \sum_{j=0}^{k-1} [\vec{\omega}]_{j} \bOne[[\vec{\theta}]_{j} \leq x]  \)
\EndProcedure
\end{algorithmic}
\end{labelalgorithm}



\subsection{Positivity by damping and the kernel polynomial method}
\label{sec:damped}

While it is clear that that the Gaussian quadrature \( \qq[g]{\Psi}{s} \) is a non-negative probability distribution function, neither \( \qq[i]{\Psi}{s} \) nor \( \qq[a]{\Psi}{s} \) are guaranteed to be weakly increasing.
We now discuss how to use damping to enforce positivity.

Towards this end, define the damping kernel,
\begin{equation*}
    P_x(y) 
    %= \sum_{i=0}^{s} \sum_{j=0}^{s} \rho_{i,j} p_i(x) p_j(y)
    = \smop{\sum_{i=0}^{s}} \rho_{i} p_i(x) p_i(y)
\end{equation*}
where \( \{ \rho_i \}_{i=0}^{s} \) are \emph{damping coefficients} as in \cref{def:damping}.
Then the damped interpolant \( \ff[d-i]{f}{s} \) and approximant \( \ff[d-a]{f}{s} \) can be written in terms of \( P_{x} \) as 
\begin{equation*}
    \ff[d-i]{f}{s}(x) = \int P_{x} f \,\d\qq[g]{\mu}{2s+1}
    \qquad \text{and} \qquad
    \ff[d-a]{f}{s}(x) = \int P_{x} f \,\d\mu.
\end{equation*}

\begin{remark}
If  \( \rho_{i} = 1 \) for all \( i \),  then \( \ff[d-i]{f}{s} = \ff[i]{f}{s} \) and \( \ff[d-a]{f}{s} = \ff[a]{f}{s} \).
\end{remark}

%We discuss the fact that interpolation at the zeros of orthogonal polynomials is equivalent to approximating the coefficients of the approximant via a Gaussian quadrature in \cref{sec:interp_gq}.

These approximations induce  \( \qq[d-i]{\Psi}{s} \) and \( \qq[d-a]{\Psi}{s} \) by
\begin{equation*}
    \int f \,\d\qq[d-i]{\Psi}{s}
    := \int \ff[d-i]{f}{s} \,\d\Psi
    \qquad\text{and}\qquad
    \int f \:  \d\qq[d-a]{\Psi}{s}
    := \int \ff[d-a]{f}{s} \,\d\Psi.
\end{equation*}
Algorithmically, this is equivalent to replacing \( m_i \) by \( \rho_{i}^{} m_i \) in the expressions for quadrature by interpolation and approximation described in \cref{sec:iq,sec:aq}.

\begin{lemma}
If \( \rho_0 = 1 \), then \( \qq[d-i]{\Psi}{s} \) and \( \qq[d-a]{\Psi}{s} \) have unit mass, and if \( P_x(y) \geq 0 \) for all \( x,y \), then \( \qq[d-i]{\Psi}{s} \) and \( \qq[d-a]{\Psi}{s} \) are weakly-increasing.
\end{lemma}

\begin{proof}
The first part of the theorem follows from the fact that
\begin{equation*}
    \int P_x \d\mu
    = \sum_{i=0}^{s} \rho_i p_i(x) \int p_i\d\mu
    = \rho_{0}.
\end{equation*}

To prove the remainder, suppose \( f \) is non-negative. 
Then clearly 
\begin{equation*}
    \ff[d-a]{f}{s}(x) = \int P_{x} f \,\d\mu
    \geq 0
\end{equation*}
so that \( \int f\,\d\qq[d-a]{\Psi}{s} \geq 0 \).
A similar argument implies \( \int f\,\d\qq[d-i]{\Psi}{s} \geq 0 \) as well.
Therefore the approximations are weakly increasing. 
\end{proof}

\begin{remark}
    While we have describe the damping procedure in terms of the damped interpolant and approximant, an equilvalent perspective is that we first smooth the distribution \( \Psi \) with the damping kernel \( P_t(y) \) and subsequently use quadrature by interpolation or approximation with this new distribution function.
\end{remark}


One particularly important damping kernel for the case \( \mu = \mu_{a,b}^T \) is given by the Jackson coefficients defined in \cref{def:jackson_coeffs}.
The associated damping kernel was used in the original proof of Jackson's theorem \cite{jackson_12} and leads to the Jackson damped KPM approximation, which is the most popular KPM variant \cite{weisse_wellein_alvermann_fehske_06,braverman_krishnan_musco_22}.
The rate of convergence of polynomial approximations using these damping coefficients is estimated in \cref{thm:damped_cheb} below.
For a discussion on other damping schemes we refer readers to \cite{weisse_wellein_alvermann_fehske_06,lin_saad_yang_16}.


\section{A priori error bounds on an interval}
\label{sec:quadrature_apriori}
%\subsection{Simple bounds from best approximation}

Most of the quadrature approximations we consider have the property that they integrate polynomials exactly. 
In this case, we can use this property to reduce error bounds to the quality of the best uniform polynomial approximations to \( f \), which we discussed in \cref{sec:uniform_poly_bounds}.


\begin{lemma}
    \label[lemma]{thm:moments_inf_norm}
    Suppose \( \Upsilon_1 \) and \( \Upsilon_2 \) are probability distribution functions whose moments are equal through degree \( s \), each constant on \( (- \infty, a) \) and \( (b,\infty) \).
    Then,
    \begin{equation*}
        \left| \int f \,\d \Upsilon_2 - \int f \,\d\Upsilon_2 \right|
         = (\TV(\Upsilon_1)+\TV(\Upsilon_2)) \min_{\deg(p)\leq s} \| f - p \|_{[a,b]}.
    \end{equation*} 
\end{lemma}

\begin{proof}%[Proof of \cref{thm:moments_inf_norm}]
    Let \( p \) be any polynomial of degree at most \( s \) and note that \( \int p \,\d\Upsilon_1 = \int p \,\d\Upsilon_2 \).
    Then, applying the triangle inequality,
    \begin{align*}
        \left| \int f \,\d \Upsilon_1 - \int f \,\d\Upsilon_2 \right|
        & = \left| \int (f - p) \,\d \Upsilon_1 - \int (f-p) \,\d\Upsilon_2 \right|
        \\& \leq \int |f - p| |\d \Upsilon_1| + \int |f - p| |\d \Upsilon_2| 
        \\& \leq \int \|f - p\|_{[a,b]} |\d \Upsilon_1| + \int \|f - p\|_{[a,b]} |\d \Upsilon_2|
        \\& = (\TV(\Upsilon_1)+\TV(\Upsilon_2)) \| f - p \|_{[a,b]}.
    \end{align*}    
    The result follows by optimizing over \( p \).
\end{proof}

\Cref{thm:moments_inf_norm} shows that the Lanczos-based Gaussian quadrature approximations always perform within a factor of two of the best polynomial approximation on \( \mathcal{I} \).
Intuitively, approaches based on explicit polynomial approximation will have error roughly equal to \( \|f-\ff{f}{s}\|_{\mathcal{I}} \) as a large portion of mass of \( \Psi \) is likely in regions where \( | f-\ff{f}{s} | \) is large.
Thus, Gaussian quadrature should not be expected to perform significantly worse than explicit polynomial methods, at least in exact arithmetic.
In fact, as we will discuss in \cref{chap:finite_precision}, even in finite preicsion arithmetic, \cref{thm:moments_inf_norm} is still morally correct.

%\subsubsection{Wasserstein approximation for damped approximant}

For some quadrature approximations we consider, polynomials are not integrated exactly.
In such cases, we turn to the following bound:
\begin{lemma}
    \label[lemma]{thm:op_inf_norm}
    Suppose \( \Upsilon_1 \) is a probability distribution function and \( \Upsilon_2 \) is defined by \( \int f \,\d\Upsilon_2 = \int \mathcal{O}[f] \d\Upsilon_1 \) for some operator \( \mathcal{O}[\:\cdot\:] \).
    Then, for any \( f \),
    \begin{equation*}
    \left| \int f \,\d\Upsilon_1
    - \int f \,\d\Upsilon_2 \right|
        = \| f - \mathcal{O}[f] \|_{[a,b]} \TV(\Upsilon_1)
    \end{equation*}
\end{lemma}

\begin{proof}%[Proof of \cref{thm:opt_inf_norm}]
The result follows by a simple application of the triangle inequality:
\begin{align*}
    \left| \int f \,\d\Upsilon_1
    - \int f \,\d\Upsilon_2 \right|
    = \left| \int ( f - \mathcal{O}[f]) \,\d\Upsilon_1 \right|
    &\leq \int \| f - \mathcal{O}[f] \|_{[a,b]} |\d\Upsilon_1|
    \\&= \| f - \mathcal{O}[f] \|_{[a,b]} \TV(\Upsilon_1).
    \tag*{\qedhere}
\end{align*}
\end{proof}

The primary downside of this bound is that it requires being able to bound \( \| f - \mathcal{O}[f] \|_{[a,b]} \).
%In some sense, this bound is less general than \cref{thm:moments_inf_norm} because it depends on the specific properties of \( \mathcal{O} \).
However, at least in the case that \( \mathcal{O} \) corresponds to undamped or Jackson damped Chebyshev interpolation or approximation, we are able to derive bounds for \( \| f - \mathcal{O}[f] \|_{[a,b]} \) directly.



\section{Qualitative comparison of algorithms}
\label{sec:QF_tradeoffs}


In \cref{sec:krylov_moments,sec:lanczos} we described \cref{alg:moments} (\ref{alg:name:moments}), \cref{alg:cheb_moments} (\ref{alg:name:cheb_moments}), \cref{alg:lanczos} (\ref{alg:name:lanczos}), \cref{alg:moments_cheb} (\ref{alg:name:moments_cheb}), and \cref{alg:moments_lanc} (\ref{alg:name:moments_lanc}) which are used to compute the information required for the quadrature approximations described in \cref{sec:quadrature}
Since \cref{alg:moments_cheb,alg:moments_lanc} respectively call \cref{alg:cheb_moments,alg:lanczos}, \cref{alg:moments,alg:cheb_moments,alg:lanczos} constitute the bulk of the computational cost of all implementations of the protoalgorithm discussed in this paper.

In each iteration, \cref{alg:moments,alg:cheb_moments,alg:lanczos} each require one matrix vector product with \( \vec{A} \) along with several scalar multiplications, vector additions, and inner products.
As such, the total computational cost of each algorithms is \( O(k(T_{\textup{mv}} + n)) \) where \( k \) is the number of iterations and \( T_{\textup{mv}} \) is the cost of a matrix-vector product with \( \vec{A} \).
Here we ignore terms depending only on \( k \) (e.g. \( k^2 \)) which are unimportant if we assume \( k\ll n \).
%\mnote{is this a safe assumption? Stefan Guettel mentioned it might not be the case in an email about a different thing}
Each of the algorithms can also be implemented using just \( O(n) \) storage; i.e. without storing the entire basis for the Krylov subspace which would cost \( O(kn) \) storage.


While the algorithms are typically quite storage efficient, there are some situations in which it may be desirable to store the whole Krylov basis.
First, \cref{alg:lanczos} is sometimes run with full reorthogonalization.
This can improve numerical stability, but increases the computation cost to \( O(k(T_{\textup{mv}} + kn)) \).
Next, by delaying all inner products to a final iteration (or using a non-blocking implementation), the number of global reductions required by \cref{alg:moments} and \cref{alg:cheb_moments} can be reduced.
Since global communication can significantly slow down Krylov subspace methods, this may speed up computation on highly parallel machines \cite{dongarra_heroux_luszczek_15,exascale}.
As mentioned earlier, there are implementations of the Lanczos algorithm which aim to decrease the number of global communications \cite{carson_demmel_15,carson_20}.
Designing Krylov subspace methods for avoiding or reducing communication costs is a large area study, but further discussion is outside the scope of this paper. 



\section{Numerical experiments}

\label{sec:numerical_experiments}

In this section, we provide a range of numerical experiments to illustrate the behavior of the algorithms described above as well as the tradeoffs between algorithms.
Our focus is primarily on quadrature approximations of the weighted CESM, as the approximation of the true CESM by the average of weighted CESMs is straightforward and well understood.

%None of the examples are meant to push the computational limits of the prototypical algorithm.

%There are many more applications than space in this paper, 


\subsection{Comparison with classical quadrature}

We begin with an example designed to illustrate some of the similarities and differences between the behavior of classical quadrature rules for continuous weight functions and the behavior of the matrix-free quadrature algorithms presented in this paper. 

\begin{figure}[htb]
    \begin{center}
        \includegraphics{imgs/ch3_runge_uniform_spec.pdf} 
    \end{center}
    \caption[{Errors for approximating \( \int f \,\d\Psi =  \vec{v}^\cT f(\vec{A}) \vec{v} \) when \( f(x) = 1/(1+16x^2) \) for a spectrum uniformly filling \( [-1,1] \).}]{%
    Errors for approximating \( \int f \,\d\Psi =  \vec{v}^\cT f(\vec{A}) \vec{v} \) when \( f(x) = 1/(1+16x^2) \) for a spectrum uniformly filling       \( [-1,1] \).
    \hspace{.25em}\emph{Legend}:
    Gaussian quadrature with 
    ({\protect\raisebox{0mm}{\protect\includegraphics[]{imgs/legend/l1.pdf}}})
    and without 
    ({\protect\raisebox{0mm}{\protect\includegraphics[]{imgs/legend/l2.pdf}}})
    reorthogonalization,
    quadrature by interpolation
    ({\protect\raisebox{0mm}{\protect\includegraphics[]{imgs/legend/l4.pdf}}}),
    and approximate quadrature by approximation
    ({\protect\raisebox{0mm}{\protect\includegraphics[]{imgs/legend/l3.pdf}}}).
    \hspace{.25em}\emph{Takeaway}: Intuition about classical approximation theory informs our understanding of algorithms for matrix-free quadrature.
    In fact, in some cases, quadrature by interpolation or approximation can provably outperform Gaussian quadrature.
    }
    \label{fig:ch3_runge_uniform_spec}
\end{figure}

\begin{figure}[htb]
    \begin{center}
        \includegraphics{imgs/ch3_runge_gap_spec.pdf} 
    \end{center}
    \caption[{Errors for approximating \( \int f \,\d\Psi =  \vec{v}^\cT f(\vec{A}) \vec{v} \) when \( f(x) = 1/(1+16x^2) \) for a spectrum uniformly filling       \( [-1,1] \) except for a gap around zero.}]{%
    Errors for approximating \( \int f \,\d\Psi =  \vec{v}^\cT f(\vec{A}) \vec{v} \) when \( f(x) = 1/(1+16x^2) \) for a spectrum uniformly filling       \( [-1,1] \) except for a gap around zero.
    \hspace{.25em}\emph{Legend}:
    Gaussian quadrature with 
    ({\protect\raisebox{0mm}{\protect\includegraphics[]{imgs/legend/l1.pdf}}})
    and without 
    ({\protect\raisebox{0mm}{\protect\includegraphics[]{imgs/legend/l2.pdf}}})
    reorthogonalization,
    quadrature by interpolation
    ({\protect\raisebox{0mm}{\protect\includegraphics[]{imgs/legend/l4.pdf}}}),
    and approximate quadrature by approximation
    ({\protect\raisebox{0mm}{\protect\includegraphics[]{imgs/legend/l3.pdf}}}).
    \hspace{.25em}\emph{Takeaway}:  
    The behavior of the algorithms are highly dependent on the eigenvalue distribution of \( \vec{A} \), and Gaussian quadrature may perform significantly better that explicit methods when the spectrum of \( \vec{A} \) has additional structure such as gaps. 
    }
    \label{fig:ch3_runge_gap_spec}
\end{figure}



Throughout this example, we use the Runge function \( f(x) = 1/(1+16x^2) \) and a vector \( \vec{v} \) with uniform weight on each eigencomponent.
%Thus, the spectral sum \( \vec{v}^\cT f[\vec{A}] \vec{v} \) can be viewed as a midpoint approximation to integrals of \( f \) over \( [-1,1] \) and \( [-1,-0.75]\cup[0.75,1] \).
We will compare the effectiveness of the Gaussian quadrature rule \( \qq[g]{\Psi}{2k-1} \), the quadrature by interpolation rule \( \qq[i]{\Psi}{2k} \) and the quadrature by approximation rule \( \qq[a]{\Psi}{2k} \). 
For the latter approximations, we set \( \mu = \mu_{-1,1}^T \), and for the quadrature by approximation rule, we use \cref{alg:aaq} with enough quadrature nodes so that the involved integrals are computed to essentially machine precision.
All three approximations can be computing using \( k \) matrix-vector products with \( \vec{A} \), and since the approaches exactly integrate polynomials of degree \( 2k-1 \) and \( 2k \) respectively, we might expect that them to behave similarly.
However, there are a variety of factors which prevent this from being the case.

In our first experiment, shown in \cref{fig:ch3_runge_uniform_spec}, the spectrum of \( \vec{A} \) uniformly fills out the interval \( [-1,1] \); i.e., \( \lambda_i = -1+(2i+1)/n \), \( i=0,1,\ldots, n-1 \).
We take \( n=10^5 \) so that \( \qq[g]{\Psi}{2k-1} \) and \( \qq[i]{\Psi}{2k} \) respectively approximate the \( k \)-point Gaussian quadrature and \( (2k-1) \)-point Fej\'{e}r quadrature rules for a uniform weight on \( [-1,1] \).
For many functions, certain quadrature by interpolation rules on \( [-1,1] \), including the Fej\'{e}r rule, behave similarly to the Gaussian quadrature rule when the same number of nodes are used \cite{trefethen_08}.
For \( f(x) = 1/(1+16x^2) \), this phenomenon is observed for some time until the convergence rate is abruptly cut in half \cite{weideman_trefethen_07}.
In our setting, a fair comparison means that the number of matrix-vector products are equal, so we see that the quadrature by interpolation approximation initially converges twice as quickly as the Gaussian quadrature approximation!
The rate of the quadrature by interpolation approximation is eventually cut in half to match the rate of the Gaussian quadrature approximation. 

\begin{figure}[htb]
    \begin{center}
        \includegraphics{imgs/ch3_model_problem.pdf} 
    \end{center}
    \caption[{Errors for approximating \( \int f \,\d\Psi =  \vec{v}^\cT f(\vec{A}) \vec{v} \) when \( f(x) = 1/x \) for model problem.}]{%
    Errors for approximating \( \int f \,\d\Psi =  \vec{v}^\cT f(\vec{A}) \vec{v} \) when \( f(x) = 1/x \) for model problem.
    \hspace{.25em}\emph{Legend}:
    Gaussian quadrature with 
    ({\protect\raisebox{0mm}{\protect\includegraphics[]{imgs/legend/l1.pdf}}})
    and without 
    ({\protect\raisebox{0mm}{\protect\includegraphics[]{imgs/legend/l2.pdf}}})
    reorthogonalization,
    quadrature by interpolation
    ({\protect\raisebox{0mm}{\protect\includegraphics[]{imgs/legend/l4.pdf}}}),
    and approximate quadrature by approximation
    ({\protect\raisebox{0mm}{\protect\includegraphics[]{imgs/legend/l3.pdf}}}).
    \hspace{.25em}\emph{Takeaway}: Without reorthogonaliztion the convergence of Gaussian quadrature is slowed. 
    However, the method still converges and can even outperform the other methods.
    }
    \label{fig:ch3_model_problem}
\end{figure}


In our second experiment, shown in \cref{fig:ch3_runge_gap_spec}, the spectrum of \( \vec{A} \) uniformly fills out the disjoint intervals \( [-1,-0.75] \cup [0.75,1] \) with the same inter-point spacing as the first example; i.e. we remove the eigenvalues in the previous example which fall between \( -0.75 \) and \( 0.75 \).
Here we observe that the Gaussian quadrature rule converges significantly faster than in the previous experiment.
This to be expected. 
Indeed, the Gaussian quadrature rule has its nodes near \( [-1,-0.75]\cup[0.75,1] \), so the union of the support of \( \Psi \) and \( \qq[i]{\Psi}{2k} \) is further from the poles of \( f \) located at \( \pm \ii /4 \).
We also note that the conditions which enabled accelerated convergence in the first experiment are no longer present, so the quadrature by interpolation approximation converges at its limiting rate \cite{trefethen_08}. 

%We conclude with the remark that \( \vec{v}^\cT f[\vec{A}] \vec{v} \) is a trapezoid rule approximation to \( \int_{-1}^{1} f(x) \,\d x \).

In the both experiments, the Lanczos based Gaussian quadrature approach behaves similar with and without reorthogonalization.
In fact, it is easily verified that the Lanczos algorithm does not lose orthogonality and behaves nearly the same regardless of whether or not reorthogonalization is used.
To the best of our knowledge, such a result has not been proved rigorously.
%Proving this phenomenon rigourously would be an interesting result.


\begin{figure}[htb]
    \begin{center}
        \includegraphics{imgs/ch3_MNIST.pdf} 
    \end{center}
    \caption[{Errors for approximating \( \int f \,\d\Psi =  \vec{v}^\cT f(\vec{A}) \vec{v} \) when \( f(x) = \bOne[x > c] \) for MNIST covariance matrix.}]{%
    Errors for approximating \( \int f \,\d\Psi =  \vec{v}^\cT f(\vec{A}) \vec{v} \) when \( f(x) = \bOne[x > c] \) for MNIST covariance matrix.
    \hspace{.25em}\emph{Legend}:
    Gaussian quadrature with 
    ({\protect\raisebox{0mm}{\protect\includegraphics[]{imgs/legend/l1.pdf}}})
    and without 
    ({\protect\raisebox{0mm}{\protect\includegraphics[]{imgs/legend/l2.pdf}}})
    reorthogonalization,
    quadrature by interpolation
    ({\protect\raisebox{0mm}{\protect\includegraphics[]{imgs/legend/l4.pdf}}}),
    and approximate quadrature by approximation
    ({\protect\raisebox{0mm}{\protect\includegraphics[]{imgs/legend/l3.pdf}}}).
    \hspace{.25em}\emph{Takeaway}: Without reorthogonaliztion the convergence of Gaussian quadrature is slowed. 
    However, the method still converges and can even outperform the other methods.
    }
    \label{fig:ch3_California}
\end{figure}



\subsection{Finite precision convergence}

In this example, we consider several experiments where orthogonality is lost and the effects of finite precision arithmetic are easily observed.
In both experiments we use diagonal matrices scaled so \( \| \vec{A} \|_2 = 1 \) and set \( \vec{v} \) to have uniform entries.
We set \( a,b \) as the largest and smallest eigenvalues respectively and again use \( \mu = \mu_{a,b}^T \) for the interpolatory and quadrature by approximations.


In the first experiment, shown in \cref{fig:ch3_model_problem}, the eigenvalues of \( \vec{A} \) are distributed according to the model problem \cref{eqn:model_problem} and \( f(x) = 1/x \).
Specifically, the eigenvalues are given by the model problem with selected parameters \( n = 300 \), \( \kappa = 10^3 \), and \( \rho = 0.85 \).
In the second experiment, shown in the right panel \cref{fig:ch3_California}, we use the \( n=9664 \) eigenvalues of the California matrix from the sparse matrix suite \cite{davis_hu_11} and the function \( f(x) = |x| \).


In both cases, the Jacobi matrices produced by Lanczos, with or without reorthogonalization, differ greatly; i.e. the difference of the matrices is on the order of \( \| \vec{A} \|_2 \). 
Even so, the modified moments for \( \mu = \mu_{a,b}^T \) obtained by \cref{alg:moments_cheb,alg:moments_lanc} differ only in the 12th digit and 14th digits respectively.
Using one approach in place of the other does not noticeably impact the convergence of the quadrature by interpolation approximations.




 

\chapter{Spectrum and spectral sum\\approximation}
\label{chap:random_quadrature}


We now turn to the tasks of spectrum and spectral sum approximation.
Specifically, we will show how the algorithms from the previous chapter can be used to produce approximations to the CESM \( \Phi \), the probability distribution function with unit mass at each eigenvalue of \( \vec{A} \) which we defined in \cref{def:CESM}.
This in turn induces approximations to \( \tr(\fA) \).

Towards this end, suppose \( \vec{v} \) is a random vector satisfying \( \EE[\vec{v}\vec{v}^\cT] = n^{-1} \vec{I} \); i.e., \( \vec{v} \) is isotropic with appropriate scale.
Then, using basic properties of the trace and expectation for any \( x\in\R \), we have 
\begin{align*}
    \EE[\Psi(x)]
    &= \EE[\vec{v}^\cT \bOne[\vec{A}\leq x] \vec{v}]
    = \EE[\tr (\vec{v}^\cT \bOne[\vec{A}\leq x] \vec{v})] 
    \\& \hspace{3em} = \EE[\tr (\bOne[\vec{A}\leq x] \vec{v} \vec{v}^\cT )]
    = \tr (\EE[ \bOne[\vec{A}\leq x] \vec{v} \vec{v}^\cT ])
    \\& \hspace{6em} = \tr( \bOne[\vec{A}\leq x] \EE[ \vec{v} \vec{v}^\cT ] )
    = n^{-1} \tr(\bOne[\vec{A}\leq x])
    = \Phi(x).
\end{align*}
That \( \Psi \) is an unbiased estimator for \( \Phi \) at every point \( x\in\R \) is illustrated in \cref{fig:ch4_CESMwCESM}.
Further, almost by definition, we see that \( \int f \,\d \Psi = \vec{v}^\cT \fA \vec{v} \) is an unbiased estimator for \( n^{-1} \tr(\fA) \).

\begin{figure}[htb]
    \includegraphics[width=\textwidth]{imgs/ch4_CESMwCESM.pdf}
    \caption[{CESM \( \Phi \) and independent 10 samples of weighted CESM \( \Psi \)}]{
        CESM \( \Phi \) 
        ({\protect\raisebox{0mm}{\protect\includegraphics[]{imgs/legend/l1nm.pdf}}})
        and 10 independent samples of weighted CESM \( \Psi \) corresponding to random \( \vec{v} \) 
        ({\protect\raisebox{0mm}{\protect\includegraphics[]{imgs/legend/l1nm_alpha.pdf}}}).
        Each copy of \( \Psi \) is an unbiased estimator for \( \Phi \) at each point \( x\in\R \).
    }
    \label{fig:ch4_CESMwCESM}
\end{figure}

Let \( \{ \Psi_\ell \}_{\ell=0}^{\nv-1} \) be independent and identically distributed (iid) copies of the weighted CESM \( \Psi \) corresponding to vectors \( \{ \vec{v}_\ell \}_{\ell=0}^{\nv-1} \) which are iid copies of \( \vec{v} \).
Then the averaged weighted CESM 
\label{eqn:ave_notation}
\begin{equation*}
    \samp{ \Psi_\ell } := \nv^{-1} \smop{\sum_{\ell=0}^{\nv-1}} \Psi_\ell 
\end{equation*}
is also an unbiased estimator for the CESM at every point \( x \).
This implies 
\begin{equation*}
    \samp[\big]{ \vec{v}_{\ell}^\cT \fA \vec{v}_\ell }
    := \nv^{-1} \smop{\sum_{\ell=0}^{\nv-1}} \vec{v}_\ell^\cT \fA \vec{v}_\ell
    = \int f \,\d \samp{\Psi_\ell},
\end{equation*}
is an unbiased estimator for \( n^{-1} \tr(\fA) \).
In both cases, the standard deviation of the averaged estimator decreases proportional to \( 1/\sqrt{\nv} \), so the averaged estimators concentrate more sharply about the mean as \( \nv \) increases.


We refer to estimators of the form \( \vec{v}^\cT \vec{B} \vec{v} \), where \( \vec{v} \) is an isotropic random vector, as quadratic trace estimators.
Thus, we see that the quadratic trace estimator for the spectral sum \( \tr(\fA) \) is an integral against the weighted CESM, which is itself a quadratic trace estimator for the CESM at every point \( x \).
Moreover, since the quadratic form \( \vec{v}^\cT \fA\vec{v} \) can be written as an integral of \( f \) against the weighted CESM, classical results about the convergence of quadrature rules for approximating this integral can be leveraged to obtain error estimates for the convergence of our Krylov subspace approximations of \( \vec{v}^\cT \fA\vec{v} \).

In order to approximate a sample of \( \Psi \), and therefore integrals against such samples, we can simply use the algorithms from the previous chapter. 
Thus, we arrive at a prototypical algorithm for spectrum and spectral sum approximation,\cref{alg:protoalg}.
The output \( \samp{\qq{\Psi_\ell}{s}} \) of \cref{alg:protoalg} is a distribution function which approximates the CESM \( \Phi \).
For any function \( f:\R\to\R \), this approximation naturally yields an approximation to the spectral sum \( \tr(\fA) \) by integration.

\begin{labelalgorithm}[H]{protoalg}{spec-approx}{Prototypical randomized spectrum and spectral sum approximation}
\begin{algorithmic}[1]
    \Procedure{\thealgorithmname}{$\vec{A}, \nv, k, \circ$}
    \For{ \( \ell = 0,1,\ldots,\nv-1 \) }
    \State define (implicitly) \( \Psi_\ell \stackrel{\text{iid}}{\sim} \Psi \) by sampling \( \vec{v}_\ell \stackrel{\text{iid}}{\sim} \vec{v} \), \( \EE[\vec{v}\vec{v}^\cT] = n^{-1} \vec{I} \) %\operatorname{Unif}(\{ \pm n^{-1/2} \}^n) \) 
    \State compute moments of \( \Psi_{\ell} \) through degree \( s \) by constructing \( \mathcal{K}_{k+1}(\vec{A},\vec{v}_\ell) \) % (using \( k \) matrix-vector products)% with \( \vec{A} \).
    \State approximate \( \Psi_{\ell} \) by \( \qq{\Psi_{\ell}}{s} \) induced by a polynomial operator \( \ff{\:\cdot\:}{s} \) 
    \EndFor
    \State \Return \( \samp{ \qq{\Psi_{\ell}}{s} } := \nv^{-1} \sum_{\ell=0}^{\nv-1} \qq{\Psi_\ell}{s} \)

\EndProcedure
\end{algorithmic}
\end{labelalgorithm}


\section{Related work and context}
\label{sec:past_work}
%\cnote{Still need to make sure this is comprehensive and representative of past work. Maybe we don't need such a big related work section.. Just add one paragraph to previous section.}

Specific implementations of the prototypical algorithm, given in \cref{alg:protoalg}, are by far the most common algorithms for spectral sum and spectrum approximation, and they have found widespread use in a range of disciplines \cite{lin_saad_yang_16,ubaru_chen_saad_17}.
As we have alluded to, the two key ingredients for such algorithms are (i) polynomial approximation and quadrature and (ii) quadratic trace estimation. 
The first of these ingredients has been studied for centuries \cite{trefethen_19}, so the popularization of the latter \cite{girard_87,hutchinson_89,skilling_89} quickly lead to a variety algorithms fitting this framework.
In this section we focus primarily on conceptual and theoretical advancements relating to the protoalgorithm.
We hope our brief review of prior work will help tie together several clusters of literature which have remained largely disjoint.


Both \cite{girard_87,hutchinson_89} focus on estimating the trace of a large implicit matrix \( \vec{B} = \vec{A}^{-1} \) for some matrix \( \vec{A} \).
While \cite{girard_87} suggests the use of the conjugate gradient algorithm, neither paper discusses in detail how to approximate products with \( \vec{B} \).
Therefore, to the best of our knowledge, \cite{skilling_89} contains the first example of an algorithm which truly fits into the form considered in this paper.
In \cite{skilling_89}, an approximation to \( \Phi \) based on an expansion in Chebyshev polynomials is described.
This approximation is then used to approximate \( \tr(\ln\A) = \ln \,\det\A \).


The Chebyshev based approach of \cite{skilling_89} was improved in \cite{silver_roder_94} where a damping Kernel was introduced to avoid Gibbs oscillations.
The connection to Jackson's damping and other classical ideas from approximation theory were subsequently considered in \cite{silver_roeder_voter_kress_96}.
The resulting algorithm is now typically called the kernel polynomial method (KPM) and is widely used in the computational physical sciences; see \cite{weisse_wellein_alvermann_fehske_06} for a review.


Essentially in parallel, stochastic trace estimation was combined with quadrature explicitly.
Typically, such approaches are based on the Lanczos algorithm which can be used in a straightforward way to compute certain quadrature rules for \( \Psi \) \cite[Chapter 6]{golub_meurant_09}, \cite{gautschi_06}.
In \cite{bai_fahey_golub_96}, Gauss, Gauss-Lobatto, and Gauss-Radau quadrature rules are used to derive upper and lower bounds for \( \int f \,\d\Psi \)
when \( f(x) = 1/x \) or \( f = \ln(x) \). 
These bounds were in turn combined with stochastic trace estimation to provide probabilistic upper and lower bounds on the traces of the corresponding matrix functions. 
The Gaussian quadrature based approach is now typically referred to as stochastic Lanczos quadrature (SLQ).


Bounds on the number of samples \( \nv \) required so that the average of iid quadratic trace estimators is within \( \epsilon \) of the true trace with at least probability \( 1-\eta \) were derived in \cite{avron_toledo_11} and subsequently improved on in \cite{roostakhorasani_ascher_14}.
These bounds enabled a range of analyses which explicitly balanced the number of samples \( \nv \) with the approximation degree \( s \). 
For instance, \cite{han_malioutov_avron_shin_17} and \cite{ubaru_chen_saad_17} respectively consider approximation of spectral sums corresponding to analytic functions by a Chebyshev based approach and SLQ.
Later, \cite{cortinovis_kressner_21} gives stronger bounds for quadratic trace estimators, and as in \cite{ubaru_chen_saad_17}, these bounds are used to analyze SLQ.

Around this time, spectrum approximation in Wasserstein distance was analyzed for KPM \cite{braverman_krishnan_musco_22} and SLQ \cite{chen_trogdon_ubaru_21}.
We remark that \cite{braverman_krishnan_musco_22,chen_trogdon_ubaru_21} both arrive at the conclusion that the number of samples required to approximate \( \Phi \) in Wasserstein distance to accuracy \( \epsilon \) actually \emph{decreases} as the matrix size \( n \) increases, provided \( \epsilon \gg n^{-1/2} \) as \( n\to\infty \).
While not stated explicitly, the analysis in \cite{cortinovis_kressner_21} implies this same fact for the number of samples required to approximate \( \int f \,\d\Phi = n^{-1} \tr(\fA) \) to additive error \( \pm \epsilon \).
This fact was already known to the physics community \cite{weisse_wellein_alvermann_fehske_06}, although, to the best of our knowledge, was not proved rigorously in the literature.


\subsection{Note on history of stochastic quadratic trace estimators and their analysis}
\label{sec:trace_hist} 

What we are calling a quadratic trace estimator is often called the \emph{Hutchinson's trace estimator}, especially when \( \vec{v} \) is chosen uniformly from the set of vectors with entries \( \pm n^{-1/2} \).
However, \cite{hutchinson_89} was not the first use of quadratic trace estimators for the task of approximating the trace of an implicit matrix; \cite{hutchinson_89} itself cites \cite{girard_87} which addresses the same task by using samples of \( \vec{v} \) drawn uniformly from the unit hypersphere.
Algorithms based on the use of random vectors back at least to the mid 1970s \cite{alben_blume_krakauer_schwartz_75,weaire_williams_76,weaire_williams_77,deraedt_devries_89}.

In fact, such estimators are a special case of the concept of \emph{typicality} in quantum physics.
%Typicality asserts that properties derived from a suitably sampled random pure state (vector) tend to be near to some typical value 
Typicality has its origins in work of Schr\"odinger \cite{schrodinger_27} and von Neumann \cite{vonneumann_29} from the late 1920s but was dismissed and/or forgotten until a resurgence in the mid 2000s \cite{gemmer_michel_mahler_09,goldstein_lebowitz_tumulka_zanghi_06,popescu_short_winter_06,reimann_07};
see \cite{goldstein_lebowitz_mastrodonato_tumulka_zanghi_10} for a historical overview and discussion in a modern context and \cite{jin_willsch_willsch_lagemann_michielsen_deraedt_21} for a review of algorithms based on typicality.

%\cite{deraedt_devries_89} DOS random state



Likewise, while the first tail bounds for quadratic trace estimators are typically attributed to \cite{avron_toledo_11,roostakhorasani_ascher_14}, quadratic trace estimators were analyzed before either of these papers.
For instance, \cite{reimann_07} provides tail bounds based on Chebyshev's inequality for quadratic trace estimators used for the specific purpose of estimating the trace of a symmetric matrix.
Sub-Gaussian concentration inequalities for quadratic trace estimators, similar to those in \cite{avron_toledo_11,roostakhorasani_ascher_14} are derived in \cite{popescu_short_winter_06} using Levy's Lemma, a general result about concentration of measure \cite{ledoux_01}; see also \cite[Theorem 2.2.2]{gogolin_10}.


There are also many earlier analyses of quadratic trace estimators outside of the specific context of trace estimation.
For instance, \cite{hanson_wright_71} provides concentration bounds for quadratic trace estimators when the entries of \( \vec{v} \) are independent symmetric sub-Gaussian random variables. 
In fact, some of the strongest bounds for quadratic trace estimators \cite{meyer_musco_musco_woodruff_21,persson_cortinovis_kressner_22} make use of so called \emph{Hanson--Wright inequalities} \cite{rudelson_vershynin_13} introduced in \cite{hanson_wright_71}.
Earlier still, \cite{grenander_pollak_slepian_59} states as fact that the expectation of such estimators, when \( \vec{v} \) has iid Gaussian entries, is the sum of the eigenvalues of the matrix in question, citing a book \cite{cramer_46} from the 1940s. 
%Similar concentration inequalities for quadratic trace estimators used in a specific task are given in \cite{popescu_short_winter_06} and are derived from L\'{e}vy's Lemma \cite{levy_51}, a result about concentration of measure  



\subsection{Other randomized trace estimation algorithms}

As a consequence of the central limit theorem, the average of iid samples of quadratic trace estimators requires \( O(\epsilon^{-2}) \) samples to reach accuracy \( \epsilon \). 
In fact, any algorithm which returns a linear combination of estimators depending on vectors drawn independently of \( \vec{A} \) requires \( O(\epsilon^{-2}) \) samples to obtain an approximation of the trace accurate to within a multiplicative factor \( 1\pm \epsilon \) \cite{wimmer_wu_zhang_14}.
A number of papers aim to avoid this dependence on the number of samples by incorporating low-rank approximation to \( \fA \) \cite{lin_17,gambhir_stathopoulos_orginos_17,saibaba_alexanderian_ipsen_17,large_matrices_density_review_18,li_zhu_21,meyer_musco_musco_woodruff_21,persson_cortinovis_kressner_22,chen_hallman_22}.


In \cite{meyer_musco_musco_woodruff_21} algorithm called Hutch++ in introduced and proved to output an estimate the trace of a positive definite matrix to relative error \( 1\pm \epsilon \) using just \( O(\epsilon^{-1}) \) matrix-vector products.
It is also shown that this \( \epsilon \) dependence is nearly optimal in certain matrix-vector query models.
The practicality of Hutch++ was improved in \cite{persson_cortinovis_kressner_22} which describes a variant which outputs an \( (\epsilon,\delta) \) approximation to the trace.
Such methods can be used to compute the trace of matrix functions by computing products with \( \fA \) (e.g. using black-box Krylov subspace methods). 

A so-called \emph{Krylov-aware} approach to estimating the trace of matrix functions was introduced in \cite{chen_hallman_22}.
Rather than treating products with \( \fA \) as a black-box, \cite{chen_hallman_22} advocates a more careful approach in which products with \( \vec{A} \) are viewed as the natural computational primative. 
This allows several efficincies not present in black-box versions of Hutch++ for matrix functions by producing better low-rank approximations. 
At least in terms of the total number of matrix-vector products used, the Krylov-aware approach always outperforms Hutch++ and related variants.

Finally, we note several more specialized techniques which may be of interest.
Variance reduction techniques based on multi-level Monte Carlo methods are studied in \cite{hallman_troester_21,frommer_khalil_ramirezhidalgo_21}.
In \cite{dharangutte_musco_21}, the problem of estimating the traces of a sequence of slowly-varying implicit matrices is studied. 
Such a setting occurs naturally in machine learning and physics. 
In physics, in order to compute important quantities for open quantum systems interacting strongly with their environment, \cite{chen_cheng_22} studies how to approximate the \emph{partial trace} of matrix functions.

\section{Analysis}

A simple approach to analyzing the protoalgorithm is to separately analyze the errors due to randomness in quadratic trace estimators from the error in approximating quadratic forms. 
Specifically, we have
\begin{align}
    \bigg| \int f \,\d\big(\Phi - \samp{\qq{\Psi_\ell}{s}}\big) \bigg|
    &\leq 
    \bigg| \int f \,\d\big(\Phi - \samp{\Psi_\ell}\big) \bigg|
    + \bigg| \int f \,\d\big(\samp{\Psi_\ell} - \samp{\qq{\Psi_\ell}{s}}\big) \bigg|
    \nonumber
    \\&=
    \bigg| \int f \,\d\big(\Phi - \samp{\Psi_\ell}\big) \bigg|
    + \bigg| \int f \,\d\samp{\Psi_\ell - \qq{\Psi_\ell}{s}} \bigg|
    \nonumber
    \\&\leq
    \bigg| \int f \,\d\big(\Phi - \samp{\Psi_\ell}\big) \bigg|
    + \samp[\bigg]{ \bigg| \int f \,\d\big(\Psi_\ell - \qq{\Psi_\ell}{s}\big) \bigg| }.
    %\nonumber
    %\\&=n^{-1} \bigg| \tr(\fA) - \samp{\vec{v}_{\ell}^\cT \fA \vec{v}_\ell} \bigg|
    %+ \samp[\bigg]{ \bigg| \int f \,\d\big(\Psi_\ell - \qq{\Psi_\ell}{s}\big) \bigg| }.
    \label{eqn:proto_triangle_f}
\end{align}
%Here we are using \( \samp{\:\cdot\:} \) to denote the average over \( \ell=1, \ldots, \nv \) and have suppressed the variable and bounds of integration for notational clarity.

The first term in \cref{eqn:proto_triangle_f} is controlled by the convergence of \( \samp{\Psi_{\ell}} \) to \( \Phi \) (as \( \nv \to \infty \)). 
Since
\begin{equation*}
    \bigg| \int f \,\d\big(\Phi - \samp{\Psi_\ell}\big) \bigg|
    =
    \left| n^{-1} \tr(\fA) - \samp{ \vec{v}_{\ell}^\cT \fA \vec{v}_\ell } \right|,
\end{equation*}
it can be analyzed in terms of bounds for quadratic trace estimators. 
Next, for each \( \ell \), the second term is controlled by the quality of the approximation of \( \Psi_{\ell} \) by \( [\Psi_{\ell}]_s^\circ \) (as \( s \to\infty \)).
Since
\begin{equation*}
    \bigg| \int f \,\d\big(\Psi_\ell - \qq{\Psi_\ell}{s}\big) \bigg|
    = 
    \bigg| \int (f - \ff{f}{s})  \,\d \Psi_\ell\bigg|,
\end{equation*}
we can analyze this term bounds for Krylov subspace methods for quadratic forms.


\subsection{Uniform unit test vectors}

\begin{definition}
    \label{eqn:hypersphere}
    The complex unit hypersphere \( \mathbb{S}^{n-1} \) is the set of unit vectors; i.e. 
    \begin{equation*}
        \mathbb{S}^{n-1} := \{ \vec{u} : \| \vec{w} \|_2 = 1\}.
    \end{equation*}
\end{definition}

In this section, we analyze the weighted CESM when \( \vec{v} \) is drawn from the uniform distribution on \( \mathbb{S}^{n-1} \).
In the case that \( \vec{A} \) is symmetric, similar results hold for uniform vectors drawn from the real unit hypersphere; see \cite{chen_trogdon_ubaru_22}.

\begin{lemma}
\label{thm:psi_beta}
Suppose \( \vec{v} \sim \operatorname{Unif}(\mathbb{S}^{n-1}) \) and, for any \( t\in \R \), define \( m(x) = n \Phi(x) \).
Then, 
\begin{align*}
    \Psi(x) \sim \operatorname{Beta} \big( m(x), n - m(x) \big).
\end{align*}
\end{lemma}

\begin{proof}
    Let \( \vec{U} = [\vec{u}_1, \ldots, \vec{u}_n] \), where \( \vec{u}_i \) is the \( i \)-th normalized eigenvector of \( \vec{A} \).
    Since \( \vec{U} \) is unitary, by the invariance of \( \operatorname{Unif}(\mathbb{S}^{n-1}) \) under orthogonal transforms, we have that \( \vec{U}^\cT \vec{v} \sim \operatorname{Unif}(\mathbb{S}^{n-1}) \).
    
    We may therefore assume \( \vec{U}^\cT\vec{v} \stackrel{\textup{dist.}}{=} \vec{x} / \| \vec{x} \|_2 \), where \( \vec{x} \sim \operatorname{ComplexNormal}(\vec{0},\vec{I}) \). % and ``\( \stackrel{\textup{dist.}}{=} \)'' denotes equality in distribution.
    Recall that the \( i \)-th weight of \( \Psi \) is given by \( w_i = |\vec{v}^\cT \vec{u}_i|^2 \).
    Thus, the \( w_i \) have joint distribution given by, 
\begin{align*}
    w_i 
    \stackrel{\textup{dist.}}{=} 
    \left|\frac{[\vec{x}]_i}{\|\vec{x}\|_2} \right|^2 
    = \frac{|[\vec{x}]_i|^2}{\sum_{i=0}^{n-1}|[\vec{x}]_i|^2},
\end{align*}
for \( i=0,1,\ldots, n-1 \).

Write, for notational convenience, \( m = m(x) = n \Phi(x) \).
Then, 
\begin{align*}
    \Psi(x)
    = \smop{\sum_{j=0}^{m-1}} w_j
    \stackrel{\textup{dist.}}{=}
    \frac{\sum_{i=0}^{m-1} |[\vec{x}]_i|^2}{\sum_{i=0}^{n-1} |[\vec{x}]_i|^2}.
\end{align*}    
    It is well known that for independent chi-square random variables \( Y \sim \chi_{\alpha}^2 \) and \( Z \sim \chi_{\beta}^2 \) (see, for example, \cite[Section 25.2]{johnson_kotz_balakrishnan_94}),
\begin{align*}
    \frac{Y}{Y+Z}\sim \operatorname{Beta} \left( \frac{\alpha}{2}, \frac{\beta}{2} \right).
\end{align*}
    Thus, since \( \sum_{i=0}^{m-1} |[\vec{x}]_i|^2 \) and \( \sum_{i=m}^{n-1} |[\vec{x}]_i|^2 \) are independent chi-square random variables with \( 2m \) and \( 2(n-m) \) degrees of freedom (because we are using complex normal random variables) respectively, \( \Psi(x) \) is a beta random variable with parameters \( m \) and \( n-m \).
\end{proof}

\begin{definition}
A random variable \( X \) is \( \sigma^2\)-sub-Gaussian if %\( \EE[|X|] < \infty \) and  
\begin{align*}
    \EE\big[ \exp(\lambda (X - \EE[X]))\big]
    \leq \exp \left( \frac{\lambda^2 \sigma^2}{2} \right)
    ,~\forall \lambda \in \mathbb{R}.
\end{align*}
%The smallest value of \( \sigma^2 \) such that \( X \) is \( \sigma^2 \) sub-Gaussian is called the optimal proxy variance of \( X \).
\end{definition}

\begin{lemma}
\label{thm:subgaussian_sum}
Suppose \( X \) is \( \sigma^2 \)-sub-Gaussian.
Let \( X_0, \ldots, X_{n_{\textup{v}}-1} \) be iid samples of \( X \).
Then for all \( \epsilon \geq 0 \),
\begin{align*}
    \PP\big[ |\samp{X_i} - \EE[X] | > \epsilon \big]
    \leq 2 \exp \left( - \frac{n_{\textup{v}}}{2 \sigma^2} \epsilon^2 \right).
\end{align*}
\end{lemma}


\begin{proof}
We follow a standard argument; see for instance \cite{vershynin_18}.
WLOG assume \( \EE[X] = 0 \).
Then,
\begin{align*}
    \PP[ \nv \samp{X_i} \geq \nv \epsilon ] 
    &= \PP[ \exp( \lambda \nv\samp{X_i} ) \geq \exp(\lambda \nv \epsilon )]
    \\&\leq \exp(- \nv \lambda \epsilon) \EE[ \exp(\lambda \nv \samp{X_i}) ] \tag{Markov}
    \\&= \exp(- \nv \lambda \epsilon) \EE[ \exp(\lambda X) ]^{\nv} \tag{iid}
    \\&\leq \exp(- \nv \lambda \epsilon) \exp( \nv \lambda^2 \sigma^2 / 2) \tag{sub-Gaussian}
    \\&= \exp(- \nv \lambda \epsilon + \nv \lambda^2 \sigma^2 / 2). 
\end{align*}
This expression is minimized when \( \lambda = t / \sigma^2 \) from which we obtain,
\begin{align*}
    \PP[ \samp{X_i} \geq t] \leq \exp \left( - \frac{\nv}{2\sigma^2} t^2 \right). \tag*{\qedhere}
\end{align*}
\end{proof}

\begin{theorem}\textup{\cite[Theorem 1]{marchal_arbel_17}}
\label{thm:beta_tails}
Suppose \( X \sim \operatorname{Beta}(\alpha,\beta) \).
Then, \( \EE[X] = \alpha / (\alpha + \beta) \), and \( X \) is \( (4(\alpha +\beta + 1))^{-1} \)-sub-Gaussian.
If \( \alpha = \beta \), then there is no smaller \( \sigma^2 \) such that \( X \) is \( \sigma^2 \)-sub-Gaussian.
\end{theorem}

With these results in place, the following theorem for spectrum approximation is straightforward.
\begin{theorem}
\label{thm:CESM_estimator_prob}
Given a positive integer \( \nv \), suppose \( \{ \vec{v}_{\ell} \}_{\ell=0}^{\nv-1} \stackrel{\textup{iid}}{\sim} \operatorname{Unif}(\mathbb{S}^{n-1}) \).
Then, for all \( \varepsilon > 0 \),
\begin{align*}
    \max_{x\in\R} \PP\left[ \left| \Phi(x) - \samp{ \Psi_{\ell}(x) } \right| > \varepsilon \right]
    &\leq
    2 \exp\left( - 2\nv (n+1) \varepsilon^2 \right).
    \\
    \PP\left[ \max_{x\in\R} \left| \Phi(x) - \samp{ \Psi_{\ell}(x) } \right| > \varepsilon \right]
    &\leq
    2 n \exp\left( - 2\nv (n+1) \varepsilon^2 \right).
\end{align*}
\end{theorem}

\begin{proof}
First note that the maximums exist because \( \Phi\) and \( \samp{\Psi_i} \) are right continuous and piecewise constant except at \( \{ \lambda_i[\vec{A}] \}_{i=1}^{n} \).

For any \( t \), let \( m  = m(x) = n \Phi(x) \).
Using \cref{thm:psi_beta,thm:beta_tails,thm:subgaussian_sum} we have that for any \( t \),
\begin{align*}
    \PP\big[ \left| \Phi(x) - \samp{ \Psi_i(x) } \right| > \epsilon \big]
    \leq 2 \exp \left( - \frac{n_{\textup{v}}}{2 (4(m + (n-m)+1) )^{-1}} \epsilon^2 \right).
\end{align*}
%Both \( \Psi \) and \( \samp{\Psi_i} \) are right continuous, so the first result follows from a limit argument. 

We also have
\begin{align*}
    \sup_{x\in\R} |\Phi(x) - \samp{\Psi_i(x)}|
    = \max_{0\leq i<n-1} | \Phi(\lambda_i[\vec{A}]) - \samp{\Psi_i(\lambda_i[\vec{A}])} |.
\end{align*}
%\note{TODO: line number}
    The second result follows by applying a union bound to the events that the maximum is attained at \( \lambda_i \) for each \( i=0,1,\ldots, n-2 \) (note since \( \|\vec{v}\|_2=1 \), \( \Phi \) and \( \Psi_{\ell} \) agree at \( \lambda_{n-1} \)).
\end{proof}

This result can be used to obtain a bound for quadratic trace estimation.
\begin{theorem}
    \label{thm:unif_QTE}
    Set \( \nv \geq 1 \) and sample \( \{ \vec{v}_{\ell} \}_{\ell=0}^{\nv-1} \stackrel{\textup{iid}}{\sim} \operatorname{Unif}(\mathbb{S}^{n-1}) \).
    Then
    \begin{equation*}
        \PP\left[ \left| n^{-1} \tr(\vec{A}) - \samp{ \vec{v}_{\ell}^\cT \vec{A} \vec{v}_\ell } \right| > \varepsilon (\lmax - \lmin)  \right]
        \leq 2 n \exp (- 2( n+1 ) \nv \varepsilon^2 ).
    \end{equation*}
\end{theorem}

\begin{proof}
    Since \( \samp{ \Psi_\ell } \) and \( \Phi\) are both constant on each of \( (-\infty , \lmin ) \) and \( ( \lmax ,\infty) \), 
\begin{equation*}
    \W(\Phi,\samp{\Psi_\ell})
    = \int | \Psi - \samp{\Psi_\ell} | \,\d{x}
    \leq (\lmax - \lmin) \| \Psi - \samp{\Psi_\ell} \|_{\R}.
\end{equation*}
Using \cref{thm:CESM_estimator_prob},
%\begin{align*}
%    \PP\left[ \max_{x\in\R}|\Phi(x) - \samp{\Psi_\ell(x)} | > \varepsilon  \right] \leq 2 n \exp (- ( n+2 ) \nv \varepsilon^2 ).
%\end{align*}
%Combining these we have that
we find that
\begin{equation*}
    \PP\big[ \W(\Phi,\samp{\Psi_\ell}) > \varepsilon (\lmax-\lmin) \big] 
    \leq 2 n \exp (- 2( n+1 ) \nv \varepsilon^2 ).
\end{equation*}
    Thus, using \cref{thm:wasserstein_1lip} and the fact that \( x \) is 1-Lipshitz, 
\begin{equation*}
    \PP\left[ \left| \int x \,\d\Phi - \int x \,\d\samp{\Psi_\ell}) \right| > \varepsilon (\lmax-\lmin) \right] 
    \leq 2 n \exp (- 2( n+1 ) \nv \varepsilon^2 ).
\end{equation*}
Next, recall that \( \int x \,\d\Phi = n^{-1} \tr(\vec{A}) \) and \( \int x \,\d\samp{\Psi_\ell} = \samp{ \vec{v}_{\ell}^\cT \vec{A} \vec{v}_\ell } \).
Thus, we obtain a bound for the quadratic trace estimator:
\begin{equation*}
    \PP\left[ \left| n^{-1} \tr(\vec{A}) - \samp{ \vec{v}_{\ell}^\cT \vec{A} \vec{v}_\ell } \right| > \varepsilon (\lmax - \lmin) \right] 
    \leq 2 n \exp (- 2( n+1 ) \nv \varepsilon^2 ). \qedhere
\end{equation*}
%The result follows by replacing \( \vec{A} \) with \( f[\vec{A}] \) in the above expression.
\end{proof}

This can be restated in terms of matrix functions.
\begin{corollary}
    \label{thm:wCESM_tr}
    Suppose \( f \) is bounded between \( f_{\textup{min}} \) and \( f_{\textup{max}} \) on the spectrum of \( \vec{A} \). 
    Set \( \nv \geq \frac{1}{2}( f_{\textup{max}} - f_{\textup{min}} )^2 (n+1)^{-1} \varepsilon^{-2} \ln(2n \eta^{-1}) \) and sample \( \{ \vec{v}_{\ell} \}_{\ell=0}^{\nv-1} \stackrel{\textup{iid}}{\sim} \operatorname{Unif}(\mathbb{S}^{n-1}) \).
    Then
    \begin{equation*}
        \PP\left[ \left| n^{-1} \tr(\fA) - \int f \,\d\samp{\Psi_\ell}) \right| > \varepsilon \right] \leq \eta.
    \end{equation*}
\end{corollary}

As we remarked in \cref{sec:past_work}, bounds similar to \cref{thm:unif_QTE} have been studied for other distributions for \( \vec{v} \). 
The best bounds are for Gaussian and Rademacher vectors, which have \emph{independent} entries. 
For such distributions, the best bounds depend on \( \|\vec{A}\|_\F^2 \) rather than \( n\|\vec{A}\|_2^2 \) and are therefore significantly stronger than \cref{thm:unif_QTE} when the stable rank \( \| \vec{A} \|_\F^2 / \| \vec{A}\|_2^2 \) is small. 
It is likely that the bounds in \cref{thm:unif_QTE} can be improved by a more careful analysis of Beta random variables. 
In particular, while the sub-Gaussian constant from \cref{thm:beta_tails} is sharp when \( \alpha = \beta \), it can be improved when \( \alpha \approx 0 \) or \( \beta\approx 0 \) \cite{zhang_zhou_20}.


\section{Numerical experiments}

\subsection{Approximating sparse spectra}

If the spectrum of \( \vec{A} \) is \( S \)-sparse; i.e., there are only \( S \) distinct eigenvalues, then the \( s \)-point Gaussian quadrature rule will be exactly equal to the weighted CESM for all \( s \geq S \), at least in exact arithmetic.
Thus, the runtime required by SLQ is determined by \( S \) and the number of samples of the weighted CESM which are required to get a good approximation to the true CESM.
The interpolation and approximation based approaches, which are based on the orthogonal polynomials of some fixed distribution function \( \mu \), are unable to take advantage of such sparsity.
Indeed, unless the eigenvalues of \( \vec{A} \) are known a priori, such methods have fixed resolution \( \sim s^{-1} \) due to the fixed locations of the zeros of the orthogonal polynomials with respect to \( \mu \).
Moreover, quadrature by approximation methods suffer from Gibbs oscillations unless a damping kernel is used, in which case the resolution is further decreased.


\begin{figure}[htb]
    \includegraphics[width=\textwidth]{imgs/ch4_Kneser_smoothed.pdf}
    \caption[{Approximations to a sparse spectrum with just \( 12 \) eigenvalues.}]{%
    Approximations to a sparse spectrum with just \( 12 \) eigenvalues.
    \hspace{.25em}\emph{Legend}:
    true spectrum  
    ({\protect\raisebox{0mm}{\protect\includegraphics[]{imgs/legend/ms4.pdf}}}).
    Gaussian quadrature approximation: \( k=12 \)
    ({\protect\raisebox{0mm}{\protect\includegraphics[]{imgs/legend/m1.pdf}}}).
    damped quadrature by approximation: \( s=500 \)
    ({\protect\raisebox{0mm}{\protect\includegraphics[]{imgs/legend/l3nm.pdf}}}).
    \hspace{.25em}\emph{Takeaway}: The Gaussian quadrature produces an extremely good approximation using just \( 12 \) matrix-vector products.
    Even with many more matrix-vector products, quadrature by approximation does not have the same resolution. 
    }
    \label{fig:Kneser}
\end{figure}


In this example, we approximate the CESM of the adjacency matrix of a Kneser graph.
The \( (N,K) \)-Kneser graph is the graph whose vertices correspond to size \( K \) subsets of \( \{1, 2, \ldots, N\} \) and whose edges connect vertices corresponding to disjoint sets.
It is not hard to see that the number of vertices is \( \binom{N}{K} \) and the number of edges is \( \frac{1}{2}\binom{N}{K}\binom{N-K}{K} \).
%Thus, for \( N < 2K \) the graph has no edges. 
The spectrum of Kneser graphs is known as well. 
Specifically, there are \( K+1 \) distinct eigenvalues whose values and multiplicities are:
\begin{equation*}
    \lambda_i = (-1)^i\binom{N-K-i}{K-i}
    ,\qquad
    m_i = \binom{N}{i} - \binom{N}{i-1}
    ,\qquad 
    i=0,1,\ldots, K.
\end{equation*}

We conduct a numerical experiment with \( N = 23 \) and \( K = 11 \), the same values used in \cite{large_matrices_density_review_18}.
This results in a graph with 1,352,078 vertices and 8,112,468 edges.
Thus, the adjacency matrix is highly sparse. 
We compare the Gaussian quadrature approximation with the damped quadrature by approximation. 
In both cases we use a single random test vector \( \vec{v} \).
For the Gaussian quadrature, we set \( k = 12 \).
For the damped quadrature by approximation we set \( s=500 \) and use Jackson damping with \( \mu = \mu_{a,b}^T \), where \( a = -11.1 \) and \( b=12.1 \).
The results are shown in \cref{fig:Kneser}.
Note that the Gaussian quadrature matches almost exactly despite having used only \( k=12 \) matrix-vector products. 
On the other hand, even after \( k=250 \) matrix-vector products, the damped quadrature by approximation has a much lower resolution.

\begin{remark}
There are sublinear time algorithms for approximate matrix-vector products with the (normalized) adjacency matrix.
Specifically, in a computational model where it is possible to 
(i) uniformly sample a random vertex in constant time, 
(ii) uniformly sample a neighbor of a vertex in constant time, and 
(iii) read off all neighbors of a vertex in linear time, 
then an \( \epsilon_{\textup{mv}} \)-accurate approximate to the a matrix-vector product with the adjacency matrix can be computed, with probability \( 1-\eta \), in time \( O(n (\epsilon_{\textup{mv}})^{-2}\ln(\eta^{-1})) \).
For dense graphs, this is sublinear in the input size \( O(n^2) \) of the adjacency matrix.
See \cite{braverman_krishnan_musco_22} for an analysis in the context of spectrum approximation.
\end{remark}

\subsection{Approximating ``smooth'' densities}
\label{sec:smooth_density}
%the limiting density of large random matrix models}

There are a range of settings in which the spectral density of \( \vec{A} \) is close to a smooth slowly varying density.
In such cases, we may hope that our approximation satisfies certain known criteria. 
For instance, that the approximation is also a slowly varying density, that the behavior of the approximation at the endpoints of the support satisfies the right growth or decay conditions, etc.
In this example, we consider how parameters in \cref{alg:protoalg} can be varied so that the resulting approximation enjoys certain desired properties. 

%\cite{large_matrices_density_review_18}
 
One setting in which \( \vec{A} \) may have a slowly varying density is when \( \vec{A} \) is a large random matrix. 
We begin this example by considering a sample covariance matrix
\begin{equation*}
    \vec{A}_n = \frac{1}{m} \vec{\Sigma}^{1/2} \vec{X} \vec{X}^\cT \vec{\Sigma}^{1/2}
\end{equation*}
where \( \vec{X} \) is random and \( \vec{\Sigma} \) is deterministic. 
Specifically, we fix constants \( \sigma > 1 \) and \( d \in (0,1) \), define \( m = n/d \), and take \( \vec{X} \) to be a \( n\times m \) matrix with iid standard normal entries and \( \vec{\Sigma} \) a diagonal matrix with \( 1/m \) as the first \( n/2 \) entries and \( \sigma/m \) as the last \( n/2 \) entries.

In the limit, as \( n\to\infty \), the spectral density \( \,\d\Phi_n /\d x \) of \( \vec{A}_n \) is convergent to a density \( \,\d\Psi_\infty/\d x \) supported on two disjoint intervals \( [a_1,b_1] \cup [a_2,b_2] \), where \( a_1 < b_1 < a_2 < b_2 \), with equal mass on each \cite{bai_silverstein_98}.
The spectral edges are equal to the values at which 
\begin{equation*}
    t\mapsto  -\frac{1}{t} + \frac{d}{2} \left( \frac{1}{t+1} + \frac{1}{t+\sigma^{-1}} \right)
\end{equation*}
attains at its local extrema.
Moreover, it is known that \( \,\d\Psi_\infty/\d x \) has square root behavior at the spectral edges.


Because we know the support of the desired density, and because we know the behavior at the spectral edges, a natural choice is to use quadrature by approximation with 
\begin{equation*}
    \mu 
    = \frac{1}{2}\mu_{a_1,b_1}^U + \frac{1}{2} \mu_{a_2,b_2}^U
\end{equation*}
where \( \mu_{a,b}^U \) is the weight function for the Chebyshev polynomials of the second kind given by
\begin{equation*}
    \frac{\d\mu_{a,b}^U}{\d{x}} 
    = \frac{4}{\pi (b-a)} \sqrt{1-\left( \frac{2}{b-a}x - \frac{b+a}{b-a}\right)} .
\end{equation*}
This will ensure that the Radon--Nikodym derivative \( \,\d\Psi_\infty / \,\d\mu \) is of order 1 at the spectral edges which seems to result in better numerical behavior than if we were to use a KPM approximation corresponding to a density which explodes at the spectral edges.


To compute the Jacobi matrix for \( \mu \), we apply the Stieltjes procedure using a slight modification of the \emph{Vandermonde with Arnoldi} approach \cite{brubeck_nakatsukasa_trefethen_21}.
In order to apply the Stieltjes procedure, we must be able to integrate polynomials against \( \mu \).
Observe that the product
\begin{equation*}
    \int p \,\d\mu
    = \frac{1}{2} \int p \,\d\mu_{a_1,b_1}^U
    + \frac{1}{2} \int p \,\d\mu_{a_2,b_2}^U
\end{equation*}
can be computed \emph{exactly} by applying a sufficiently high degree quadrature rule to each of the right hand side integrals.
If we aim to compute the \( s\times s \) Jacobi matrix associated with \( \mu \) the maximum degree polynomial we will integrate will be of degree \( 2s-1 \) when we orthogonalize \( x p_{s-1} \) against \( p_{s-1} \).
Therefore, it suffices to use the degree \( s \) Gaussian quadrature rules for \( \mu_{a_1,b_1}^U \) and \( \mu_{a_2,b_2}^U \) for all of the first \( s \) iterations of the Stieltjes procedure.

One simple approach to running the Stieltjes procedure in this manner is to place the quadrature nodes on the diagonal of a matrix \( \vec{N} \) and the corresponding weights on a vector \( \vec{w} \).
Then the weighted CESM corresponding to \( \vec{N} \) and \( \vec{w} \) is a quadrature rule which integrate polynomials of degree up to \( 2s-1 \) against \( \mu \) exactly.
The tridiagonal matrix obtained by the Lanczos algorithm run for \( s \) iterations will be exactly the upper \( s\times s \) block of the Jacobi matrix \( \vec{M}(\mu) \).
Some potentially more computationally efficient approaches are outlined in \cite{fischer_golub_91}.

\begin{figure}[htb]
    \includegraphics[width=\textwidth]{imgs/ch4_RM_split_AQ.pdf}
    \caption[{Approximations to a ``smooth'' spectrum using quadrature by approximation with various choices of \( \mu \).}]{%
        Approximations to a ``smooth'' spectrum using quadrature by approximation with various choices of \( \mu \).
        \hspace{.25em}\emph{Legend}:
        \( \mu = \mu_{a_1,b_2}^U \) 
        ({\protect\raisebox{0mm}{\protect\includegraphics[]{imgs/legend/l3nm.pdf}}}).
        \( \mu = \frac{1}{2} \mu_{a_1,b_1}^U + \frac{1}{2} \mu_{a_2,b_2}^U \)
        ({\protect\raisebox{0mm}{\protect\includegraphics[]{imgs/legend/l1nm.pdf}}}).
        \hspace{.25em}\emph{Takeaway}: A priori knowledge about the spectrum allows for better choices of parameters such as \( \mu \). 
    }
    \label{fig:RM_split_AQ}
\end{figure}

We conduct a numerical experiment with \( n = 10^4 \) and \( d=0.3 \).
We use \( s=60 \) and average over 10 trials, resampling \( \vec{A}_n \) in each trial.
To generate an approximation to the density we expand the support of the limiting density by \( 0.001 \) on endpoint to avoid eigenvalues of \( \vec{A}_n \) lying outside the support of \( \mu \).
In \cref{fig:RM_split_AQ} we show the approximations with \( \mu = \mu_{a_1,b_2}^U \) and \( \mu = \frac{1}{2} \mu_{a_1,b_1}^U + \frac{1}{2} \mu_{a_2,b_2}^U \).
As shown in the inset image of \cref{fig:RM_split_AQ}, we observe that the approximation with \( \mu = \frac{1}{2} \mu_{a_1,b_1}^U + \frac{1}{2} \mu_{a_2,b_2}^U \) exhibits the correct square root behavior at the endpoints as well as fewer oscillations throughout the interior of the support of the density.


\begin{remark}
In recent work \cite{ding_trogdon_21} it was shown how Lanczos performs on such a sample covariance matrix. 
In particular, one sample from stochastic Lanczos quadrature will converge almost surely, as \( n\to\infty \), to the desired distribution.
In this same work another density approximation scheme was proposed based on Stieltjes transform inversion. 
Analysis and comparison for this method is an interesting open problem.
\end{remark}


\subsubsection{Smoothing by convolution}
\label{sec:smoothing}


The Gaussian quadrature approximation is the sum of weighted Dirac delta functions.
A simple approach to obtain a density function from a distribution function involving point masses is to approximate each point masses with some concentrated probability density function; e.g. Gaussians with a small variance \cite{lin_saad_yang_16,ghorbani_krishnan_xiao_19}.
This is simply convolution with this distribution, and if the smoothing distribution has small enough variance, the Wasserstein distance between the original and smoothed distributions will be small. 
Specifically, we have the following standard lemma:
\begin{lemma}
    \label{thm:smoothing_wass}
    Given a smooth positive probability distribution function \( G_\sigma \), define the smoothed approximation \( \Upsilon_\sigma \) to \( \Upsilon \) by the convolution
    \begin{equation*}
        \Upsilon_\sigma(x)  := \int_{-\infty}^{\infty} G_\sigma(t-y) \,\d\Upsilon(y).
    \end{equation*}
    Then, \( \W(\Upsilon,\Upsilon_\sigma) \leq \W(\bOne[ x < 0], G_\sigma) \TV(\Upsilon) \).
    Moreover, if \( G_\sigma \) has median zero and standard deviation \( \sigma \), then \( \W(\Upsilon,\Upsilon_\sigma)  \leq \sigma \TV(\Upsilon) \). 
\end{lemma}

\begin{figure}[htb]
\includegraphics[width=\textwidth]{imgs/ch4_RM_split_GQ.pdf}
    \caption[{Approximations to a ``smooth'' spectrum using smoothed Gaussian quadrature for various smoothing parameters \( \sigma \).}]{%
    Approximations to a ``smooth'' spectrum using smoothed Gaussian quadrature for various smoothing parameters \( \sigma \).
    \hspace{.25em}\emph{Legend}:
    \( \sigma = 3/k \) 
        ({\protect\raisebox{0mm}{\protect\includegraphics[]{imgs/legend/l3nm.pdf}}}).
    \( \sigma = 8/k \) 
        ({\protect\raisebox{0mm}{\protect\includegraphics[]{imgs/legend/l2nm.pdf}}}).
    \( \sigma = 15/k \) 
        ({\protect\raisebox{0mm}{\protect\includegraphics[]{imgs/legend/l1nm.pdf}}}).
    \hspace{.25em}\emph{Takeaway}: Gaussian quadrature is not always the best choice of algorithm. 
    Here we observe that it is difficult to produce a density approximation using the specified smoothing scheme.
}
\label{fig:GQ_conv_smoothing}
\end{figure}


It is well known that if \( G_\sigma \) is differentiable then the smoothed distribution function \( \Upsilon_\sigma \) will also be differentiable.
Thus, we can obtained a density function \( \,\d\Upsilon_\sigma/\d{x} \) even if \( \Upsilon \) has discontinuities.
Moreover, the bounds obtained earlier can easily be extended to smoothed spectral density approximations obtained by convolution using the triangle inequality.

While the smoothing based approach has a simple theoretical guarantee in Wasserstein distance, it does not need to provide a good approximation to the density. 
Indeed, if the variance of the smoothing kernel is too small, then the smoothed distribution will still look somewhat discrete. 
On the other hand, if the variance of the smoothing kernel is too large, then the smooth distribution will become blurred out and lose resolution.
As shown in \cref{fig:GQ_conv_smoothing}, this is particularly problematic if different parts of the spectrum would naturally require different amounts of smoothing.

There are of course many different smoothing schemes that could be used. 
These include adaptively choosing the variance parameter based on the position in the spectrum, using a piecewise constant approximation to the density, interpolating the distribution function with a low degree polynomial or splines, etc. 
Further exploration of these approaches is beyond the scope of this thesis since they would likely be context dependent. 
For instance, in random matrix theory, it may be desirable to enforce square root behavior at endpoints whereas in other applications it may be desirable to have smooth tails.


We conclude with the remark that alternate metrics of closeness, such as the total variation distance, are likely better suited for measuring the quality of approximations to ``smooth'' densities.
However, since the actual spectral density \( \,\d\Psi/\d{x} \) is itself the sum of Dirac deltas, some sort of regularization is required to obtain a proper density \cite{lin_saad_yang_16} which of course relates closely to what it actually means to be ``close to a smooth slowly varying density''.
A rigorous exploration of this topic would be of interest. 


\subsubsection{Handling isolated spikes}

In some situations one may encounter spectra which are nearly ``smooth'' except at a few points at which there are large jumps in the CESM (for instance, low rank matrices may have many repeated zero eigenvalues). 

To model such a situation, we consider a matrix 
\begin{equation*}
    \vec{A}_n := 
    \begin{bmatrix}
        m^{-1} \vec{X}\vec{X}^\cT & \vec{0} \\
        \vec{0} & z\vec{I} + \sigma \vec{D}
    \end{bmatrix}
\end{equation*}
where \( \vec{X} \) is a \( n'\times m \) matrix standard normal entries and \( \vec{D} \) is a \( (n-n')\times (n-n') \) diagonal matrix with standard normal entries. 
In both cases,  \( m = n'/d \) for some fixed \( d\in(0,1) \). 
While this particular matrix is block diagonal, the protoalgorithm is mostly oblivious to this structure and would work similarly well if the matrix were conjugated by an arbitrary unitary matrix so that the block diagonal structure is lost.

When \( n\to\infty \) and \( \sigma \to 0 \), the spectral density \( \,\d\Phi_n/\d{x} \) is convergent to a density \( \,\d\Phi_\infty/\d{x} \) equal to the sum of a scaled Marchenko--Pastur distribution and a weighted Dirac delta distribution. 
Thus, a natural approach would be to use quadrature by approximation with
\begin{equation*}
    \mu = (1-p) \mu_{a,b}^{U} + p \,\delta(x-z).
\end{equation*}
As above, we can use a modified version of the Vandermonde with Arnoldi approach to compute the orthogonal polynomials with respect to \( \mu \). 
The resulting approximation to the ``smooth'' part of the density \( \,\d\Phi/\d{x} \) is shown in \cref{fig:RM_spike}.

\begin{figure}[htb]
\includegraphics[width=\textwidth]{imgs/ch4_RM_AQ_spiked.pdf}
    \caption[{Approximations to a ``smooth'' spectrum with a spike using quadrature by approximation with various choices of \( \mu \).}]{%
    Approximations to a ``smooth'' spectrum with a spike using quadrature by approximation with various choices of \( \mu \).
    \hspace{.25em}\emph{Legend}:
    absolutely continuous part of true limiting density
    ({\protect\raisebox{0mm}{\protect\includegraphics[]{imgs/legend/l4nm.pdf}}}).
    quadrature by approximation: \( \mu = (1-p) \mu_{a,b}^{U} + p \,\delta(x -z) \)
    ({\protect\raisebox{0mm}{\protect\includegraphics[]{imgs/legend/l1nm.pdf}}}).
    quadrature by approximation: \( \mu = \mu_{a,b}^U \) 
    ({\protect\raisebox{0mm}{\protect\includegraphics[]{imgs/legend/l3nm.pdf}}}).
    \hspace{.25em}\emph{Takeaway}: A priori knowledge of the location of a singularity allows for a better approximation to the absolutely continuous part of the spectrum. 
    }
\label{fig:RM_spike}
\end{figure}



We set \( n = 10^6 \), \( n' = n/10 \), \( d = 0.3 \), \( z=1.5 \), and \( \sigma = 10^{-10} \).
As before, we average of \( 10 \) trials where \( \vec{A}_n \) is resampled in each trial.
For each sample, we compute the quadrature by approximation with \( s=200 \) for \( \mu = (1-p) \mu_{a,b}^{U} + p \,\delta(x-z) \) with \( p = 0.2 \) and \( \mu = \mu_{a,b}^U \).
The results are show in \cref{fig:RM_spike}.

Clearly, accounting for the spike explicitly results in a far better approximation to the density.
Note that this approach does not require that the mass of the spike is accurately matched. 
For instance, in our example, we estimate the spike mass to be 0.2 while the actual mass is 0.1.
On the other hand, if the location of the spike is misestimated, then the approximation to the density may have massive oscillations.
In our example the spike has width roughly \( 10^{-10} \) which does not cause issues for the value of \( s \) used. 
However, if \( s \) is increased, the width of the spike is increased, or the location of the estimate of the spike is offset significantly, then the existing oscillations become large.
Approaches for adaptively finding the location of spikes would an interesting area of further study. 




\subsection{Energy spectra of small spin systems}

The quantum Heisenberg model can be used to study observables of magnetic systems \cite{weisse_wellein_alvermann_fehske_06,schnalle_schnack_10,schnack_richter_steinigeweg_20,schulter_gayk_schmidt_honecker_schnack_21,schlter_richter_schnack_22}.
For a system with \( N \) spins of spin number \( S \), the Heisenberg spin Hamiltonian is an operator on a Hilbert space of dimension \( (2S+1)^N \) given by
\begin{equation*}
    \vec{H} = \sum_{i=0}^{N-1} \sum_{j=0}^{N-1} \left( 
     [\vec{J}^{\textup{x}}]_{i,j} \vec{s}^{\textup{x}}_i  \vec{s}^{\textup{x}}_j 
    +[\vec{J}^{\textup{y}}]_{i,j}  \vec{s}^{\textup{y}}_i  \vec{s}^{\textup{y}}_j
    +[\vec{J}^{\textup{z}}]_{i,j} \vec{s}^{\textup{z}}_i  \vec{s}^{\textup{z}}_j
    \right).
\end{equation*}
Here \( \vec{s}^\sigma_i \) gives the component spin operator for the \( i \)-th spin site and acts trivially on the Hilbert spaces associated with other spin sites but as the \( (2S+1)\times (2S+1) \) component spin matrix \( \vec{s}^{\sigma} \) on the \( i \)-th spin site.
Thus, \( \vec{s}^\sigma_i \) can be represented in matrix form as
\begin{equation*}
    \vec{s}^\sigma_i
    = \underbrace{\vec{I} \otimes \cdots \otimes \vec{I}}_{i\text{ terms}} 
    \otimes ~ \vec{s}^\sigma \otimes 
    \underbrace{\vec{I} \otimes \cdots \otimes \vec{I}}_{N-i-1\text{ terms}}.
\end{equation*}
%Each of the terms in the Kronecker product representation of \( \vec{s}_i^\sigma \) are dimension \( 2s+1 \) operators.% so that \( \vec{s}^\sigma_i \) is a \( (2s+1)^N \) dimension operator.

\begin{figure}[htb]
    \includegraphics[width=\textwidth]{imgs/ch4_spin_heat_capacity.pdf}
    \caption[Heat capacity as a function of temperature for a small spin system.]{%
    Heat capacity as a function of temperature for a small spin system.
    \hspace{.25em}\emph{Legend}:
    exact diagonalization 
    ({\protect\raisebox{0mm}{\protect\includegraphics[]{imgs/legend/l4nm.pdf}}}),
    Gaussian quadrature
    ({\protect\raisebox{0mm}{\protect\includegraphics[]{imgs/legend/l1nm.pdf}}}),
    quadrature by approximation 
    ({\protect\raisebox{0mm}{\protect\includegraphics[]{imgs/legend/l2nm.pdf}}}), and
    damped quadrature by approximation 
    ({\protect\raisebox{0mm}{\protect\includegraphics[]{imgs/legend/l3nm.pdf}}}).
    \hspace{.25em}\emph{Takeaway}: While damping produces a physical result, the resulting ghost bump may be more difficult to identify than the nonphysical ghost dip obtained without damping.
    }
    \label{fig:spin_heat_capacity}
\end{figure}

The CESM of \( \vec{H} \) gives the energy spectrum of the system and can be used to compute many important quantities.
For instance, given an observable \( \vec{O} \) (i.e. a Hermitian matrix), the corresponding thermodynamic expectation of the observable in thermal equilibrium at inverse temperature \( \beta \)  is given by
\begin{equation*}
    \frac{\tr(\vec{O}\exp(-\beta \vec{H}))}{\tr(\exp(-\beta\vec{H}))}.
\end{equation*}

Quantities depending on observables which are matrix functions \( \vec{H} \) can be written entirely in terms of matrix functions of \( \vec{H} \).
For instance, the system heat capacity is given by
\begin{equation*}
    \frac{C(T)}{k_B}
    = 
    \frac{\tr\left( (\beta\vec{H})^2 \exp(-\beta \vec{H}) \right)}{\tr \left( \exp(-\beta \vec{H}) \right)}
    -\left[ \frac{\tr\left( \beta \vec{H} \exp(-\beta \vec{H}) \right)}{\tr \left( \exp(-\beta \vec{H}) \right)}\right]^2.
\end{equation*}
Thus, for fixed finite temperature, evaluating the heat capacity amounts to evaluating several matrix functions. 

In some cases, symmetries of the system can be exploited to diagonalize or block diagonalize \( \vec{H} \) \cite{schnalle_schnack_10}. 
Numerical diagonalization can be applied to blocks to obtain a full diagonalization.
Even so, the exponential dependence of the size of \( \vec{H} \) on the number of spin sites \( N \) limits the size of systems which can be treated in this way. 
Moreover, such techniques are not applicable to all systems. 
Thus, approaches based on \cref{alg:protoalg} are widely used; see \cite{schnack_richter_steinigeweg_20} for examples using a Lanczos based approach and \cite{schulter_gayk_schmidt_honecker_schnack_21} for examples using a Chebyshev based approach.

In this example, we consider a Heisenberg ring (\( [\vec{J}^{\textup{x}}]_{i,j} = [\vec{J}^{\textup{y}}]_{i,j} = [\vec{J}^{\textup{z}}]_{i,j} = \bOne[|i-j| = 1 \pmod{N}] \)) with \( N = 12 \) and \( S = 1/2 \).
Similar examples, with further discussion in the context of the underlying physics, are considered in \cite{schnack_richter_steinigeweg_20,schulter_gayk_schmidt_honecker_schnack_21}.
We take \( k = 50 \) and \( \nv = 300 \) and compute approximations to the heat capacity at many temperatures using Gaussian quadrature, quadrature by interpolation, and damped quadrature by interpolation. 
For the latter two approximations we use \( \mu = \mu_{a,b}^T \) where \( a \) and \( b \) are chosen based on the nodes of the Gaussian quadrature. 
Note that averages over random vectors are computed for each trace rather than for \( C(\beta) \), and that we use the same vectors for all four traces. 
The results are shown in \cref{fig:spin_heat_capacity}.


We note the presence of an nonphysical ``ghost dip'' in the quadrature by interpolation approximation. 
If the approximation to the CESM is non-decreasing, the Cauchy--Schwarz inequality guarantees a positive heat capacity.
Thus, when we use Jackson's damping, the heat capacity remains positive for all temperatures.
However, as noted in \cite{schulter_gayk_schmidt_honecker_schnack_21}, this is not necessarily desirable as the ghost dip is easily identifiable while the ghost peak may be harder to identify.

We conclude with the remark that it would be interesting to provide bounds for the accuracy of the approximations to the quantity \( \tr(\vec{O} \exp(-\beta \vec{H})) / \tr(\exp(-\beta \vec{H}))  \) for all \( \beta > 0 \).


 

\chapter{Optimal rational matrix function\\ approximation}
\label{chap:opt_algs}

We now shift gears a bit and consider methods for approximating \( \fA\vec{v} \) from Krylov subspace \( \mathcal{K}_k = \operatorname{span}\{\vec{v}, \vec{A}\vec{v}, \ldots, \vec{A}^{k-1}\vec{v} \} \).
The most ideal Krylov subspace method for this task would output iterates \( \mathsf{alg}_k \) satisfying
\begin{equation*}
    \mathsf{alg}_k = \argmin_{\vec{x}\in\mathcal{K}_k} \| \fA \vec{v} - \vec{x} \|.
\end{equation*}
The above condition guarantees the algorithm produces approximations with smaller error (in the given norm) than any other Krylov subspace method. 
Moreover, under the assumption \( \|\cdot\| \) is induced by a positive definite matrix with the same eigenvectors as \( \vec{A} \), \cref{thm:norm_exchange} implies the iterates satisfy a bound
\begin{equation*}
    \| \fA\vec{v} - \mathsf{alg}_k \| 
    \leq \min_{\deg(p) < k} \| \fA \vec{v} - \pA \vec{v} \|
    \leq \min_{\deg(p) < k} \| f - p \|_\Lambda \| \vec{v} \|.
\end{equation*}
In other words, the iterates satisfy a minimax bound \emph{on the eigenvalues of \( \vec{A} \)}.
As we saw in the introduction, a bound on the eigenvalues can be substantially stronger than a bound on an interval containing the eigenvalues.

A number of well-known algorithms, including the conjugate gradient (CG), minimum residual (MINRES), and quasi-minimum residual (QMR) algorithms, are standard methods for solving linear systems of equations $\vec{A}\vec{x} = \vec{v}$; i.e. for approximating $\vec{A}^{-1} \vec{v}$.
Each of these methods is \emph{optimal} for a certain norm and for certain classes of matrix \( \vec{A} \).
Moreover, such methods can be implemented in such a way that the amount of storage they use does not grow with the number of iterations \( k \). 

\iffalse
For general functions $f$ the situation is murkier. 
Presently, it is unknown whether general purpose KSMs, like the well-known Lanczos method for matrix function approximation (Lanczos-FA) \cite{saad_92}, which we discuss in \cref{chap:general_matfunc}, return near-optimal approximation to $\fA\vec{v}$ from $\mathcal{K}_k$, except in the case of particular functions such as the inverse or exponential \cite{druskin_greenbaum_knizhnerman_98}. 
% Instead, we must settle for weaker spectrum independent approximation guarantees \cite{musco_musco_sidford_18} \Chris{what's the right citation for these sorts of bounds?}. 
Work on more general spectrum dependent bounds remains somewhat ad-hoc, even for the basic case of rational functions \cite{eshof_frommer_lippert_schilling_van_der_vorst_02,ilic_turner_simpson_09,frommer_guttel_schweitzer_14,frommer_guttel_schweitzer_14a,chen_greenbaum_musco_musco_22a}.
Moreover, in terms of computational cost, to return an approximation to $\fA\vec{v}$ from $\mathcal{K}_k$ methods for general functions like Lanczos-FA require either (i) storing $k$ vectors of length $n$, or 
(ii) storing a constant number of vectors of length $n$, but increasing the number of matrix-vector products by a factor of two.
\fi

In this chapter, we describe the Lanczos method for optimal rational matrix function approximation (Lanczos-OR).
When $f$ is a rational function, Lanczos-OR outputs the \emph{optimal} (in a certain norm) approximation to $\fA\vec{v}$ from $\mathcal{K}_{k+1}$ using slightly more than $k$ matrix-vector products.
We provide a practical implementation of Lanczos-OR that
only requires storing a number of vectors of length $n$ proportional to the degree of the denominator in the rational function.
Therefore, for a fixed rational function, the storage costs do not grow with the iteration $k$. The approach used to derive this implementation of Lanczos-OR can also be used for computing the Lanczos-FA approximations to rational matrix functions, avoiding storage costs growing with $k$ in that widely used method.


Lanczos-OR is closely related to existing methods for linear systems. 
In particular, if $f = 1/(x-z)$, then, depending on the choice of $z$, the CG, MINRES, and QMR iterates can all be obtained as special cases of Lanczos-OR.
The Lanczos-OR iterate is mathematically equivalent to an optimal Galerkin projection method as described in \cite[Section 4]{lopez_simoncini_06}, provided the denominator matrix is positive definite. 
However, this method was mostly viewed as of theoretical interest since it could be used to help explain the behavior of Lanczos-FA. 
On the other hand, we show that such approximations can be computed efficiently.
Our approach is also somewhat more general in that it works with \emph{any} rational function.



\section{A bit of notation}


To simplify analysis, it will be useful to consider the recurrence that would obtained if the Lanczos algorithm were run to completion. 
In exact arithmetic, for some $K\leq n$, $\beta_{K-1} = 0$ in which case the algorithm terminates. 
Then, the final basis \( \Qhat := [ \vec{q}_0 , \ldots , \vec{q}_{K-1} ]\) and symmetric tridiagonal $\widehat{\vec{T}}$ with  diagonals $[\alpha_0, \ldots, \alpha_{K-1}]$ and off diagonals $[\beta_0, \ldots, \beta_{K-2}]$
satisfy
 a three-term recurrence 
\begin{equation}
    \label{eqn:lanczos_three_term_complete}
    \vec{A} \Qhat = \Qhat \widehat{\vec{T}}.
\end{equation}
%where \( \widehat{\vec{T}} \) is a real symmetric tridiagonal matrix with entries \Cam{I feel like we should move the picture up to where we first define $T$ and then just say what  $\hat T$ is without showing picture.} 
%\begin{equation*}
%    \widehat{\vec{T}}
 %   := \begin{bmatrix}
 %       \alpha_0 & \beta_0 \\
 %       \beta_0 & \alpha_1 & \ddots \\
 %       & \ddots & \ddots & \beta_{K-2}\\
 %       & & \beta_{K-2} & \alpha_{K-1}
 %   \end{bmatrix}.
%\end{equation*}
We emphasize that the algorithms we discuss \emph{do not} require Lanczos to be run to completion; the introduction of $\Qhat$ and $\widehat{\vec{T}}$ is for analysis purposes only.
% We often consider approximations from $\mathcal{K}_k$, and 
We note that $\vec{Q} = [\Qhat]_{:,:k}$ and $\vec{T} =[\widehat{\vec{T}}]_{:k,:k}$.
Since the columns of $\Qhat$ are orthogonal, we have that
\begin{equation*}
    \widehat{\vec{T}}
    = \Qhat^\cT \vec{A} \Qhat,
\end{equation*}
from which we easily see that, after any number of iterations $k$, $\vec{T} = \vec{Q}^\cT \vec{A} \vec{Q}$.
Note also that, for any shift $z\in\mathbb{C}$, 
\begin{equation*}
    (\vec{A} - z\vec{I}) \Qhat = \Qhat (\widehat{\vec{T}} - z\vec{I}).
\end{equation*}
In other words, the Krylov subspaces generated by $(\vec{A},\vec{v})$ and $(\vec{A}-z\vec{I},\vec{v})$ coincide, and the associated tridiagonal matrices are easily related by a diagonal shift. 
This shift invariance of Krylov subspaces is critical in a number of Krylov subspace methods.



\section{Existing algorithms}
\label{sec:linear_solvers}


It is well known that when $\vec{A}$ is positive definite, CG minimizes the $\vec{A}$-norm of the error over the Krylov subspace; i.e., the CG approximation $\mathsf{cg}_k$ is given by
\begin{equation*}
    \mathsf{cg}_k := \argmin_{\vec{x}\in\mathcal{K}_k}
    \| \vec{A}^{-1} \vec{v} - \vec{x} \|_{\vec{A}}.
\end{equation*}
%Recall that for a positive definite matrix $\vec{M}$, the norm $\|\vec{x}\|_{\vec{M}}$ denotes $\|\vec{x}\|_{\vec{M}} = \sqrt{\vec{x}^\cT\vec{M}\vec{x}}$. 
Since $\vec{Q}$ is a basis for $\mathcal{K}_k$, we can equivalently write
\label{eqn:cg_def}
\begin{equation*}
    \mathsf{cg}_k = \vec{Q} \argmin_{\vec{c}\in\R^k}
    \| \vec{A}^{1/2}(\vec{A}^{-1} \vec{v} - \vec{Q}\vec{c}) \|_2
    = \vec{Q} \argmin_{\vec{c}\in\R^k}
    \| \vec{A}^{-1/2} \vec{v} - \vec{A}^{1/2}\vec{Q}\vec{c} \|_2.
\end{equation*}
The solution to this least squares problem is
%Thus, by solving the normal equations \Anne{[This will be misleading to many people because it is the normal equations based on $\vec{A}^{1/2}$ not $A$.  Maybe just say ``The solution is ...''.]}, we arrive at the expression
\begin{equation*}
    \mathsf{cg}_k
    = \vec{Q} ( \vec{Q}^\cT \vec{A}^{1/2} \vec{A}^{1/2} \vec{Q})^{-1} \vec{Q}^\cT \vec{A}^{1/2} \vec{A}^{-1/2} \vec{v}
    = \vec{Q} \vec{T}^{-1} \vec{e}_0.
\end{equation*}
Here we have used that $\vec{Q}^\cT\vec{A}\vec{Q} = \vec{T}$ and that $\vec{Q}^\cT \vec{v} = \| \vec{v}\|_2\vec{e}_0 = \vec{e}_0$ (since we are assuming $\|\vec{v}\|_2=1$).

If $\vec{A}$ is indefinite, then the $\vec{A}$-norm of the error is not well defined and the CG iterates need not be optimal.
%\Anne{[This is not quite accurate since there is no norm defined analogously to the $A$-norm
%for positive definite $A$.  Maybe just omit this sentence?]}
A common alternative to CG for indefinite systems is MINRES, which minimizes the $\vec{A}^2$-norm of the error (the 2-norm of the residual) over the Krylov subspace; i.e., the MINRES approximation $\mathsf{mr}_k$ is given by
\begin{equation*}
    \mathsf{mr}_k
    := \argmin_{\vec{x}\in\mathcal{K}_k}
    \| \vec{A}^{-1} \vec{v} - \vec{x} \|_{\vec{A}^2}
    =
    \argmin_{\vec{x}\in\mathcal{K}_k}
    \| \vec{v} - \vec{A}\vec{x} \|_{2}
\end{equation*}
Now note that
\begin{equation*}
    \vec{A} \vec{Q} 
    = [\vec{A} \Qhat]_{:,:k}
    =[\Qhat\widehat{\vec{T}}]_{:,:k}
    =[\Qhat]_{:,:k+1} [\widehat{\vec{T}}]_{:k+1,:k}.
\end{equation*}
Therefore, since $\widehat{\vec{T}}$ is symmetric and $\Qhat$ has orthonormal columns,
%\begin{equation*}
%    \mathsf{mr}_k = \vec{Q} \argmin_{\vec{c}\in\R^k}
%    \| \vec{v} - \vec{A}\vec{Q} \vec{c} \|_{2}
%\end{equation*}
%so
\label{eqn:mr_def}
\begin{equation*}
    \mathsf{mr}_k = \vec{Q} (\vec{Q}^\cT \vec{A} \vec{A} \vec{Q})^{-1} \vec{Q}^\cT \vec{A} \vec{v}
    = \vec{Q} ([\widehat{\vec{T}}]_{:k,:k+1}[\widehat{\vec{T}}]_{:k+1,:k})^{-1} \vec{T} \vec{e}_0.
\end{equation*}

More generally,  suppose $z$ is an arbitrary complex number.
Then a special case of the quasi-minimum residual method (QMR) \cite{freund_92} can be used to compute the optimal $(\vec{A}^2+|z|^2 \vec{I})$-norm
approximation to $(\vec{A}-z \vec{I})^{-1}\vec{v}$; i.e., the QMR approximation $\mathsf{qmr}_k(z)$ is given by
\begin{align}
    \mathsf{qmr}_k(z)
    &:= \argmin_{\vec{x}\in\mathcal{K}_k}
    \| (\vec{A}-z \vec{I})^{-1} \vec{v} - \vec{x} \|_{(\vec{A}^2+|z|^2 \vec{I})}
    \nonumber\\&=
    \argmin_{\vec{x}\in\mathcal{K}_k}
    \| (\vec{A}-\overline z \vec{I})^{1/2}(\vec{A}-z \vec{I})^{-1/2}\vec{v} - (\vec{A}-\overline z \vec{I})^{1/2}(\vec{A}-z \vec{I})^{1/2}\vec{x} \|_{2}.
\end{align}
Here we have used that $\vec{A}^2+|z|^2 \vec{I} = (\vec{A}-\overline z \vec{I})(\vec{A}-z \vec{I})$.
Next, using the shift invariance of Krylov subspace, we see that
\label{eqn:qmr_def}
\begin{align*}
    \mathsf{qmr}_k(z)
    &:=  \vec{Q} (\vec{Q}^\cT (\vec{A}-\overline z \vec{I})(\vec{A}-z \vec{I}) \vec{Q})^{-1} \vec{Q}^\cT (\vec{A}- z \vec{I})\vec{v}.
    \\&=  \vec{Q} ( \vec{Q}( \vec{A}^2+|z|^2 \vec{I})\vec{Q}^\cT)^{-1} \vec{Q}^\cT (\vec{A}- \overline{z} \vec{I})\vec{v}
    \\&=  \vec{Q} ([\widehat{\vec{T}}]_{:k,:k+1}[\widehat{\vec{T}}]_{:k+1,:k}+|z|^2 \vec{I})^{-1} (\vec{T} - \overline{z} \vec{I})\vec{e}_0.
\end{align*}

%((\vec{A}-\Re(\theta)\vec{I})^2 + \Im(\theta)^2 \vec{I})

At first glance, CG, MINRES, and QMR all require the matrix \( \vec{Q} \) which is of size \( nk \).
However, by taking advantage of the tridiagonal structure of \( \widehat{\vec{T}} \), each of these algorithms can be implemented in a way which require storing just a few vectors of length $n$. 
In particular, the algorithms work by implicitly forming $\vec{L}\vec{D}\vec{L}^\cT$ factorization of \( \widehat{\vec{T}} \) \cite{paige_saunders_75,freund_92,liesen_strakos_13,simonova_tichy_21}.
The low-memory implementations for Lanczos-OR and Lanczos-FA that we derive in \cref{sec:low_mem} are based on this idea. 


\section{Optimal rational function approximation}


We now describe an optimal iterate for approximating $r(\vec{A})\vec{b}$ when $r$ is an arbitrary fixed rational function.
We will describe a low-memory implementation of this algorithm in \cref{sec:low_mem}.
Our low-memory implementation can also be used to compute the Lanczos-FA approximations to $r(\vec{A})\vec{b}$ without doubling the number of matrix-vector products.

\begin{definition}
\label{def:lanczos_or}
    Let \( r :\R\to\R \) be a rational function decomposed as \( r = M/N \), where \( N : \R\to\R \) is a polynomial with leading coefficient one and \( M:\R\to\R \) is a polynomial which does not divide \( N \).
For any polynomial \( R : \R\to\R \), define \( \tilde{M} = MR \) and \( \tilde{N} = NR \).
Then the Lanczos-OR iterate is defined as
\begin{equation*}
    \lanopt_k(r,R)
    :=  \vec{Q}  ([\tilde{N}(\widehat{\vec{T}})]_{:k,:k})^{-1} [\tilde{M}(\widehat{\vec{T}})]_{:k,:k} \vec{e}_0.
 \end{equation*}
\end{definition}

Those familiar with CG, MINRES, and the version of QMR for shifted Hermitian systems will note that these optimal algorithms are each obtained as special cases of Lanczos-OR.
Specifically, when $\vec{A}$ is positive definite, CG is obtained with \( r(x) = 1/x \) and \( R(x) = 1 \), MINRES is obtained with $r(x) = 1/x$ and \( R(x) = x \), and QMR is obtained if $r(x) = 1/(x-z)$ and \( R(x) = (x-\overline{z}) \).
In fact, we have a more general optimality result for Lanczos-OR:
\begin{theorem}
\label{thm:lanczos_opt}
    Given a rational function $r(x) = M(x)/N(x)$ as in \cref{def:lanczos_or}, choose a polynomial $R$ so that $\vec{H} = \tilde{N}(\vec{A}) = N(\vec{A})R(\vec{A})$ is positive definite.
Then $\lanopt_k(r,R)$ is the $\vec{H}$-norm optimal approximation to $r(\vec{A})\vec{v}$ from $\mathcal{K}_k$.
\end{theorem}


A simple way to ensure $\vec{H}$ is positive definite is to take \( R(x) = N(x) \) so that $\vec{H} = N(\vec{A})^2$.
However, in some situations, we may be able to get away with a lower degree choice for \( R \).
For instance, in the case of symmetric linear systems, while one can always use MINRES ($r(x)=1/xi$, $R(x)=x$), if \( \vec{A} \) is positive definite, then one may hope to use CG ($r(x)=1/x$, $R(x)=1$).
A simple way to obtain a lower degree choice of \( R \) is to only take the terms in \( N \) which are indefinite.
\begin{definition}
    \label{def:Rstar}
    Given a rational function $r = M/N$ as in \cref{def:lanczos_or}, factor 
    \begin{equation*}
        N(x) = \bigg( {\prod_{i=0}^{d_1-1}} (x-z_i) \bigg) \bigg( {\prod_{i=0}^{d_2-1}} (x-{z_i'})(x-\overline{z_i'}) \bigg)
    \end{equation*}
    where $z_i\neq \overline{z}_j$ for all $i,j = 0,1,\ldots, d_1-1$ with $j \neq i$.
    Then \( R^* \) is defined by
    \begin{equation*}
        R^*(x) = \xi {\prod_{i=0}^{d_1-1}} (x-\overline{z}_i)^{\alpha_i}   
    \end{equation*}    
    where, for \( i=0,1,\ldots,q-1 \), $\alpha_i = 0$ if $z_i\in \mathbb{R} \setminus \mathcal{I}$  and $\alpha_i = 1$ otherwise and $\xi \in \{ \pm 1 \}$ is chosen so that $R^*(\lambda_{\textup{min}})N(\lambda_{\textup{min}}) \geq 0$.  
\end{definition}

\begin{lemma}
    Given a rational function $r = M/N$ as in \cref{def:lanczos_or}, choose \( R^* \) as in \cref{def:Rstar}. 
    Then \( \vec{H} = N(\vec{A}) R^*(\vec{A}) \) is positive definite.
\end{lemma}

\begin{proof}
    Clearly $\prod_{i=0}^{q'-1} (x-{z'}_i)(x-\overline{z'}_i) \geq 0$ for all $x \in \mathbb{R}$, and for $z_i \in \mathbb{R} \setminus \mathcal{I}$, $(x-z_i)$ does not change signs over $\mathcal{I}$. 
The choice of $\xi$ ensures that $\tilde{N}(\lambda_{\textup{min}})>0$, and by assumption $N(\lambda) \neq 0$ for all $\lambda \in \Lambda$ so that that $\tilde{N}(\lambda) \neq 0$ for all $\lambda\in\Lambda$. 
It follows that $\tilde{N}(\lambda) > 0$ for all $\lambda\in\Lambda$; i.e. that $\vec{H} = \tilde{N}(\vec{A})$ is positive definite.
\end{proof}



\begin{proof}[Proof of \cref{thm:lanczos_opt}]
Since $\vec{y}\in\mathcal{K}_k$, we have $\vec{y} = \vec{Q} \vec{c}$ for some vector $\vec{c}$.
Thus, we can consider the problem
\begin{equation*}
    \argmin_{\vec{c}\in\R^k} \| r(\vec{A})\vec{v} - \vec{Q} \vec{c} \|_{\vec{H}}
    =
    \argmin_{\vec{c}\in\R^k} \| \vec{H}^{1/2} r(\vec{A})\vec{v} - \vec{H}^{1/2} \vec{Q} \vec{c} \|_2.
\end{equation*}
But this is just a standard least squares problem which has solution
\begin{equation*}
    \vec{c} = ( (\vec{H}^{1/2} \vec{Q})^\cT (\vec{H}^{1/2} \vec{Q}) )^{-1} (\vec{H}^{1/2} \vec{Q})^\cT (\vec{H}^{1/2} r\A \vec{v}).
\end{equation*}
Thus, we see that the optimal iterate has the form
\begin{equation*}
    \vec{Q}( \vec{Q}^\cT \vec{H} \vec{Q})^{-1} \vec{Q}^\cT \vec{H} r(\vec{A}) \vec{v}.
\end{equation*}
By definition, $\vec{H} = N(\vec{A})R(\vec{A}) $ so $\vec{H}r(\vec{A}) = M(\vec{A})R(\vec{A}) = \tilde{M}(\vec{A})$. Thus,
%By \cref{thm:lanczos_exact_poly}, if $k > \deg(\tilde{M})$, then
\begin{equation*}
    \vec{Q}^\cT \vec{H} r(\vec{A}) \vec{v}
    = \vec{Q}^\cT M(\vec{A}) \vec{v}
    = \vec{Q}^\cT \Qhat\tilde{M}(\widehat{\vec{T}}) \Qhat^\cT \vec{v}
    = [\tilde{M}(\widehat{\vec{T}})]_{:k,:k} \vec{e}_0.
\end{equation*}
Next, since $\vec Q$ consists of the first $k$ columns $\Qhat$,
\begin{equation*}
    \vec{Q}^\cT \vec{A}^q \vec{Q} = [\Qhat^\cT \vec{A}^q \Qhat]_{:k,:k}  
    % = (\widehat{\vec{I}}_k)^\cT \Qhat^\cT \vec{A}^q \Qhat \widehat{\vec{I}}_k
    % = (\widehat{\vec{I}}_k)^\cT  \widehat{\vec{T}}^q \widehat{\vec{I}}_k
    = [ \widehat{\vec{T}}^{q} ]_{:k,:k}
\end{equation*}
so since $\vec{H}$ is a linear combination of powers of $\vec{A}$ we obtain
\begin{equation*}
    \vec{Q}^\cT \vec{H} \vec{Q}
    = \vec{Q}^\cT \tilde{N}(\vec{A}) \vec{Q}
    =  [\tilde{N}(\widehat{\vec{T}}) ]_{:k,:k}.
\end{equation*}
The result follows by combining and rearranging the above expressions. 
\end{proof}

As we noted at the beginning of this chapter, the optimality of Lanczos-OR implies an a priori scalar polynomial error bound on the eigenvalues of $\vec{A}$ analogous to the well known minimax bounds for CG, MINRES, and QMR \cite{greenbaum_97}.
\begin{theorem}
Given a rational function $r = M/N$ as in \cref{def:lanczos_or} and a polynomial \( R \) so that $\vec{H} = \tilde{N}\A = N\A R\A$ is positive definite, 
\begin{equation*}
    \frac{\| r(\vec{A}) \vec{v} - \lanopt_k(r,R) \|_{\vec{H}}}{\| \vec{v} \|_{\vec{H}}}
    \leq \min_{\deg(p) < k} \max_{\lambda \in \Lambda} | r(\lambda) - p(\lambda) |.
\end{equation*}
\end{theorem}

\begin{proof}
Since $\lanopt_k(r,R)$ is the $\vec{H}$-norm optimal approximation over the Krylov subspace, we have
\begin{equation*}
    \| r(\vec{A})\vec{v} - \lanopt_k(r,R) \|_{\vec{H}}
    = \min_{\vec{x}\in \mathcal{K}_k} \| r(\vec{A})\vec{v} - \vec{x} \|_{\vec{H}}
    = \min_{\deg(p)<k} \| r(\vec{A})\vec{v} - p(\vec{A})\vec{v} \|_{\vec{H}}
\end{equation*}
Next, using the fact that $\vec{A}$ and $\vec{H}^{1/2}$ commute, we note that,
\begin{equation*}
    \| r(\vec{A})\vec{v} - p(\vec{A})\vec{v} \|_{\vec{H}}
    = \| (r(\vec{A})- p(\vec{A}))\vec{H}^{1/2}\vec{v} \|
    \leq \| r(\vec{A})- p(\vec{A}) \| \| \vec{v} \|_{\vec{H}}.
\end{equation*}
Finally, using the definition of the spectral norm,
\begin{equation*}
    \| r(\vec{A})- p(\vec{A}) \|
    = \| (r - p)(\vec{A}) \|
    = \max_{\lambda \in \Lambda} | r(\lambda)-p(\lambda) |.
\end{equation*}
The result follows.
\end{proof}


It is not yet apparent that the iterate can be computed efficiently. 
Indeed the expression involves the terms  $[\tilde{N}(\widehat{\vec{T}})]_{:k,:k}$ and $[\tilde{M}(\widehat{\vec{T}})]_{:k,:k}$, and computing $\widehat{\vec{T}}$ requires running Lanczos to completion.
However, since \( \widehat{T} \) is tridiagonal, these terms can be computed without much more information than is in $\vec{T}$ and the Lanczos-OR iterate can be computed with only a few more matrix vector products than required to compute the Lanczos-FA iterate.

\begin{lemma}
\label{thm:poly_That}
Suppose $p$ is a polynomial with $q := \deg(p) > 0$. 
Then $[p(\widehat{\vec{T}})]_{:k,:k}$ can be computed using the coefficients generated by $k+\lfloor (q-1)/2 \rfloor$ iterations of Lanczos.
Moreover, if $k':=k+\lfloor q/2 \rfloor$, then
\begin{equation*}
    [p(\widehat{\vec{T}})]_{:k,:k} = [p([\widehat{\vec{T}}]_{:k',:k'})]_{:k,:k}.
\end{equation*}
\end{lemma}

\begin{proof}
    Note that after \( k+d \) iterations of Lanczos, one obtains \( [\widehat{\vec{T}}]_{:k+d+1,:k+d} \).
    The result then follows immediately as a special case of \cref{thm:tridiag_power_dependence}.
\end{proof}

\iffalse
\begin{proof}%[Proof of \cref{thm:poly_That}]
It suffices to consider the case $p = x^q$.
Let $\widehat{\vec{I}}_\ell$ be the $K\times \ell$ matrix with the $\ell\times \ell$ identity in the top $\ell$ rows and zeros in the bottom $K-\ell$ rows.
Then, we have that 
\begin{equation*}
    \widehat{\vec{T}} \widehat{\vec{I}}_\ell
    = [\widehat{\vec{T}}]_{:,:\ell}
    = \widehat{\vec{I}}_{\ell+1} [\widehat{\vec{T}}]_{:\ell+1,:\ell}.
\end{equation*}
Repeatedly applying this relation,
\begin{equation*}
    \widehat{\vec{T}}^j \widehat{\vec{I}}_k
    = 
    \widehat{\vec{I}}_{k+j}
    [\widehat{\vec{T}}]_{:k+j,k+j-1}
    \cdots
    [\widehat{\vec{T}}]_{:k+2,k+1}
    [\widehat{\vec{T}}]_{:k+1,k}
    = \widehat{\vec{I}}_{k+j} \vec{B}(k+j,k)
\end{equation*}
where we have defined
\begin{equation*}
    \vec{B}(k+j,k) := [\widehat{\vec{T}}]_{:k+j,k+j-1}
    \cdots
    [\widehat{\vec{T}}]_{:k+2,k+1}
    [\widehat{\vec{T}}]_{:k+1,k}.
\end{equation*}
Therefore, since $(\widehat{\vec{I}}_{k+j})^\cT \widehat{\vec{I}}_{k+j}$ is the $(k+j)\times (k+j)$ identity,
\begin{equation*}
    [\widehat{\vec{T}}^{2j}]_{:k,:k}
    = \vec{B}(k+j,k)^\cT \vec{B}(k+j,k)
\end{equation*}
and, since $(\widehat{\vec{I}}_{k+j-1})^\cT \widehat{\vec{T}} \widehat{\vec{I}}_{k+j-1} = [\vec{T}]_{:k+j-1,:k+j-1}$,
\begin{equation*}
    [\widehat{\vec{T}}^{2j-1}]_{:k,:k}
    = \vec{B}(k+j-1,k)^\cT  [\widehat{\vec{T}}]_{:k+j-1,:k+j-1} \vec{B}(k+j-1,k).
\end{equation*}
In both cases, we see that these expressions depend only on $[\widehat{\vec{T}}]_{:k+j,:k+j-1}$ which can be obtained using $k+j-1$ matrix-vector products.
The result follows by noting that $\lfloor (2j-1)/2 \rfloor(x) = \lfloor(2j-1-1)/2 \rfloor(x) = j-1$.

To complete the lemma, note that for $\ell+1\leq k'$,
\begin{equation*}
    [\widehat{\vec{T}}]_{:k',:k'} \widetilde{\vec{I}}_\ell
    = [\widehat{\vec{T}}]_{:k',:\ell}
    = \widetilde{\vec{I}}_{\ell+1} [\widehat{\vec{T}}]_{:\ell+1,:\ell}
\end{equation*}
where $\widetilde{\vec{I}}_\ell$ is defined like $\widehat{\vec{I}}_\ell$ but is of size $k' \times \ell$.
The same argument as above then implies that
\begin{equation*}
    ([\widehat{\vec{T}}]_{:k',:k'})^j \widetilde{\vec{I}}_k
    =\widetilde{\vec{I}}_{k+j}\vec{B}(k+j,k)
\end{equation*}
provided that \( k+j \leq k' \).
We therefore have that
\begin{equation*}
    [\widehat{\vec{T}}^{q}]_{:k,:k} = [([\widehat{\vec{T}}]_{:k',:k'})^{q}]_{:k,:k}. 
\end{equation*}
\end{proof}
\fi

\subsection{Relation of Lanczos-OR to QMR on inverse quadratics}

The Lanczos-OR approximation to \( r(x) = 1/(x^2+|z|^2) \) is closely related to the approximation of two linear systems by QMR. 
\begin{lemma}
Suppose $z\in\R$ and define $r= 1/(x^2+|z|^2)$, $R^\pm = x\pm \ii z$, and $r^\pm = 1/R^\pm$.
Then, for all \( k\geq 1 \),
\begin{equation*}
    \lanopt_k(r,1)
    %= \frac{1}{2\ii z} \left( \lanopt_k(r^-,R^+) - \lanopt_k(r^+,R^-) \right).
    = \frac{1}{2\ii z} \left( \mathsf{qmr}_k(-\ii z) - \mathsf{qmr}_k(\ii z) \right).
\end{equation*}
\end{lemma}

\begin{proof}
Observe that
\begin{equation*}
    \lanopt_k(r^\pm,R^\mp) = \vec{Q}([\widehat{\vec{T}}^2+|z|^2\vec{I}]_{:k,:k})^{-1} [\vec{T} \mp \ii z \vec{I}]_{:k,:k} \vec{e}_0
\end{equation*}
so that 
\begin{equation*}
    \mathsf{qmr}_k(-\ii z) - \mathsf{qmr}_k(\ii z) 
    = 2 \ii z  \vec{Q}([\widehat{\vec{T}}^2+|z|^2\vec{I}]_{:k,:k})^{-1} \vec{e}_0
    = 2 \ii z \lanopt_k(r,1).
\end{equation*}
The result is then obtained by rearranging the above expression.
\end{proof}

Whether it is better to use Lanczos-OR with $r$ and $R=1$ or with $r^\pm$ and $R^\pm$ (i.e. QMR) is somewhat unclear.
The Lanczos-OR based approach avoids the need for complex arithmetic, which simplifies implementation slightly. 
However, since QMR has been studied longer, it is likely to have more practical low-memory implementations.


\section{Error estimates for Lanczos-OR}
\label{sec:error_OR}

We now describe an approach for estimating the Lanczos-OR error.
This approach has been widely studied for estimating the \( \vec{A} \)-norm of the error in CG, and we refer to \cite{strakos_tichy_02,meurant_tichy_18,estrin_orban_saunders_19,meurant_papez_tichy_21} and the references within for more details.
Note that several of these works also study whether such estimates are still reasonable in finite precision arithmetic as well as how to derive estimates for other norms such as the 2-norm.


\begin{theorem}
    \label{thm:OR_error}
Let \( r \) be a rational function and \( R \) a polynomial.
Write \( r(x) = M/N \) where \( N \) has leading coefficient one and \( M \) does not divide \( N \) and define \( \tilde{M} = MR \) and \( \tilde{N} = N R \).
Suppose \( \vec{H} = \tilde{N}\A \) is positive definite. 
Then the Lanczos-OR iterates satisfy,
\begin{equation*}
    \| r\A\vec{v} - \lanopt_k(r,R) \|_{\vec{H}}^2
    = \sum_{i=k}^{n-1}
    \| \lanopt_{i}(r,R) - \lanopt_{i+1}(r,R) \|_{\vec{H}}^2.
\end{equation*}

\end{theorem}

\begin{proof}
Let \( \{ \vec{p}_i \}_{i=0}^{n-1} \) be an \( \vec{H} \)-orthonormal set satisfying
\begin{equation*}
    \operatorname{span}\{\vec{p}_0, \ldots, \vec{p}_{i} \} 
    = \mathcal{K}_{i+1} 
\end{equation*}
for all \( i=0,1, \ldots, n-1 \), and decompose \( r\A\vec{v} \) as 
\begin{equation*}
    r\A\vec{v} = \sum_{i=0}^{n-1} c_i \vec{p}_i.
\end{equation*}
\Cref{thm:lanczos_opt} asserts that \( \lanopt_k(r,R) \) is the \( \vec{H} \)-norm optimal approximation to \( r\A\vec{v} \) from \( \mathcal{K}_k \).
Thus, for all \( j=0,1,\ldots, n-1 \), 
\begin{equation*}
    \lanopt_j(r,R) = \sum_{i=0}^{j-1} c_i \vec{p}_i.
\end{equation*}
This implies that
\begin{equation*}
    r\A\vec{v} - \lanopt_k(r,R) 
    = r\A\vec{v} - \sum_{i=0}^{k-1} c_i \vec{p}_i
    = \sum_{i=k}^{n-1} c_i \vec{p}_i.
\end{equation*}
    so that, by the \( \vec{H} \)-orthonormality of the \( \{ \vec{p}_i \}_{i=0}^{n-1} \),
\begin{equation*}
    \| r\A\vec{v} - \lanopt_k(r,R) \|_{\vec{H}}^2
    = \sum_{i=k}^{n-1} c_i^2.
\end{equation*}
    But we also have
    \begin{equation*}
        \lanopt_{i}(r,R) - \lanopt_{i+1}(r,R)
        = - c_i \vec{p}_i
    \end{equation*}
    so that
    \begin{equation*}
        \| \lanopt_{i}(r,R) - \lanopt_{i+1}(r,R) \|_{\vec{H}}^2 = c_i^2. \qedhere
    \end{equation*}
\end{proof}

Note that \cref{thm:OR_error} implies that
\begin{equation}
    \label{eqn:OR_est}
    \| r\A\vec{v} - \lanopt_k(r,R) \|_{\vec{H}}^2
    \geq \smop{\sum_{i=k}^{k+d-1}}
    \| \lanopt_{i}(r,R) - \lanopt_{i+1}(r,R) \|_{\vec{H}}^2.
\end{equation}
While this is a lower bound, if we assume that
\begin{equation*}
    \| r\A\vec{v} - \lanopt_{k+d}(r,R) \|_{\vec{H}}^2 = 
    \smop{\sum_{i=k+d}^{n-1}}
    \| \lanopt_{i}(r,R) - \lanopt_{i+1}(r,R) \|_{\vec{H}}^2
\end{equation*}
is negligible compared to \( \| r\A\vec{v} - \lanopt_k(r,R) \|_{\vec{H}}^2 \), then \cref{eqn:OR_est} becomes an approximate equality.

Typically $d$ can be taken as a small constant, say $d=5$, so the extra work required to obtain these estimate is not too large.

\begin{remark}
As shown in \cite{eshghi_reichel_21}, a similar approach can be used for general matrix functions. In particular, 
\begin{equation*}
    \| r\A\vec{v} - \mathsf{alg}_k \|
    \leq \| r\A\vec{v} - \mathsf{alg}_{k+d} \| + \| \mathsf{alg}_{k+d} - \mathsf{alg}_k \|
    \approx \| \mathsf{alg}_{k+d} - \mathsf{alg}_k \|
\end{equation*}
provided that \( \| r\A\vec{v} - \mathsf{alg}_{k+d} \| \ll \| r\A\vec{v} - \mathsf{alg}_k \| \).
    Note, however, that in situations where the convergence of \( \mathsf{alg}_k \) is oscillatory, it may be hard to guarantee \( \| r\A\vec{v} - \mathsf{alg}_{k+d} \| \ll \| r\A\vec{v} - \mathsf{alg}_k \| \), even if \( d \) is large.
\end{remark}



\subsection{Numerical experiment}

We choose \( \vec{A} \) with eigenvalues from the model problem \cref{eqn:model_problem} with \( n=300 \), \( \rho = 0.8 \), and \( \kappa = 1000 \) and \( \vec{b} \) with equal projection onto each eigencomponent.
We set \( r(x) = 1/(x^2+1) \), and run Lanczos-OR using \( R=1 \) with and without reorthgonalization.
In each case, we compute \cref{eqn:OR_est} with \( d=4 \).

The resulting estimates, shown in \cref{fig:ch5_OR_error_est}, are accurate for most iterations, with larger error in initial iterations where the true Lanczos-OR error is not decreasing as quickly as in later iterations.
Interestingly, the bound seems to work well in finite precision arithmetic.
Understanding this further is of interest.
In particular, a unified analysis of Lanczos-OR could provide more information about MINRES, for which existing bounds in finite precision arithmetic are somewhat weaker than CG.


\begin{figure}[htb]
    \begin{center}
        \includegraphics{imgs/ch5_OR_error_est.pdf} 
    \end{center}
    \caption[{Error estimates for Lanczos-OR for \( r(x) = 1/(x^2+1) \) and \( R(x) = 1 \).}]{%
    Error estimates for Lanczos-OR for \( r(x) = 1/(x^2+1) \) and \( R(x) = 1 \).
    \hspace{.25em}\emph{Legend}:
    Lanczos-OR error with reorthogonalization 
    ({\protect\raisebox{0mm}{\protect\includegraphics[]{imgs/legend/l3.pdf}}}) 
    and corresponding estimate \cref{eqn:OR_est} with \( d=4 \)
    ({\protect\raisebox{0mm}{\protect\includegraphics[]{imgs/legend/l4.pdf}}}).
    Lanczos-OR error without reorthogonalization 
    ({\protect\raisebox{0mm}{\protect\includegraphics[]{imgs/legend/l1.pdf}}}) 
    and corresponding estimate \cref{eqn:OR_est} with \( d=4 \)
    ({\protect\raisebox{0mm}{\protect\includegraphics[]{imgs/legend/l2.pdf}}}).
    \hspace{.25em}\emph{Takeaway}: The error estimates are remarkably accurate, even in finite precision arithmetic. 
    }
    \label{fig:ch5_OR_error_est}
\end{figure}



\section{Implementing Lanczos-OR using low memory}
\label{sec:low_mem}


We now describe a low-memory implementation of Lanczos-OR which is similar in spirit to CG, MINRES, and QMR. %In  particular, it is inspired by the LDL based version of CG described in \cite{simonova_tichy_21}.

For convenience, we will denote $\vec{M} := [\tilde{M}\mf{\widehat{\vec{T}}}]_{:k,:k}$ and $\vec{N} := [\tilde{N}\mf{\widehat{\vec{T}}}]_{:k,:k}$ so that the Lanczos-OR output is given by $\vec{Q} \vec{N}^{-1} \vec{M} \vec{e}_0$.
Then, at a high level, our approach is to:
\begin{itemize}
    \item Take one iteration of Lanczos to generate one more column of $\Qhat$ and $\widehat{\vec{T}}$
    \item Compute one more column of each of $\vec{M}$ and $\vec{N}$
    \item Compute one more factor of $\vec{L}^{-1} = \vec{L}_{k-1}\cdots \vec{L}_1 \vec{L}_0$ and one more entry of $\vec{D}$ where $\vec{L}$ and $\vec{D}$ are defined by the LDL factorization $\vec{N} =\vec{L}\vec{D}\vec{L}^\cT$
    \item Compute one more term of the sum:
    \begin{equation*}
        \vec{Q}\vec{N}^{-1}\vec{M}\vec{e}_0
        = \vec{Q} \vec{L}^{-1} \vec{D}^{-1} \vec{L}^{-\cT} \vec{M} \vec{e}_0 
        = \smop{\sum_{i=0}^{k-1}} \frac{[\vec{L}^{-\cT}\vec{M}\vec{e}_0]_i}{[\vec{D}]_{i,i}} [\vec{Q}\vec{L}^{-1}]_{:,i}
    \end{equation*}
\end{itemize}



There are two critical observations which must be made in order to see that this gives a memory-efficient implementation.
The first is that, since $\widehat{\vec{T}}$ is tridiagonal, $\vec{M}$, $\vec{N}$, and therefore $\vec{L}$ are all of half-bandwidth $q := \max(\deg(\tilde{M}),\deg(\tilde{N}))$.
This means that it is possible to compute the entries of $\vec{D}$ and the factors of $\vec{L}^{-1} = \vec{L}_{k-1}\cdots \vec{L}_1 \vec{L}_0$ one by one as we get the entries of $\widehat{\vec{T}}$.
The second is that because $\vec{L}$ is of bandwidth $q$, we can compute $[\vec{Q}\vec{L}^{-1}]_{:,i}$ without saving all of $\vec{Q}$. 
More specifically, $[\vec{L}^{-1} \vec{M} \vec{e}_0]_i$ and $[\vec{Q} \vec{L}^{-1}]_{:,i}$ can respectively be computed from $\vec{L}_{j-1}\cdots \vec{L}_1 \vec{L}_0 \vec{M} \vec{e}_0$ and $\vec{Q} \vec{L}_0^\cT \vec{L}_1^\cT \cdots \vec{L}_{k-1}^\cT$ and can therefore be maintained iteratively as the factors of $\vec{L}^{-1}$ are computed. 
Moreover, because of the banded structure of the factors $\vec{L}_i$, we need only maintain a sliding window of the columns of $\vec{Q}\vec{L}^{-1}$ which will allow us to access the relevant columns when we need them and discard them afterwards.

For clarity of exposition, we only describe how to compute $\vec{M}$ and $\vec{N}$ in the case that $\tilde{M}$ and $\tilde{N}$ are degree at most two.
The rest of the subroutines are fully described for any degree.
The syntax we use follows Python and other object oriented languages closely.

%, the intricacies of stably computing matrix polynomials in a manner compatible with our streaming approach, although interesting, is beyond the scope of this paper.
%This unlikely to be a major practical concern.
%Indeed, as we discussed in \cref{sec:rational_func}, when using a rational function obtained by discretizing an integral representation via a numerical quadrature, it is seemingly more natural to compute the term-wise optimal approximation.


\subsection{Computing LDL factorization}

For the time being, we will assume that we can sequentially access the rows of $\vec{M}$ and $\vec{N}$.
Our first step is to compute an LDL factorization of $\vec{N}$.
A LDL factorization can be computed via a symmetrized version of Gaussian elimination and is guaranteed to exist if \( \vec{N} \) is positive definite \cite{higham_02}.
Gaussian elimination can be viewed as transforming the starting matrix \( \vec{N}_0 = \vec{N} \) to a diagonal matrix \( \vec{N}_{k-1} = \vec{D} \) via a sequence of row and column operations
\begin{equation*}
    \vec{N}_{i+1} = \vec{L}_i \vec{N}_i \vec{L}_i^\cT
\end{equation*}
where
\begin{equation*}
    \vec{L}_i := 
    \vec{I}_k + \vec{l}_i\vec{e}_i^\cT
    ,\qquad
    \vec{l}_i := 
        \bigg[ \underbrace{\vphantom{\bigg|}0 , \cdots , 0}_{i+1\text{ zeros}} , -\frac{[\vec{N}_{i}]_{i+1,i}}{[\vec{N}_i]_{i,i}}  , \cdots , -\frac{[\vec{N}_i]_{k-1,i}}{[\vec{N}_i]_{i,i}}  \bigg]^\cT.
\end{equation*}
Note that the entries of \( \vec{L}_i \) are chosen to introduce zeros to the \( i \)-th row and column of \( \vec{N}_i \) such that \( [\vec{N}_{i+1}]_{:i+1,:i+1} \) is diagonal.
Therefore, if the algorithm terminates successfully, we will have obtained a factorization 
\begin{equation*}
    \vec{D} = (\vec{L}_{k-1} \cdots \vec{L}_1 \vec{L}_0) \vec{N} (\vec{L}_0^\cT \vec{L}_1^\cT \cdots \vec{L}_{n-1}^\cT)
\end{equation*}
where \( \vec{D} \) is diagonal and each \( \vec{L}_i \) is unit lower triangular.
To obtain the factorization \( \vec{N} = \vec{L} \vec{D} \vec{L}^\cT \), simply define  \( \vec{L} := (\vec{L}_{k-1} \cdots \vec{L}_1 \vec{L}_0)^{-1} \) and note that 
\begin{equation*}
    \vec{L} = \vec{I}_k - \smop{\sum_{i=0}^{k-1}} \vec{l}_i \vec{e}_i^\cT.
\end{equation*}
We remark that that $\vec{l}_{k-1}$ is the zeros vector and is only included in sums for ease of indexing later on. 
For further details on LDL factorizations, we refer readers to \cite{higham_02}.

To implement a LDL factorization, observe that the procedure above defines a recurrence
\begin{align*}
    [\vec{D}]_{j,j} &= [\vec{N}]_{j,j} - \smop{\sum_{\ell=0}^{j-1}} [\vec{L}_{j,\ell}]^2 [\vec{D}]_{\ell,\ell}
    \\
    [\vec{L}]_{i,j} &= \frac{1}{[\vec{D}]_{j,j}} \left( [\vec{N}]_{i,j} - \smop{\sum_{\ell=0}^{j-1}}[\vec{L}]_{j,\ell} [\vec{L}]_{i,\ell} [\vec{D}]_{\ell,\ell} \right)
    ,\qquad i > j.
\end{align*}
We therefore have \cref{alg:LDL}.

\begin{labelalgorithm}[htb]{LDL}{LDL}{LDL factorization}
\begin{algorithmic}[1]
    \Procedure{\thealgorithmname}{$\vec{N}$}
    \For{\(j=0,1,\ldots, k-1\)}
    \State \( [\vec{D}]_{j,j} = [\vec{N}]_{j,j} - \sum_{\ell=0}^{j-1} [\vec{L}_{j,\ell}]^2 [\vec{D}]_{\ell,\ell} \)
    \State \( [\vec{L}]_{j,j} = 1 \)
    \For{\(i=j+1, j+2, \ldots, k-1\)}
    \State \( [\vec{L}]_{i,j} = (1/[\vec{D}]_{j,j})( [\vec{N}]_{i,j} - \sum_{\ell=0}^{j-1}[\vec{L}]_{j,\ell} [\vec{L}]_{i,\ell} [\vec{D}]_{\ell,\ell} ) \)
    \EndFor
    \EndFor
    \State \Return \( \vec{L}, \vec{D} \)
\EndProcedure
\end{algorithmic}
\end{labelalgorithm}

\subsubsection{Streaming version}

%\Cref{thm:banded_LDL}
The fact that $\vec{L}$ has the same half bandwidth as $\vec{N}$ allows a more efficient implementation of \cref{alg:LDL} where terms which are known to be zero are not computed and only the important diagonals of $\vec{L}$ are stored.
Moreover, \cref{alg:LDL} already only accesses $\vec{N}$ one column at a time so it is easily converted to a streaming algorithm.
Making these changes gives the implementation \cref{alg:streaming_LDL} which is fed a stream of the columns of $\vec{N}$ in order, as shown in \cref{fig:streaming_LDL}.
Here the diagonal of $\vec{D}$ is stored as $\ms{d}$ and the $(j+1)$-st diagonal of $\vec{L}$ is stored as $[\ms{L}]_{j,:}$.
Thus, $\vec{L}_{i,j} = [\ms{L}]_{i-j-1,j}$ as long as $i-j \in 0:q+1$.


   
\begin{labelalgorithm}[htb]{streaming_LDL}{streaming-LDL}{Streaming LDL factorization}
\begin{algorithmic}[1]
\Class{\thealgorithmname}{$q,k$}
\State{stream: $[\vec{N}]_{0,0:q+1}, [\vec{N}]_{1,1:q+2}, \ldots, [\vec{N}]_{k-1,k-1:q+k-1}$}
\State $\ms{L} = \textsc{zeros}(q,k)$
\State $\ms{d} = \textsc{zeros}(k)$
\State $\ms{j}\gets0$
\Procedure{read-stream}{$\vec{n}$}
    \State $[\ms{d}]_{\ms{j}} \gets [\vec{n}]_0 - \sum_{\ell=\max(0,\ms{j}-q)}^{\ms{j}-1} [\ms{L}]_{\ms{j}-\ell-1,\ell}^2 [\ms{d}]_{\ell}$
%\State $[\ms{L}]_{0,\ms{j}} \gets 1$
    \For{$i=\ms{j}+1, \ms{j}+2, \ldots,  \min(\ms{j}-q,n-1)$}
    \State $[\ms{L}]_{i-\ms{j}-1,\ms{j}} \gets (1/[\ms{d}]_{\ms{j}})([\vec{n}]_{i-\ms{j}} - \sum_{\ell=\max(0,i-q)}^{i-1} [\ms{L}]_{i-\ell-1,\ell} [\ms{L}]_{\ms{j}-\ell-1,\ell} [\ms{d}]_{\ell} )$
\EndFor
\State $\ms{j}\gets \ms{j}+1$
\EndProcedure
\EndClass
\end{algorithmic}
\end{labelalgorithm}

\begin{figure}
    \centering
    \begin{subfigure}{.48\textwidth}\centering
    \includegraphics[height=4.5cm]{imgs/ch5_streaming_LDL.pdf}
    \caption{Pattern for $\vec{N}$ in \cref{alg:streaming_LDL}.}
    \label{fig:streaming_LDL}
    \end{subfigure}
    \hfill
    \begin{subfigure}{.48\textwidth}\centering
    \includegraphics[height=4.5cm]{imgs/ch5_streaming_tridiag_square.pdf}
    \caption{Pattern for $\vec{T}$ in \cref{alg:streaming_tridiag_square}.}
    \label{fig:streaming_tridiag_square}
    \end{subfigure}
    \\
    \begin{subfigure}{\textwidth}\centering
    \includegraphics[height=4.5cm]{imgs/ch5_streaming_binv_Q.pdf}
    \hfill   
    \includegraphics[height=4.5cm]{imgs/ch5_streaming_binv_L.pdf}
    \hfill    
    \includegraphics[height=4.5cm]{imgs/ch5_streaming_binv_d.pdf}
    \hfill
    \includegraphics[height=4.5cm]{imgs/ch5_streaming_binv_M.pdf}
    \caption{Pattern for $\vec{Q}$, $\vec{L}$, $\vec{d}$, $\vec{M}$ in \cref{alg:streaming_banded_prod}.}
    \label{fig:streaming_banded_prod}
    \end{subfigure}
    \caption{Access patterns for inputs to streaming functions used in low-memory implementations of Lanczos-OR and Lanczos-FA.
    Indices indicate what information should be streamed into the algorithm at the given iteration.}
%    \label{}
\end{figure}


\begin{labelalgorithm}[htb]{streaming_banded_prod}{streaming-banded-prod}{Streaming banded product}
\begin{algorithmic}[1]
\Class{\thealgorithmname}{$n,k,q$}
\State{stream:$ $}
\State $\ms{X\_} \gets \textsc{zeros}(n,q+1)$
\State $\ms{y\_} \gets \textsc{zeros}(q+1)$
\State $\ms{out} \gets \textsc{zeros}(n)$
\State $\ms{j}\gets-1$
\Procedure{read-stream}{$\vec{v},\vec{l},d,\vec{y}_0$}
\If{$\ms{j}=-1$}
\State $[\ms{X\_}]_{:,:q} = \vec{v}$
\Else
\If{$\ms{i}=-1$}
\State $\ms{y\_} \gets \vec{y}_0$
\EndIf
\State $\ms{out} \gets \ms{out} + ([\ms{y\_}]_{0}/d)[\ms{X\_}]_{:,0}$
\State $[\ms{y\_}]_{:q} \gets [\ms{y\_}]_{1:} - [\ms{y\_}]_{0} \vec{l}$
\State $[\ms{y\_}]_{-1} \gets 0$
\State $[\ms{X\_}]_{:,-1} \gets \vec{v}$
\State $[\ms{X\_}]_{:,:q} \gets [\ms{X\_}]_{:,1:} + [\ms{X\_}]_{:,0} \vec{l}^\cT$
\EndIf
\State $\ms{j}\gets \ms{j}+1$
\EndProcedure
\EndClass
\end{algorithmic}
\end{labelalgorithm}


\subsection{Inverting the LDL factorization}

Once we have computed a a factorization \( \vec{N} = \vec{L}\vec{D}\vec{L}^\cT \), we can easily evaluate $\vec{Q}\vec{L}^{-1}\vec{D}^{-1} \vec{L}^{-\cT}\vec{M}\vec{e}_1$ using the fact that \( \vec{L}^{-1} = \vec{L}_{k-1} \cdots \vec{L}_1 \vec{L}_0 \).
Moreover, because the \( \vec{L}_j \) can be computed one at a time, there is hope that we can derive a memory efficient implementation.

Towards this end, define \( \vec{y}_j := \vec{L}_{j-1} \cdots \vec{L}_1 \vec{L}_0 \vec{M}\vec{e}_0 \) and \( \vec{X}_j := \vec{Q} \vec{L}_0^\cT \vec{L}_1^\cT \cdots \vec{L}_{j-1}^\cT  \).
Then, setting \( \vec{y}_{0} = \vec{M} \vec{e}_1 \) we have that
\begin{equation*}
    \vec{y}_{j+1} 
    = \vec{L}_{j} \vec{y}_{j}
    = (\vec{I} + \vec{l}_j \vec{e}_j^\cT ) \vec{y}_{j}
    = \vec{y}_{j} +  (\vec{e}_{j}^\cT \vec{y}_{j}) \vec{l}_j.
\end{equation*}
Similarly, setting \( \vec{X}_{0} = \vec{Q} \) we have that
\begin{equation*}
    \vec{X}_{j+1} 
    = \vec{X}_{j} \vec{L}_j^\cT
    = \vec{X}_{j} ( \vec{I} + \vec{e}_j \vec{l}_j^\cT)
    = \vec{X}_{j} + \vec{X}_{j} \vec{e}_j \vec{l}_j^\cT.
\end{equation*}
Then \( \vec{Q}\vec{L}^{-1}\vec{D}^{-1} \vec{L}^{-\cT}\vec{M}\vec{e}_1 = \vec{X}_{k} \vec{D}^{-1} \vec{y}_{k} \) can be computed accessing \( \vec{L} \), and therefore \( \vec{N} \), column by column.

\subsubsection{Streaming version}
\label{sec:streaming_banded_solve}

Recall that $[\vec{l}_i]_{:,\ell}$ is zero if $\ell \leq i$ or $\ell > i +q$.
Since $[\vec{l}_i]_{:i}$ is zero, we have
\begin{equation*}
    [\vec{y}_j]_j
    = [\vec{y}_{j} +  (\vec{e}_{j}^\cT \vec{y}_{j}) \vec{l}_j]_j
%    = [\vec{L}_j \vec{y}_j]_j
    = [\vec{y}_{j+1}]_j
    =\cdots
    = [\vec{y}_{k}]_j
\end{equation*}
and
\begin{equation*}
    [\vec{X}_{j}]_{:,j}
    = [\vec{X}_{j} + \vec{X}_{j} \vec{e}_j \vec{l}_j^\cT]_{:,j}
%    = [\vec{X}_{j}\vec{L}_j^\cT]_{:,j}
    = [\vec{X}_{j+1}]_{:,j}
    = \cdots = [\vec{X}_{k}]_{:,j}.
\end{equation*}
We therefore have that
\begin{equation*}
    \vec{X}_{k} \vec{D}^{-1} \vec{y}_{k}
    = \smop{\sum_{j=0}^{k-1}} \frac{[\vec{y}_{k}]_j}{[\vec{D}]_{j,j}} [\vec{X}_{k}]_{:,j}
    = \smop{\sum_{j=0}^{k-1}} \frac{[\vec{y}_{j}]_j}{[\vec{D}]_{j,j}} [\vec{X}_{j}]_{:,j}.
\end{equation*}
Similarly, since $[\vec{l}_i]_{i+q+1:}$ is zero,
\begin{equation*}
    [\vec{y}_j]_{j+q:}
    = [\vec{y}_{j-1} + (\vec{e}_{j-1}^\cT\vec{y}_{j-1}) \vec{l}_{j-1}]_{j+q:}
    = [\vec{y}_{j-1}]_{j+q:}
    = \cdots
    = [\vec{y}_0]_{j+q:}
\end{equation*}
and
\begin{equation*}
    [\vec{X}_{j}]_{:,j+q:}
    = [\vec{X}_{j-1}+\vec{X}_{j-1}\vec{e}_{j-1}\vec{l}_{j-1}^\cT]_{:,j+q:}
    = [\vec{X}_{j-1}]_{:,j+q:}
    =\cdots
    = [\vec{X}_{0}]_{:,j+q:}.
\end{equation*}
By definition, $\vec{y}_0 = \vec{M}\vec{e}_0$ and $\vec{X}_0 = \vec{Q}$.
Thus, we see that it is not necessary to know the later columns of $\vec{X}_j$ immediately. 

We can define a streaming algorithm by maintaining only the relevant portions of the $\vec{X}_i$ and $\vec{y}_i$.
Towards this end, define the length \( q+1 \) vector \( \bar{\vec{y}}_j := [\vec{y}_j]_{j:j+q+1} \) the \( n\times (q+1) \) matrix \( \bar{\vec{X}}_j = [\vec{X}_j]_{:,j:j+q+1} \).
Using the above observations, we see that these quantities can be maintained by the recurrences
\begin{equation*}
    \bar{\vec{y}}_j =  
    \begin{bmatrix}
        [\bar{\vec{y}}_{j-1}]_{1:}\\0
    \end{bmatrix}
    +  [\bar{\vec{y}}_{j-1}]_0 [\vec{l}_j]_{j+1:j+q+1} 
\end{equation*}
and
\begin{equation*}
    [\bar{\vec{X}}_j]_{:,:q}
    = [\bar{\vec{X}}_{j-1}]_{:,1:} 
        + ( [\bar{\vec{X}}_{j-1}]_{:,1} ) ([\vec{l}_j]_{j+1:j+q+1})^\cT
    ,\qquad
    [\bar{\vec{X}}_j]_{:,q} = [\vec{Q}]_{j+q}.
\end{equation*}
Note then that,
\begin{equation*}
    \vec{X}_{k-1} \vec{D}^{-1} \vec{y}_{k-1}
    = \smop{\sum_{j=0}^{k-1}}  \frac{[\bar{\vec{y}}_j]_{1}}{[\vec{D}]_{j+1,j+1}} [\bar{\vec{X}}_j]_{:,1}.
\end{equation*}

This results in \cref{alg:streaming_banded_prod} whose streaming pattern is outlined in \cref{fig:streaming_banded_prod}.


\begin{labelalgorithm}[htb]{streaming_banded_inv}{streaming-banded-inv}{Streaming banded inverse}
\begin{algorithmic}[1]
\Class{\thealgorithmname}{$n,k,q$}
\State{stream:$ $}
\State $\ms{LDL} \gets \textsc{streaming-LDL}(k,q)$
\State $\ms{Q0} \gets \textsc{zeros}(n,q)$
\State $\ms{j}\gets0$
\Procedure{read-stream}{$\vec{q},\vec{n},\vec{y}_0$}
\If{$\ms{j}<q$}
\State $[\ms{Q0}]_{:,j} \gets \vec{v}$
\If{$\ms{j}=q-1$}
    \State $\ms{b-prod} \gets \textsc{streaming-banded-prod}(n,k,q)$
    \State $\ms{b-prod.}\textsc{read-stream}(\ms{V0},\ms{none},\ms{none},\ms{none})$
\EndIf
\Else
\State $\ms{LDL.}\textsc{read-stream}(\vec{n})$
\State $\ms{b-inv.}\textsc{read-stream}(\vec{q},-[\ms{LDL.}\ms{L}]_{:,j-q},[\ms{LDL.}\ms{d}]_{j-q},\vec{y}_0)$
\EndIf
\State $\ms{j}\gets \ms{j}+1$
\EndProcedure
\EndClass
\end{algorithmic}
\end{labelalgorithm}



\subsection{Computing polynomials in $\vec{T}$}

The last major remaining piece is to construct $\vec{M} = \tilde{M}\mf{\vec{T}}$ and $\vec{N} = [\tilde{N}\mf{\widehat{\vec{T}}}]_{:k,:k}$.
Recall that we have assumed $\tilde{M}$ and $\tilde{N}$ are of degree at most two for convenience.
In iteration $\ell$ of Lanczos, we obtain $\alpha_{\ell}$ and $\beta_{\ell}$.
Observe that $\vec{T}^2$ is symmetric and that, defining $\beta_{-1} = \beta_{k} = 0$, the lower triangle is given by
\begin{equation*}
    [\vec{T}^2]_{i,j}
    = 
    \begin{cases}
    \beta_{j-1}^2 + \alpha_j^2 + \beta_j^2 & j=i \\
    (\alpha_j + \alpha_{j+1}) \beta_i & j=i-1 \\
    \beta_j \beta_{j+1} & j = i-2 \\ 
    0 & \text{o.w.}
    \end{cases}
\end{equation*}
%Thus, assuming we have stored $\{ \alpha_i \}_{i=1}^{\ell-1}$ and $\{ \beta_i \}_{i=1}^{\ell-1}$ we can compute $[\vec{M}]_{:\ell-1,:\ell-1}$ and $[\vec{N}]_{:\ell-1,:\ell-1}$.



We can use this to implement the streaming algorithm, \cref{alg:streaming_tridiag_square}, for computing the entries of $\vec{T}^2$.
Rather than being fed the entire tridiagonal matrix $\vec{T}$, \cref{alg:streaming_tridiag_square} is fed a stream of the columns of $\vec{T}$ in order, as shown in \cref{fig:streaming_tridiag_square}.
The algorithm respectively stores the $j$-th diagonals of $\vec{T}$ and $\vec{T}^2$  as $[\ms{T}]_{j,:}$ and $[\ms{Tp2}]_{j,:}$.

\begin{labelalgorithm}[htb]{streaming_tridiag_square}{streaming-tridiag-square}{Streaming tridiagonal square}
\begin{algorithmic}[1]
\Class{\thealgorithmname}{$k$}
\State{stream: $(\alpha_0, \beta_0), \ldots, (\alpha_{k-1},\beta_{k-1})$}
\State $\ms{T} \gets \textsc{zeros}(2,k)$
\State $\ms{Tp2} \gets \textsc{zeros}(3,k)$
\State $\ms{j}\gets0$
\Procedure{read-stream}{$\alpha,\beta$}
\State $[\ms{T}]_{0,\ms{j}} = \alpha$
\State $[\ms{T}]_{1,\ms{j}} = \beta$
\If{$\ms{i}=0$}
\State $[\ms{Tp2}]_{0,\ms{j}} \gets [\ms{T}]_{0,\ms{j}}^2+[\ms{T}]_{1,\ms{j}}^2$
\Else
\State $[\ms{Tp2}]_{0,\ms{j}} \gets [\ms{T}]_{0,\ms{j}}^2+[\ms{T}]_{1,\ms{j}}^2+[\ms{T}]_{1,\ms{j}-1}^2$
\State $[\ms{Tp2}]_{1,\ms{j}} \gets ([\ms{T}]_{0,\ms{j}}+[\ms{T}]_{0,\ms{j}-1})[\ms{T}]_{1,\ms{j}-1}$
\State $[\ms{Tp2}]_{2,\ms{j}} \gets [\ms{T}]_{1,\ms{j}}[\ms{T}]_{1,\ms{j}-1}$
\EndIf
\State $\ms{j}\gets \ms{j}+1$
\EndProcedure
\EndClass
\end{algorithmic}
\end{labelalgorithm}




Since we maintain the columns of $\vec{T}^2$ with \cref{alg:streaming_tridiag_square}, we can easily compute $\vec{M}$ and $\vec{N}$ using \cref{alg:get_poly}.

\begin{labelalgorithm}[htb]{get_poly}{get-poly}{Get polynomial of tridiagonal matrix}
\begin{algorithmic}[1]
\Procedure{\thealgorithmname}{$P,\ms{STp2},k,j$}
\State $a,b,c = P(0), P'(0), P''(0)$
\State $\vec{p} \gets \textsc{zeros}(3)$
\State $[\vec{p}]_{:3} \gets a [\ms{STp2.}\ms{Tp2}]_{:,j}$
\State $[\vec{p}]_{:2} \gets b [\ms{STp2.}\ms{T}]_{:,j}$
\State $[\vec{p}]_{:1} \gets c $
\EndProcedure
\end{algorithmic}
\end{labelalgorithm}

\subsection{Putting it all together}

With this algorithm in place, putting everything together is straightforward, and the full implementation is shown in \cref{alg:rational_lanczos}.

\begin{labelalgorithm}[htb]{rational_lanczos}{banded-rational}{Streaming banded rational inverse}
\begin{algorithmic}[1]
\Class{\thealgorithmname}{$n,k,\tilde{M},\tilde{N}$}
\State $\ms{b-inv} \gets \textsc{banded-inv}(n,k,2)$
\State $\ms{STp2} \gets \textsc{streaming-tridiag-square}(k)$
\State $\ms{j}\gets0$
\Procedure{read-stream}{$\vec{q},\alpha,\beta$}
\If{$\ms{j}<k$}
    \State $\ms{STp2.}\textsc{read-stream}(\alpha,\beta)$
    \State $\ms{b-inv.}\textsc{read-stream}($
    \State \hspace{1em} $\vec{q},$
    \State \hspace{1em} $\textsc{getpoly}(\tilde{N},\ms{STp2},k,\ms{j}-1)~\textbf{if}~\ms{j}\geq 2~\textbf{else}~\ms{none},$
    \State \hspace{1em} $\textsc{getpoly}(\tilde{M},\ms{STp2},k,\ms{j}-1)~\textbf{if}~\ms{j}= 2~\textbf{else}~\ms{none},$
    \State$)$
\EndIf
\State $\ms{LDL.}\textsc{read-stream}(\vec{n})$
\State $\ms{j}\gets \ms{j}+1$
\EndProcedure
\Procedure{finish-up}{$ $}
\For{$i=k,k+1$}
    \State $\ms{b-inv.}\textsc{read-stream}( \ms{none}, \textsc{getpoly}(\tilde{N},\ms{STp2},k,\ms{j}-1), \ms{none} )$
\EndFor
\EndProcedure
\Procedure{get-output}{$ $}
\State \Return \ms{b-inv.}\ms{b-prod.}\ms{out}
\EndProcedure
\EndClass 
\end{algorithmic}
\end{labelalgorithm}

This can be incorporated into any Lanczos implementation and used to compute the Lanczoz-OR iterates.
For concreteness, we show this with the implementation of Lanczos from \cref{alg:lanczos}.
We call the resulting implementation Lanczos-OR-lm.

\begin{labelalgorithm}[htb]{lanczos_OR}{Lanczos-OR-lm}{Lanczos-OR (low memory)}
\begin{algorithmic}[1]
\Procedure{\thealgorithmname}{$\vec{A},\vec{v},k,M,N$}
\State \( \vec{q}_{-1} = \vec{0} \),
\( \beta_{-1} = 0 \),
    \( \vec{q}_0  = \vec{v} \)
\State Set $\tilde{M}$ and $\tilde{N}$ as in \cref{def:lanczos_or}
\State $\ms{lam-lm} \gets \textsc{banded-rational}(n,k,\tilde{M},\tilde{N})$
\For {\( j=0,1,\ldots, k-1 \)}
    \State \( \tilde{\vec{q}}_{j+1} = \vec{A} \vec{q}_{j} - \beta_{j-1} \vec{q}_{j-1} \)
    \State \( \alpha_j = \langle \tilde{\vec{q}}_{j+1}, \vec{q}_j \rangle \)
    \State \( \tilde{\vec{q}}_{j+1} = \tilde{\vec{q}}_{j+1} - \alpha_j \vec{q}_j \)
%    \If{reorthogonalize}
        \State optionally, reorthogonalize \( \tilde{\vec{q}}_{j+1} \) against \( \{\vec{q}_i\}_{i=0}^{j-1} \)
%    \EndIf
    \State \( \beta_{j} = \| \tilde{\vec{q}}_{j+1} \| \)
    \State \( \vec{q}_{j+1} = \tilde{\vec{q}}_{j+1} / \beta_{j} \)
    %\If {Lanczos-FA-lm~\textbf{and}~$j=k-1$}
    %\State $\ms{lam-lm.}\textsc{read-stream}(\vec{q}_{k-1},\alpha_{k-1},0)$
    %\Else%If {Lanczos-OR-lm}
    \State $\ms{lam-lm.}\textsc{read-stream}(\vec{q}_{j},\alpha_{j},\beta_{j})$
    %\EndIf
\EndFor
\State $\ms{lam-lm.}\textsc{finish-up}()$
\EndProcedure 
\end{algorithmic}
\end{labelalgorithm}

We can easily obtain an implementation of Lanczos-FA, which we call Lanczos-FA-lm, by replacing $\beta_{k-1}$ with $0$ in the final iteration of the loop. 

\subsection{Some comments on implementation}

Our main goal is to describe how to implement Lanczos-FA and Lanczos-OR in a way that requires $k$ matrix-vector products and $O(n)$ storage, when $M$ and $N$ are each at most degree two. As mentioned, the approach can be extended to any constant degree.
There are a range of improvements to our implementation which may be useful in practice.

First, the amount of storage used can be reduced somewhat.
Indeed, the implementation described above saves $\vec{T}$, $\vec{T}^2$, $\vec{L}$, and $\vec{d}$, but only accesses a sliding window of these quantities.
We have chosen to save them for convenience since they require only $O(k)$ storage.
However, storing only the relevant information from these quantities would result in an implementation with storage costs independent of the number of iterations $k$.
In this vein, a practical implementation would likely determine $k$ adaptively by monitoring the residual or other measures of the error.
Improvements to the number of vectors of length $n$ may be possible as well.
For example, storage could possibly be reduced somewhat by incorporating the Lanczos iteration more explicitly with the inversion of the LDL facorization, much like the classical Hestenes and Stiefel implementation of CG \cite{hestenes_stiefel_52}.


As with other short-recurrence based Krylov subspace methods, the behavior of Lanczos-FA-lm and Lanzos-OR-lm in finite precision arithmetic may be different than in exact arithmetic.
However, with the exception of the Lanczos algorithm, the other aspects of our algorithm are essentially backwards stable.
It is therefore more or less clear that Lanczos-FA-lm and Lanczos-OR-lm will accurately compute the expressions $\vec{Q} N\mf{\vec{T}}^{-1} M\mf{\vec{T}} \vec{e}_0$ and $\vec{Q}  ([\tilde{N}\mf{\widehat{\vec{T}}}]_{:k,:k})^{-1} \tilde{M}\mf{\vec{T}} \vec{e}_0$ provided that $M\mf{\vec{T}}$ $N\mf{\vec{T}}$, $[\tilde{N}\mf{\widehat{\vec{T}}}]_{:k,:k}$ are reasonably well conditioned.
Indeed, in practice solving linear systems by symmetric Gaussian elimination is accurate; see for instance \cite[Chapter 10]{higham_02}.
Thus, such bounds and techniques can be applied to Lanczos-FA-lm and Lanczos-OR-lm.

\iffalse
However, our implementation is equivalent, even in finite precision arithmetic, to \begin{itemize}
    \item running Lanczos to obtain $\vec{Q}$ and $\vec{T}$
    \item computing $\vec{M}$ and $\vec{N}$ via certain standard approachs
    \item computing an LDL factorizaton of $\vec{N}$ via a certain standard approach
    \item computing $\vec{Q} \vec{L}^{-1} \vec{D}^{-1} \vec{L}^{-\cT}\vec{M}\vec{e}_0$ in an essentially standard way.
\end{itemize}
\fi




\chapter{General matrix function approximation}
\label{chap:general_matfunc}

We now consider Krylov subspace methods for approximating \( \fA \vec{v} \) for more general \( f \).
In \cref{sec:lanczos_FA} we discuss the Lanczos method for matrix function approximation (Lanczos-FA), which is probably the most widely used algorithm for this task.
Then, in \cref{sec:lanczos_OR_gen}, we discuss how Lanczos-OR iterates can be used to generate approximations to a wide range of functions.


\section{Explicit polynomial methods}

A simple approach to approximating \( \fA\vec{v} \) is to compute 
\begin{equation*}
    \ff{f}{k-1}\mf{\vec{A}} \vec{v} \in \mathcal{K}_{k}
\end{equation*}
where \( \ff{f}{k-1} \) is some polynomial chosen to approximate \( f \).
For instance, the Chebyshev semi-iterative method mentioned in the example from \cref{chap:intro} falls into this category of algorithms.

If \( |f - \ff{f}{k-1}| \) is not strongly correlated with the eigenvalues of \( \vec{A} \), we might expect 
\begin{equation*}
    \| \fA \vec{v} - \ff{f}{k-1}\mf{\vec{A}} \vec{v} \|
    > c \| f - \ff{f}{k-1} \|_{\mathcal{I}} \| \vec{v} \|
\end{equation*}
for some reasonable constant \( c \).

Like mentioned in our discussion of algorithms for quadratic forms in \cref{sec:QF_tradeoffs}, one nice property of explicit polynomial methods is the lack of inner products, which can be expensive on supercomputers. 
In addition, explicit polynomial methods for \( \fA\vec{v} \) do not require storing a basis for Krylov subspace.

\section{Lanczos-FA}
\label{sec:lanczos_FA}

Early uses of Lanczos-FA were focused primarily on computing matrix exponentials applied to a vector; i.e. \( f = \exp(tx) \).
As far as we can tell, Lanczos-FA was introduced in \cite{nauts_wyatt_83} and first used for general \( f \) in \cite{vandervorst_87}.
Soon after Lanczos-FA was first used, a number of papers studying the algorithm and its convergence properties were published \cite{park_light_86,druskin_knizhnerman_88,druskin_knizhnerman_89,gallopoulos_saad_92,saad_92}.
These early works were followed by a number of papers demonstrating the effectiveness of Lanczos-FA in finite precision arithmetic \cite{druskin_knizhnerman_91,druskin_knizhnerman_95,druskin_greenbaum_knizhnerman_98}, a topic we discuss further in \cref{chap:finite_precision}.


\begin{definition}
    \label{def:lanczos_fa}
The \( k \)-th Lanczos-FA approximation to \( f\A\vec{v} \) is
\begin{equation*}
    \lan_k(f) 
    := \vec{Q} f\mf{\vec{T}} \vec{e}_0.
\end{equation*}
\end{definition}

\begin{remark}
\label{rem:CG_MINRES}
If \( \vec{A} \) is positive definite, then the Lanczos-FA approximation to \( f = 1/x \) coincides with the CG iterate defined in \cref{sec:linear_solvers}.
 Even when $\vec{A}$ is indefinite, we can use the Lanczos-FA iterate as an approximation to \( \vec{A}^{-1}\vec{v} \).
However, the resulting algorithm is not guaranteed to be optimal, and if $\vec{T}$ has an eigenvalue near to or at zero then $\vec{T}^{-1}$ will be poorly conditioned or even undefined. 
Even so, the \emph{overall} convergence of the Lanczos-FA iterate is closely related to the convergence of MINRES.
    We described this phenomenon in detail in \cref{sec:CG_indefinite}. \qedhere
\end{remark}


A basic property of the Lanczos-FA iterate is that polynomials are applied exactly. 
More precisely, we have the following, well known, theorem.
\begin{theorem}
    \label{thm:lanczosFA_exact}
    Suppose \( \deg(p) < k \). 
    Then,
    \begin{equation*}
        \lan_k(p) = \pA \vec{v}.
    \end{equation*}
\end{theorem}

\begin{proof}
    Using \cref{thm:tridiag_power} we have
    \begin{align*}
        \vec{A}^{q} \vec{v} = \vec{A}^q \Qhat \vec{e}_0
        = \Qhat \widehat{\vec{T}}^q \vec{e}_0
        = \vec{Q} \vec{T}^q \vec{e}_0.
    \end{align*}
\end{proof}


This theorem implies that, when \( \mu = \Psi \), 
\begin{equation*}
    \lan_k(f) = \ff[i]{f}{k-1}\mf{\vec{A}} \vec{v}.
\end{equation*}
In other words, the Lanczos-FA iterate is obtained by interpolating \( f \) at the eigenvalues of \( \vec{T} \) with a degree \( k-1 \) polynomial.



\subsection{A priori error bounds on an interval}

Akin to the bounds we saw in \cref{sec:quadrature_apriori}, we can derive a bound based on best approximation on an interval.
\begin{theorem}
\label{thm:poly_unif}
    The Lanczos-FA iterate satisfies
\begin{equation*}
    \frac{\| \fA \vec{v} - \lan_k(f) \|}{\|\vec{v}\|}
    %&\leq  2 \left( \min_{\deg(p)<k} \max_{x \in \mathcal{I}\A} | f(x) - p(x) | \right) \| \vec{v} \|_2.
    \leq  2 \min_{\deg(p)<k} \| f-p \|_{\mathcal{I}} .
\end{equation*}   
\end{theorem}
\begin{proof}
For any polynomial \( p \) with \( \deg(p) < k \),
\begin{align*}
\| \fA \vec{v} - \lan_k(f) \|
& \leq \| \fA \vec{v} - \pA \vec{v} \| + \| \lan_k (p) - \lan_k (f) \| 
\\& = \| (\fA - \pA )\vec{v} \| + \| \vec{Q} ( p\mf{\vec{T}} - f\mf{\vec{T}} ) \vec{Q}^{\cT} \vec{v} \| 
\\& \le \| \fA - \pA \|_2  \| \vec{v} \| + \| \vec{Q} ( p\mf{\vec{T}} - f\mf{\vec{T}} ) \vec{Q}^{\cT} \|_2  \|\vec{v} \|
\\& \le ( \| \fA - \pA \|_2 + \| p(\vec{T} ) - f\mf{\vec{T}} \|_2 )  \| \vec{v} \| . 
    \\& = ( \| f-p \|_{\Lambda} + \| f-p \|_{\Lambda\mf{\vec{T}}} )  \| \vec{v} \| . 
\end{align*}
Then, optimizing over polynomials of degree less than \( k \),
\begin{equation}
\| \fA \vec{v} - \lan_k(f) \|
    \leq \min_{\deg(p)< k} \left( \| f-p \|_{\Lambda} + \| f - p \|_{\Lambda\mf{\vec{T}}} \right) \| \vec{v} \|. \label[ineq]{eqn:triangle_ineq1}
\end{equation}
Finally, using that \( \Lambda,\Lambda\mf{\vec{T}} \subset \mathcal{I} \), we obtain the result.
\end{proof}

As we will discuss in \cref{chap:finite_precision}, bounds for Lanczos-FA based on polynomial approximation on \( \mathcal{I} \) still hold, to close approximation, in finite precision arithmetic.

%\Cref{thm:poly_unif} asserts that, except for a possible factor of \( 2 \), the error of the Lanczos-FA approximation to \( \fA \vec{v} \) is at least as good as the best polynomial approximation to \( f \) on the interval containing the eigenvalues of \( \vec{A} \).
%For arbitrary \( f \), \cref{thm:poly_unif} remains the standard bound for Lanczos-FA. 
%As we discuss in \cref{chap:finiite_precision}, it has been studied carefully and is known to hold to a close degree in finite precision arithmetic \cite{musco_musco_sidford_18}.
%However, as we have noted throughout this thesis, such bounds are insufficient to truly describe the behavior of many methods, including Lanczos-FA.

\subsection{Two-pass Lanczos-FA}

A major downside of Lanczos-FA compared with explicit polynomial approaches is that a simple implementation requires that \( \vec{Q} \) be stored.
Fortunately, this storage cost can be avoided by incurring additional computational cost. 
In particular, we can use an implementation called two pass Lanczos-FA \cite{boricci_00a,frommer_simoncini_08}.
On the first pass, the tridiagonal matrix $\vec{T}$ is computed using the short-recurrence version of Lanczos; i.e., without storing all of $\vec{Q}$.
Once $\vec{T}$ has been computed, $f(\vec{T}) \vec{e}_0$ can be evaluated using $O(k^2)$ storage.
Lanczos is then run again and the product $\vec{Q} f(\vec{T}) \vec{e}_0$ is computed as the columns of $\vec{Q}$ become available.
Note that on the second run, the \emph{exact same} Lanczos vectors (even in finite precision arithmetic) can be computed without any inner products by using the values computed in the first run and stored in \( \vec{T} \).


Such an approach can be generalized by re-generating the Lanczos recurrence from multiple points simultaneously on the second pass \cite{li_22}.
Specifically, on the first pass, vectors \( \vec{q}_j \) and \( \vec{q}_{j-1} \) can be saved for \( j=0,d,2d,\ldots \). 
Then, on the second pass, the rest of the Lanczos vectors can be constructed by continuing the three-term Lanczos recurrence \cref{eqn:lanczos_three_term} from each of the roughly \( n/d \) start points in parallel. 
Thus, the number of matrix-loads is reduced by a factor of roughly \( d \) at the cost of storing roughly \( 2n/d \) vectors.
The case \( d=n \) gives the original two-pass approach.

\section{Lanczos-OR based methods}
\label{sec:lanczos_OR_gen}

\iffalse
As an alternate motivation for the Lanczos-FA iterate, we note that many functions can be written in the form
\begin{equation*}
    f(x) = \int_{\Gamma} \frac{g(z)}{x-z} \d{z}
\end{equation*}
for some contour \( \Gamma \) and function \( g \).
Indeed, such functions include Stieltjes functions, analytic functions, and rational functions.
In such cases, it is not hard to see that
\begin{equation*}
    \fA \vec{v} = \int_{\Gamma} g(z) (\vec{A} - z\vec{I})^{-1} \vec{v} \d{z}.
\end{equation*}
As discussed in \cref{rem:CG_MINRES} and justified in \cref{sec:CG_MINRES}, for a fixed value \( z\in\mathbb{C} \), we can approximate \( (\vec{A} - z\vec{I})^{-1} \vec{v} \) by means of the formal expression for the conjugate gradient algorithm; i.e. by 
\begin{equation*}
    (\vec{A} - z\vec{I})^{-1} \vec{v} \approx \vec{Q} (\vec{T} - z\vec{I})^{-1} \vec{e}_0.
\end{equation*}
Integrating the resulting approximation we obtain
\begin{equation*}
    \int_{\Gamma} g(z) \vec{Q} (\vec{T} - z\vec{I})^{-1} \vec{e}_0 = \vec{Q} f\mf{\vec{T}} \vec{e}_0 = \lan_k(f).
\end{equation*}
We will use the integral representation to derive error bounds for Lanczos-FA in \cref{chap:CIF}.

\fi

We can use integral representations of functions to derive algorithms based on Lanczos-OR iterates.
For concreteness, we consider the case of the matrix-sign function and rational functions in partial fraction form.
It's clear that a similar approach can be applied to other functions, and further study of the resulting algorithms would be interesting.


\subsection{The matrix sign function}
\label{sec:sign_function}

We begin by noting that, for any $a>0$,
\begin{equation*}
    \frac{1}{\sqrt{a}}
    = \frac{2}{\pi} \int_0^\infty \frac{1}{a + z^2} \d{z}.
\end{equation*}
%\Anne{[Is it standard to write this as $\frac{1}{\sqrt{x^2}}$ instead of  $\frac{1}{|x|}$?]}
Thus, if $f = \operatorname{sign}=x/|x| = x/\sqrt{x^2}$, we have 
\begin{equation*}
    \fA\vec{v} = \frac{2}{\pi} \int_0^\infty  \vec{A}(\vec{A}^2 + z^2\vec{I})^{-1} \vec{v} \: \d{z}.
\end{equation*}
The Lanczos-OR approximation to $\vec{A}(\vec{A}^2 + z^2\vec{I})^{-1} \vec{v}$ is $\vec{Q} ([\widehat{\vec{T}}^2]_{:k,:k} + z^2\vec{I})^{-1} \vec{T} \vec{e}_0$, which is optimal over Krylov subspace in the $(\vec{A}^2 + z^2\vec{I})$-norm.
This yields the approximation 
\begin{equation*}
    \frac{2}{\pi} \int_0^\infty \vec{Q} ([\widehat{\vec{T}}^2]_{:k,:k} + z^2\vec{I})^{-1} \vec{T} \vec{e}_0 \: \d{z}
    = \vec{Q} \left([\widehat{\vec{T}}^2]_{:k,:k} \right)^{-1/2} \vec{T}  \vec{e}_0.
\end{equation*}
Thus, we can define the induced iterate as
\begin{equation}
    \label{def:sign_or}
\mathsf{sign-OR}_k := \vec{Q} \left([\widehat{\vec{T}}^2]_{:k,:k} \right)^{-1/2} \vec{T}  \vec{e}_0
=
\vec{Q} \left( [\widehat{\vec{T}}]_{:k,:k+1}[\widehat{\vec{T}}]_{:k+1,:k} \right)^{-1/2} \vec{T}  \vec{e}_0.
\end{equation}


\subsubsection{Relation to Lanczos-FA}

The Lanczos-OR and Lanczos-FA iterates for a given rational matrix function are clearly related. 
In particular, $\tilde{N}\mf{\vec{T}}$ and $[\tilde{N}(\widehat{\vec{T}})]_{:k,:k}$ differ only in the bottom rightmost $(q-1)\times (q-1)$ principle submatrix, where $q = \deg(\tilde{N})$.
Using this fact, it can be shown that the Lanczos-OR and Lanczos-FA iterates ``tend to coalesce as convergence takes place'' \cite[Proposition 5.1]{lopez_simoncini_06}.
We now show that a similar phenomenon occurs with the induced Lanczos-OR approximation to the sign function and the Lanczos-FA approximation.

\begin{theorem}
\label{thm:sign_approx_compare}
The Lanczos-FA and induced Lanczos-OR approximations to the matrix sign function satisfy
\begin{equation*}
    \| \lan_k(\operatorname{sign}) - \mathsf{sign-OR}_k \|_2
    \leq
    \beta_{k-1}^2 
    \frac{\sqrt{\alpha_0^2+\beta_0^2}}{2 \sigma_{\textup{min}}\mf{\vec{T}}^3}.
\end{equation*}
where $\sigma_{\textup{max}}\mf{\vec{T}}$ and $\sigma_{\textup{min}}\mf{\vec{T}}$ are the largest and smallest singular values of $\vec{T}$ respectively.
\end{theorem}


\begin{proof}%[Proof of \cref{thm:sign_approx_compare}]
Let ${N} = x^2 + z^2$ and ${M} = x$.
Note that ${N}\mf{\vec{T}} = [{N}(\widehat{\vec{T}})]_{:k,:k} - \beta_{k-1}^2 \vec{e}_{k-1}\vec{e}_{k-1}^\cT$ so,
\begin{equation*}
    {M}\mf{\vec{T}} \vec{e}_0
%    = {N}\mf{\vec{T}} {N}\mf{\vec{T}} ^{-1} {M}\mf{\vec{T}} \vec{e}_0
    = ([{N}(\widehat{\vec{T}})]_{:k,:k} - \beta_{k-1}^2 \vec{e}_{k-1}\vec{e}_{k-1}^\cT){N}\mf{\vec{T}}^{-1} {M}\mf{\vec{T}} \vec{e}_0.
\end{equation*}
    Thus, rearranging terms and multiplying by $([{N}\mf{\widehat{\vec{T}}}]_{:k,:k})^{-1}$ we find that
\begin{align*}
    &{N}\mf{\vec{T}}^{-1} {M}\mf{\vec{T}} \vec{e}_0 - ([{N}(\widehat{\vec{T}})]_{:k,:k})^{-1} {M}\mf{\vec{T}} \vec{e}_0
    \\&\hspace{6em}= \beta_{k-1}^2 ([{N}(\widehat{\vec{T}})]_{:k,:k})^{-1} \vec{e}_{k-1}\vec{e}_{k-1}^\cT {N}\mf{\vec{T}}^{-1}{M}\mf{\vec{T}} \vec{e}_0
\end{align*}
After left multiplying with $\vec{Q}$, this implies that
\begin{equation*}
    \lan_k(r) - \lanopt_k(r,1)
    = \beta_{k-1}^2 \vec{Q} ([{N}(\widehat{\vec{T}})]_{:k,:k})^{-1} \vec{e}_{k-1}\vec{e}_{k-1}^\cT  {N}\mf{\vec{T}}^{-1}{M}\mf{\vec{T}} \vec{e}_0.
\end{equation*}
Now, suppose that $f = \operatorname{sign}$ and set $ \mathsf{diff}_k := \lan_k(\operatorname{sign}) - \mathsf{sign-OR}_k$ as above.
Then, since the Lanczos-FA approximation can also be induced by an integral over $z\in[0,\infty)$, we have that,
\begin{equation*}
    \mathsf{diff}_k
    = \beta_{k-1}^2 
    \frac{2 }{\pi} \int_{0}^{\infty}
    \vec{Q} ([\widehat{\vec{T}}^2]_{:k,:k} + z^2 \vec{I})^{-1} \vec{e}_{k-1}\vec{e}_{k-1}^\cT (\vec{T}^2 + z^2 \vec{I})^{-1} \vec{T} \vec{e}_0 
    \d{z}.
\end{equation*}
Note that $[\widehat{\vec{T}}^2]_{:k,:k} - \vec{T}^2 = \beta_{k-1}^2 \vec{e}_{k-1} \vec{e}_{k-1}^\cT$ is positive semidefinite.
Therefore, using that $\sigma_{\textup{min}}([\widehat{\vec{T}}^2]_{:k,:k}) \geq \sigma_{\textup{min}}(\vec{T}^2) = \sigma_{\textup{min}}\mf{\vec{T}}^2$,
\begin{align*}
    \| \mathsf{diff}_k \|_2
    &
    = \beta_{k-1}^2 
     \bigg\| \vec{Q} \bigg( \frac{2 }{\pi} \int_{0}^{\infty}
    ([\widehat{\vec{T}}^2]_{:k,:k} + z^2 \vec{I})^{-1} \vec{e}_{k-1}\vec{e}_{k-1}^\cT (\vec{T}^2 + z^2 \vec{I})^{-1} \d{z} \bigg)  \vec{T} \vec{e}_0 
     \bigg\|_2
    \\&
    \leq \beta_{k-1}^2 
     \bigg( \frac{2 }{\pi} \int_{0}^{\infty}
    \| ([\widehat{\vec{T}}^2]_{:k,:k} + z^2 \vec{I})^{-1} \|_2 \| (\vec{T}^2 + z^2 \vec{I})^{-1} \|_2 \d{z} \bigg)  \| \vec{T}\vec{e}_0 \|_2
    \\& \leq
    \beta_{k-1}^2 
     \bigg( \frac{2 }{\pi} \int_{0}^{\infty}
    | (\sigma_{\textup{min}}\mf{\vec{T}}^2 + z^2 )^{-1} |~ | (\sigma_{\textup{min}}\mf{\vec{T}}^2 + z^2 )^{-1} | \d{z} \bigg)  \sqrt{\alpha_0^2+\beta_0^2}
    \\&= \beta_{k-1}^2 
    \frac{\sqrt{\alpha_0^2+\beta_0^2}}{2 \sigma_{\textup{min}}\mf{\vec{T}}^3}.
    \tag*{\qedhere}
\end{align*}
\end{proof}

Since $|\beta_{k-1}|$ tends to decrease as the Lanczos method converges, this seemingly implies that the induced Lanczos-OR iterate and the Lanczos-FA iterate tend to converge in this limit.
However, recall that $\vec{T} = [\widehat{\vec{T}}]_{:k,:k}$ changes at each iteration $k$.
In particular, there is the difficulty that $\vec{T}$ may have an eigenvalue near zero, in which case the preceding bound could be useless.

It is known that \( \vec{T} \) cannot have small eigenvalues in two consecutive iterations, provided the eigenvalues of \( \vec{A} \) are not small \cite{greenbaum_druskin_knizhnerman_99}, a result we will recall in \cref{thm:Tk_invertible_bd}.
Since $\beta_{k-1}$ has little to do with the minimum magnitude eigenvalue of $\vec{T}$ (recall that the Lanczos recurrence is shift invariant), we expect that the ``overall'' convergence of the induced Lanczos-OR iterate and the Lanczos-FA iterate will be similar as Lanczos converges.


\subsection{Rational function approximation}
\label{sec:rational_func}


We can use a similar approach to derive (non-optimal) approximations to rational matrix functions $r(\vec{A})\vec{b}$. 
In many settings, particularly when the rational function $r$ is used as a proxy for a function $f$, this approach is more natural than computing the Lanczos-OR approximation to the rational function directly.
The convergence of such methods is closely related to the quality of the (scalar) rational function approximation as well as the quality of the approximation to the rational matrix function.
Specifically, for any output $\mathsf{alg}(r)$ meant to approximate $r(\vec{A})\vec{b}$, we have the following bound:
\begin{align}
\| f(\vec{A}) \vec{b} - \mathsf{alg}(r) \|
&\leq 
\nonumber
\| f(\vec{A}) \vec{b} - r(\vec{A}) \vec{b} \| +  \| r(\vec{A}) \vec{b} - \mathsf{alg}(r) \|
\\&\leq \nonumber
\| f(\vec{A}) - r(\vec{A}) \|_2 \|\vec{b} \| +  \| r(\vec{A}) \vec{b} - \mathsf{alg}(r) \|
\\&\leq \nonumber
\|\vec{b}\|  \max_{\lambda \in \Lambda} | f(\lambda) - r(\lambda) | +  \| r(\vec{A}) \vec{b} - \mathsf{alg}(r) \|
\\& \leq \underbrace{\|\vec{b}\|  \max_{\lambda \in \mathcal{I}} | f(\lambda) - r(\lambda) |}_{\text{approximation error}} 
    + \underbrace{\vphantom{\max_{\lambda\in\mathcal{I}}}\| r(\vec{A}) \vec{b} - \mathsf{alg}(r) \|}_{\text{application error}}.
    \label{eqn:triangle_ineq}
\end{align}
In many cases, very good or even optimal scalar rational function approximations to a given function on a single interval are known or can be easily computed. 
Thus, the approximation error term can typically be made small with a rational function of relatively low degree \cite{trefethen_19,nakatsukasa_sete_trefethen_18}.

Of course, this bound is only meaningful if the approximation error term is small relative to the application error.
Indeed, we also have
\begin{align}
\| f(\vec{A}) \vec{b} - \mathsf{alg}(r) \|
&\geq 
\big| \| f(\vec{A}) \vec{b} - r(\vec{A}) \vec{b} \| -  \| r(\vec{A}) \vec{b} - \mathsf{alg}(r) \| \big|.
\label{eqn:reverse_triangle_ineq}
\end{align}
This shows that the size of $\| f(\vec{A}) \vec{b} - \mathsf{alg}(r) \|$ is rouhgly the size of $\| f(\vec{A}) \vec{b} - r(\vec{A}) \vec{b} \|$ when $\mathsf{alg}(r)$ is a good approximation to $r(\vec{A}) \vec{b}$.


As we noted, rational function approximations commonly are obtained by discretizing an integral representation using a numerical quadrature approximation. 
For instance, the matrix sign function may be approximated as 
\begin{equation}
    \label{eqn:rational_func}
    r_q\A\vec{v} = \sum_{i=0}^{q-1} \omega_i \vec{A} (\vec{A}^2 + z_i^2 \vec{I})^{-1} \vec{v}
\end{equation}
where $z_i$ and $\omega_i$ are appropriately chosen quadrature nodes and weights \cite{hale_higham_trefethen_08}.

We can of course write $r_q = M_q / N_q$, so it's tempting to set $R_q = 1$ and $\vec{H}_q = N_q\A$ and then use Lanczos-OR to compute the $\vec{H}_q$-norm optimal approximation. 
However, while $r_q$ is convergent to $f$ as $q\to\infty$, $N_q := \prod_{i=0}^{q-1} (x^2
+ z_i^2 )$ is not convergent to any fixed function.
In fact $N_q$ will increase in degree and $\vec{H}_q$ will be increasingly poorly conditioned.
This presents a numerical difficulty in computing the Lanczos-OR iterate in this limit.
More importantly, it is not clear that it is meaningful to approximate a function in this way. 
Indeed, it seems reasonable to expect that, for fixed $k$, as $q\to\infty$,  our approximation should be convergent to something. 
However, we cannot guarantee $\lanopt_k(M_q,N_q)$ is convergent in this limit.

On the other hand, akin to our approach for approximating the integral described in the previous subsection, we can compute the \emph{term-wise optimal} approximations to each term in the sum representation of $r_q$ and output
\begin{equation*}
    \sum_{i=0}^{q-1} \omega_i \vec{Q} ([\widehat{\vec{T}}^2]_{:k,:k} + z_i^2\vec{I})^{-1} \vec{T} \vec{e}_0.
\end{equation*}
In this case, as $q\to\infty$, the approximation is convergent to the integral output. 
It would be interesting to understand when a term-wise optimal approximation behaves nearly optimally.


\section{Numerical experiments}
\label{sec:comparison}


\subsection{The matrix sign function}

We now provide several examples which illustrate various aspects of the convergence properties of Lanczos-OR and Lanczos-OR based algorithms, and show when these new methods can outperform more standard techniques like the classic Lanczos-FA.

As we noted in \cref{sec:sign_function}, Lanczos-OR can be used to obtain an approximation to the matrix sign function. 
A related approach, which interpolates the sign function at the so called ``harmonic Ritz values'', is described in \cite[Section 4.3]{eshof_frommer_lippert_schilling_van_der_vorst_02}.
The harmonic Ritz values are characterized by the generalized eigenvalue problem
\begin{equation*}
     [\widehat{\vec{T}}^2]_{:k,:k} \vec{y} = \theta \vec{T} \vec{y}
\end{equation*}
and are closely related to MINRES in the sense that MINRES produces a polynomial interpolating $1/x$ at the harmonic Ritz values \cite{paige_parlett_vandervorst_95}.
Finally, a standard approach is using Lanczos-FA (or equivalently, Gaussian quadrature) as described in \cref{chap:QF,chap:random_quadrature}.


\subsubsection{Spectrum approximation}

In this example, we show the spectrum approximations induced by the algorithms described above.
We now set $\vec{A}$ to be a diagonal matrix with $1000$ eigenvalues set to the quantiles of a Chi-squared distribution with parameters $\alpha=1$ and $\beta=10$. 
We set $k=10$ and consider approximations to the function $c\mapsto \vec{v}^\cT \bOne [\vec{A} \leq c] \vec{v}$ for a range of values $c$. 
Here $\bOne[x\leq c] = (1-\operatorname{sign}(x-c))/2$ is one if $x\leq c$ and zero otherwise.
We pick $\vec{v}$ as a unit vector with equal projection onto each eigencomponent so that $\vec{v}^\cT \bOne [\vec{A} \leq c] \vec{v}$ gives the fraction of eigenvalues of $\vec{A}$ below $c$.
In the $n\to\infty$ limit, this function will converge pointwise to the cumulative distribution of a Chi-squared random distribution with parameters $\alpha=1$ and $\beta=10$. 
The results are shown in \cref{fig:ch6_sign_spec}.

Note that the Lanczos-FA based approach is piecewise constant with jumps at each eigenvalue of $\vec{T}$.
On the other hand, the harmonic Ritz value and Lanczos-OR based approaches produce continuous approximations to the spectrum.
In this particular example, the spectrum of $\vec{A}$ is near to a smooth limiting density, so the harmonic Ritz value and Lanczos-OR based approaches seem to produce better approximations.
Note that these approximations differ from the KPM approximations in \cref{chap:random_quadrature} in that, like Gaussian quadrature, they adapt automatically to the spectrum of \( \vec{A} \).


We note that it is not typically possible to pick $\vec{v}$ with equal projection onto each eigencomponent since the eigenvectors of $\vec{A}$ are unknown. 
However, by choosing $\vec{v}$ from a suitable distribution, it can be guaranteed that $\vec{v}$ has roughly equal projection onto each eigencomponent.
This is discussed thoroughly in \cref{chap:random_quadrature}.


\begin{figure}[ht]
    \begin{center}
        \includegraphics{imgs/ch6_sign_spec.pdf} 
    \end{center}
    \caption[{Comparison of Lanczos-based spectrum approximation algorithms.}]{%
    Comparison of Lanczos-based spectrum approximation algorithms.
    \hspace{.25em}\textit{Legend}:
    Lanczos-OR induced approximation
    ({\protect\raisebox{0mm}{\protect\includegraphics[]{imgs/legend/l1nm.pdf}}}), 
    Lanczos-FA (GQ) ({\protect\raisebox{0mm}{\protect\includegraphics[]{imgs/legend/l2nm.pdf}}}),
    harmonic Ritz values based approximation
    ({\protect\raisebox{0mm}{\protect\includegraphics[]{imgs/legend/l3nm.pdf}}}), and
    limiting density
    ({\protect\raisebox{0mm}{\protect\includegraphics[]{imgs/legend/l4nm.pdf}}}).
    \hspace{.25em}\emph{Takeaway}: The Lancos-OR and harmonic Ritz value based approximations produce smooth approximations to the spectral density.
    }
    \label{fig:ch6_sign_spec}
\end{figure}

\subsubsection{Quality of approximation}

We now study how the number of matrix vector products impact the quality of approximation for a fixed sign function.

We construct a matrix with $400$ eigenvalues, 100 of which are the negatives of the values of a model problem \cref{eqn:model_problem} with parameters $\kappa = 10^2$, $\rho=0.9$, and $n=100$ and 300 of which are the values of a model problem with parameters $\kappa = 10^3$, $\rho=0.8$, $n=300$.
We then compute the Lanczos-OR induced approximation, the Lanczos-FA approximation, the harmonic Ritz value based approximation from \cite{eshof_frommer_lippert_schilling_van_der_vorst_02}, and the optimal $\vec{A}^2$-norm approximation to the matrix sign function.
The results are shown in \cref{fig:ch6_sign}.
In all cases, we use the Lanczos algorithm with full reorthogonalization.
Because eigenvalues of $\vec{T}$ may be near to zero, Lanczos-FA exhibits oscillatory behavior
On the other hand, the Lanczos-OR based approach and the harmonic Ritz value based approach have much smoother convergence.
Note that the Lanczos-OR induced approximation is not optimal, although it seems to perform close to optimally after a few iterations.

\begin{figure}[ht]
    \begin{center}
        \includegraphics{imgs/ch6_sign.pdf} 
    \end{center}
    \caption[{Optimality ratio for $\vec{A}^2$-norm errors for approximating $\operatorname{sign}\A\vec{v}$.}]{%
    Optimality ratio for $\vec{A}^2$-norm errors for approximating $\operatorname{sign}\A\vec{v}$.
    \hspace{.25em}\textit{Legend}:
    Lanczos-OR induced approximation
    ({\protect\raisebox{0mm}{\protect\includegraphics[]{imgs/legend/l1.pdf}}}), 
    Lanczos-FA (GQ) ({\protect\raisebox{0mm}{\protect\includegraphics[]{imgs/legend/l2.pdf}}}),
    harmonic Ritz values based approximation
    ({\protect\raisebox{0mm}{\protect\includegraphics[]{imgs/legend/l3.pdf}}})
    optimal
    ({\protect\raisebox{0mm}{\protect\includegraphics[]{imgs/legend/l4.pdf}}}). 
    \hspace{.25em}\textit{Takeaway}: The Lanczos-OR induced approximation to the matrix sign function performs well.
    }
    \label{fig:ch6_sign}
\end{figure}



\subsection{Rational matrix functions}

We now illustrate the effectiveness of the Lanczos-OR based approach to approximating rational matrix functions described in \cref{sec:rational_func}.


\subsubsection{Sign function}

In this example, we use the same spectrum as in the first example.
However, rather than approximating the sign function directly, we instead use Lanczos-OR to approximate each term of a proxy rational function of the form \cref{eqn:rat_form}.
In particular, we consider the best uniform approximation\footnote{Note that the eigenvalues of $\vec{A}$ live in $[-10^2,-1]\cup[1,10^3]$, so we could have used an asymmetric approximation to the sign function.
This would reduce the degree of the rational function required to obtain an approximation of given accuracy, but the qualitative behavior of Lanczos-OR-lm would not change substantially.} of degree $(39,40)$ to the sign function on $[-10^3,1]\cup[1,10^3]$.
Such an approximation is due to Zolotarev \cite{zolotarev_77}, an can be derived from the more well known Zolotarev approximation to the inverse square root function on $[1,10^6]$.
Our implementation follows the partial fractions implementation in the Rational Krylov Toolbox \cite{berljafa_elsworth_guttel_20} and involves computing the sum of $20$ terms of degree $(1,2)$.
The/ results are shown in \cref{fig:ch6_sign_rat}.

\begin{figure}[ht]
    \begin{center}
        \includegraphics{imgs/ch6_sign_rat.pdf} 
    \end{center}
    \caption[{$\vec{A}^2$-norm error in Lanczos-OR-lm based rational approximation to matrix sign function.}]{%
    $\vec{A}^2$-norm error in Lanczos-OR-lm based rational approximation to matrix sign function.
    \hspace{.25em}\textit{Legend}:
    Lanczos-OR-lm based approximation of matrix sign function with
    ({\protect\raisebox{0mm}{\protect\includegraphics[]{imgs/legend/l3.pdf}}})
    and without 
    ({\protect\raisebox{0mm}{\protect\includegraphics[]{imgs/legend/l1.pdf}}}), 
    reorthogonalization.
    Lanczos-OR-lm based approximation of proxy rational matrix function with
    ({\protect\raisebox{0mm}{\protect\includegraphics[]{imgs/legend/l4.pdf}}})
    and without 
    ({\protect\raisebox{0mm}{\protect\includegraphics[]{imgs/legend/l2.pdf}}})
    reorthogonalization
    \hspace{.25em}\textit{Takeaway}: Lanczos-OR-lm can be applied to each term of a proxy rational function approximation of the sign function.
    }
    \label{fig:ch6_sign_rat}
\end{figure}


At least while while the application error for the Lanczos-OR approximation to the proxy rational matrix function is large relative to the approximation error, then as seen in \cref{eqn:triangle_ineq}, the error in approximating the matrix sign function is similar to the error in approximating the proxy rational matrix function.
However, as seen in \cref{eqn:reverse_triangle_ineq}, the final accuracy of approximating the matrix sign function is limited by the quality of the scalar approximation.

We also note that it really only makes sense to use Lanczos-OR-lm with a short-recurrence version of Lanczos, in which case the effects of a perturbed Lanczos recurrence are prevalent. 
In particular, as we noted in \cref{chap:intro} and discuss in detail in \cref{chap:finite_precision}, the algorithm encounters a delay of convergence as compared to what would happen with reorthogonalization.


\subsection{Lanczos-FA vs Lanczos-OR vs CG}
\label{ex:lanczos_msCG_comparison}

%
%All three approaches are use a rational function approximation to $f$, and we will 

%If the matrix $(a_i\vec{A}^2 + b_i \vec{A} + c_i \vec{I})$ is positive definite, as is often the case, then $\vec{x}_i = (a_i\vec{A}^2 + b_i \vec{A} + c_i \vec{I})^{-1} \vec{v}$ can be approximated using CG.


We now illustrate the effectiveness of the Lanczos-OR based approach to approximating rational matrix functions described in \cref{sec:rational_func}.
Then we compare an existing low-memory approach, called multishift CG, to the analogous approaches based on Lanczos-OR-lm and Lanczos-FA-lm.

Throughout this example, we will assume that $r$ is a rational function of the form
\begin{align}
\label{eqn:rat_form}
r = \sum_{i=1}^{m} \frac{A_ix^2+B_ix+C_i}{a_ix^2+b_ix+c_i}.
\end{align}
so that $r(\vec{A})\vec{b}$ has the form
\begin{align*}
    r(\vec{A})\vec{b} = \sum_{i=1}^{m} (A_i\vec{A}^2 + B_i\vec{A} + C_i\vec{I}) \vec{x}_i,
\end{align*}
where $\vec{x}_i$ is obtained by solving the linear system of equations $(a_i\vec{A}^2 + b_i \vec{A} + c_i \vec{I}) \vec{x}_i = \vec{b}$.
This is relatively general since any real valued rational function $r:\R\to\R$ with numerator degree smaller than denominator degree and only simple poles can be written in this form (in fact, this would be true even if $A_i=0$).
A range of rational functions of this form appear naturally; for instance by a quadrature approximation to a Cauchy integral formula representation of $f$ \cite{hale_higham_trefethen_08}.
Similar rational functions are seen in \cite{eshof_frommer_lippert_schilling_van_der_vorst_02,frommer_simoncini_09} and in the rational approximation tot the sign function given described in \cref{eqn:rational_func}.


In certain cases, the shift invariance of Krylov subspace can be used to simultaneously compute all of the $\vec{x}_i$ using the same number of matrix-vector products as would be required to approximate a single $\vec{x}_i$.
Specifically, suppose $(a_i\vec{A}^2 + b_i \vec{A} + c_i \vec{I})^{-1} \vec{v}$ can be written as $\vec{B}+z_i \vec{I}$ for all $i=1,\ldots, m$.
Then $\mathcal{K}_k(\vec{B}-z_i\vec{I},\vec{v})) = \mathcal{K}_k(\vec{B},\vec{v})$, so by constructing a single Krylov subspace $\mathcal{K}_k(\vec{B},\vec{v})$, one can compute all of the $\vec{x}_i$ and therefore $r\A\vec{v}$.
The resulting algorithms are typically called multishift-CG or multishift-MINRES \cite{eshof_frommer_lippert_schilling_van_der_vorst_02,frommer_simoncini_08,frommer_simoncini_08,guttel_schweitzer_21,pleiss_jankowiak_eriksson_damle_gardner_20}.
However, such an approach only works when $a_i = 0$ or $a_i = a$ and $b_i = b$, and in the latter case, matrix-vector products with $\vec{B}$ require two matrix-vector products with $\vec{A}$ and convergence depends only on the properties of $\vec{A}^2$ rather than $\vec{A}$.

Instead, we might apply Lanczos-FA-lm to compute individual terms of $r\A\vec{v}$.
However, if $(a_i\vec{A}^2 + b_i \vec{A} + c_i \vec{I})$ is indefinite, then Lanczos-FA may exhibit oscillatory behavior due to eigenvalues of $(a_i\vec{T}^2 + b_i \vec{T} + c_i \vec{I})$ near zero.
This may result in a breakdown of Lanczos-FA-lm similar to the breakdown which may be encountered by standard implementations of CG on indefinite linear systems.
Lanczos-OR-lm avoids such issues.

To highlight some of the tradeoffs between the algorithms, we construct several test problems by placing eigenvalues uniformly throughout the specified intervals.
In all cases, $\vec{v}$ has uniform weight onto each eigencomponent.
The outputs are computed using standard Lanczos, but we note that the spectrum and number of iterations are such that the behavior is quite similar to if full reorthgonalization were used. 
In particular, orthogonality is not lost since no Ritz value converges. 
The results of our experiments are shown in \cref{fig:msCG_squared}.

\begin{figure}[htb]
    \centering
    \includegraphics[]{imgs/ch6_lanczos_msCG_squared.pdf}
    \caption[{Comparison of $(\vec{A}^2+c\vec{I})$-norm errors for CG and Lanczos-FA for computing $(\vec{A}^{2}+c\vec{I})^{-1}\vec{v}$ with \( c=0.05 \).}]{%  
    Comparison of $(\vec{A}^2+c\vec{I})$-norm errors for CG and Lanczos-FA for computing $(\vec{A}^{2}+c\vec{I})^{-1}\vec{v}$ with \( c=0.05 \).  
    Here CG works with $\vec{A}^2+c\vec{I}$ and requires two matrix-vector products per iteration whereas Lanczos-FA works with $\vec{A}$ and requires just one.
    \hspace{.25em}\emph{Legend}: Lanczos-OR 
    ({\protect\raisebox{0mm}{\protect\includegraphics[]{imgs/legend/l1.pdf}}}), 
    Lanczos-FA 
    ({\protect\raisebox{0mm}{\protect\includegraphics[]{imgs/legend/l2.pdf}}}), and
    CG on squared system 
    ({\protect\raisebox{0mm}{\protect\includegraphics[]{imgs/legend/l3.pdf}}}). 
    \hspace{.5em}
    \emph{Left}: eigenvalues on $[1,10]$. 
    \emph{Middle}: eigenvalues on $[-1.5,-1]\cup[1,10]$.
    \emph{Right}: eigenvalues on $[-10,-1]\cup[1,10]$. 
    \hspace{.25em}\emph{Legend}: Optimal algorithms have many nice convegence properties.
    }
    \label{fig:msCG_squared}
\end{figure}

We consider approximations to $r = 1/(x^2+0.05)$ with eigenvalues spaced with increments of $0.005$ throughout $[1,10]$, $[-1.5,-1]\cup[1,10]$, and $[-10,-1]\cup[1,10]$ respectively.
For each  example, the condition number of $\vec{A}^2+0.05\vec{I}$ is roughly 100 and the eigenvalues of $\vec{A}^2+0.05 \vec{I}$ fill out the interval $[1,100.05]$.
As such we observe that multishift CG converges at a rate (in terms of matrix products with $\vec{A}$) of roughly $\exp(-k/\sqrt{\kappa(\vec{A}^2)}) = \exp(-k/\sqrt{100})$ on all of the examples.

In the first example, $\vec{A}$ is positive definite.
Here Lanczos-FA and Lanczos-OR converge similarly to CG on $\vec{A}$ at a rate of roughly $\exp(-2k/\sqrt{10})$, where $k$ is the number of matrix-vector products with $\vec{A}$.

In the next example $\vec{A}$ is indefinite.
The convergence of CG is unchanged, because CG acts on $\vec{A}^2 + c\vec{I}$, it is unable to ``see'' the asymmetry in the eigenvalues of $\vec{A}$.
While the convergence of Lanczos-FA and Lanczos-OR is slowed considerably, both methods converges more quickly than CG due to the asymmetry in the intervals to the left and the right of the origin.
The convergence of these methods is at a rate of roughly $\exp(-k/\sqrt{15})$, although the exact rate is more complicated to compute \cite{fischer_96,schiefermayr_11}.
We also note the the emergence of oscillations in the error curve of Lanczos-FA.

In the third example, the asymmetry in the eigenvalue distribution about the origin is removed, and Lanczos-FA and Lanczos-OR converge at a rate very similar to that of multishift CG.
Note that Lanczos-FA displays larger oscillations, since the symmetry of the eigenvalue distribution of $\vec{A}$ ensures that $\vec{T}$ has an eigenvalue at zero whenever $k$ is odd.
However, the size of the oscillations is regularized by the fact that $c>0$.





\chapter{Spectrum dependent bounds and \\a posteriori error estimates}
\label{chap:CIF}


Thus far, we have not rigorously justified why Lanczos-FA and other non-optimal Lanczos-based methods typically outperform explicit polynomial methods. 
In \cref{sec:cif_bound,sec:cif_results,sec:cif_examples}, we describe a general technique for bounding the error of Lanczos-based methods for matrix functions via a reduction to the error of Lanczos-FA used to solve a certain linear system of equations.
Since the error of Lanazos-FA on linear systems is well studied, this approach can be used to derive a priori error bounds as well as a posteriori error bounds and estimates for general functions.
The effectiveness of our approach is demonstrated by a range of numerical experiments.
Finally, in \cref{sec:CG_indefinite}, we discuss the error of Lanczos-FA on \emph{indefinite} linear systems where Lanczos-FA is not optimal.
These bounds explain why Lanczos-FA performs well in theory.
%In practice, however, such bounds are not computable unless the spectrum of \( \vec{A} \) is known.
%Fortunately, it is possible to compute a posteriori error bounds and estimates for the Lanczos-FA error on linear systems, for instance, by computing the residual norm.


\iffalse
\subsection*{A bound based on interlacing}
\label{ex:unif_discrete}


Recall \cref{thm:poly_unif} asserts
\begin{equation}
\label{eqn:poly_unif}
    \| f(\vec{A}) \vec{b} - \lan_k(f) \|_2 
    \leq  2 \min_{\deg p<k} \left( \| f-p \|_{\mathcal{I}(\vec{A})}  \right) \| \vec{b} \|_2.
\end{equation}
This bound depends only on the range of eigenvalues \(\mathcal{I}(\vec{A})\) and not on more fine-grained information like the presence of eigenvalue clusters or isolated eigenvalues, which are known to lead to faster convergence.
As we have seen many times up to this point, such bounds are too pessimistic to capture the true behavior of Lanczos based methods.

From the proof of \cref{thm:poly_unif} we also have that
\begin{equation}
    \label{eqn:triangle_ineq1}
    \| f(\vec{A}) \vec{b} - \lan_k(f) \|
    \leq \min_{\deg p< k} \left( \| f-p \|_{\Lambda(\vec{A})} + \| f - p \|_{\Lambda(\vec{T})} \right) \| \vec{b} \|. 
\end{equation}
This expression is more accurate but less useful; it cannot be used as an \emph{a priori} bound since it involves the eigenvalues of the tridiagonal matrix $\vec{T}$, which depend on $\vec{b}$ nor can it be used as a practical \emph{a posteriori} bound since it involves all  eigenvalues of \( \vec{A} \).


It is well known that the eigenvalues of \( \vec{T} \) are interlaced by those of \( \vec{A} \); that is, \( \Lambda(\vec{T}) \subset \mathcal{I}(\vec{A}) \) and between each pair of eigenvalues of \( \vec{T} \) is at least one eigenvalue of \( \vec{A} \).
With this property in mind, define \( \mathcal{J}_k(\vec{A}) \) as the set of all \( k \)-tuples \( \bm{\mu} = ( \mu_1 , \ldots , \mu_k ) \in \mathbb{R}^k \) that are interlaced by the eigenvalues of \( \vec{A} \). %; that is, \( \lambda_{\text{min}} ( \vec{A} ) < \min_j \mu_j \), \( \lambda_{\text{max}} ( \vec{A} ) > \max_j \mu_j \), and between each consecutive pair \( \mu_i , \mu_j \) there is at least one eigenvalue of \( \vec{A} \).
Then, we can use \cref{eqn:triangle_ineq1} to write
\begin{equation}
    \label{eqn:triangle_ineq2} 
\| f(\vec{A}) \vec{b} - \lan_k(f) \| 
 & \leq  \max_{\bm{\mu}\in \mathcal{J}_k(\vec{A})} \min_{\deg p<k} \left( \| f-p \|_{\Lambda(\vec{A})} + \| f - p \|_{\bm{\mu}} \right) \| \vec{b} \|. 
\end{equation}
The bound \cref{eqn:triangle_ineq2} is an {\em a priori} error bound, and at least in some special cases, provides more insight than \cref{eqn:poly_unif} in situations where the eigenvalues of $\vec{A}$ are clustered.


Consider \( \vec{A} \) with many eigenvalues uniformly spaced through the interval \( [0,1] \) and a single isolated eigenvalue at \( \kappa > 1 \).
Since the eigenvalues of \( \vec{T}_k \) are interlaced by those of \( \vec A \), there is at most one eigenvalue of \( \vec{T} \) between \( 1 \) and \( \kappa \); that is, \( \Lambda( \vec{A} )\cup\Lambda( \vec{T} ) \) is contained in \( [0,1] \cup \{ \mu , \kappa \} \) for some \( \mu\in[1,\kappa] \).
We then have
\begin{align}
\label{eqn:ex_minimax_bd}
    \| f(\vec{A}) \vec{b} - \lan_k(f) \| 
%\leq \max_{\mu\in[1,\kappa]} \min_{\deg p<k} \left( \max_{\lambda\in[0,1]\cup\{\mu,\kappa\}} | f(\lambda)-p(\lambda)| \right) \| \vec{b} \|.
\leq 2 \max_{\mu\in[1,\kappa]} \min_{\deg p<k} \left( \| f-p \|_{[0,1]\cup\{\mu,\kappa\}} \right) \| \vec{b} \|.
\end{align}
For \( \kappa = 5\), \( f(x) = \exp(-x) \), and \( k=6 \), we use a numerical optimizer to determine that the value  maximizing the right hand side of \cref{eqn:ex_minimax_bd} is  \( \mu^*\approx 4.96\).
In \cref{fig:lanc_poly_approx} we show the error of the Lanczos-FA polynomial along with the optimal uniform polynomial approximations to \( f \) on \( [0,5] \), which contains \( [0,1]\cup\{\mu^*,5\} \).
As expected, the bound from \cref{eqn:ex_minimax_bd}  is significantly better than that from the uniform approximation.

\begin{figure}[ht]
    \includegraphics[width=.97\textwidth]{CIF_imgs/lanc_poly_approx_diff.pdf}

   \caption{
    Comparison of errors of degree 5 polynomial approximations to \( f(x) = \exp(-x) \).
    \emph{Legend}:
    \emph{Remarks}: 
    Note that the Lanczos-FA approximation becomes very inaccurate on \( (1,5) \) which allows a smaller error on the eigenvalues of $\vec A$, which is the only error that impacts our approximation to $f(\vec A)\vec b$.
    As a result, the uniform approximation on \( [0,1] \cup \{\mu^*,5\} \) is a much better bound for the Lanczos-FA error than the uniform approximation on \( [0,5] \), which remains equally accurate over the entire interval \([0,5]\). 
    }
    \label{fig:lanc_poly_approx}
\end{figure}
\fi

\section{An integral representation of the Lanczos-FA error}%and the Cauchy integral formula}
\label{sec:cif_bound}

Assuming \( f : \mathbb{C} \rightarrow \mathbb{C} \) is analytic in a neighborhood of the eigenvalues of \( \vec{A} \) and \( \Gamma \) is a simple closed curve or union of simple closed curves inside that neighborhood and enclosing the eigenvalues of \( \vec{A} \), the Cauchy integral formula states that
\begin{equation}
    \label{eqn:fofACIF1}
    \fA\vec{v} = - \frac{1}{2\pi \ii} \oint_{\Gamma} f(z) (\vec{A} - z \vec{I} )^{-1} \vec{v} \, \d{z}.
\end{equation}
If \( \Gamma \) also encloses the eigenvalues of \( \vec{T} \) we can similarly write the Lanczos-FA approximation as
\begin{equation}
\label{eqn:fofT_CIF}
%    \vec{Q} \fT \vec{Q}^\cT \vec{v}
    \lan_k(f)
    = - \frac{1}{2\pi \ii} \oint_{\Gamma} f(z) \vec{Q} (\vec{T} - z \vec{I} )^{-1} \vec{Q}^\cT \vec{v} \, \d{z} .
\end{equation}
Observing that the integrand of \cref{eqn:fofACIF1} contains the solution to the shifted linear system \( (\vec{A}-z\vec{I}) \vec{x} = \vec{v} \) while  \cref{eqn:fofT_CIF} contains the Lanczos-FA approximation to the solution, we make the following definition.
\begin{definition}\label{def:err}
For \( z \in \mathbb{C} \), %let \( h_z (x) = 1/(x-z) \).
define the \( k \)-th Lanczos-FA error and residual for the linear system \( (\vec{A}-z\vec{I}) \vec{x} = \vec{v} \) as,
\begin{align*}
    \err_k(z,\vec{A},\vec{v}) &:= (\vec{A} - z \vec{I})^{-1} \vec{v} - \vec{Q}(\vec{T}-z\vec{I})^{-1}\vec{Q}^\cT \vec{v}%\lan_k ( h_z )
    ,\\
    \Res_k(z,\vec{A},\vec{v}) &:= \vec{v} - (\vec{A} - z \vec{I}) \vec{Q}(\vec{T}-z\vec{I})^{-1}\vec{Q}^\cT \vec{v}.%\,\lan_k ( h_z ).
\end{align*}
As with the Lanczos-FA approximation, we will typically omit the arguments \( \vec{A} \) and \( \vec{v} \), and in the case \( z = 0 \), we will often write \( \err_k \) and \( \Res_k \).
\end{definition}
%\Tyler{Should we mention CG at some point?} \Anne{No need.}
With \cref{def:err} in place, the error of the Lanczos-FA approximation to \( \fA\vec{v} \) can be written as
\begin{equation}
    \label{eqn:lanczos_FA_error}
    \fA \vec{v} - \lan_k(f)%\vec{Q} f( \vec{T} ) \vec{Q}^{\cT} \vec{v}
    = - \frac{1}{2\pi \ii} \oint_{\Gamma} f(z) \, \err_k(z) \, \d{z}.
\end{equation}
Therefore, if for every $z\in\Gamma$ we are able to understand the convergence of Lanczos-FA on the linear system \( (\vec{A}-z\vec{I}) \vec{x} = \vec{v} \), then this formula lets us understand the convergence of Lanczos-FA for \( \fA\vec{v} \).
To simplify bounding \cref{eqn:lanczos_FA_error}, we will write $\err_k(z)$ for all  \( z\in\Gamma \) in terms of the error in solving a single shifted linear system.

%We now turn our attention to the shifted linear system \( (\vec{A}-z \vec{I}) \vec{y} = \vec{v} \), where \( z \) is arbitrary and possibly complex.
To do this, we use the fact that the Lanczos factorization \cref{eqn:lanczos_three_term} can be shifted, even for complex \( z \), to obtain
\begin{equation}
    \label{eqn:shifted_lanczos_factorization}
    ( \vec{A} - z \vec{I} ) \vec{Q}
    = \vec{Q} ( \vec{T} - z \vec{I} ) + \beta_{k-1} \vec{q}_{k} \vec{e}_{k-1}^\rT.
\end{equation}
That is, Lanczos applied to \( (\vec{A},\vec{v}) \) for \( k \) steps produces output \( \vec{Q} \) and \( \vec{T} \) satisfying \cref{eqn:lanczos_three_term} while Lanczos applied to \( (\vec{A} - z \vec{I}, \vec{v}) \) for \( k \) steps produces output \( \vec{Q} \) and \( \vec{T}-z \vec I \) satisfying \cref{eqn:shifted_lanczos_factorization}. %, where \( \vec{Q} \) and \( \vec{T} \) are the same in both factorizations. \Cam{Are these the ones produced by iterating on $\vec A$ right? We should say explicitly.}
%In the case that \( z \) is not real, we interpret the inner products in the Lanczos method as complex inner products linear in the first argument; i.e. \( \langle a \vec{x} , \vec{y} \rangle = a \langle \vec{x}, \vec{y} \rangle \) for all \( a \in \mathbb{C} \).
Using this fact, we have the following well known lemma.
\begin{lemma}
    \label[lemma]{thm:shifted_lanczos_equivalence}
For all \( z \) where \( \vec{T} - z \vec{I} \) is invertible, 
\begin{equation*}
    \Res_k(z) 
    = \| \vec{v} \|_2 \left( \frac{(-1)^{k}}{\det(\vec{T} -z \vec{I}) }\prod_{j=0}^{k-1} \beta_j \right) \vec{q}_{k}.
\end{equation*}
\end{lemma}
\begin{proof}
    From \cref{eqn:shifted_lanczos_factorization}, and the fact that $\vec{Q}$'s first column is $\vec{v}/\|\vec{v}\|_2$, it is clear that,
    \begin{align*}
        (\vec{A} - z\vec{I}) \vec{Q}(\vec{T}-z\vec{I})^{-1}\vec{Q}^\cT \vec{v}
        &= (\vec{A}-z \vec{I}) \vec{Q} (\vec{T} - z \vec{I})^{-1} \| \vec{v} \|_2 \vec{e}_0
        \\&= \vec{Q} \| \vec{v} \|_2 \vec{e}_0 + \beta_k \vec{q}_{k} \vec{e}_{k-1}^\cT (\vec{T} - z \vec{I})^{-1} \| \vec{v} \|_2 \vec{e}_0
        \\&= \vec{v} + \beta_k \vec{q}_{k} \vec{e}_{k-1}^\cT (\vec{T} - z \vec{I})^{-1} \| \vec{v} \|_2 \vec{e}_0.
    \end{align*}

    Using the formula \( (\vec{T} - z \vec{I} )^{-1} = (1/\!\det(\vec{T} - z \vec{I})) \operatorname{adj}( \vec{T} - z \vec{I} ) \), we see that
    \begin{equation*}
        \vec{e}_{k-1}^\cT (\vec{T} - z \vec{I})^{-1} \vec{e}_0
        = \frac{(-1)^{k-1}}{\det( \vec{T} - z \vec{I} )}\prod_{j=0}^{k-2} \beta_j.
    \end{equation*}
    The result then follows by combining these expressions.
\end{proof}
We use \cref{thm:shifted_lanczos_equivalence} to relate \( \err_k(z) \) to \( \err_k(w) \) for any $z,w \in \mathbb{C}$.


\begin{definition}
    \label{def:hwz}
For $w,z\in\mathbb{C}$ define $h_{w,z}:\R\to\mathbb{C}$ and $h_z:\R\to\mathbb{C}$ by
\begin{equation*}
    h_{w,z}(x) := \frac{x-w}{x-z}
    ,\qquad
    h_z(x):= \frac{1}{x-z}.
\end{equation*}
\end{definition}

\begin{corollary} 
    \label[corollary]{thm:shifted_linear_system_error}
For all \( z , w \in \mathbb{C} \), where \( \vec{A} - z \vec{I} \) and \( \vec{A} - w \vec{I} \) are both invertible,
%\( \vec{A} - b \vec{I} \) is invertible. Then,
\begin{align*}
    \err_k(z)
    &= \det(h_{w,z}(\vec{T})) \,h_{w,z}(\vec{A}) 
    \,\err_k(w)
    \\
    \Res_k(z)
    &= \det(h_{w,z}(\vec{T}))
    \,\Res_k(w).
\end{align*}
\end{corollary}

\begin{proof}
By \cref{thm:shifted_lanczos_equivalence},
\begin{equation*}
    \det(\vec{T} - z\vec{I}) \,\Res_k (z) = \det(\vec{T} - w \vec{I}) \,\Res_k (w). 
\end{equation*}
Thus,
\begin{equation*}
    \Res_k (z) = \frac{\det(\vec{T} - w\vec{I})}{\det(\vec{T} - z\vec{I})} \,\Res_k (w) = \det(h_{w,z}(\vec{T})) \,\Res_k (w) .
%    \frac{\det(\vec{T}-w\vec{I})}{\det(\vec{T} - z\vec{I})} (\vec{A} - w\vec{I}) \err_k (w)
%    = (\vec{A} - z \vec{I}) \err_k(z)
\end{equation*}
Noting that \( \Res_k (z) = (\vec{A}-z\vec{I}) \,\err_k(z) \) and 
\( \Res_k (w) = (\vec{A}-w\vec{I}) \,\err_k(w) \), we obtain the relation between the
errors,
%Thus,
\begin{align*}
    \err_k(z)
    &= \det(h_{w,z}(\vec{T})) (\vec{A} - z \vec{I})^{-1} ( \vec{A} - w \vec{I})\, \err_k ( w ) 
    \\&= \det(h_{w,z}(\vec{T}))  \,h_{w,z}(\vec{A}) \,\err_k(w) .
    \tag*{\qedhere}
\end{align*}
\end{proof}

In summary, combining \cref{eqn:lanczos_FA_error} and \cref{thm:shifted_linear_system_error} we have the following corollary.
This result is by no means new, and appears throughout the literature; see for instance \cite{frommer_simoncini_09} and \cite[Theorem 3.4]{frommer_guttel_schweitzer_14}.
\begin{corollary}
\label[corollary]{thm:integral_error_vec} 
Suppose \( \vec{A} \) is a Hermitian matrix and \( f : \mathbb{C} \rightarrow \mathbb{C} \) is a function analytic in a neighborhood of the eigenvalues of \( \vec{A} \) and \( \vec{T} \), where  \( \vec{T} \) is the tridiagonal matrix output by Lanczos run on \( \vec{A},\vec{v} \) for \( k \) steps.
Then, if\/ \( \Gamma \) is a simple closed curve or union of simple closed curves inside this neighborhood and enclosing the eigenvalues of \( \vec{A} \) and \( \vec{T} \) and \( w\in\mathbb{C} \) is such that \( w \not\in \Lambda(\vec{T})\cup\Lambda \),
\begin{equation*}
    \fA \vec{v} - \lan_k(f)
    = \left( - \frac{1}{2\pi \ii} \oint_{\Gamma} f(z) \det(h_{w,z}(\vec{T})) \,h_{w,z}(\vec{A}) \,\d{z} \right) \, \err_k(w).
\end{equation*}
\end{corollary}
%\Cam{Should we make the above equation a corrloary?}

\subsection{A reduction to linear system error}

Our main result is a flexible bound for the  Lanczos-FA error, obtained by bounding the integral in the right-hand side of \cref{thm:integral_error_vec}. As we will see in \cref{sec:cif_results}, we can instantiate this theorem to obtain effective a priori and a posteriori error bounds in many settings.
% the Cauchy integral formula \cref{eqn:lanczos_FA_error}, and bounding each $\err_k(z)$ term with $\err_k(w)$ for a single choice of $w$.
%Given a scalar function \( h:\R\to\R \) and set \( S\subset \R \), we denote the infinity norm of \( h \) on \( S \) by \( \|h\|_S \); i.e. \( \|h\|_S := \sup_{x\in S}| h(x)| \).
%Using this notation, we then have the following theorem.
\begin{theorem}
\label{thm:err_int}
In the setting of \cref{thm:integral_error_vec}, if for some $S,S_0,\ldots,S_{k-1} \subset \mathbb{R}$ we have \( \Lambda\subset S_0\ \) and \( \lambda_i(\vec{T}) \in S_i \) for  \( i=0,\ldots, k-1 \), then
\begin{equation*}
    \| \fA\vec{v} - \lan_k(f) \|
    %\leq \underbrace{\vphantom{ \bigg| }\left( \frac{1}{2\pi}\oint_{\Gamma} |f(z)| \cdot \left(\prod_{i=0}^{k-1} \|h_{w,z}\|_{S_i}\right) \cdot \|h_{w,z}\|_{S_0} \cdot | \d{z} | \right)}_{\text{integral term}}  \hspace{-1.2 em}\underbrace{\vphantom{ \Bigg| } \| \err_k(w) \| . \hspace{-.4em} }_{\text{linear system error}} \hspace{-.5em}
    \leq \underbrace{\vphantom{ \Bigg| }\Bigg( \frac{1}{2\pi}\oint_{\Gamma} |f(z)| \left(\prod_{i=0}^{k-1} \| h_{w,z}\|_{S_i}\right) \|h_{w,z}\|_{S} | \d{z} | \bigg)}_{\text{integral term}} \hspace{-.2em}\underbrace{\vphantom{ \Bigg| } \| \err_k(w) \|.}_{\text{linear system error}}
\end{equation*}
\end{theorem}
Analogously, we have a bound for Gaussian quadrature
\begin{theorem}
\label{thm:err_int_GQ}
In the setting of \cref{thm:integral_error_vec}, if for some $S,S_0,\ldots,S_{k-1} \subset \mathbb{R}$ we have \( \Lambda\subset S_0\ \) and \( \lambda_i(\vec{T}) \in S_i \) for  \( i=0,\ldots, k-1 \), then
\begin{equation*}
    | \vec{v}^\cT\fA\vec{v} - \int f \,\d\qq[g]{\Psi}{2k-1} |
    \leq  \underbrace{\vphantom{\Bigg|}\Bigg( \frac{1}{2\pi} \oint_{\Gamma} | f(z) |  \left( \prod_{i=0}^{k-1} \|h_{w,z}\|_{S_i}^2  \right) \| h_z \|_{S} |\d{z}| \Bigg)}_{\text{integral term}} \hspace{-.2em}\underbrace{\vphantom{\Bigg|} \| \Res_k(w) \|_2^2.}_{\text{linear system error}}  \label{eqn:err_quad_int}
\end{equation*}
\end{theorem}


The above bounds depend on our choices of $\Gamma$, $w$, and the sets $S,S_0,\ldots,S_{k-1}$, which must contain the eigenvalues of $\vec A$ and $\vec T_k$. 
The sets $S,S_0,\ldots,S_{k-1}$ should be chosen based on the information we have about $\vec A$ and $\vec T_k$. 
For example, we could take all these sets to be the eigenvalue range $\mathcal{I}( \vec A)$. 
If we have more information a priori about the eigenvalues of $\vec A$, we can obtain a tighter bound by choosing smaller $S$, with correspondingly lower $\|h_{w,z}\|_{S}$. 
For an a posteriori bound, we can simply set $S_i = \{ \lambda_i(\vec T_k) \}$, for $i = 0,\ldots,k-1$. 
This gives an optimal value for $\|h_{w,z}\|_{S}$. Both approaches are detailed in \cref{sec:cif_results}.

We emphasize that, in both bounds, the integral term and linear system error term in the theorem are entirely decoupled.
Thus, once the integral term is computed, bounding the error of Lanczos-FA for \( \fA\vec{v} \) is reduced to bounding \( \| \err_k(w) \| \), and if the integral term can be bounded independently of \( k \), \cref{thm:err_int} implies that, up to a constant factor, the Lanczos-FA approximation to \( \fA\vec{v} \) converges at least as fast as \( \| \err_k(w) \| \).

Note that \cref{thm:err_int} depends on $\|\err_k(w)\|$ whereas \cref{thm:err_int_GQ} depends on $\|\Res_k(w)\|_2^2$. 
Thus, heuristically, we can expect the quadratic form to converge at a rate twice that of the norm of the error of the matrix function.
This is exacted as Gaussian quadrature is exact for polynomials of degree \( 2k-1 \) whereas Lanczos-FA is exact for polynomials of degree \( k-1 \). 
In the case that the contour $\Gamma$ does not pass through \( \mathcal{I} \), the bound of \cref{eqn:err_quad_int} is essentially  as easy to compute as that of  \cref{thm:err_int}.
However, if the contour passes through \( \mathcal{I} \) at \( w \), to ensure that \( S \) does not contain points in the contour, it must be chosen as a set other than $\mathcal{I}$. This set must contain all of $\vec A$'s eigenvalues and we must bound its distance to the contour (in particular, to $w$).

\begin{proof}[Proof of \cref{thm:err_int}]
    We begin by taking the norm on both sides of \cref{thm:integral_error_vec}.
    Applying the triangle inequality for integrals and using the fact that \( \|\,\cdot\,\| \) is induced by a matrix with the same eigenvectors as \( \vec{A} \) (see \cref{thm:norm_exchange}) we have
\begin{equation}
    \label{eqn:integral_error}
    \| \fA \vec{v} - \lan_k(f) \|
    \leq \left( \frac{1}{2\pi}\oint_{\Gamma} |f(z)|\mkern1mu  |\!\det(h_{w,z}(\vec{T}))|  \| h_{w,z}(\vec{A}) \|_2  |\d{z} | \right) \| \err_k(w) \| .
\end{equation}
Next, since %as an immediate consequence of the definition of \( \| h_{w,z} (\vec{A}) \|_2 \), if 
\( \Lambda \subseteq S \) then
\begin{equation*}
%    \label{eqn:hwz_qwz}
    \| h_{w,z}(\vec{A}) \|_2 = \max_{i=0,\ldots, n-1} | h_{w,z}(\lambda_i(\vec{A})) | \leq \|h_{w,z}\|_{S},
\end{equation*}
and similarly, if  \( \lambda_i(\vec{T}) \in S_i \) for \( i=0, \ldots , k-1 \), then 
\begin{equation}
    \label{eqn:dkwz_qwz}
    |\!\det(h_{w,z}(\vec{T}))| = \left| \prod_{i=0}^{k-1} h_{w,z}(\lambda_i(\vec{T})) \right| \leq
    \prod_{i=0}^{k-1} \|h_{w,z}\|_{S_i}.
\end{equation}
Combining these inequalities yields the result.
\end{proof}
%root-condition number bound on favorable spectra. 


\begin{proof}[Proof of \cref{thm:err_int_GQ}]
Recall
    \begin{equation*}
    \vec{v}^\cT \lan_k(f) 
    = \vec{v}^\cT \vec{Q} \fT \vec{Q}^\cT \vec{v}
    = \| \vec{v} \|_2^2\i: \vec{e}_0^\cT \fT \vec{e}_0
    = \int f \,\d\qq[g]{f}{2k-1}.
\end{equation*}

Since \( \vec{A} \) is Hermitian, \( (\vec{A} - z\vec{I})^\cT = \vec{A} - \overline{z} \vec{I} \). 
Thus, since 
\begin{equation*}
    \vec{v}^\cT (\vec{A} - z\vec{I})^{-1}
    = ( ( \vec{A} - \overline{z} \vec{I}   )^{-1} \vec{v} )^\cT
    = ( \lan_k(h_{\overline{z}}) + \err_k(\overline{z}) ) \vec{v} )^\cT
\end{equation*}
we can expand the quadratic form error as
\begin{equation*}
    \vec{v}^\cT \err_k(z)
    = \vec{v}^\cT ( \vec{A} - z \vec{I})^{-1} \Res_k(z)
    = \left( \lan_k ( h_{\overline{z}} ) ) + \err_k(\overline{z}) \right)^\cT \Res_k(z).
\end{equation*}
Now, by definition, \( \lan_k ( h_{\overline{z}}(x) )  = \vec{Q} h_{\overline{z}}(\vec{T}) \vec{Q}^\cT \vec{v} \) and by \cref{thm:shifted_lanczos_equivalence} \( \Res_k(z) \) is proportional to \( \vec{q}_{k+1} \).
Thus, since, at least in exact arithmetic, \( \vec{q}_{k+1} \) is orthogonal to $\vec{Q}$,
\begin{equation*}
    \vec{v}^\cT \err_k(z)
    = \err_k(\overline{z})^\cT \Res_k(z)
    = (( \vec{A} - \overline{z}\vec{I})^{-1}  \Res_k(\overline{z}))^\cT \Res_k(z).
\end{equation*}
Next, using \cref{thm:shifted_linear_system_error} and the fact that \( h_{w,z}(x) h_{w,\overline{z}}(x) = |h_{w,z}(x)|^2 \) for \( w,x\in\R \), 
\begin{equation*}
    \vec{v}^\cT \err_k(z)
    = |\!\det(h_{w,z}(\vec{T}))|^2 \Res_k(w)^\cT ( \vec{A} - z\vec{I})^{-1} \Res_k(w).
\end{equation*}

We then have,
\begin{equation*}
    | \vec{v}^\cT \err_k(z) |
    \leq | \! \det(h_{w,z}(\vec{T})) |^2  \| ( \vec{A} - z\vec{I})^{-1} \|_2  \| \Res_k(w) \|_2^2.
\end{equation*}
\iffalse
which gives
\begin{equation*}
    | \vec{v}^\cT \fA \vec{v} - \vec{v}^\cT \lan_k(f) |
    &\leq \left( \frac{1}{2\pi} \oint_{\Gamma} | f(z) | \cdot |\det(h_{w,z}(\vec{T}))|^2 \cdot \| h_z(\vec{A}) \|_2 \cdot |\d{z}| \right) \| \Res_k(w) \|_2^2.
\end{equation*}
\fi
Applying the Cauchy integral formula we therefore obtain a bound for the quadratic form error analogous to \cref{thm:err_int} we obtain the result.
\end{proof}


\subsection{Comparison with previous work}
%\Tyler{Is this okay here? If we should move it back to after section 2, I think it should be its own section and perhaps beefed up a bit.}

%\Cam{I think before this section we need to at least say what our framework is, outside the abstract. Should we put this at the end of section 2? Otherwise we keep comparing things prior work does in comparision to what we do, but we haven't talked about our approach at all yet.}

Our framework for analyzing Lanczos-FA has four properties which differentiate it from past work:
(i) it is applicable to a wide range of functions, 
(ii) it yields a priori bounds dependent on fine-grained properties of the spectrum of \( \vec{A} \) such as clustered or isolated eigenvalues, 
(iii) it can be used a posteriori as a practical stopping criterion,
and (iv) it is applicable when computations are carried out in finite precision arithmetic. %when used with a perturbed Lanczos recurrence.
To the best of our knowledge, no existing analysis satisfies more than two of these properties simultaneously.
In this section, we provide a brief overview of the most relevant past work.

Most directly related to our framework is a series of works which also make use of the shift-invariance of Krylov subspaces when \( f \) is a Stieltjes function\footnote{A function \( f \) defined on the positive real axis is a Stieltjes function if and only if \( f(x) \geq 0 \) for all $x\in\R$ and \( f \) has an analytic extension to the cut plane \( \mathbb{C} \setminus (-\infty,0] \) satisfying \( \Im(f(x)) \leq 0 \) for all \( x \) in the upper half plane \cite[Theorem 3.2]{berg_07} \cite[p. 127 attributed to Krein]{aheizer_65}.} \cite{frommer_guttel_schweitzer_14a,frommer_schweitzer_15,ilic_turner_simpson_09} or a certain type of rational function \cite{frommer_kahl_lippert_rittich_13,frommer_simoncini_08b,frommer_simoncini_09}.
These analyses are applicable a priori and a posteriori and in fact allow for corresponding error \emph{lower bounds} as well.
However, these bounds cannot be applied to more general functions, and the impact of a perturbed Lanczos recurrence in finite precision is not considered.

The most detailed generally applicable analysis is \cite{musco_musco_sidford_18}, which extends \cite{druskin_knizhnerman_91,druskin_knizhnerman_95} and studies \cref{thm:poly_unif}, the classical bound for Lanczos-FA based on polynomial approximation on \( \mathcal{I} \), when Lanczos is run in finite precision arithmetic.
However, as we have seen throughout this thesis, \cref{thm:poly_unif} is often too pessimistic in practice as it does not depend on the fine-grained properties about the distribution of eigenvalues.
Another generally applicable analysis is \cite{hochbruck_lubich_selhofer_98}, which suggests replacing \( \err_k(z) \) with \( \Res_k(z) \) in \cref{eqn:lanczos_FA_error}.
Since \( \Res_k(z) \) can be computed once the outputs of Lanczos have been obtained, the resulting integral can be computed (or at least approximated by a quadrature rule).
However, this approach does not take into account the actual relationship between \( \Res_k(z) \) and \( \err_k(z) \), and therefore gives only an estimate of the error, not a true bound.
Another Cauchy integral formula based approach is \cite{hochbruck_lubich_97} which shows that Lanczos-FA exhibits superlinear convergence for the matrix exponential and certain other specific analytic functions.

There are a variety of other bounds specialized to individual functions.
For example, it is known that if \( \vec{A} \) is nonnegative definite and \( t > 0 \), then the error in the Lanczos-FA approximation for the matrix exponential \( \exp( t \vec{A}) \vec{v} \) can be related to the maximum over \( s \in [0,t] \) of the error in the  optimal approximation to \( \exp( s \vec{A}) \vec{v} \) over a Krylov space of slightly lower dimension \cite{druskin_greenbaum_knizhnerman_98}.
More recent work involving the matrix exponential are \cite{jia_lv_14,jawecki_auzinger_koch_19,jawecki_21}.
%in exact and finite precision arithmetic \cite{druskin_greenbaum_knizhnerman_98}.
There is also a range of work which analyzes the convergence of Lanczos-FA and related methods for computing the square root and sign functions \cite{boricci_99,borici_03,eshof_frommer_lippert_schilling_van_der_vorst_02}.
%Further results for the matrix exponential \cite{orecchia_sachdeva_nisheeth_12}.

\section{Applying our framework}
\label{sec:cif_results}

We proceed to show how to effectively bound the integral term of \cref{thm:err_int}, to give  a priori and a posteriori  bounds on the Lanczos-FA error, assuming accurate bounds on \( \| \err_k(w) \| \) are available.
Throughout this chapter, we assume $w\in\mathbb{R}$ and we do not discuss in detail how to bound this linear system error  -- there are many known approaches, both a priori and a posteriori, and the best bounds to use are often context dependent. 
Some of these approaches are similar to those used for Lanczos-OR in \cref{sec:error_OR}.


To use \cref{thm:err_int}, we must evaluate or bound \( \|h_{w,z}\|_{S_i} \).
Towards this end, we introduce the following lemmas, which apply when  $S_i$ is an interval.
These lemmas are also useful when $S_i$ is a union of intervals -- in that case $\|h_{w,z}\|_{S_i}$ is bounded by the maximum bound on any of these intervals.
i
\begin{lemma}
    \label[lemma]{thm:Q_wz_value}
    For any interval \( [a,b] \subset \mathbb{R} \), if \( z \in \mathbb{C} \setminus [a,b] \) and \( w\in\mathbb{R} \), we have
\begin{equation*}
    \|h_{w,z}\|_{[a,b]}
    %Q_{w,z} = \max_{x \in \mathcal{I} } \left| \frac{x-w}{x-z} \right| 
    = \max \left\{ 
    \left| \frac{a-w}{a-z} \right|, 
    \left| \frac{b-w}{b-z} \right|,
    \left( \left| \frac{z-w}{\Im(z)} \right|~\text{ if } x^{*} \in [a,b] ~\text{else}~0 \right)
    \right\}
\end{equation*}
where 
\begin{equation*}
    x^* := \frac{\Re(z)^2 + \Im(z)^2 - \Re(z) w}{\Re(z)-w}.
\end{equation*}

\end{lemma}
\begin{proof}
    Note that for $x\in\mathbb{R}$, 
\begin{equation*}
    | h_{w,z} (x) |^2 = \left| \frac{x-w}{x-z} \right|^2 = \frac{(x-w )^2}{( x - \Re(z) )^2 + \Im(z)^2} ,
\end{equation*}
and
\begin{equation*}
    \frac{\d}{\d{x}} \left( | h_{w,z} (x) |^2 \right) = \frac{[ (x - \Re(z) )^2 + \Im(z )^2 ]2(x-w) - (x-w )^2 2 ( x - \Re(z) )}{[ (x - \Re(z) )^2 + \Im(z )^2 ]^2} .
\end{equation*}
%Define 
%\begin{equation*}
%    g(x) = \frac{x-w}{\sqrt{(\Re(z)-x)^2 + \Im(z)^2}}.
%\end{equation*}
%    g'(x) = \frac{\Im(z)^2 + (\Re(z)-w)(\Re(z)-x)}{(\Re(z)-x)^2 + \Im(z)^2}.
Aside from \( x=w \), where \( h_{w,z} (x) = 0 \), the only value \( x \in \R \) for which \( \frac{\d}{\d{x}} \left( | h_{w,z} (x) |^2 \right) = 0 \) is \( x^* \).
%is the only value \( x\in \R \) for which \( g'(x) = 0 \).
This implies that the only possible local extrema of  \( | h_{w,z} (x)|  \) on \( [a,b] \) are  \( a \), \( b \), and \( x^* \) if \( x^* \in [a,b] \).
Substituting the expression for \( x^{*} \) into that for \( | h_{w,z} ( x^{*} ) | \), one finds, after some algebra, that \( | h_{w,z} ( x^{*} ) | = | z-w | / | \Im (z) | \).
%We have that \( | h_{w,z}(x) | = |g(x)| \)
%Moreover, \( g(x) \) only changes sign at \( x=w \) and \( g(w) = 0 \), 
%so the only possible local maxima of \( | h_{w,z}(x) | \) on \( [a,b] \) are also  \( %a \), \( b \), and \( x^* \) if \( x^* \in [a,b] \).
%The result follows from the fact that \( | g(x^*) | = |z-w|/|\Im(z)| \).
\iffalse
\Tyler{the last sentence could be replaced by the full derivation:
It remains to show that \( | h_{w,z} (x^*) | = |z-w|/|\Im(z)| \).
Indeed, this holds since
\begin{align*}
    g(x^*)
    &= \frac{\Re(z) + \frac{\Im(z)^2}{\Re(z)-w} - w}{\sqrt{ \left( \Re(z) - \left( \Re(z) + \frac{\Im(z)^2}{\Re(z)-w} \right)\right)^2 + \Im(z)^2 }}
    \\&= \operatorname{sign}(\Re(z)-w)\frac{(\Re(z)-w)^2 + \Im(z)^2}{|\Im(z)|\sqrt{(\Re(z)-w)^2 + \Im(z)^2}}
    \\&= \operatorname{sign}(\Re(z)-w) \left| \frac{z-w}{\Im(z)} \right|.
\end{align*}
}
\fi
\end{proof}

\begin{lemma}
    \label[lemma]{thm:level_sets}
Fix \( r >0 \), let  \( \mathcal{D}(c,t) \) be the disc in the complex plane centered at \( c \) with radius \( t \geq 0\), and define
\begin{equation*}
    X_r = \bigcup_{ x \in [a,b]  } \mathcal{D}\left( x, \frac{|x-w|}{r} \right).
\end{equation*}
%where \( \mathcal{D}(c,t) \) is the disc centered at \( c \) with radius \( t \geq 0\).
Then for \( z \in \mathbb{C} \setminus X_r \), we have
\begin{equation*}
    \|h_{w,z}\|_{[a,b]} 
    %= \max_{x \in \mathcal{I}} \left| \frac{x-w}{x-z} \right| 
    \leq r.
\end{equation*}
In particular, if \( z \) is on the boundary of \( X_r \), then \( \|h_{w,z}\|_{[a,b]}  = r \).
\end{lemma}

\begin{proof}
Let \( z \in \mathbb{C}\setminus X_r \) and pick any \( x \in [a,b] \).
 Since \( z \not\in \mathcal{D}(x,|x-w|/r) \) it follows that \( |z-x| > |x-w|/r \) and therefore \( |h_{w,z}(x)| = |x-w|/|x-z| < r \).
Maximizing over \( x \) yields the result.

If \( z\) is on the boundary of \( X_r \),  then for some \( x \in [a,b] \), \( |z-x| = |x-w|/r \), which means that for this \( x \), \(|h_{w,z}(x)| = r \).
\end{proof}
Note that if \( r \leq 1 \) and \( w \in \R\setminus[a,b] \), then the region described in \cref{thm:level_sets} is simply a disc about \( b \) if \( w < a\) or a disc about \( a \) if \( w > b \).  If \( r > 1 \) and \( w \) is real, then the region described is that in the discs about \( a \) and \( b \) and between the two external tangents to these two discs.


Similar to \cref{thm:Q_wz_value} we have the following bound on $\|h_z\|_{S_i}$ when $S_0$ is an  interval.
This allows a bound on \cref{eqn:err_quad_int} analogous to \cref{eqn:integral_error}.
\begin{lemma}
\label[lemma]{thm:Qz}
For any interval $[a,b]\subset \R$, if $z\in\mathbb{C}\setminus[a,b]$, we have    
    %Fix \( z \in \mathbb{C} \setminus [a,b] \) and \( w\in\mathbb{R} %\).
    \begin{equation*}
        \| h_z \|_{[a,b]}
    =
    \begin{cases}
        1/|\Im(z)| &  \Re(z) \in \mathcal{I} \\
        1/|a-z| & \Re(z) < a \\
        1/|b-z| & \Re(z) > b
    \end{cases}
\end{equation*}   
\end{lemma}



\subsection{A priori bounds}

We can use \cref{thm:err_int} to give a priori bounds, as long as we choose \( S \) and \( S_i \), \( i = 0, \ldots, k-1 \) independently of $\vec b$ (and in turn $\vec{T}$). %\Cam{In intro we complain that certain bounds depend on the spectrum making them bad a priori bounds. Should we reword this a bit?}

The simplest possibility is to take \( S = S_i = \mathcal{I} \).
In this case, as an immediate consequence of \cref{thm:err_int,thm:level_sets} we have the following a priori bound,
\begin{corollary}
    \label[corollary]{thm:disk}
    Suppose that for some \( w < \lmin \), \( f \) is analytic in a neighborhood of \( \mathcal{D}(\lmax,\lmax-w) \).
%on \( X_1 = \mathcal{D}(\lmax, \lmax) \).
    Then, taking \( \Gamma \) to be the boundary of this disk, 
%\( \Gamma = \{ r \left( \exp(i t) + r \right): r = \lambda_{\text{max}}, t \in [0,\pi] \} \),
   \begin{align*}
        \| \fA\vec{v} - \lan_k(f) \|
        &\leq \left( \frac{1}{2\pi} \oint_{\Gamma} |f(z)| |\d{z}| \right) \| \err_k(w) \|
        \\&\leq \left( (\lmax-w) \: \max_{z\in\Gamma}|f(z)|\right) \| \err_k(w) \|.
    \end{align*}
\end{corollary}
\begin{proof}
To obtain the first inequality  observe that \cref{thm:level_sets} with \( [a,b] = \mathcal{I} \) implies \( \|h_{w,z}\|_{\mathcal{I}} = 1 \) on this contour.
The second inequality follows since the length of \( \Gamma \) is \( 2 \pi (\lmax-w) \).
\end{proof}
%We remark that taking \( \Gamma \) as the \( \|h_{w,z}\|_{\mathcal{I}} = 1 \) level curve gives a somewhat shorter contour which can be used in place of the disk used in \cref{thm:disk}.
%This will often yield a slightly better bound, albeit at the cost of a more complicated parameteriztion of the contour.
This bound is closely related to \cite[Theorem 6.6]{frommer_guttel_schweitzer_14a} which bounds the error in Lanczos-FA for Stieltjes
functions in terms of the error in the Lanczos approximation for a certain linear system.

\iffalse
Using that \( \err_0(w) = (\vec{A}-w\vec{I})^{-1} \vec{v} \), we can rewrite \cref{thm:disk} as
\begin{equation*}
    \frac{ \| \fA \vec b - \lan_k(f) \|_2}{\| \fA\vec b \|_2} \le \max_{z \in \Gamma} |f(z)| \cdot  \frac{(\lmax- w) \|(\vec{A} - w\vec{I})^{-1} \vec b \|_2}{\| \fA\vec b \|_2} \cdot \frac{\|\err_k(w)\|_2}{\| \err_0(w) \|_2}. 
\end{equation*}
This can be used to obtain simple relative error bounds for many functions. 
For instance, suppose $\vec{A}$ is positive definite,  $f(x) = x^{-q}$ for \( q > 1 \), and $w = c \lambda_{\text{min}}$ for $c\in (0,1)$.
Then $\max_{z\in\Gamma} |z^{-q}| = w^{-q} = c^{-q} \lmin^{-q}$, \( \| (\vec{A} - w\vec{I})^{-1} \vec{v} \|_2 \leq (\lmin-w)^{-1} \| \vec{v} \| \) and \( \| \vec{A}^{-q} \vec{v} \|_2 \geq \lmax^{-q} \| \vec{v} \| \).
We then have the bound\footnote{Slightly stronger bounds can be obtained by bounding \( \| (\vec{A}-w\vec{I})^{-1} \vec{v} \|_2 / \| \vec{A}^{-q} \vec{v} \|_2 \) directly, rather than bounding the numerator and denominator separately.} 
\begin{equation*}
    \frac{ \| \vec{A}^{-q} \vec b - \lan_k(f) \|_2}{\| \vec{A}^{-q}\vec b \|_2} 
    %&\le \max_{z \in \Gamma} |z^{-q}| \cdot  \frac{(\lmax- w) \cdot \|(\vec{A} - w\vec{I})^{-1} \vec b \|_2}{\| \vec{A}^{-q} \vec b \|_2} \cdot \frac{\|\err_k(w)\|_2}{\| \err_0(w) \|_2} 
    \leq 
    c^{-q} \kappa(\vec{A})^{q} \kappa( \vec{A} - w\vec{I})  \frac{\|\err_k(w)\|_2}{\| \err_0(w) \|_2} .
\end{equation*}
\fi


\Cref{thm:disk} provides simple reductions to the error of solving a positive definite linear system involving \( \vec{A} - w\vec{I} \) using Lanczos. 
However, these bounds may be a significant overestimate in practice.
In particular, for any \( k>1 \), \cref{eqn:dkwz_qwz} cannot be sharp due to the fact that \( \|h_{w,z}\|_{\mathcal{I}} = \sup_{x \in \mathcal{I}} |h_{w,z}(x)| \) cannot be attained at every eigenvalue of \( \vec{T} \).
In fact, for most values \( \lambda_i ( \vec{T} ) \) and most points \( z \in \Gamma \), we expect \( | h_{w,z}( \lambda_i ( \vec{T} ) ) | \ll \|h_{w,z}\|_{\mathcal{I}} \).
\Cref{fig:ch7_hwz} shows sample level curves for \( \|h_{w,z}\|_{\mathcal{I}} / |\det(h_{w,z}(\vec{T}))|^{1/k} \) which illustrate  the slackness in the bound.


To derive sharper a priori bounds, there are several approaches. 

First, if more information is known about the eigenvalue distribution of \( \vec{A} \), then the \( S_i \) can be chosen based on this information.
For example, it is possible to exploit the interlacing property of the eigenvalues of \( \vec{T} \).
\begin{example}
Suppose \( \vec{A} \) has eigenvalues in \( [0,1] \) with a single eigenvalue at \( \kappa > 1 \).  Assume \( w \leq 0 \).
Then there is at most one eigenvalue of \( \vec{T} \) in \( [1,\kappa] \) so in \cref{thm:err_int} we can pick $S_i = [0,1]$ for $i = 0,\ldots, k-2$ and $S_{k-1} = [0,\kappa]$. We have
\begin{equation*}
    |\!\det(h_{w,z}(\vec{T}))| = \left| \prod_{i=0}^{k-1} h_{w,z}(\lambda_i(\vec{T})) \right| \leq
      \left(\| h_{w,z} \|_{[0,1]}\right)^{k-1} \|h_{w,z}\|_{[0,\kappa]}.
\end{equation*}
If \( z \) is near to \( \kappa \) then \( \| h_{w,z} \|_{[0,1]} \) may be much smaller than \( \| h_{w,z} \|_{[0,\kappa]} \).
\end{example}

\begin{figure}
    \begin{center}
        \includegraphics{imgs/ch7_hwz1.pdf} 
        \includegraphics{imgs/ch7_hwz2.pdf} 
    \end{center}
     \caption[{Contour plot of \( \|h_{w,z}\|_{\mathcal{I}} / |\!\det(h_{w,z}(\vec{T}))|^{1/k} \) as a function of $z \in \mathbb{C}$}]{%
         Contour plot of \( \|h_{w,z}\|_{\mathcal{I}} / |\!\det(h_{w,z}(\vec{T}))|^{1/k} \) as a function of $z \in \mathbb{C}$ for a synthetic example with \( w=0 \) (top) and \( w=1 \) (bottom), \( \mathcal{I} = [0.5,3] \), and \( \Lambda(\vec{T}) = \{ 0.5,0.8,1.2,1.5,3 \} \) ($k=5$).
    Larger slackness in \cref{eqn:dkwz_qwz} corresponds to darker regions. 
    \hspace{.25em}\emph{Legend}:
    Here \( w \) is indicated by the white diamond 
    ({\protect\raisebox{0mm}{\protect\includegraphics[]{imgs/legend/mD.pdf}}}). 
    and the eigenvalues of \( \vec{T} \) are indicated by white x'is 
    ({\protect\raisebox{0mm}{\protect\includegraphics[]{imgs/legend/mX.pdf}}}). 
    \hspace{.25em}\emph{Takeaway}: Slackness exchibits structure; in particular, it is lower far from \( \Lambda \). 
    }
    \label{fig:ch7_hwz}
\end{figure}


Second, the contour \( \Gamma \) can be chosen to try to reduce the slackness in \cref{eqn:dkwz_qwz}.
Intuitively, the slackness is exacerbated when \( z \in \Gamma \) is close to \( S_i \) but far from \( \lambda_i(\vec{T}) \).
For instance, for any \( k>1 \),
\begin{equation*}
\lim_{|z|\to\infty} \frac{\|h_{w,z}\|_{\mathcal{I}}^k}{|\det(h_{w,z}(\vec{T}))|} \to 1 
,\qquad\text{and }\qquad
\forall \lambda \in \mathcal{I},~ \lim_{z\to\lambda} \frac{\|h_{w,z}\|_{\mathcal{I}}^k}{|\!\det(h_{w,z}(\vec{T}))|} \to \infty.
\end{equation*}
This behavior is also observed in \cref{fig:ch7_hwz}.

These observations suggest that we should pick \( \Gamma \) to be far from the spectrum of \( \vec{A} \).
Of course, we are constrained by properties of \( f \) such as branch cuts and singularities.
Moreover, certain contours may increase the slackness in \cref{thm:err_int} itself. 
These considerations are discussed further in \cref{ex:sqrt_contours}.


\subsection{A posteriori error bounds}

After the Lanczos factorization \cref{eqn:lanczos_three_term} has been computed, \( \vec{T} \) is known and \( \Lambda(\vec{T}) \) can be cheaply computed. 
Thus, in \cref{thm:err_int} we can take \( S_i = \{ \lambda_i(\vec{T}) \} \) for \( i =0,\ldots,k-1 \), which is the best possible choice. 
In this case \cref{eqn:dkwz_qwz} is an equality and \( \det(h_{w,z}(\vec{T})) = \det(\vec{T}-w)/\det(\vec{T}-z) \) can be computed via tridiagonal determinant formulas rather than using the eigenvalues of \( \vec{T} \).

If \( \mathcal{I} \) is not known, the extreme Ritz values \( \lambda_{\text{min}}(\vec{T}) \) and \( \lambda_{\text{max}}(\vec{T}) \) can be used to estimate the extreme eigenvalues of \( \vec{A} \) \cite{kuczyski_wozniakowski_92,parlett_simon_stringer_82}.
All together, this means that it is not difficult to efficiently obtain accurate estimates of the bound from \cref{thm:err_int}.




\subsection{Numerical computation of integrals}
\label{sec:numerical_integral}
Typically, to produce an a priori or a posteriori error bound, the integral term in \cref{thm:err_int} must be computed numerically.
Consider a discretization of the integral
\begin{equation*}
    \fA = -\frac{1}{2\pi \ii} \oint_{\Gamma} f(z) (\vec{A}-z\vec{I})^{-1} \d{z}
\end{equation*}
using nodes \( z_i \) and weights \( w_i \), \( i=0,1,\ldots,q-1 \). 
This yields a rational matrix function
\begin{equation*}
    r_q(\vec{A}) := -\frac{1}{2\pi \ii} \sum_{i=0}^{q-1}  w_i f(z_i) (\vec{A}-z_i\vec{I})^{-1}. 
\end{equation*}
Using the triangle inequality, we can write
\begin{align}
\hspace{1em}&\hspace{-1em}\nonumber
\| \fA \vec{v} - \lan_k(f) \|
\\&\leq 
\nonumber
\| \fA \vec{v} - r_q(\vec{A}) \vec{v} \| + \| r_q(\vec{A}) \vec{v} - \lan_k(r_q)) \| + \| \lan_k(r_q) - \lan_k(f) \|
\\&\leq 2  \left( \max_{x \in \Lambda \cup \Lambda(\vec{T})} | f(x) - r_q(x) | \right) \|\vec{v}\|  + \| r_q(\vec{A}) \vec{v} - \lan_k(r_q) \|. \label{eqn:triangle_ineq_rational}
%\nonumber
%\\& \leq \underbrace{2 \|\vec{v}\| \left( \max_{x \in \mathcal{I}} | f(x) - r(x) | \right)}_{\text{approximation error}} 
%    + \underbrace{\vphantom{ \bigg| } \| r(\vec{A}) \vec{v} - \lan_k(r) \|}_{\text{application error}}.
%    \label{eqn:triangle_ineq}
\end{align}
Now, observe that analogous to \cref{thm:err_int},
\begin{align}
    \label{eqn:err_int_rational}
    \hspace{3em}&\hspace{-3em}\nonumber
    \| r_q(\vec{A})\vec{v} - \lan_k(r_q) \| 
    \\&\leq \left( \frac{1}{2\pi} \sum_{i=0}^{q-1} w_i \, |f(z_i)|  \left( \prod_{i=0}^{k-1} \|h_{w,z}\|_{S_i} \right)  \|h_{w,z}\|_{S_0} \right) \| \err_k(w) \|.
\end{align}
If we use the same nodes and weights to evaluate the integral term in \cref{thm:err_int}, we obtain exactly the expression on the right hand side of \cref{eqn:err_int_rational}.
Thus, this discretization of \cref{thm:err_int} is a true upper bound for the Lanczos-FA error to within an additive error of size equal to twice the approximation error of \( r(x) \) to \( f(x) \) on \( \Lambda \cup \Lambda(\vec{T}) \) times \( \| \vec{v} \| \).
In many cases, we expect exponential convergence of \( r_q \) to \( f \), which implies that this term can be made less than any desired value $\epsilon > 0$ using a number of quadrature nodes that grows only as the logarithm of \( \epsilon^{-1} \) \cite{hale_higham_trefethen_08,trefethen_weideman_14}.

We note that fast convergence of \( r_q \) to \( f \) suggests that, instead of applying Lanczos-FA, we can approximate \( \fA \vec{v} \) by first finding \( r_q \) and then solving a small number of linear systems \( (\vec{A} - z_i \vec{I}) \vec{x}_i = \vec{v} \) to compute \( r_q(\vec{A})\vec{v} \).
Solving these systems with any fast linear system solver yields an algorithm for approximating \( \fA\vec{v} \) inheriting, up to logarithmic factors in the error tolerance, the same convergence guarantees as the linear system solvers used.
A recent example of this approach is found in \cite{jin_sidford_19} which uses a modified version of stochastic variance reduced gradient (SVRG) to obtain a nearly input sparsity time algorithm for \( \fA\vec{v} \) when \( f \) corresponds to principal component projection or regression.

A range of work suggests using a Krylov subspace method and the shift invariance of the Krylov subspace to solve these systems and compute $r_q(\vec A)\vec b$ explicitly.
This was studied in \cite{frommer_kahl_lippert_rittich_13,frommer_simoncini_09} for the Lanczos method, and in \cite{pleiss_jankowiak_eriksson_damle_gardner_20} for MINRES, the latter of which uses the results of \cite{hale_higham_trefethen_08} to determine the quadrature nodes and weights. 
However, as the above argument demonstrates, the limit of the Lanczos-based approximation as the discretization becomes finer is simply the Lanczos-FA approximation to \( \fA \vec{v} \).
Therefore, there is no clear advantage to such an approach over Lanczos-FA in terms of the convergence properties, unless preconditioning is used.
.

On the other hand, there are some advantages to these approaches in terms of computation.
Indeed, Krylov solvers for symmetric/Hermitian linear systems require just $O(n)$ storage; i.e. they do not require more storage as more iterations are taken. 
A naive implementation of Lanczos-FA requires $O(kn)$ storage, and while Lanczos-FA can be implemented to use \( O(n) \) storage by taking two passes, this has the effect of doubling the number of matrix-vector products required.
See \cite{guttel_schweitzer_21} for a recent overview of limited-memory Krylov subspace methods.


\section{Examples and numerical verification}
\label{sec:cif_examples}

We next present examples in which we apply  \cref{thm:err_int} to give a posteriori and a priori error bounds for approximating common matrix functions with Lanczos-FA. These examples illustrate the general approaches to applying \cref{thm:err_int} described in \cref{sec:cif_results}. 
All integrals are computed either analytically or using SciPy's \texttt{integrate.quad} which is a wrapper for QUADPACK routines. 

In all cases, we exactly compute the $\| \err_k(w) \|$ term in the bounds. 
In practice, one would bound this quantity a priori or a posteriori using existing results on bounding the Lanczos error for linear system solves.
By computing the error exactly, we separate any looseness due to our bounds from any looseness due to an applied bound on  $\| \err_k(w) \|$. 



\subsection{Choice of contour}
\label{ex:sqrt_contours}
Let \( \vec{A} \) be positive definite and  \( f(x) = \sqrt{x} \).
Perhaps the simplest bound is obtained by using \cref{thm:err_int} with $w = 0$, \( S_i = \mathcal{I} \) and \( \Gamma \) chosen as the boundary of the disk \( \mathcal{D}(\lmax, \lmax ) \).%or as the \( \|h_{w,z}\|_{\mathcal{I}} = 1 \) level curve of \( \|h_{w,z}\|_{\mathcal{I}} \).
We then obtain a bound via \cref{thm:disk}. 
However, this bound may be loose -- note that except through $\|\err_k(w)\|$, it does not depend on the number of iterations $k$.  Thus it cannot establish convergence at a rate faster than that of solving a linear system with coefficient matrix \( \vec{A} \).


\begin{figure}[h]
\centering
\begin{subfigure}{.25\textwidth}\centering
\includegraphics[scale=.7]{imgs/ch7_circle_contour.pdf}
\caption{circle contour}
\label{fig:circle}
\end{subfigure}
\hfill
\begin{subfigure}{.25\textwidth}\centering
\includegraphics[scale=.7]{imgs/ch7_pacman_contour.pdf}
\caption{Pac-Man contour}
\label{fig:pacman}
\end{subfigure}
\hfill
\begin{subfigure}{.35\textwidth}\centering
\includegraphics[scale=.7]{imgs/ch7_double_circle_contour.pdf}
\caption{double circle contour}
\label{fig:double_circle}
\end{subfigure}
\caption[Circle, Pac-Man and double circle contours]{%
Circle, Pac-Man and double circle contours described in \cref{ex:sqrt_contours,ex:cif_pw} respectively.
    All three figures show \( \mathcal{I} \) 
    ({\protect\raisebox{0mm}{\protect\includegraphics{imgs/legend/spec.pdf}}})
    and \( w \) 
    ({\protect\raisebox{0mm}{\protect\includegraphics{imgs/legend/mD.pdf}}}).
}
\label{fig:contours}
\end{figure}

Keeping $w = 0$, we can obtain tighter bounds by letting $\Gamma$ be a ``Pac-Man'' like contour that consists of a large circle about the origin of radius \( R \) with a small circular cutout of radius \( r \) that excludes the origin and a small strip cutout to exclude the negative real axis. %, as pictured in \cref{fig:Pac-man}.
    That is, as shown in \cref{fig:pacman}, the boundary of the set,
\begin{equation*}
    \mathcal{D}(0,R) \setminus ( \{ z : \Re(z) \leq 0, |\Im(z)| < r \} \cup \mathcal{D}(0,r) ).
\end{equation*}

As the outer radius \( R \rightarrow \infty \), the integral over the large circular arc goes to \( 0 \) since \( \|h_{w,z}\|_{\mathcal{I}} = O(R^{-1}) \), \( | f(z) | = O(R^{1/2}) \), and the length of the circular arc is on the order of \( R \). 
Thus, the product \( f(z) (\| h_{w,z} \|_{\mathcal{I}})^{k+1} \) goes to \( 0 \) as \( R \rightarrow \infty \), for all \( k \geq 1 \).  
Similarly, as \( r\to 0 \), the length of the small arc goes to zero.
Therefore, we need only consider the contributions to the integral on \( [ -R \pm ir, \pm ir] \) in the limit \( R\to \infty, r\to 0 \).


In this case, when $S_i = \mathcal{I}$ for all $i$, we can  compute the value of the integral term in \cref{thm:err_int} analytically. We have
\begin{align*}
\| f( \vec{A} ) \vec{v} - \lan_k (f) \| 
&\leq \left(\frac{1}{2 \pi} \int_{-\infty}^0 | (x \pm 0 i)^{1/2} | \:  \| h_{w,x \pm 0 i} \|_{\mathcal{I}}^{k+1} \,\d{x}  \right) \| \err_k \| 
\\&= \left(\frac{1}{2 \pi} \int_{-\infty}^0 | x \pm 0 i|^{1/2}\frac{ \lambda_{\text{max}} ( \vec{A})^{k+1}}{( \lambda_{\text{max}} ( \vec{A} ) - x )^{k+1}}\,\d{x}   \right) \| \err_k \| 
\\&= \left(\frac{1}{\pi} \lambda_{\text{max}} ( \vec{A} )^{k+1}  \int_0^{\infty} \frac{y^{1/2}}{( \lambda_{\text{max}} ( \vec{A} ) + y )^{k+1}}\d{y} \right) \| \err_k \| 
\\&= \left( \frac{\lmax^{3/2}}{2\sqrt{\pi}} \frac{\Gamma(k-1/2)}{\Gamma(k+1)}  \right) \| \err_k \|,
%\frac{1}{2 \pi} \int_{-R}^0 | (x \pm ir )^{1/2} |\cdot|\! \det ( h_{w,x \pm ir} (\vec{T})) |\cdot\| h_{w,x \pm ir} ( \vec{A} ) \|\,\d{x} 
%\frac{1}{2 \pi} \int_{-R}^0 | (x \pm ir )^{1/2} |\cdot|\! \det ( h_{w,x \pm ir} (\vec{T})) |\cdot\| h_{w,x \pm ir} ( \vec{A} ) \|\,\d{x} 
\end{align*}
where we have made the change of variable \( y = -x \). 
Note that
\begin{equation*}
    \lim_{k \to \infty }k^{3/2}  \frac{\Gamma(k-1/2)}{\Gamma(k+1)} = 1.
\end{equation*}
This proves that \( \lan_k ( \sqrt{\cdot}) \) converges somewhat faster than the Lanczos algorithm applied to the corresponding linear system \( \vec{A} \vec{x} = \vec{v} \).


\begin{figure}[htb]
    \begin{center}
        \includegraphics{imgs/ch7_sqrt_Anorm.pdf} 
    \end{center}
    \caption[{\( \vec{A} \)-norm error bounds for \( f(x) = \sqrt{x} \) where \( \vec{A} \) has \( n=1000 \) eigenvalues spaced uniformly in \( [10^{-2},10^2] \) and \( \Gamma \) is a circular contour (left) or Pac-Man contour (right).}]{%
    \( \vec{A} \)-norm error bounds for \( f(x) = \sqrt{x} \) where \( \vec{A} \) has \( n=1000 \) eigenvalues spaced uniformly in \( [10^{-2},10^2] \) and \( \Gamma \) is a circular contour (left) or Pac-Man contour (right).
    \hspace{.25em}\emph{Legend}:
    Lanczos-FA error 
    ({\protect\raisebox{0mm}{\protect\includegraphics[]{imgs/legend/l1.pdf}}}), a priori bounds obtained by using \cref{thm:err_int} with \( S = S_i = \mathcal{I}  \)
    ({\protect\raisebox{0mm}{\protect\includegraphics[]{imgs/legend/l4.pdf}}}) 
    and \( S = S_i = \tilde{\mathcal{I}}(\vec{A}) = [ \lmin/2, 2 \lmax] \)
    ({\protect\raisebox{0mm}{\protect\includegraphics[]{imgs/legend/l3.pdf}}}), 
    a posteriori bounds obtained by using \cref{thm:err_int} with \( S = \tilde{\mathcal{I}} \)
    ({\protect\raisebox{0mm}{\protect\includegraphics[]{imgs/legend/l2.pdf}}}).
    \hspace{.25em}\emph{Takeaway}: The a posteriori bounds tend to be quite accurate. The choice of contour impacts the quality of the bounds, particuarly the a priori bounds.
    }
    \label{fig:ch7_sqrt_Anorm}
\end{figure}


In \cref{fig:ch7_sqrt_Anorm} we plot the bounds from \cref{thm:err_int} for the circular and Pac-Man contours described above.
For both contours we consider \( S_i = \mathcal{I} \) for all \( i \), as well as bounds based on an overestimate of this interval,  \( S_i = \tilde{\mathcal{I}}(\vec{A}) \) where 
\(
    \tilde{\mathcal{I}}(\vec{A}) = [ \lmin/2, 2 \lmax].
\)
This provides some sense of how sensitive the bounds are to the choice of \( S_i \) when $S_i$ is a single interval. 
For a posteriori bounds, we set \( S_i \) to $\{\lambda_i(\vec T_k)\}$ for \( i > 0 \).

We remark that the bounds from \cref{thm:err_int} are upper bounds for \cref{eqn:integral_error} which implies that the slackness of \cref{eqn:integral_error} is relatively small.
This suggests that the roughly 2 orders of magnitude improvement in \cref{thm:err_int} when moving from the circular contour to the Pac-Man contour is primarily due to reducing the slackness in \cref{eqn:dkwz_qwz}, aligning with our intuition.

\subsection{Piecewise analytic functions}
\label{ex:cif_pw}

We now discuss the application of \cref{thm:err_int} to piecewise analytic functions.
Functions of this class have found widespread use throughout scientific computing and data science but have proven particularly difficult to analyze using existing approaches \cite{dinapoli_polizzi_saad_16,frostig_musco_musco_sidford_16,jin_sidford_19,eshof_frommer_lippert_schilling_van_der_vorst_02}.

Let \( f(x) \) be one of \( |x-a| \), \( \operatorname{step}(x-a) \), or \( \operatorname{step}(x-a)/x \) for \( a \in \mathcal{I} \), where, for \( z \in \mathbb{C} \) we define $\operatorname{step}(z) := 0$ for $\text{Re}(z) < 0$ and $\operatorname{step}(z) := 1$ for $\text{Re}(z) \ge 0$.
Note that the latter two functions correspond to principle component projection and principle component regression respectively.
In the case of principle component regression, we assume $\vec{A}$ is positive semi-definite.
The step function is also closely related to the sign function, which is widely used in quantum chromodynamics to compute the overlap operator \cite{eshof_frommer_lippert_schilling_van_der_vorst_02}.

We take \( w = a \) and define \( \Gamma_1 \) and \( \Gamma_2 \) as the boundaries of the disks 
\begin{equation*}
    \mathcal{D}_1 := \mathcal{D} ( \lmin, w - \lmin - \varepsilon) 
    \quad \text{and}\quad
    \mathcal{D}_2 := \mathcal{D} ( \lmax , \lmax - w - \varepsilon ),
\end{equation*}
for some sufficiently small \( \varepsilon > 0 \). 
To extend \( |x-a| \) to the complex plane, we replace \( |x-a| \) by \( z-a \) if \( \text{Re}(z) > a \) and by \( a-z \) if \( \text{Re}(z) \leq a \). 
Then \( f \) is analytic in a neighborhood of the union of these two disks, so assuming none of the eigenvalues of $\vec{A}$ or $\vec{T}$ are equal to $a$, we can apply \cref{thm:level_sets}.
\begin{table}[htb]
 \begin{align*}
\begin{array}{rlll}
    \toprule
    f(x) & f(z), z \in \Omega_1 & f(z), z \in \Omega_2 & \frac{1}{2 \pi} \sum_{j=1}^{2} | \Gamma_j | \max_{z\in\Gamma_j} | f(z) | \\\midrule
    |x-a| & a-z & z-a &  2(a- \lambda_{\text{min}} )^2 + 2(\lambda_{\text{max}}-a)^2  \\
    \operatorname{step}(x-a) & 0 & 1 &  (\lambda_{\text{max}}-a) \\
    \operatorname{step}(x-a)/x & 0 & 1/z & (\lambda_{\text{max}}-a)/a\\
    \bottomrule
\end{array}
\end{align*}
\caption{Values of the factor in parentheses on the right-hand side of \cref{eqn:ex_cif_pw_a_priori2} (ignoring $\varepsilon$) for several common piecewise analytic functions.}
\label{tab:cif_piecewise}
\end{table}


Note that \( \|h_{w,z}\|_{\mathcal{I}} \to 1 \) as \( z \to w \) from outside \( [a,b] \), avoiding a potential singularity which would occur if the contour $\Gamma$ passed through $\mathcal{I}$ at any other points.
In fact, ignoring the contribution of $\epsilon$, \( \| h_{w,z} \|_{\mathcal{I}} = 1 \) for all \( z\in \Gamma_1 \) and for all \( z\in\Gamma_2 \).
Thus, \cref{thm:disk} can be written as
\begin{equation}
    \| \fA \vec{v} - \lan_k(f) \|
%    &\leq \left( \frac{1}{2\pi} \sum_{j=1}^{2} \int_{\Gamma_j} | f(z) | \d{z}| \right) \| \err_k(w) \| \label[ineq]{eqn:ex_cif_pw_a_priori1}
    \leq \left(\frac{1}{2\pi} \sum_{j=1}^{2}  |\Gamma_j| \max_{z\in\Gamma_j} | f(z) | \right) \| \err_k(w) \| \label{eqn:ex_cif_pw_a_priori2}.
\end{equation}
The values of this bound for all three functions are summarized in \cref{tab:cif_piecewise}.

If \( w \in \mathcal{I} \)  we note that \( \| \err_k(w) \| \) corresponds to the indefinite linear system \( (\vec{A} - w \vec{I}) \vec{x} = \vec{v} \), so standard results for the Conjugate Gradient algorithm are not applicable.
However, the residual of this system can still be computed exactly once the Lanczos factorization \cref{eqn:lanczos_three_term} has been obtained, and as we discuss in \cref{sec:CG_indefinite}, a priori bounds for the convergence of MINRES \cite{cullum_greenbaum_96} can be extended to the Lanczos algorithm for indefinite systems.
It is also clear that, at the cost of having to compare against the error of multiple different linear systems, functions which are piecewise analytic on more than two regions can be handled.

\begin{figure}[ht]
    \begin{center}
        \includegraphics{imgs/ch7_pcr.pdf} 
    \end{center}
    \caption[{\( (\vec{A}-w\vec{I})^2 \)-norm error bounds for \( f(x) = \operatorname{step}(x-a)/x \) where \( \vec{A} \) is a random matrix whose limiting density is supported on \( [a_1,b_1]\cup[a_2,b_2] \), $a = (b_1+a_2)/2$, and \( \Gamma \) is a double circle contour.}]{% 
    \( (\vec{A}-w\vec{I})^2 \)-norm error bounds for \( f(x) = \operatorname{step}(x-a)/x \) where \( \vec{A} \) is a random matrix whose limiting density is supported on \( [a_1,b_1]\cup[a_2,b_2] \), $a = (b_1+a_2)/2$, and \( \Gamma \) is a double circle contour. 
    \hspace{.25em}\emph{Legend}:
    Lanczos-FA error 
    ({\protect\raisebox{0mm}{\protect\includegraphics[]{imgs/legend/l1.pdf}}}), 
    a priori bounds obtained by using \cref{thm:err_int} with \( S = S_i = [a_1-0.1,b_2+0.1]  \)
    ({\protect\raisebox{0mm}{\protect\includegraphics[]{imgs/legend/l3.pdf}}}) and \cref{eqn:ex_cif_pw_a_priori2} with the values from \cref{tab:cif_piecewise}
    ({\protect\raisebox{0mm}{\protect\includegraphics[]{imgs/legend/l4.pdf}}}) 
    a posteriori bounds obtained by using \cref{thm:err_int} with \( S = [a_1-0.1,b_2+0.1] \)
    ({\protect\raisebox{0mm}{\protect\includegraphics[]{imgs/legend/l2.pdf}}}). 
    \hspace{.25em}\emph{Takeaway}: The bounds work well, even for pieceiwise analytic functions.
    }
    \label{fig:ch7_pcr}
\end{figure}


In \cref{fig:ch7_pcr}, we plot the bounds from \cref{thm:err_int} for the contour described above for principle component regression with \( f(x) = \operatorname{step}(x-a)/x \).
Here we use the same model as in \cref{sec:smooth_density}. 
In particular, we set \( n=2000 \), \( d=0.3 \), and \( \sigma = 8 \) and take
\begin{equation*}
    \vec{A} = \frac{1}{m} \vec{\Sigma}^{1/2} \vec{X} \vec{X}^\cT \vec{\Sigma}^{1/2},
\end{equation*}
where \( m = n/d \), \( \vec{X} \) is a \( n\times m \) matrix with iid standard normal entries, and \( \vec{\Sigma} \) a diagonal matrix with \( 1/m \) as the first \( n/2 \) entries and \( \sigma/m \) as the last \( n/2 \) entries.
As discussed in \cref{sec:smooth_density}, in the large \( n \) limit, the spectrum of such matrices is supported on intervals \( [a_1,b_1]\cup[a_2,b_2] \), so we take \( a = (b_1+a_2)/2 \) and \( S = S_i = [a_1-0.1,b_2+0.1]\) for a priori bounds and \( S = [a_1-0.1,b_2+0.1 \) for a posteriori bounds.
Thus, \( f(x) \) corresponds to solving a linear system involving the eigenmodes of the right cluster of eigenvalues supported on \( [a_2,b_2] \).



%For reference, we also plot the bound \cref{eqn:ex_cif_pw_a_priori2} which shows that there is not much loss from using \cref{eqn:ex_cif_pw_a_priori2} in place of \cref{thm:err_a_priori}.
%Again, all integrals are computed using SciPy's \texttt{integrate.quad}.


\subsection{Quadratic forms}
\label{ex:quadform_step}

Let \( f(x) = \operatorname{step}(x-a) \) for  \( a \in \mathcal{I} \), and set $w = a$.
Similarly to the previous example we use \cref{eqn:err_quad_int} to obtain a bound for the quadratic form error $| \vec{v}^\cT \fA\vec{v} - \vec{v}^\cT \lan_k(f)|$. 
However, since \( \| h_z \|_{S_i} \) has singularities at each point in \( S_i \), we must have $S_i$ avoid where $\Gamma$ crosses the real axis.

\begin{figure}[ht]
    \begin{center}
        \includegraphics{imgs/ch7_step_qf.pdf} 
    \end{center}
    \caption[{Quadratic form error bounds for \( f(x) = \operatorname{step}(x-a) \) where \( \vec{A} \) is a random matrix whose limiting density is supported on \( [a_1,b_1]\cup[a_2,b_2] \), $a = (b_1+a_2)/2$, and \( \Gamma \) is a double circle contour.}]{%
    Quadratic form error bounds for \( f(x) = \operatorname{step}(x-a) \) where \( \vec{A} \) is a random matrix whose limiting density is supported on \( [a_1,b_1]\cup[a_2,b_2] \), $a = (b_1+a_2)/2$, and \( \Gamma \) is a double circle contour. 
    \hspace{.25em}\emph{Legend}:
    Lanczos-FA quadratic form error 
    ({\protect\raisebox{0mm}{\protect\includegraphics[]{imgs/legend/l1.pdf}}}), 
    squared Lanczos-FA 2-norm ({\protect\raisebox{0mm}{\protect\includegraphics[]{imgs/legend/l4.pdf}}}),
    a priori bounds obtained by using \cref{thm:err_int} with \( S = S_i = [a_1-0.1,b_1+0.1]\cup[a_2-0.1,b_2+0.1]  \)
    ({\protect\raisebox{0mm}{\protect\includegraphics[]{imgs/legend/l3.pdf}}})
    a posteriori bounds obtained by using \cref{thm:err_int} with \( S = [a_1-0.1,b_1+0.1]\cup[a_2-0.1,b_2+0.1] \)
    ({\protect\raisebox{0mm}{\protect\includegraphics[]{imgs/legend/l2.pdf}}}). 
    \hspace{.25em}\emph{Legend}: The bounds are applicable to quadratic forms.
    }
    \label{fig:ch7_pcr_qf}
\end{figure}



Suppose \( \lambda_{\text{max}}^{\text{l},w}(\vec{A}) \) and \( \lambda_{\text{min}}^{\text{r},w}(\vec{A}) \) are consecutive eigenvalues of \( \vec{A} \) so that  \( \lambda_{\text{max}}^{\text{l},w}(\vec{A}) < w < \lambda_{\text{min}}^{\text{r},w}(\vec{A}) \).
Then we can define
\begin{equation*}
\mathcal{I}_w(\vec{A}) := [ \lmin , \lambda_{\text{max}}^{\text{l},w}(\vec{A}) ] \cup [ \lambda_{\text{min}}^{\text{r},w}(\vec{A}) , \lmax ].
\end{equation*}
In this case, \( \|h_{z}\|_{\mathcal{I}_w(\vec{A})} = \max\{ \|h_{z}\|_{ [\lambda_{\text{min}},\lambda_{\text{max}}^{\text{l},w} ] }, \|h_{z}\|_{ [\lambda_{\text{min}}^{\text{r},w},\lambda_{\text{max}} ] }) \} \) can be computed using \cref{thm:Qz}. 
We can then apply \cref{eqn:err_quad_int} to obtain a bound for the quadratic form error $| \vec{v}^\cT \fA\vec{v} - \vec{v}^\cT \lan_k(f)|$.
A priori bounds are obtained with $S_0, S_i = \mathcal{I}_w(\vec A)$ while a posteriori bounds are obtained with $S = \mathcal{I}_w(\vec A)$ and $S_i = \{ \lambda_i(\vec{T}) \}$.

In \cref{fig:ch7_pcr_qf} we use the same matrix as in \cref{ex:cif_pw}.
This time, however, we use \( S = S_i = [a_1-0.1,b_1+0.1]\cup[a_2-0.1,b_2+0.1]  \) for a priori bounds and \( S = [a_1-0.1,b_1+0.1]\cup[a_2-0.1,b_2+0.1]  \) for a posteriori bounds.
Note that the squared error \( \| \fA\vec{v} - \lan_k(f) \|^2 \) is close to that of the quadratic form error \( | \vec{v}^\cT \fA\vec{v} - \int f \,\d\qq[g]{\Psi}{2k-1} | \).

\section{Error bounds for Lanczos-FA on indefinite systems}
\label{sec:CG_indefinite}

In this section, we review several results which rigorously justify the claim that, for any choice of \( w \) with \( \vec{A} - w\vec{I} \) invertible, \( \|\err_k(w)\| \) satisfies a spectrum-dependent error bound.


Note that on indefinite problems, the standard implementation of CG (or the LDL version on which Lanczos-OR is based) may fail in such situations since \( \vec{T} \) can be singular in which case the inversion will break down.
As a result, in such situations, it is standard practice to use MINRES or other related algorithms.
On the other hand, the Lanczos algorithm does not break down if \( \vec{T} \) is singular, and so the Lanczos-FA approximation to \( \vec{A}^{-1} \vec{v} \) can be computed whenever \( \vec{T} \) is non-singular, even if it was singular at earlier iterations. 
Interestingly, however, the ``overall'' convergence of the algorithm tends to be comparable to MINRES in the sense that at many iterations the error is quite similar.


The first result we remark on was first proved in \cite{cullum_greenbaum_96} compares the residual norms of Lanczos-FA and MINRES.
\begin{theorem}
    %\cite[Ex. 5.1]{greenbaum_97}
\label{thm:minres_CG_residuals1}
Let \( \vec{A} \) be a nonsingular Hermitian matrix and define \( \vec{r}_k^M \) as the MINRES residual at step $k$; i.e. 
\begin{equation*}
    \vec{r}_k^M := \vec{v} - \vec{A} \hat{\vec{y}}
    ,\qquad
    \hat{\vec{y}} = \argmin_{\vec{y}\in\mathcal{K}_k} \| \vec{v} - \vec{A} \vec{y} \|_2 .
\end{equation*}
Then, assuming that the initial residuals in the two procedures are the same,
\begin{equation*}
    \frac{\| \Res_k  \|_2}{\|\Res_0\|_2} 
    = \frac{\| \vec{r}_{k}^M \|_2/\| \vec{r}_{0}^M\|_{2}}{\sqrt{1- \left( \| \vec{r}_{k}^M \|_2 / \| \vec{r}_{k-1}^M \|_2 \right)^2}}.
\end{equation*}
\end{theorem}

Therefore, we see that if MINRES makes good progress at step \( k \) (i.e. \( \| \vec{r}_k^M \|_2 / \| \vec{r}_{k-1}^M \|_2 \) is small), then \cref{thm:minres_CG_residuals1} implies \( \| \Res_k \|_2/\|\Res_0\|_2 \approx \| \vec{r}_k^M \|_2/\|\vec{r}_0^M\|_2 \).
%Loosely, we therefore expect that Lanczos-FA to converge similarly to MINRES in the sense that
Thus, since MINRES converges at a linear rate,
there must be iterations where the MINRES residual norm decreases enough that the Lanczos-FA residual norm is similarly small.
This is made precise in \cite[Corollary A.2]{chen_greenbaum_musco_musco_22a} which demonstrates the iteration complexity of Lanczos-FA on indefinite systems is nearly the same as that of MINRES.

\iffalse
\begin{corollary}
\label{thm:indefinite_CG}
Suppose $\Lambda\subset [a,b]\cup[c,d]$, where $a < b < 0 < c < d$ with $b-a = d-c$, and define $\gamma = \sqrt{|ad|/|bc|}$.
Then, for any $\epsilon < \gamma/4$ there exsits \( k \leq 2 \gamma \log(\sqrt{2} \gamma/\epsilon) \) so that \( \| \Res_k \|_2 / \| \Res_0 \|_2 < \epsilon \).
\end{corollary}

While the proof is not difficult, we do not include it as the corollary only improves the bound from ... by a constant factor.
\fi

In fact, stronger results are known. 
In particular, it is known that Lanczos-FA process iterates whose residuals satisfy a minimax bound on the eigenvalues of \( \vec{A} \), at least at every other iteration \cite{greenbaum_druskin_knizhnerman_99}.
While not well known, the argument proving this claim is amazingly simple, so we provide proofs for the exact arithmetic case.
These results hold to close approximation in finite precision arithmetic; see \cite{greenbaum_druskin_knizhnerman_99} for the statements and proofs in this setting. 

We begin with several lemmas. 
\begin{lemma}
    \label[lemma]{thm:TkTk1_rvbd}
    Suppose \( \theta \) is an eigenvalue of \(\,\vec{T}\) with eigenvector \( \vec{s} \)  and \( \mu \) is an eigenvalue of \( \,[\vec{T}]_{:k-1,:k-1} \) with eigenvector \( \vec{v} \).
Then,
\begin{equation*}
    \theta - \mu = \frac{\beta_{k-2}[\vec{s}]_{k-1}[\vec{v}]_{k-2}}{\vec{s}^\cT \hat{\vec{v}}}
\end{equation*}
    where \( \hat{\vec{v}} \) is \( \vec{v} \) with a zero appended at the bottom.
\end{lemma}

\begin{proof}
Observe that
\begin{equation*}
    \theta \vec{s}^\cT \hat{\vec{v}} 
    = \vec{s}^\cT \vec{T} \hat{\vec{v}} 
    = \vec{s}^\cT ( \mu \hat{\vec{v}} + \beta_{k-2}[\vec{v}]_{k-2} \vec{e}_{k-1} )
    = \mu \vec{s}^\cT \hat{\vec{v}} + \beta_{k-2} [\vec{s}]_{k-1}[\vec{v}]_{k-2}.
\end{equation*}
The result follows by rearranging the above expression.
\end{proof}

\begin{lemma}
    \label[lemma]{thm:rv_bound}
Suppose \( \theta \) is an eigenvalue of \( \,\vec{T} \) with eigenvector \( \vec{s} \). 
    Then
\begin{equation*}
    \min_{0\leq i<n} |\theta - \lambda_i| \leq | \beta_{k-1} [\vec{s}]_{k-1}|.
\end{equation*}
\end{lemma}

\begin{proof}
Using the Lanczos recurrence \cref{eqn:lanczos_three_term} and the eigendecomposition \( \vec{A} = \vec{U} \vec{\Lambda}\vec{U}^\cT \)  we find,
\begin{equation*}
    \beta_{k-1} \vec{q}_{k} \vec{e}_{k-1}^\cT \vec{s} = \vec{A}\vec{Q}\vec{s} - \vec{Q}\vec{T}\vec{s}
    = \vec{A} \vec{Q}\vec{s} - \theta \vec{Q}\vec{s}
    = \vec{U}(\vec{\Lambda} - \theta \vec{I}) \vec{U}^\cT \vec{Q} \vec{s}.
\end{equation*}
%
Rearranging and taking norms on both sides we then have
\begin{equation*}
    1 = \| \vec{Q}\vec{s} \| 
    = \| \beta_{k-1} \vec{U}(\Lambda - \theta \vec{I})^{-1} \vec{q}_{k} \vec{e}_{k-1}^\cT \vec{s} \|_2
    \leq |\beta_{k-1} [\vec{s}]_{k-1}| \| (\vec{\Lambda} - \theta \vec{I})^{-1} \|_2 
\end{equation*}
where we have used that \( \|\vec{U}\|_2 = 1 \) and \( \|\vec{Q}\vec{s}\|_2 = 1 \).
The result then follows from the fact that
\begin{equation*}
   \|(\vec{\Lambda} - \theta \vec{I})^{-1}\|_2^{-1}
    = \big(\max_{0\leq i<n} |\lambda_i - \theta|^{-1} \big)^{-1}
    = \min_{0\leq i<n} |\lambda_i - \theta|.
    \qedhere
\end{equation*}
\end{proof}

The first main result of \cite{greenbaum_druskin_knizhnerman_99} asserts that eigenvalues of \( \vec{T} \) are bounded away from zero at least at every other iteration, provided \( \vec{A} \) is not singular. 
\begin{theorem}
\label{thm:Tk_invertible_bd}
Suppose \( \theta \) is an eigenvalue of \(\,\vec{T}\) and \( \mu \) is an eigenvalue of \( \,[\vec{T}]_{:k-1,:k-1} \).
Then
\begin{equation*}
    \frac{\max\{|\theta|,|\mu|\}}{\|\vec{A}\|} \geq \frac{\kappa^2}{2+\sqrt{3}}.
\end{equation*}
\end{theorem}

\begin{proof}
Applying \cref{thm:rv_bound} to $\vec{T}$ and $[\vec{T}]_{:k-1,:k-1}$ and then using \cref{thm:TkTk1_rvbd} and the fact that $\|\vec{s}\|_2 = \|\vec{v}\|_2 = 1$ we have
\begin{equation*}
    |\theta||\mu| \leq \beta_{k-2}\beta_{k-1} [\vec{s}]_{k-1} [\vec{v}]_{k-2}
    = \beta_{k-1} |\theta - \mu| \vec{s}^\cT \hat{\vec{v}}
    \leq \beta_{k-1} |\theta - \mu|.
\end{equation*}

Let $\tau := \max\{|\theta|,|\mu|\}$. 
Then $|\theta-\mu|\leq 2\tau$, $|\theta| \geq \sigma_{\textup{min}}-\tau$,  and $|\mu| \geq \sigma_{\textup{min}}-\tau$.
This implies that
\begin{equation*}
    (\sigma_{\textup{min}} - \tau)^2
    \leq |\theta||\mu|
    \leq \beta_{k-1} | \theta - \mu |
    \leq 2 \beta_{k-1} \tau.
\end{equation*}
Solving for $\tau$ we find
\begin{equation*}
    \tau \geq \frac{\sigma_{\textup{min}}^2}{\sigma_{\textup{min}} + \beta_{k-1} + \sqrt{\beta_{k-1}^2 + 2\beta_{k-1} \sigma_{\textup{min}}}}
    \geq \frac{\kappa}{2+\sqrt{3}}
\end{equation*}
where, in the final inequality, we have use the fact that both $\beta_{k-1} $ and $\sigma_{\textup{min}}$ are bounded above by $\| \vec{A} \|_2$.
\end{proof}

We are nearly ready to show that the Lanczos-FA iterate satisfies a minimax bound on $\Lambda$. 
First, however, we require the following lemma relating the bottom left entry of \( \vec{T}^{-1} \) to the norm of \( \vec{T}^{-1} \).
\begin{lemma}
    \label[lemma]{thm:Tinv_poly}
\begin{equation*}
    | \vec{e}_{k-1}^\cT\vec{T}^{-1} \vec{e}_0 |
    \leq \| \vec{T}^{-1} \|_2 \min_{\deg(p)<k, p(0)=1} \| p(\vec{T}) \vec{e}_1 \|_2.
\end{equation*}
\end{lemma}

\begin{proof}
Write $p = 1-xq$ for some $q$ with $\deg(q)<k-1$. 
Then $\vec{e}_{k-1}^\cT q(\vec{T}) \vec{e}_0 = 0$ so
\begin{equation*}
    \vec{e}_{k-1}^\cT\vec{T}^{-1} \vec{e}_0
    = \vec{e}_{k-1}^\cT(\vec{T}^{-1}-q(\vec{T})) \vec{e}_0
    = \vec{e}_{k-1}^\cT\vec{T}^{-1} p(\vec{T}) \vec{e}_0.
\end{equation*}
Now, applying a submultiplicative bound, we find
\begin{equation*}
    |\vec{e}_{k-1}^\cT\vec{T}^{-1} \vec{e}_0|
    \leq \| \vec{T}^{-1} \|_2 \| p(\vec{T}) \vec{e}_1 \|_2.
\end{equation*}
The result follows by optimizing over $p$.
\end{proof}

Proving a minimax bound for the Lanczos-FA residual is now straightforward.
\begin{theorem}
At least at every other iteration,
\begin{equation*}
    \| \res_k \|_2 
    \leq \frac{\kappa^2}{2+\sqrt{3}}l\min_{\deg(p)<k, p(0)=1} \| p \|_{\Lambda}
\end{equation*}
\end{theorem}

\begin{proof}
From \cref{thm:shifted_lanczos_equivalence} we see have that
\begin{equation*}
    \| \res_k \|_2
    = |\beta_{k-1} \vec{e}_0^\cT \vec{T}^{-1} \vec{e}_0|.
\end{equation*}
The result then follows immediately from \cref{thm:Tinv_poly,thm:Tk_invertible_bd}.
\end{proof}



Interestingly, it is known that $\vec{T}$ cannot have eigenvalues near zero in two successive iterations, at least assuming that the eigenvalues of $\vec{A}$ are not too close to zero.
Specifically,  \cite[Equation 3.10]{greenbaum_druskin_knizhnerman_99} asserts that
\begin{equation}
    \label{eqn:nodoublezeroT}
\max\{ \sigma_{\textup{min}}([\vec{T}]_{:k-1,:k-1}),\sigma_{\textup{min}}(\vec{T})\} > \frac{\sigma_{\textup{min}}\A^2}{(2+\sqrt{3})\|\vec{A}\|_2}.
\end{equation}
Thus, as noted in \cref{rem:CG_MINRES}, we might still use the Lanczos-FA approximation to \( \vec{A}^{-1} \vec{v} \).





\chapter{Finite precision arithmetic}
\label{chap:finite_precision}

As mentioned in the introduction and observed in the many numerical experiments throughout this thesis, even in finite precision arithmetic, Lanczos-based methods for matrix functions tend to perform at least as well as their explicit polynomial counterparts.
In fact, they often perform significantly better.
This is in direct conflict with the widespread notion that costly reorthogonalization schemes are necessary for Lanczos-based methods \cite{jaklic_prelovsek_94,aichhorn_daghofer_evertz_vondelinden_03,weisse_wellein_alvermann_fehske_06,ubaru_chen_saad_17,granziol_wan_garipov_19}.

In this chapter, we provide an overview of several existing theoretically rigorous results which explain this phenomenon.
We also prove that the reduction technique from \cref{chap:CIF} still works in finite precision arithmetic.
It is our hope that our treatment of this topic will provide an accessible starting point for those outside of numerical analysis to better understand the impact of finite precision arithmetic on Lanczos-based methods.
Thus, while this chapter consists mostly of exposition on existing results, we believe it to be it to be one of the more important contributions of this thesis.

In this chapter, we will use \( \| \cdot \| \) rather than \( \| \cdot \|_2 \) to denote the operator 2-norm of matrices and Euclidean norm of vectors.
We will also make the simplifying assumption that \( \vec{A} \) is real symmetric.


\section{Preliminaries}

Almost all of modern scientific computing involves computations in finite precision arithmetic, and specifically, floating point arithmetic.
The introduction of rounding errors can have potentially large impacts on the output of an algorithm, so accounting for these errors is important. 
In fact, this is one of the main goals of the field of numerical analysis.

There are many possible implementations of floating point arithmetic and other low-level math kernels.
For instance, the number of bits of precision may vary, the rounding scheme may vary, the way basic functions such as the square root and logarithm are implemented may vary, etc.
To avoid the need for a separate analysis of each implementation, it is standard to work in a model of computation which captures the essential qualities of a broad number of basic math routines.

Perhaps the most commonly studied model of finite precision computing assumes that basic operations are carried out to relative accuracy \( \epsilon_{\textup{mach}} \), a constant referred to as the machine precision.
For floating point numbers \( \alpha \) and \( \beta \) and standard binary arithmetic operations \( \circ \in \{ +, -, \times, \div \} \), these assumptions take the form
\begin{equation*}
%    \label{eqn:fp}
    | \!\operatorname{fp}( \alpha \circ \beta ) - \alpha \circ \beta \,|
    \leq \epsilon_{\textup{mach}} \, | \alpha \circ \beta \,|.
\end{equation*}
Similar assumptions are also made for unary operations such as the square root,
\begin{equation*}
     | \!\operatorname{fp}(\sqrt{\alpha}) - \sqrt{\alpha} | \leq \epsilon_{\textup{mach}} | \sqrt{\alpha} \,|.
\end{equation*}
Assuming overflow and underflow do not occur, the above assumptions are satisfied for IEEE 754 floating point arithmetic \cite{ieee_19}.
Since the vast majority of modern computers use IEEE 754 floating point arithmetic, such assumptions are relatively safe.\footnote{It is worth noting, however, that the above bounds do not necessarily hold for other number systems.
Notable examples include Cray supercomputers prior to the mid 1990s as well as a number of other early computers \cite{higham_02}.
In fact, with the recent rise in low precision number formats and custom hardware acceleration methods, the above bounds cannot be universally assumed \cite{fasi_higham_mikaitis_pranesh_21}.}

Under the above assumptions, the accuracy of basic linear algebraic primitives can be bounded. 
For instance, for floating point vectors \( \vec{x} \), \( \vec{y} \), floating point number \( \alpha \), and floating point matrix \( \vec{A} \), reasonable implementations of basic linear algebraic tasks \cite{higham_02} will satisfy
\begin{align*}
%    \label{eqn:axpy}
    \| \!\operatorname{fp}(\vec{x} + \alpha \vec{y}) - (\vec{x} + \alpha \vec{y}) \| 
    &\leq \epsilon_{\textup{mach}} \, ( \|\vec{x}\| + 2 |\alpha | \|\vec{y}\| ) \\
%    \label{eqn:ip}
    \| \!\operatorname{fp}(\langle \vec{x}, \vec{y} \rangle ) - \langle \vec{x}, \vec{y} \rangle \| 
    &\leq \epsilon_{\textup{mach}} \,  n \, \|\vec{x}\| \|\vec{y}\| \\
%    \label{eqn:Ax}
    \| \!\operatorname{fp}(\vec{A} \vec{x}) - \vec{A} \vec{x} \| 
    &\leq \epsilon_{\textup{mach}}\, c \, \|\vec{A}\| \|\vec{x}\|.
\end{align*}
Here \( c \leq n^{3/2} \) is a dimensional constant depending on the method of matrix multiplication and the sparsity of \( \vec{A} \) which is often written in terms of the ratio of the norm of the absolute value of \( \vec{A} \) and the norm of \( \vec{A} \).
These results can then be applied to analyze linear algebra routines.

\section{Three term recurrences}

The Lanczos algorithm as well as the explicit polynomial methods \cref{alg:cheb_moments,alg:lanczos} from \cref{chap:QF} compute \( \vec{q}_{i+1} \) by a symmetric three term recurrence of the form
\begin{equation*}
    \vec{q}_{i+1} = \frac{1}{\beta_{i}} \left( \vec{A} \vec{q}_i - \alpha_i \vec{q}_i - \beta_{i-1} \vec{q}_{i-1}  \right).
\end{equation*}
In \cref{alg:moments,alg:cheb_moments} the coefficients are predetermined, whereas Lanczos chooses the coefficients adaptively in order to enforce orthogonality.
Regardless, in finite precision arithmetic, we will instead have a perturbed recurrence 
\begin{equation*}
    \vec{q}_{i+1} = \frac{1}{\beta_{i}} \left( \vec{A} \vec{q}_i - \alpha_i \vec{q}_i - \beta_{i-1} \vec{q}_{i-1}  \right) + \vec{f}_{i+1}
\end{equation*}
where \( \vec{f}_{i+1} \) accounts for local rounding errors made in the computation of \( \vec{q}_{i+1} \).
These simple arithmetic computations are all stable in the sense described above, so \( \vec{f}_{i+1} \) is small (on the order of \( \epsilon_{\textup{mach}} \|\vec{A}\| \)) relative to the involved quantities.


While \( \vec{q}_{i+1} = p_{i+1}\A \vec{v} \) in exact arithmetic, this is no longer the case in finite precision arithmetic. 
Indeed, the difference between \( \vec{q}_{i+1} \) and \( p_{i+1}\A\vec{v} \) depends on the \emph{associated polynomials} of the recurrence applied to the \( \vec{f}_{j} \)'s; see for instance \cite{meurant_06}.
In the case of the Chebyshev polynomials of the first kind, the associated polynomials are well behaved on \( [-1,1] \)\footnote{In fact, the associated polynomials are the Chebyshev polynomials of the second kind.}, so it can be shown that \( \vec{q}_{i+1} \approx p_{i+1}\A \vec{v} \).
Here ``\(\approx\)'' means the error has a polynomial dependence on \( k \) and a linear dependence on the machine precision, along with other reasonable dependence on the dimension and matrix norm.
This can be easily seen by performing an analysis similar to that in the proof of \cref{thm:LF_fp} found later in this section.

As such, the computed modified moments for \( \mu = \mu_{a,b}^T \) can be expected to be near to the true modified moments and \cref{thm:moments_inf_norm} can be expected to hold to close degree as long as \( \mathcal{I} \subset [a,b] \).
On the other hand, for different \( \{ \alpha_i \}_{i=0}^{\infty} \) and \( \{ \beta_i \}_{i=0}^{\infty} \), for instance those generated by the Lanczos algorithm, the associated polynomials may grow exponentially in \( \mathcal{I}  \) and the modified moments obtained from the finite precision computation may differ greatly from their exact arithmetic counterparts unless very high precision is used.
In fact, this situation includes \( \mu_{a,b}^T \) if \( a \) and \( b \) are not chosen so that \( \mathcal{I} \subset [a,b] \).

\subsection{The Lanczos algorithm}

In the case of Lanczos, the coefficients are computed adaptively and therefore depend on \( \vec{q}_{i-1} \), \( \vec{q}_i \), and \( \vec{q}_{i+1} \).
It is well known that even if the \( \vec{f}_j \)'s are small, the coefficients produced by Lanczos run in finite precision arithmetic may differ \emph{greatly} from what would be obtained in exact arithmetic and the Lanczos vectors \( \{ \vec{q}_j \}_{j=1}^{k+1} \) need not be orthogonal. 
Moreover, the tridiagonal matrix \( \vec{T} \) from the finite precision computation may have multiple ``ghost'' eigenvalues near certain eigenvalues of $\vec{A}$, even when the eigenvalues of \( \vec{T} \) would have been well separated in exact arithmetic.
In this sense, the algorithm is notoriously unstable, and such instabilities can appear even after only a few iterations.


A great deal is known about the Lanczos algorithm in finite precision arithmetic; see for instance \cite{greenbaum_97,meurant_strakos_06,meurant_06}.
In this section, we summarize the content needed to understanding why Lanczos-based methods for matrix functions still work in finite precision arithmetic. 
A fully rigorous and self-contained treatment of the topic would be very long and exceedingly technical, and such a treatment would not improve understanding of the big-picture (in fact, it might do exactly the opposite).
As such, we aim to provide an overview of existing theory which provides an intuitive understanding.

Throughout the next several sections the symbol ``\( \preceq \)'' suppresses absolute constants (e.g. \( 5 \), \( 12 \), and \( 1/10 \)) and higher order terms in the machine precision, \( \epsilon_{\textup{mach}} \) (e.g. \( \epsilon_{\textup{mach}}^2 \)).
The results in the literature typically state explicit constants for terms linear in \( \epsilon_{\textup{mach}} \), but such a precise analysis is not necessary to understand the intuition we wish to convey.

We will write the matrix-form of the perturbed Lanczos recurrence as
\begin{equation}
    \label{eqn:lanczos_factorization_fp}
    \vec{A} \vec{Q} = \vec{Q} \vec{T} + \beta_{k-1} \vec{q}_{k} \vec{e}_{k-1}^\rT + \vec{F}.
\end{equation}
We denote by \( \vec{R} \) the strictly upper triangular part of \( \vec{Q}^\rT\vec{Q} \); i.e. 
\begin{equation*}
    \vec{Q}^\rT \vec{Q} = \vec{R} + \vec{R}^\rT + \vec{D} .
\end{equation*}
It can then be shown that 
\begin{equation*}
    \label{eqn:lanczos_R_factorization_fp}
    \vec{T} \vec{R} =  \vec{R} \vec{T}  - \beta_{k-1}\vec{Q}^\rT \vec{q}_k\vec{e}_{k-1}^\rT + \vec{H}
\end{equation*}
for some upper triangular perturbation term \( \vec{H} \) which is \( \vec{0} \) in exact arithmetic; see for instance \cite[Equation 41]{paige_76}.
We will also denote by \( \eta \) the minimum value so that
\begin{equation*}
    \Lambda\mf{\vec{T}} \subset [\lmin - \eta, \lmax+\eta].
\end{equation*}
    

The first work which truly explained why the Lanczos algorithm was still useful as an eigenvalue algorithm was the PhD thesis \cite{paige_71} of Paige (and the technical reports leading up to the thesis).
The results in \cite{paige_71} were subsequently simplified and extended in \cite{paige_72,paige_76,paige_80}.
The main result we need is \cite[Theorem 1]{paige_76} which we have simplified to the needs of this chapter.
\begin{theorem}
    \label{thm:paige}
    Suppose the implementation of the Lanczos algorithm given in \cref{alg:lanczos} is run in finite precision arithmetic with machine precision \( \epsilon_{\textup{mach}} \). 
    Then, for \( i<k \), under some mild technical assumptions on \( \epsilon_{\textup{mach}} \),
    the following quantities 
    \begin{equation*}
        \| \vec{D} - \vec{I} \|
        ,\qquad
        \| \vec{F} \| 
        ,\qquad
        \| \vec{H} \| 
        ,\qquad
        \eta
    \end{equation*}
    are bounded by \( O(k^\alpha n^\beta \| \vec{A} \| \epsilon_{\textup{mach}}) \) for small constants \( \alpha,\beta \).
\end{theorem}

\begin{remark}
    The full version of \cref{thm:paige} from \cite{paige_76} is significantly more precise than our above statement might suggest.
    In particular, the above quantities (and several others) are each explicitly bounded to first order in \( \epsilon_{\textup{mach}} \), with the constants and the dependence on \( k \), \( n \), and \( \| \vec{A} \| \) stated explicitly for each quantity.
\end{remark}

In exact arithmetic, Lanczos-based approaches such as Gaussian quadrature and Lanczos-FA apply sufficiently low degree polynomials exactly.
We will now show that, even in finite precision arithmetic, these approaches apply (appropriately scaled) Chebyshev polynomials accurately. 
For functions which have a Chebyshev expansion with bounded coefficients, this implies that \cref{thm:moments_inf_norm,thm:poly_unif} holds to close degree.
Thus, Lanczos-based approaches should be expected to perform at least as well as explicit polynomial approaches for most reasonable functions \( f \).

For convenience, from this point to the end of the chapter, we will assume that \( \vec{A} \) has been shifted and scaled so that \( \Lambda \) and \( \Lambda\mf{\vec{T}} \) are each contained in \( [-1,1] \).

Recall that the Chebyshev polynomials of the first kind satisfy the recurrence
\begin{align*}
    \hspace{2em}T_i &= 2x T_{i-1} - T_{i-2}, & T_1&=x, & T_0&=1.\hspace{2em}
\intertext{and that the Chebyshev polynomials of the second kind satisfy the recurrence}
    U_i &= 2x U_{i-1} - U_{i-2}, & U_1&=2x, & U_0&=1.
\end{align*}
We have the following, well known, bound.
\begin{lemma}\label{thm:cheb_bounded}
    For all \( j \geq 0 \),
    \( \| T_j \|_{[-1,1]} \leq 1 \) and \( \| U_j \|_{[-1,1]} \leq j+1 \).
\end{lemma}

For notational brevity, for \( i\geq 1 \), introduce the vectors
\begin{equation*}
    \vec{t}_i = T_i\mf{\vec{A}}\vec{v}
    ,\qquad \tilde{\vec{t}}_i = T_i\mf{\vec{T}} \vec{e}_0
    ,\qquad \vec{\psi}_i = \vec{t}_i - \vec{Q} \tilde{\vec{t}}_i.
\end{equation*}


\section{Lanczos-FA}

To the best of our knowledge, the result in this section first appeared in \cite[Section 4]{druskin_knizhnerman_91}; see also \cite[Lemma 10]{musco_musco_sidford_18}.

\begin{theorem}\label{thm:LF_fp}
For all \( i =0,1,\ldots, k-1 \), 
\begin{equation*}
    \| T_i\mf{\vec{A}} \vec{v} - \vec{Q}T_i\mf{\vec{T}} \vec{e}_0 \|_2
    \leq k^2 \| \vec{F} \|.
%    \preceq k^{3/2} \epsilon_{\textup{mach}} \| \vec{A}\|_2.
\end{equation*}
\end{theorem}

\begin{proof}
Using \cref{eqn:lanczos_factorization_fp} and \cref{thm:tridiag_power}, observe that for \( i>1 \), the \( \vec{d}_i \) satisfy a perturbed three term recurrence
\begin{align*}
    \vec{\psi}_i 
    &= (2\vec{A} \vec{t}_{i-1} - \vec{t}_{i-2}) - (2\vec{Q}\vec{T} \tilde{\vec{t}}_{i-1} - \vec{Q}\tilde{\vec{t}}_{i-2})
    \\&= 2(\vec{A} \vec{t}_{i-1} - (\vec{A} \vec{Q} \tilde{\vec{t}}_{i-1} + \beta_k\vec{q}_{k-1} \vec{e}_{k-1}^\rT \tilde{\vec{t}}_{i-1} + \vec{F} \tilde{\vec{t}}_{i-1})) - \vec{\psi}_{i-2}
    \\&= 2\vec{A} \vec{\psi}_{i-1} - \vec{\psi}_{i-2} + 2 \vec{F} \tilde{\vec{t}}_{i-1}.
\end{align*}
%
By direct computation, we also have \( \vec{\psi}_0 = \vec{0} \) and \( \vec{\psi}_1 = \vec{F} \tilde{\vec{t}}_{0} \).
Then, it's easy to see that
\begin{equation*}
    \vec{d}_i = U_{i-1}\mf{\vec{A}} \vec{F} \tilde{\vec{t}}_{0} + 2 \sum_{j=1}^{i} U_{i-j-1}\mf{\vec{A}} \vec{F} \tilde{\vec{t}}_{j}.
\end{equation*}
    We now use \cref{thm:cheb_bounded} to obtain
\begin{equation*}
    \| \vec{d}_i \| \leq 2 \sum_{j=0}^{i-1} \| U_{i-j-1}\mf{\vec{A}} \| \|\vec{F}\| \| \tilde{\vec{t}}_{j} \|
    \leq 2 \sum_{j=0}^{i-1} (i-j) \| \vec{F} \|
    \leq 2i^2 \|\vec{F} \|. \qedhere
\end{equation*}
%The result follows from \cref{thm:paige}.
\end{proof}

\section{Gaussian quadrature}
\label{sec:GQ_fp}

The results in this section are summarized from \cite{knizhnerman_96}.
Interestingly, this work seems to be relatively unknown, despite its significance given the widespread use of Lanczos-based quadrature methods.
It is our hope that the resurfacing of these results will help assuage some of the hesitancy to use Lanczos-based quadrature methods without reorthogonalization.

We begin with several lemmas.
\begin{lemma}\label{thm:RTe0}
    For all \( i=0,1,\ldots,k-1 \),
    \begin{equation*}
        \| \vec{R} T_i\mf{\vec{T}} \vec{e}_0 \| 
        \leq k^2 \| \vec{H} \|.
%        \preceq k^3 \epsilon_{\textup{mach}} \|\vec{A}\|.
    \end{equation*}
\end{lemma}

\begin{proof}
    Write \( \vec{\Delta}_i =  \vec{R} T_i\mf{\vec{T}} \vec{e}_0 \).
    Using \cref{eqn:lanczos_R_factorization_fp} and \cref{thm:tridiag_power}, for \( i>1 \), analogous to \cref{thm:LF_fp}, the \( \vec{\Delta}_i \)    satisfy the perturbed three term recurrence
    \begin{align*}
        \vec{\Delta}_i
        &= 2 \vec{R}\vec{T}\tilde{\vec{t}}_{i-1} - \vec{R} \tilde{\vec{t}}_{i-2}
        \\&= 2(\vec{T}\vec{R} \tilde{\vec{t}}_{i-1} + (\beta_{k-1} \vec{Q}^\rT \vec{q}_k \vec{e}_{k-1}^\rT \tilde{\vec{t}}_{i-1} + \vec{H}  \tilde{\vec{t}}_{i-1}) - \vec{R} \vec{\Delta}_{i-2}
        \\&= 2 \vec{T} \vec{\Delta}_{i-1} - \vec{\Delta}_{i-2}  + 2\vec{H} \tilde{\vec{t}}_{j-1}
    \end{align*}
    Since \( \vec{R} \) is strictly upper triangular we have \( \vec{\Delta}_0 = \vec{0} \) and, by direct computation, \( \vec{\Delta}_1 = \vec{H} \vec{e}_0 \). 
    % 
    This implies that
    \begin{equation*}
        \vec{\Delta}_i 
        = U_{i-1}\mf{\vec{T}}  \vec{H} \vec{e}_0 + 2 \sum_{j=1}^{i-1} U_{i-j-1}\mf{\vec{T}} \vec{H} \tilde{\vec{t}}_{j}.
    \end{equation*}
    We again use \cref{thm:cheb_bounded} to obtain
    \begin{equation*}
        \| \vec{\Delta}_i \|
        \leq 2\sum_{j=0}^{i-1} \| U_{i-j-1}\mf{\vec{T}} \| \| \vec{H} \| \| \tilde{\vec{t}}_{j} \|
        \leq 2\sum_{j=0}^{i-1} (i-j) \| \vec{H} \|
        \leq 2 i^2 \| \vec{H} \|.
        \qedhere
    \end{equation*}
%    and the result follows from \cref{thm:paige}.
\end{proof}

We now state the main result.
\begin{theorem}
    \label{thm:QF_fp}
    For all \( i=0,1,\ldots, 2k-2 \)
    \begin{equation*}
        | \vec{v}^\rT T_i\mf{A} - \vec{e}_0^\rT T_i\mf{\vec{T}} \vec{e}_0 | \preceq k^2( \| \vec{F}\| + \| \vec{H} \|) + \| \vec{D} - \vec{I} \|.
    \end{equation*}
\end{theorem}

\begin{proof}
    Recall that the Chebyshev polynomial satisfy the identities
    \begin{equation*}
        T_{2i} = 2 (T_i)^2 - 1 
        ,\qquad
        T_{2i+1} = 2 T_i T_{i+1} - x. 
    \end{equation*}
    It therefore suffices to bound \( | \vec{v}^\rT T_i\mf{\vec{A}} T_j\mf{\vec{A}} \vec{v} - \vec{e}_0^\rT T_i\mf{\vec{T}} T_j\mf{\vec{T}} \vec{e}_0| \).

    By definition,
    \begin{equation*}
        \vec{t}_i^\rT \vec{t}_j
        = ( \tilde{\vec{t}}_i\vec{Q}^\rT + \vec{\psi}_i^\rT)(\vec{\psi}_j + \vec{Q} \tilde{\vec{t}}_j )
    \end{equation*}
    Thus,
    \begin{align*}
        | \vec{t}_i^\rT \vec{t}_j - \tilde{\vec{t}}_i^\rT \tilde{\vec{t}}_j |
%        | \vec{v}^\rT T_i\mf{\vec{A}} T_j\mf{\vec{A}} \vec{v}
 %       - \vec{e}_0^\rT T_i\mf{\vec{T}} T_j\mf{\vec{T}} \vec{e}_0|
        &\leq 
        | \tilde{\vec{t}}_i \vec{Q}^\rT \vec{Q} \tilde{\vec{t}}_j 
        - \tilde{\vec{t}}_i^\rT \tilde{\vec{t}}_j |
        + \|\vec{\psi}_i\| \|\vec{Q}\tilde{\vec{t}}_j\| + \|\vec{\psi}_j\|\|\vec{Q} \tilde{\vec{t}}_i\| + \|\vec{\psi}_i\| \| \vec{\psi}_j \|
    \end{align*}
    By definition of \( \vec{R} \),
    \begin{equation*}
        \tilde{\vec{t}}_i^\rT \vec{Q}^\rT \vec{Q} \tilde{\vec{t}}_j
        =  \tilde{\vec{t}}_i^\rT (\vec{R}+\vec{R}^\rT + \vec{I} + (\vec{D} - \vec{I}))  \tilde{\vec{t}}_j.
    \end{equation*}
    Thus, applying \cref{thm:RTe0,thm:paige,thm:cheb_bounded} we have 
    \begin{equation*}
        | \tilde{\vec{t}}_i \vec{Q}^\rT \vec{Q} \tilde{\vec{t}}_j 
        - \tilde{\vec{t}}_i^\rT \tilde{\vec{t}}_j |
        \leq \| \tilde{\vec{t}}_j \| \| \vec{R} \tilde{\vec{t}}_i \|
            + \| \tilde{\vec{t}}_i \| \| \vec{R} \tilde{\vec{t}}_j \|
            + \| \vec{D} - \vec{I} \|\| \tilde{\vec{t}}_i \|\| \tilde{\vec{t}}_j \|
        \preceq k^2 \| \vec{H} \| + \| \vec{D} - \vec{I} \|.
    \end{equation*}
    Now, observe that by \cref{thm:cheb_bounded}
    \begin{equation*}
        \| \vec{Q} \tilde{\vec{t}}_j \|
        = \| \vec{\psi}_j  - \vec{t}_j \|
        \leq \| \vec{\psi}_j \| + \| \vec{t}_j \|
        \preceq 1 + k^2 \| \vec{F} \|.%\epsilon_{\textup{mach}} \| \vec{A} \|.
    \end{equation*}
    Then, dropping higher order terms in \( \epsilon_{\textup{mach}} \),
    \begin{equation*}
        \| \tilde{\vec{t}}_i \|\| \vec{Q} \tilde{\vec{t}}_j \|
        \preceq k^2 \| \vec{F} \|
        \quad\text{and}\quad
        \| \tilde{\vec{t}}_j \|\| \vec{Q} \tilde{\vec{t}}_i \|
        \preceq k^2 \| \vec{F} \|.
    \end{equation*}
    The result follows by combining the above expressions.
\end{proof}

\begin{remark}
    In fact, when Lanczos is run for \( k \) iterations, \cref{thm:QF_fp} holds for \( i=0,1,\ldots, 2k-1 \).
    However, in the context of this thesis, the additional work required to prove this more general statement is not warranted.
    See \cite{knizhnerman_96} for details.
\end{remark}


\section{Backwards stability of the Lanczos algorithm}

The Lanczos method is clearly far from forward stable in the sense that the \( \vec{Q} \) and \( \vec{T} \) output are far from what would be produced in exact arithmetic.
The work of Greenbaum \cite{greenbaum_89} shows that the matrix \( \vec{T} \) in a perturbed Lanczos recurrence of the form \cref{eqn:lanczos_factorization_fp} can be viewed as the output of the Lanczos algorithm run in exact arithmetic on a certain ``nearby'' problem, provided the conclusions of \cref{thm:paige} are satisfied; i.e. it shows that Lanczos is backwards stable.
In particular, \cite{greenbaum_89} shows that if these conditions are satisfied, there exists a \( N\times N \) matrix \( \bar{\vec{A}} \) and vector \( \bar{\vec{v}} \) such that Lanczos run on \( \bar{\vec{A}} \), \( \bar{\vec{v}} \) in exact arithmetic for \( k \) steps produces \( \vec{T} \) (i.e., in the notation from \cref{chap:QF}, that \( \Psi_{\vec{T},\vec{e}_0} = \qq[g]{\Psi_{\bar{\vec{A}},\bar{\vec{v}}}}{2k-1} \)) and , that
\begin{enumerate}[label=(\roman*),nolistsep]
    \item  
        Eigenvalues of \( \bar{\vec{A}} \) are clustered near to those of \( \vec{A} \):
        for any \( j \in 0,1,\ldots, N-1 \), there exists \( i\in 0,1,\ldots, n-1   \) such that 
    \begin{align*}
        \lambda_j(\bar{\vec{A}}) \approx \lambda_i(\vec{A}).
    \end{align*}

    \item 
        The sum of squares of first components of eigenvectors corresponding to eigenvalues or clusters of eigenvalues of \( \bar{\vec{A}} \) approximately equal tot he squares of the projections of \( \vec{v} \) onto the eigenvectors of \( \vec{A} \):
        for an eigenvalue \( \lambda_i(\vec{A}) \)
        \begin{align*}
            w_i \approx \sum_{j\in S} \bar{w}_j
        \end{align*}
        where \( S_i \) is the set of indices such that \( \lambda_j(\bar{\vec{A}}) \approx \lambda_i(\vec{A}) \) for all \( j\in S \).
\end{enumerate}

Together, these conditions imply that
\begin{align}
    \label{eqn:finprec_wCESM}
    \Psi_{\vec{A},\vec{v}}
    \approx \Psi_{\bar{\vec{A}},\bar{\vec{v}}}.
\end{align}
%Of course, what the ``\( \approx \)'' in \labelcref{eqn:finprec_wCESM} means depend on the details of both the analysis in \cite{greenbaum_89} and implementation of the Lanczos algorithm used. 


\subsection{A new approach}

While the analysis of \cite{greenbaum_89} provides a backwards stability result, the proofs are highly technical.
We now show how the result of \cite{knizhnerman_96} implies a simple backwards stability result.

Towards this end, denote by $\{\sigma_m\}_{m=0}^{\infty}$ and $\{\rho_m\}_{m=0}^{\infty}$ the Chebyshev moments of $\Psi_{\vec{A},\vec{v}}$ and $\Psi_{\vec{T},\vec{e}_0}$ respectively; i.e. 
\begin{equation*}
    \sigma_m = \int T_m \,\d\Psi_{\vec{A},\vec{v}}
    \quad\text{and}\quad
    \rho_m = \int T_m \,\d\Psi_{\vec{T},\vec{e}_0}.
\end{equation*}
Now define the distribution function
\begin{equation*}
    \widehat{\Psi}
    = \Psi_{\vec{A},\vec{v}} + \Xi
    ,\quad \text{where} \quad
    \frac{\d\Xi}{\d{x}} = \frac{\d\mu_{-1,1}^T}{\d{x}} \sum_{m=0}^{2k-1} (\rho_m - \sigma_m) T_m.
\end{equation*}
By construction, for $m=1,2,\ldots, 2k-1$,
\begin{align*}
    \int_{-1}^{1} T_m \widehat{\Psi}
    &= \int_{-1}^{1} T_m \,\d\Psi_{\vec{A},\vec{v}} + \int T_m \, \d\Xi
    \\&= \sigma_m + \sum_{\ell=0}^{2k-1}(\rho_\ell-\sigma_\ell) \int T_m T_\ell \,\d \mu_{-1,1}^T
    \\&= \sigma_m + (\rho_m - \sigma_m).
\end{align*}
Thus, the modified moments $\{\widehat{\sigma}_m\}_{m=0}^{\infty}$ of $\widehat{\Psi}$ satisfy
\begin{equation*}
    \widehat{\sigma}_m 
    = \int_{-1}^{1} T_m \,\d\widehat{\Psi}
    = 
    \begin{cases}
    \rho_m & m = 0,1,\ldots, 2k-1 \\
    \sigma_m & m = 2k, 2k+1, \ldots
    \end{cases}.
\end{equation*}


Now, observe that
\begin{equation*}
    \W(\Psi,\widehat{\Psi})
    = \int_{-1}^{1} |\Xi(x) | \,\d x
    \leq \sum_{m=0}^{2k-1} |\rho_m-\sigma_m| \int_{-1}^{1} \left| \int_{-1}^{x} T_m \,\d\mu_{-1,1}^T \right| \d{x}.
\end{equation*}
The Chebyshev polynomials are bounded by one on $[-1,1]$, so
\begin{align*}
    \left| \int_{-1}^{x} T_m \,\d\mu_{-1,1}^T \right|
    \leq \int_{-1}^{x} |T_m \,\d\mu_{-1,1}^T |
    \leq \int_{-1}^{1} \d\mu_{-1,1}^T = 1.
\end{align*}
Thus, assuming \( |\rho_m - \sigma_m | \leq \epsilon(k) \),
\begin{equation*}
    \W(\Psi,\widehat{\Psi})
    \leq \sum_{m=0}^{2k-1} 2 |\rho_m-\sigma_m|
    \leq 4k \epsilon(k).
\end{equation*}

The distribution function $\widehat{\Psi}$ is near to $\Psi_{\vec{A},\vec{v}}$ in the sense of Wasserstein distance, and $\vec{T}_k$ is produced when the Stieltjes algorithm is applied. 
However, there are two shortcomings.
First, the support of $\widehat{\Psi}$ is all of $[-1,1]$ as $\Xi$ is absolutely continuous on $(-1,1)$.
Second $\Xi$ is not necessarily increasing, and so $\widehat{\Psi}$ is not necessarily increasing. 



\section{CIF bounds Finite precision}
\label{sec:finite_precision}

While the tridiagonal matrix \( \vec{T} \) and the matrix \( \vec{Q} \) of Lanczos vectors produced in finite precision arithmetic may be very different from those produced in exact arithmetic, we now show that our error bounds, based on the \( \vec{T} \) and \( \vec{Q} \) actually produced by the finite precision computation, still hold to a close approximation.
First, we argue that \cref{thm:shifted_lanczos_equivalence} holds to a close degree provided \( \vec{F} \) is not too large.
Towards this end, note that we have the shifted perturbed recurrence,
\begin{equation}
    \label{eqn:shifted_lanczos_factorization_fp}
    ( \vec{A} - z \vec{I} ) \vec{Q}
    = \vec{Q} ( \vec{T} - z \vec{I} ) + \beta_{k-1} \vec{q}_{k} \vec{e}_{k-1}^\cT + \vec{F}.
\end{equation}

From \cref{eqn:shifted_lanczos_factorization_fp}, it is then clear that,
    \begin{equation*}
        (\vec{A}-z \vec{I}) \vec{Q} (\vec{T} - z \vec{I})^{-1} \vec{e}_0
        = \vec{Q} \vec{e}_1 + \beta_{k-1} \vec{q}_{k} \vec{e}_{k-1}^\cT (\vec{T} - z \vec{I})^{-1} \vec{e}_0 + \vec{F} (\vec{T} - z \vec{I})^{-1} \vec{e}_0.
\end{equation*}
This implies that \cref{thm:shifted_linear_system_error} also holds closely.
More specifically,
\begin{align*}
    \Res_k(z)
    &= \det(h_{w,z}(\vec{T})) \Res_k(w)
    + \vec{f}_k(w,z)
\\
    \err_k(z)
    &= \det(h_{w,z}(\vec{T}))  h_{w,z}(\vec{A}) 
    \err_k(w)
    + (\vec{A}-z\vec{I})^{-1} \vec{f}_k(w,z)
\end{align*}
where
\begin{equation*}
    \vec{f}_k(w,z) := 
    \vec{F} \left( (\vec{T} - z\vec{I})^{-1} - \det(h_{w,z}(\vec{T})) (\vec{T}-w\vec{I)}^{-1} \right)\vec{e}_0.
\end{equation*}


Using this we have,
\begin{equation*}
    \fA\vec{v} - \lan_k(f)
    = - \frac{1}{2\pi \ii} \oint_{\Gamma} f(z) \err_k(z) \d{z} - \frac{1}{2\pi \ii} \oint_{\Gamma} f(z) (\vec{A} - z\vec{I})^{-1} \vec{f}_k(w,z) \,\d{z}
\end{equation*}
which we may bound using the triangle inequality as
\begin{equation*}
    \| \fA\vec{v} - \lan_k(f) \|
    \leq \frac{1}{2\pi} \left\| \oint_{\Gamma} f(z) \err_k(z) \,\d{z} \right \| +  \frac{1}{2\pi} \left\|  \oint_{\Gamma} f(z) (\vec{A} - z\vec{I})^{-1} \vec{f}_k(w,z) \,\d{z} \right\|.
\end{equation*}
This expression differs from \cref{thm:err_int} only by the presence of the term involving \( \vec{f}_k(w,z) \) (and, of course, by the fact that \( \err_k(z) \) now denotes the error in the finite precision computation).
If we take \( \| \cdot \| \) as the \( (\vec{A}-w\vec{I})^2 \)-norm, then this additional term can be bounded by using that
\begin{align}
    \left\|  \oint_{\Gamma} f(z) (\vec{A} - z\vec{I})^{-1} \vec{f}_k(w,z) \,\d{z} \right\|
    &\leq \oint_{\Gamma} |f(z)|   \| (\vec{A}-w\vec{I})(\vec{A} - z\vec{I})^{-1} \|_2 \| \vec{f}_k(w,z) \|_2   |\d{z}| \nonumber
    \\*&\leq \oint_{\Gamma} |f(z)| \|h_{w,z}\|_{S_0} \| \vec{f}_k(w,z) \|_2  |\d{z}|. \label{eqn:a_posteriori_fp}
\end{align}

\begin{figure}[htb]
    \begin{center}
        \includegraphics{imgs/ch8_CIF.pdf} 
    \end{center}
    \caption[{\( \vec{A}^2 \)-norm error bounds for \( f(x) = \sqrt{x} \) where \( \vec{A} \) has \( n=50 \) eigenvalues spaced according to the model problem with \( \rho = 0.8 \) and \( \kappa = 10^{3} \)}]{% 
        \( \vec{A}^2 \)-norm error bounds for \( f(x) = \sqrt{x} \) where \( \vec{A} \) has \( n=50 \) eigenvalues spaced according to the model problem with \( \rho = 0.8 \) and \( \kappa = 10^{3} \). We take \( \Gamma \) as a circular contour of radius \( \lmax \) centered at \( \lmax \) and \( w=0 \). Computations without reorthogonalization are run in single precision arithmetic, and error bounds are computed using \cref{thm:err_int} using the finite precision quantities.
    \hspace{.25em}\emph{Legend}:
    Lanczos-FA error with  
    ({\protect\raisebox{0mm}{\protect\includegraphics[]{imgs/legend/l2.pdf}}}) and without 
    ({\protect\raisebox{0mm}{\protect\includegraphics[]{imgs/legend/l1.pdf}}}) reorthogonalization, 
    a priori bounds with \( S = S_i = \mathcal{I} \)
    ({\protect\raisebox{0mm}{\protect\includegraphics[]{imgs/legend/l3.pdf}}}), and
    a posteriori bounds obtained by using \cref{thm:err_int} with \( S = \mathcal{I} \)
    ({\protect\raisebox{0mm}{\protect\includegraphics[]{imgs/legend/l4.pdf}}}).
    \hspace{.25em}\emph{Takeaway}: The bounds in finite precision arithmetic are accurate until near the ultimately attainable accuracy. 
    }
    \label{fig:ch8_CIF}
\end{figure}

Note that \( 1/(2\pi) \) times \cref{eqn:a_posteriori_fp} can be viewed as an upper bound of the ultimate obtainable  accuracy of Lanczos-FA in finite precision after convergence.
 If the inequalities do not introduce too much slack, this upper bound will also produce a reasonable estimate.
Since \( \| \vec{F} \| \) is small, one may simply ignore the contribution of \cref{eqn:a_posteriori_fp}, at least provided the Lanczos-FA error is not near the final accuracy.
Finally, we have worked in the \( (\vec{A}-w\vec{I})^2 \) norm as it simplifies some of the analysis, but in principle, a similar approach could be used with other norms.
This is straightforward, but would involve bounding something other than \( \|h_{w,z}\|_{S_0} \).

\subsection{Numerical experiment}



To illustrate the point of above analysis, we use a setup similar to what was used to produce \cref{fig:ch7_sqrt_Anorm}.
However, we now use the model problem and run Lanczos without reorthgonalization.
We use the \( \vec{T} \) produced in finite precision arithmetic in our computation of the error bounds from \cref{thm:err_int} and report the results in \cref{fig:ch8_CIF}.
Note that we use \cref{thm:err_int} and therefore do not account for the roundoff term analyzed above. 
However, since this term can be expected to be on the order of machine precision, the absence of this term in our computed bound does not significantly impact the bounds until near the ultimately attainable accuracy of Lanczos-FA. 
In other words, \cref{thm:err_int} still holds until the Lanczos-FA error is small.




\iffalse
\TODO{

\section{Lanczos-OR error}
\note{TODO if possible}
Recall that, in exact arithmetic,
\begin{align*}
    \| \err_k \|_{\tilde{N}\A}^2
    = \vec{e}_0^\cT \tilde{M}(\widehat{\vec{T}}) \tilde{N}(\widehat{\vec{T}})^{-1} \tilde{M}(\widehat{\vec{T}}) \vec{e}_0
    - \vec{e}_0^\cT [\tilde{M}(\widehat{\vec{T}})]_{:k,:k} ([\tilde{N}(\widehat{\vec{T}})]_{:k,:k})^{-1} [\tilde{M}(\widehat{\vec{T}})]_{:k,:k} \vec{e}_0.
\end{align*}

The expression on the right depends entirely on \( \Psi_{\vec{A},\vec{v}} \) and  the entries of \( \widehat{\vec{T}} \) required to compute \( [\tilde{M}(\widehat{\vec{T}})]_{:k,:k} \) and \( [\tilde{N}(\widehat{\vec{T}})]_{:k,:k} \).
Thus, we see that this expression will be nearly identical for the finite precision computaiton with \( \vec{A} \) and \( \vec{v} \) and the exact arithmetic computation with \( \bar{\vec{A}} \) and \( \bar{\vec{v}} \).

That the expressions on the left and right are near requires a bit more work.
In particular, we have made use of the identities \( \vec{Q}^\cT \tilde{M}(\vec{A}) \vec{v} = \tilde{M}(\widehat{\vec{T}}) \vec{e}_0 \) and \( \vec{Q}^\cT \tilde{N}\A \vec{Q} = [\tilde{N}(\widehat{\vec{T}})]_{:k,:k} \).
}

\fi


\chapter{Outlook}


So, what's next for algorithms for computing expressions involving matrix functions?
While this is far too vague a question to provide any sort of definitive answer, I will gladly discuss several directions which I hope will be pursued further in the near future.
These topics are simply a collection of directions for future work I personally find interesting, and they should not be viewed as any sort of statement regarding the direction the field as a whole should move in.
Indeed, many important topics, such as algorithms for high performance computing and the use of mixed precision arithmetic \cite{mixed_21} are not discussed.


\section{Randomization}

It is now widely recognized that randomization is an extremely powerful algorithmic tool in numerical linear algebra \cite{halko_martinsson_tropp_11,martinsson_tropp_20}.
While a number of topics have ``matured'', the use of randomization in Krylov subspace methods and related algorithms remains ripe for further study.

A big question is how to compute a low-rank approximation to \( \fA \), given access to products with \( \vec{A} \).
Some progress on this question has been made, primarily with the end goal of estimating the spectral sum \( \tr(\fA) \) \cite{lin_17,gambhir_stathopoulos_orginos_17,saibaba_alexanderian_ipsen_17,li_zhu_21,chen_hallman_22}, but a general theoretical understanding of randomized Krylov subspace methods for approximating \( \fA \) is an open problem.
A natural starting point is the analysis of block Krylov subspace methods applied to a set of random vectors \cite{musco_musco_15,martinsson_tropp_20,tropp_21}.
However, block Lanczos methods are perhaps even more susceptible to the effects of finite precision arithmetic than the standard Lanczos methods, and not much is known theoretically about their behavior in finite precision arithmetic.


Another interesting question is how randomization can be used to speed up computations of \( \fA\vec{v} \).

For overdetermined linear systems \( \vec{A} \vec{x} = \vec{v} \), methods such as the randomized Kaczmarz algorithm \cite{strohmer_vershynin_08,needell_srebro_ward_14}, accelerated coordinate descent \cite{lee_sidford_13,allenzhu_qu_richtarik_yuan_16}, and stochastic heavy ball momentum \cite{bollapragada_chen_ward_22} can all outperform applying CG to the normal equations \( \vec{A}^\cT \vec{A} \vec{x} = \vec{A}^\cT \vec{v} \).
In fact, for positive definite systems, accelerated coordinate descent methods can outperform CG applied directly to the system of interest.

A natural way to extend these fast linear system solvers to matrix functions is by applying them to a proxy rational function whose individual terms are each positive definite linear systems.
This technique was used in \cite{jin_sidford_19} to obtain a fast algorithm for approximating products with the matrix sign function and related quantities. 
However, this approach treats each term in the proxy rational function as independent, despite the fact that there is significant shared structure.
From a theoretical perspective this is acceptable as long as the number of terms in the proxy rational function is logarithmic in the accuracy tolerance, which is typically the case \cite{gopal_trefethen_19}.
However, this is likely somewhat wasteful in practice, so it would be worthwhile to study how such ideas can be implemented more efficiently.

%Finally, it would be interesting to see whether randomization can be used for indefinite linear systems.
%A concrete goal would be to develop a randomized algorithm which can match the convergence bound for MINRES based on polynomial approximation on the union of two intervals.
%Such an algorithm could outperform the above methods for general least squares problems on such systems.



\section{Typicality}


Recall that \emph{typicality}, discussed in \cref{sec:trace_hist}, is essentially the physics version of concentration of quadratic trace estimators.
I find typicality fascinating for a number of reasons. 
First, typicality provides a physical meaning to quadratic trace estimators, which have become one of the most widely studied methods in randomized numerical linear algebra.
Second, the literature on typicality has a rich history, with the earliest works dating back nearly a century.
This not only means the popular opinion on typicality has evolved, but it makes typicality an interesting case study in the fragmentation of knowledge between disciplines.
While several review papers have been published recently in the physics literature \cite{goldstein_lebowitz_mastrodonato_tumulka_zanghi_10,jin_willsch_willsch_lagemann_michielsen_deraedt_21}, I believe a review from the perspective of numerical linear algebra would yield many interesting historical insights.
Indeed, applied mathematicians have seemingly overlooked several important lines of literature on this topic.

\section{Accessibility to non-experts}

As mentioned in the introduction, it is my sense that practitioner knowledge of Lanczos based methods for matrix functions is limited by the lack of resources providing easy to understand background for such methods. 
While I hope that this thesis provides a more accessible introduction to the topic, by nature, a thesis emphasizes the author's own work and only touches on the important work of others.
A more balanced treatment of methods for matrix functions, with a treatment of methods for non-symmetric problems as well as a further emphasis on the important case of linear systems would be of general interest.

Separately, I hope that easy-to-use black-box versions of some of the algorithms studied in this thesis are eventually implemented.
A natural starting point would be implementing the integral based bounds from \cref{chap:CIF} in such a way that they could be easily integrate into existing codes.
In order for such a tool to be truly black-box would require additional study into how to choose parameters such as the contour of integration.
However, even if some user input is required, such a tool would help ensure more efficient resource allocation.




\chapter{Notation and other reference sheets}
\label{chap:notation}

\section{Basic notation}

Here we provide a reference for some common notation.
The page number is the first page on which the notation is used.
In many cases, some of the parameters will be suppressed for notational convince.
for instance, while the \( i \)-th eigenvalue of a matrix \( \vec{B} \) is denoted \( \lambda_i(\vec{B}) \), we will often write \( \lambda_i \) for the \( i \)-th eigenvalue of \( \vec{A} \) 

\begin{center}
    \begin{tabularx}{\textwidth}{p{1.9cm}Xr}\toprule
     notation & description & page \\\midrule\endfirsthead
    \toprule 
    notation & description & page \\\midrule\endhead
    \midrule\multicolumn{3}{r}{\itshape continues on next page}\\\midrule\endfoot
    \bottomrule\endlastfoot
%        \( x \) & identity function \( x : t\mapsto t \) & \pageref{eqn:A} \\
%    \midrule
        \( \vec{A} \) & \( n\times n \) Hermitian matrix & \pageref{eqn:A} \\
        \( \lambda_i \), \( \lambda_i(\vec{A}) \) & \( i \)-th eigenvalue of \( \vec{A} \) & \pageref{eqn:A} \\
        \( \Lambda \), \( \Lambda(\vec{A}) \) & set of eigenvalues of \( \vec{A} \) & \pageref{eqn:A} \\
        \( \mathcal{I} \), \( \mathcal{I}(\vec{A}) \) & smallest interval containing \( \Lambda(\vec{A} \) & \pageref{eqn:unif_bound} \\
    \midrule
    %=================================
    \( \fA \) & matrix function & \pageref{def:fA} \\
    \( \tr(\fA) \) & spectral sum & \pageref{eqn:spectral_sum} \\
        \( \Phi \), \(\Phi_{\vec{A}} \) & cumulative empirical spectral measure (CESM) & \pageref{def:CESM} \\
        \( \Psi \), \( \Psi_{\vec{A},\vec{v}} \) & weighted CESM & \pageref{def:wCESM} \\
        \( \mathcal{K}_k \), \( \mathcal{K}_k(\vec{A},\vec{v}) \) & dimension \( k \) Krylov subspace & \pageref{def:Kk} \\
    \( \vec{Q},\vec{T} \) & Lanczos vectors and coefficients after \( k \) iterations & \pageref{eqn:lanczos_three_term} \\
        \( \Qhat,\widehat{\vec{T}} \) & Lanczos vectors and coefficients after completion & \pageref{eqn:lanczos_three_term_complete} \\
    % ================================
    \midrule
%        \( \EE \) & expectation \\
%        \( \PP \) & probability \\
%        \( \operatorname{Unif}(\cdot) \) & uniform probability distribution \\
        \( \bOne \) & indicator function & \pageref{eqn:A} \\
        \( \| g \|_S \) & supremum of \( g:\mathbb{C}\to\mathbb{C} \) on \( S\subset \mathbb{C} \) & \pageref{def:norm_sup} \\
        \( \| \cdot \| \) & norm induced by matrix with same eigenvectors as \( \vec{A} \) & \pageref{def:norm} \\
    \midrule
        \( \mu \) & non-negative unit-mass distribution function & \pageref{def:mu} \\
        \( \langle \cdot, \cdot \rangle_{\mu} \) & inner product induced by \( \mu \) & \pageref{eqn:mu_ip} \\
        \( p_i \) & degree \( i \) orthogonal polynomial of \( \mu \) & \pageref{def:OP}\\ 
        \( \theta_j^{(s+1)} \) & \( j \)-th zero of \( p_s \) & \pageref{thm:jacobi_eigen}\\
    \( \vec{M} \), \( \vec{M}(\nu) \) & Jacobi matrix for \( \mu \), \( \nu \) & \pageref{def:jacobi}\\
        \( \mu_{a,b}^T \) & Chebyshev distribution function on \( [a,b] \) & \pageref{def:muT}\\
        \( m_i \) & degree \( i \) modified moment with respect to \( \mu \) &\pageref{def:modified_moments}\\ 
    % ================================
    \midrule
        \( \ff[a]{f}{s} \) & degree \( s \) projection of \( f \) in \( \langle \cdot, \cdot \rangle_{\mu} \) & \pageref{def:fffs}\\
        \( \ff[i]{f}{s} \) & degree \( s \) interpolation of \( f \) at zeros of \( p_s \) & \pageref{def:fffs} \\
        \( \ff[d-a]{f}{s} \) & degree \( s \) damped projection of \( f \) in \( \langle \cdot, \cdot, \rangle_{\mu} \) & \pageref{def:fffs} \\
        \( \ff[d-i]{f}{s} \) & degree \( s \) damped interpolation of \( f \) at zeros of \( p_s \) & \pageref{def:fffs} \\
        \( \{ \rho_i \}_{i=0}^{s} \) & damping coefficients & \pageref{def:damping} \\
        \( \{ \rho_i^J \}_{i=0}^{s} \) & Jackson's damping coefficients & \pageref{def:jackson_coeffs} \\
    \midrule
        \( \vec{C}_{\mu\to\nu} \) & connection coefficient matrix & \pageref{def:connection_coeffs} \\
        \( \langle \,\cdot\, \rangle \) & average over \( \ell = 0, \ldots, \nv-1 \) & \pageref{eqn:ave_notation} \\
        \( \mathbb{S}^{n-1} \) & unit hypersphere on \( \mathbb{C}^n \) & \pageref{eqn:hypersphere} \\ 
    \midrule
        \( \mathsf{cg}_k \) & CG iterate at step \( k \) & \pageref{eqn:cg_def} \\
        \( \mathsf{mr}_k \) & MINRES iterate at step \( k \) & \pageref{eqn:mr_def} \\
        \( \mathsf{qmr}_k \) & QMR iterate at step \( k \) & \pageref{eqn:qmr_def} \\
        \( \lanopt_k(r,R) \) & Lanczos-OR iterate & \pageref{def:lanczos_or} \\
        \( \lan_k(f) \) & Lanczos-FA iterate & \pageref{def:lanczos_fa} \\
        \( \mathsf{sign-OR} \) & Lanczos-OR induced iterate to matrix sign & \pageref{def:sign_or} \\ 
    \midrule
        \( \err_k(z) \) & Lanczos-FA error at step \( k \) for for \( \vec{A}^{-1} \vec{v} \) & \pageref{def:err} \\ 
        \( \res_k(z) \) & Lanczos-FA residual at step \( k \) for for \( \vec{A}^{-1} \vec{v} \) & \pageref{def:err} \\ 
        \( h_{w,z} \) & \( h_{w,z} = (x-w)/(x-z) \) & \pageref{def:hwz} \\
        \( h_{z} \) & \( h_{z} = 1/(x-z) \) & \pageref{def:hwz} \\
        \( S, S_0, \ldots, S_{k-1} \) & \( \Lambda \subset  S \), \( \lambda_i(\vec{T}) \subset S_i \) & \pageref{thm:err_int} \\
\end{tabularx}
\end{center}

\section{Indexing for matrices}

Given a matrix \( \vec{B} \), we use \( [\vec{B}]_{r:r',c:c'} \) to denote the submatrix matrix of \( \vec{B} \) consisting of rows \( r \) up to (but not including) row \( r' \) and columns \( c \) up to  (but not including) \( c' \).
Thus, the dimension of \( [\vec{B}]_{r:r',c:c'} \) is \( (r'-r)\times (c'-c) \).
\emph{Indexing of matrices starts at zero}.
If any of these indices are equal to \( 0 \) or the corresponding max dimension of \( \vec{B} \), they may be omitted.
If \( r' = r+1  \) or \( c' = c+1 \), then we will simply write \( r \) or \( c \).

As an example, suppose
\begin{equation*}
    \vec{B} = 
    \begin{bmatrix}
        1 & 2 & 3 & 4 \\
        5 & 6 & 7 & 8 \\
        9 & 10 & 11 & 12
    \end{bmatrix}
\end{equation*}
Then
\begin{equation*}
    [\vec{B}]_{:,:2}
    = 
    \begin{bmatrix}
        1 & 2 \\
        5 & 6 \\
        9 & 10
    \end{bmatrix}
    ,\quad
    [\vec{B}]_{1,:} 
    = 
    \begin{bmatrix}5&6&7&8\end{bmatrix},
    \quad\text{and}\quad
    [\vec{B}]_{0,3} = 4.
\end{equation*}

I am sure many readers are wondering why I would use such a notation. 
Perhaps some are even thinking of \href{https://xkcd.com/927/}{xkcd number 927, Standards}.

\begin{center}
    \includegraphics[width=8cm]{imgs/standards.png}
\end{center}

While I do not expect this notation to become standard in linear algebra, it was chosen intentionally after much consideration. 
The two primary motivations are as follows:
\begin{itemize}
    \item Much of this thesis relies on the theory of orthogonal polynomials, and it is natural for orthogonal polynomials to be indexed by their degree (which starts at zero). 
    It then makes sense that matrices involving orthogonal polynomials are indexed in a way which matches the polynomials.
        
    \item This thesis is written to be as accessible as possible to practitioners, particularly those in physics and data science.
        Zero-indexed programming languages are more common than one-indexed languages in these fields. 
        In such languages, non-inclusive endpoints are idiomatic, so that the number of objects in a range is equal to the difference of the endpoints.\footnote{The reason many languages use such conventions is perhaps due in part to Dijkstra's 1982 letter, \emph{Why numbering should start at zero}, which advocates for indexing to start at zero and for ranges to include the start point but not the end point.}
        In fact, our notation is identical to that of Python/NumPy, the language of choice in many disciplines from these fields.
\end{itemize}

In my opinion, it would have been nice if sums were also indexed in a similar way. 
However, using a notation like \( \sum_{r\leq i<r'} \) was deemed too verbose, and modifying the standard notation \( \sum_{i=r}^{r'-1} \) would have caused too much confusion.


\section{The model problem}

The model problem \cite{strakos_91,strakos_greenbaum_92} is a standard class of problems used in the analysis of the finite precision behavior of Lanczos based algorithms, especially in the context of solving linear systems of equations.
This is because the exponential spacing of the eigenvalues is favorable to Lanczos based linear system solvers in exact arithmetic yet simultaneously causes the Lanczos algorithm to rapidly lose orthogonality in finite precision arithmetic.
The model problem is parameterized by the dimension \( n \), the condition number \( \kappa \), and a parameter \( \rho \) controlling the rate of growth of eigenvalues.
Specifically,
\begin{equation}
    \label{eqn:model_problem}
    \lambda_0 = 1
    ,\quad \lambda_{n-1} = \kappa
    ,\quad \lambda_i = \lambda_1 + \left( \frac{i}{n-1} \right) \cdot (\kappa -1) \cdot \rho^{n-i-1}
    ,\qquad i=1,\ldots, n-1.
\end{equation}


\section{Some basic properties}


\begin{lemma}
    \label[lemma]{thm:norm_exchange}
    Let \( \|\cdot\| \) be a norm induced by a positive definite matrix with the same eigenvectors as \( \vec{A} \). 
    Then
    \begin{equation*}
        \| g\mf{\vec{A}} \vec{v} \| \leq \| g\mf{\vec{A}} \|_2 \| \vec{v} \| .
    \end{equation*}
\end{lemma}

\begin{proof}
    Let \( \vec{B}^2 \) be the matrix inducing \( \| \cdot \| \).
    By assumption \( \vec{A} \) and \( \vec{B} \) commute, so
    \begin{equation*}
        \| g\mf{\vec{A}} \vec{v} \|
        = \| \vec{B} g\mf{\vec{A}} \vec{v}\|_2
        = \| g\mf{\vec{A}} \vec{B} \vec{v}\|_2
        \leq \| g\mf{\vec{A}} \|_2 \|\vec{B} \vec{v}\|_2
        \leq \| g\mf{\vec{A}} \|_2 \|\vec{v}\|.
        \qedhere
    \end{equation*}
\end{proof}

We now provide a number of useful facts about powers of tridiagonal matrices. 
To simplify our proofs, we recall the following fact.

\begin{lemma}
For any \( q>0 \) and \( k_0,k_q=0,1,\ldots, n-1 \),
\begin{equation*}
    [\vec{A}^q]_{k_0,k_q}
    = \sum_{k_1=0}^{n-1} \sum_{k_2=0}^{n-1} \cdots \sum_{k_{q-1}=0}^{n-1}
    [\vec{A}]_{k_0,k_1} [\vec{A}]_{k_1,k_2} \cdots [\vec{A}]_{k_{q-1},k_q}.
\end{equation*}
\end{lemma}

\begin{proof}
This is the definition of matrix multiplication applied \( q \) times.
\end{proof}

Note that if \( \vec{T} \) is tridiagonal, then \( [\vec{T}]_{k_\ell,k_{\ell+1}} = 0 \) whenever \( |k_{\ell}-k_{\ell+1}|>1 \).
Thus, the product 
\begin{equation}
    \label{eqn:tridiag_power_term}
    [\vec{T}]_{k_0,k_1} [\vec{T}]_{k_1,k_2} \cdots [\vec{T}]_{k_{q-1},k_q}
\end{equation}
is nonzero, if and only if \( |k_{\ell} - k_{\ell+1}|\leq 1 \) for all \( \ell \). 
Thus, assuming \cref{eqn:tridiag_power_term} is nonzero, we can view \( \{k_\ell\} \) as a walk on \( \{0,1,\ldots, n-1\} \), starting from \( k_0 \) and ending at \( k_q \), where, at each iteration \( \ell \), we stay put or move to an adjacent index. 
Clearly
\begin{equation*}
    |k_0-k_1| + |k_1-k_2| + \cdots + |k_{q-1} - k_{q}| \leq q.
\end{equation*}
In other words, the total distance moved during the walk is at most \( q \).

Using this perspective, we immediately find that powers of tridiagonal matrices are banded.
\begin{corollary}
    \label[corollary]{thm:tridiag_power}
    Suppose \( \vec{T} \) is a tridiagonal matrix and \( q \geq 0 \) an integer.
    Then, for all \( i,j \) with \( |i-j|>q \), 
    \begin{equation*}
        [\vec{T}^q]_{i,j} = 0.
    \end{equation*}
\end{corollary}

\begin{proof}
    Consider a nonzero term \cref{eqn:tridiag_power_term}.
    If \( |i-j|>q \), then it is not possible to move from \( i \) to \( j \) in \( q \) steps.
\end{proof}


More generally, we find that the entries of powers of a tridiagonal matrix depend only on nearby entries of the base tridiagonal matrix.
We consider the symmetric case for simplicity.
\begin{corollary}
    \label[corollary]{thm:tridiag_power_dependence}
    Suppose \( \vec{T} \) is a symmetric tridiagonal matrix and \( q \geq 0 \) an integer.
    Then, for any \( i,j \) with  \( |j-i|\leq q \),
    \begin{equation*}
        [\vec{T}^q]_{i,j}
    \end{equation*}
    is determined entirely by \( [\vec{T}]_{:k,:k} \), where \( k = \max(i,j) + \lfloor (q-|j-i|)/2 \rfloor \).
    In fact, if \( q-|j-i| \) is even, then there is no dependence on \( [\vec{T}]_{k-1,k-1} \). 
\end{corollary}

\begin{proof}
    Without loss of generality, we may assume \( i\leq j \).

    Consider a nonzero term \cref{eqn:tridiag_power_term}.
    We require at least \( j-i \) of our allocated \( q \) movements to move from \( i \) to \( j \).
    Since we must end at \( j \), we could move past \( j \) at most \( \lfloor (q-(j-i))/2 \rfloor \) indices before returning.

    If \( (q-(j-i))/2 \) is an integer, then when we reach the maximum point \( j+(q-(j-i))/2 \) we must immediately return towards \( j \). 
    Thus, \( k_\ell \neq k_{\ell+1} \) for any \( \ell \), and in particular, for \( \ell = j+(q-(j-i))/2 \).
\end{proof}


%\listofalgorithms
\begin{landscape}
\section{List of algorithms}
\begin{center}
\fontsize{10}{14}\selectfont
\renewcommand{\arraystretch}{.9}
    \begin{tabularx}{\linewidth}{p{5.3cm}Xr}
    \toprule
    name & description & reference \\
    \midrule
    \ref{alg:name:lanczos} & \nameref{alg:lanczos}  & \cref{alg:lanczos} \\
    \midrule
    \ref{alg:name:stieltjes_naive} & \nameref{alg:stieltjes_naive}  & \cref{alg:stieltjes_naive} \\
    \ref{alg:name:stieltjes} & \nameref{alg:stieltjes}  & \cref{alg:stieltjes} \\
    \midrule
    \ref{alg:name:moments} & \nameref{alg:moments}  & \cref{alg:moments} \\
    \ref{alg:name:cheb_moments} & \nameref{alg:cheb_moments}  & \cref{alg:cheb_moments} \\
    \ref{alg:name:connection_coeffs} & \nameref{alg:connection_coeffs}  & \cref{alg:connection_coeffs} \\
    \ref{alg:name:moments_cheb} & \nameref{alg:moments_cheb}  & \cref{alg:moments_cheb} \\
    \ref{alg:name:moments_lanc} & \nameref{alg:moments_lanc}  & \cref{alg:moments_lanc} \\
    \ref{alg:name:iq} & \nameref{alg:iq}  & \cref{alg:iq} \\
    \ref{alg:name:gq} & \nameref{alg:gq}  & \cref{alg:gq} \\
    \ref{alg:name:aq} & \nameref{alg:aq}  & \cref{alg:aq} \\
    \ref{alg:name:aaq} & \nameref{alg:aaq}  & \cref{alg:aaq} \\
    \ref{alg:name:protoalg} & \nameref{alg:protoalg}  & \cref{alg:protoalg} \\
    \midrule
    \ref{alg:name:LDL} & \nameref{alg:LDL}  & \cref{alg:LDL} \\
    \ref{alg:name:streaming_LDL} & \nameref{alg:streaming_LDL}  & \cref{alg:streaming_LDL} \\
    \ref{alg:name:streaming_banded_prod} & \nameref{alg:streaming_banded_prod}  & \cref{alg:streaming_banded_prod} \\
    \ref{alg:name:streaming_banded_inv} & \nameref{alg:streaming_banded_inv}  & \cref{alg:streaming_banded_inv} \\
    \ref{alg:name:streaming_tridiag_square} & \nameref{alg:streaming_tridiag_square}  & \cref{alg:streaming_tridiag_square} \\
    \ref{alg:name:get_poly} & \nameref{alg:get_poly}  & \cref{alg:get_poly} \\
    \ref{alg:name:rational_lanczos} & \nameref{alg:rational_lanczos}  & \cref{alg:rational_lanczos} \\
    \ref{alg:name:lanczos_OR} & \nameref{alg:lanczos_OR}  & \cref{alg:lanczos_OR} \\
%    \ref{alg:name:} & \ref{alg:}  & \cref{alg:} \\
    \bottomrule
\end{tabularx}
\end{center}
\end{landscape}


\section{List of figures}
%\listoffiguresp
\makeatletter
{
\fontsize{11}{13}\selectfont
\@starttoc{lof}% Print List of Figures
}
\makeatother



\printbibliography[heading=bibnumbered]

\end{document}
